% 鉄緑風プリント テンプレート（非公式）  LuaLaTeX で2回コンパイルする
\documentclass[paper=b5j, fontsize=8pt, fleqn, twoside]{jlreq}

%% ---- フォント：ヒラギノがあれば使い、なければ原ノ味 ----
\usepackage[no-math]{fontspec}
\newif\ifTRhiragino
\IfFontExistsTF{HiraMinProN-W3}{\TRhiraginotrue}{\TRhiraginofalse}
\ifdefined\TRforceharano \TRhiraginofalse \fi % 原ノ味で確かめたいとき
\ifTRhiragino
  \usepackage[hiragino-pron,deluxe,match]{luatexja-preset}
  \setsansjfont{HiraginoSans-W3}[BoldFont=HiraginoSans-W6]
  \setsansfont{HiraginoSans-W3}[BoldFont=HiraginoSans-W6]
\else
  \usepackage[haranoaji,deluxe,match]{luatexja-preset}
  \setsansfont{HaranoAjiGothic-Regular}[BoldFont=HaranoAjiGothic-Bold]
\fi
\DeclareRobustCommand{\TRgt}{\sffamily\gtfamily}
\DeclareTextFontCommand{\textgt}{\TRgt} % 和文も英数字もゴシックにする
% ゴシックの太さを 0〜9 で選ぶ（ヒラギノ角ゴシック W0〜W9 相当）
\makeatletter
\NewDocumentCommand{\HiraginoKakuGothic}{m}{%
  \ifTRhiragino
    \csname TRhira@\number#1\endcsname
  \else
    \csname TRharano@\ifcase#1 xl\or xl\or l\or r\or r\or m\or b\or h\or h\or h\else r\fi\endcsname
  \fi}
\ifTRhiragino
  \@for\TR@w:=0,1,2,3,4,5,6,7,8,9\do{%
    \expandafter\newjfontfamily\csname TRhiraJ@\TR@w\endcsname{HiraginoSans-W\TR@w}%
    \expandafter\newfontfamily\csname TRhiraL@\TR@w\endcsname{HiraginoSans-W\TR@w}%
    \expandafter\edef\csname TRhira@\TR@w\endcsname{%
      \expandafter\noexpand\csname TRhiraL@\TR@w\endcsname
      \expandafter\noexpand\csname TRhiraJ@\TR@w\endcsname}}
\else
  \def\TR@harano#1#2{%
    \expandafter\newjfontfamily\csname TRharanoJ@#1\endcsname{HaranoAjiGothic-#2}%
    \expandafter\newfontfamily\csname TRharanoL@#1\endcsname{HaranoAjiGothic-#2}%
    \expandafter\edef\csname TRharano@#1\endcsname{%
      \expandafter\noexpand\csname TRharanoL@#1\endcsname
      \expandafter\noexpand\csname TRharanoJ@#1\endcsname}}
  \TR@harano{xl}{ExtraLight}\TR@harano{l}{Light}\TR@harano{r}{Regular}
  \TR@harano{m}{Medium}\TR@harano{b}{Bold}\TR@harano{h}{Heavy}
\fi
\makeatother

%% ---- パッケージ ----
\usepackage{
  luwa-ul,     % \namiKKAuto（答えの波線）、\underLineKK
  KKsymbols,   % \kakko{答}
  keisennote,  % \notefill（ノート罫）
  multicol,
  amsmath, amssymb, mathtools,
  fancyhdr, enumitem,
  tikz,
  tabularray,
}
\usetikzlibrary{calc,arrows.meta}
\makeatletter
\def\VDNT@xscaler{1-1.5/\linewidth}% ノート罫の端の点を段の内側に収める
\makeatother
\UseTblrLibrary{booktabs}
\usepackage[tikz]{multicolrule}
\SetMCRule{line-style=dense-solid-circles, width=0.8pt, extend-bot=-.69pt}
\usepackage[most]{tcolorbox}
\usepackage{calc, varwidth}

\makeatletter
\def\ascb@gtfamily{\TRgt}
\def\ascb@zw#1#2{#1\zw}
\newdimen\ascb@parindent@dimen
\ascb@parindent@dimen=\zw
\newcommand{\ascb@parindent}[1]{\setlength{\parindent}{#1}\relax}

%% 図・表：中央に置き、上下に空きを取る（段組の中では center の空きが小さいため）
\NewDocumentEnvironment{zuhyo}{}{\par\addvspace{.8\zw}\centering}{\par\addvspace{.8\zw}}

%% 参考問題：【参考問題】〈出典〉の白抜き罫枠
\DeclareTColorBox{simple}{ o m O{.5} O{} }%
{empty, left=2mm, right=2mm, top=-1mm, grow to left by=-#3pt/2, grow to right by=-#3pt/2,
attach boxed title to top left={xshift=\ascb@zw{1.2}{11pt}},
boxed title style={empty,left=-.5mm,right=-.5mm},
colframe=black, coltitle=black, coltext=black, breakable,
before upper={\ascb@parindent{\ascb@parindent@dimen}},
underlay unbroken={\draw[black,line width=#3pt](title.east) -- (title.east-|frame.east) -- (frame.south east) -- (frame.south west) -- (title.west-|frame.west) -- (title.west); },
underlay first={\draw[black,line width=#3pt](title.east) -- (title.east-|frame.east) -- (frame.south east) ;
\draw[black,line width=#3pt] (frame.south west) -- (title.west-|frame.west) -- (title.west); },
underlay middle={\draw[black,line width=#3pt](frame.north east) -- (frame.south east) ;
\draw[black,line width=#3pt](frame.south west) -- (frame.north west) ;},
underlay last={\draw[black,line width=#3pt](frame.north east) -- (frame.south east) -- (frame.south west) -- (frame.north west) ;},
fonttitle=\ascb@gtfamily, IfValueTF={#1}{title=【#2】〈#1〉}{title=【#2】}, #4}

%% KEY ボックス：黒ラベル＋灰色の見出し帯＋淡灰の本文（見出しの改行に対応）
\newlength{\len@ptbs@kk@D}\newlength{\lenn@ptbs@kk@D}\newlength{\myfontsize@ptbs@kk@D}
\newcommand{\titlelength@ptbs@kk@D}[1]{%
  \setlength{\myfontsize@ptbs@kk@D}{\f@size pt}%
  \def\titletext@ptbs@D{\TRgt\normalsize\selectfont#1}%
  \settowidth{\lenn@ptbs@kk@D}{\titletext@ptbs@D}%
  \setlength{\len@ptbs@kk@D}{\linewidth}%
  \addtolength{\len@ptbs@kk@D}{-\lenn@ptbs@kk@D}%
  \addtolength{\len@ptbs@kk@D}{-.3em}%
  \addtolength{\len@ptbs@kk@D}{-1pt}%
  \addtolength{\len@ptbs@kk@D}{4mm}%
}
\DeclareTColorBox{ptbs}{ m O{\phantom{A}} O{} }
{enhanced, breakable, boxsep=0mm,
arc=.5mm, bottom=2mm, top=2mm, leftupper=4mm, rightupper=4mm,
colbacktitle=black, colframe=black,
coltitle=white, colback=black!10!white, boxrule=1pt,
lefttitle=.3em,
before upper={\ascb@parindent{\ascb@parindent@dimen}},
fonttitle=\TRgt\normalsize,
fontupper=\TRgt\normalsize,
title={\titlelength@ptbs@kk@D{#1}#1\kern.3\zw\kern1pt},
after title={\rlap{% 帯は右の余白まで伸ばす（Overfull を出さない）
  \tcbox[on line, boxsep=.25\myfontsize@ptbs@kk@D, boxrule=0pt,
  top=.1\myfontsize@ptbs@kk@D, bottom=.1\myfontsize@ptbs@kk@D,
  left=.5mm, right=.5mm,
  width=\len@ptbs@kk@D, colback=black!30!white, arc=.5mm]
  {\raisebox{.2ex}{\parbox{\len@ptbs@kk@D-.5\myfontsize@ptbs@kk@D-1mm}{\renewcommand{\baselinestretch}{.5}\selectfont#2}}}}%
}, #3}

%% 例題：点線＋角丸の内枠と、四隅の灰色の扇形
\DeclareTColorBox{ascolorbox4A}{ o m O{3} O{} }%
{enhanced, colback=white, colframe=white,
attach boxed title to top left={xshift=1cm,yshift=-\tcboxedtitleheight/2},
varwidth boxed title=0.85\linewidth, coltitle=black,
fonttitle=\ascb@gtfamily,
before skip=.5mm, after skip=.8mm,
before upper={\ascb@parindent{\ascb@parindent@dimen}},
enlarge top by=2mm, enlarge bottom by=2mm, grow to left by=-.25pt, grow to right by=-.25pt, breakable, sharp corners,
boxed title style={colback=white,left=-.6em,right=-.6em},
borderline={.75pt}{#3pt}{black,dotted},
underlay unbroken={\draw[black,line width=.5pt]
  (title.east|-frame.north east)--([xshift=-#3*4pt]frame.north east)
  arc [start angle=180, end angle=270, radius=#3*4pt]
  -- ([yshift=#3*4pt]frame.south east)
  arc [start angle=90, end angle=180, radius=#3*4pt]
  -- ([xshift=#3*4pt]frame.south west)
  arc [start angle=0, end angle=90, radius=#3*4pt]
  -- ([yshift=-#3*4pt]frame.north west)
  arc [start angle=270, end angle=360, radius=#3*4pt]
  -- (frame.north west-|title.west) ;
  \filldraw[fill=gray,draw=gray] (frame.north east) -- ++(0,-#3*3pt) arc [start angle=270, end angle=180, radius=#3*3pt] -- cycle ;
  \filldraw[fill=gray,draw=gray] (frame.north west) -- ++(#3*3pt,0) arc [start angle=0, end angle=-90, radius=#3*3pt] -- cycle ;
  \filldraw[fill=gray,draw=gray] (frame.south west) -- ++(0,#3*3pt) arc [start angle=90, end angle=0, radius=#3*3pt] -- cycle ;
  \filldraw[fill=gray,draw=gray] (frame.south east) -- ++(-#3*3pt,0) arc [start angle=180, end angle=90, radius=#3*3pt] -- cycle ;
},
underlay first={\draw[black,line width=.5pt]
  (title.east|-frame.north east)--([xshift=-#3*4pt]frame.north east)
  arc [start angle=180, end angle=270, radius=#3*4pt] -- (frame.south east) ;
  \draw[black,line width=.5pt]
  (frame.south west) -- ([yshift=-#3*4pt]frame.north west)
  arc [start angle=270, end angle=360, radius=#3*4pt]
  -- (frame.north west-|title.west) ;
  \filldraw[fill=gray,draw=gray] (frame.north east) -- ++(0,-#3*3pt) arc [start angle=270, end angle=180, radius=#3*3pt] -- cycle ;
  \filldraw[fill=gray,draw=gray] (frame.north west) -- ++(#3*3pt,0) arc [start angle=0, end angle=-90, radius=#3*3pt] -- cycle ;
},
underlay middle={\draw[black,line width=.5pt] (frame.north east)--(frame.south east) ;
  \draw[black,line width=.5pt] (frame.south west)--(frame.north west) ;
},
underlay last={\draw[black,line width=.5pt]
  (frame.north east) -- ([yshift=#3*4pt]frame.south east)
  arc [start angle=90, end angle=180, radius=#3*4pt]
  -- ([xshift=#3*4pt]frame.south west)
  arc [start angle=0, end angle=90, radius=#3*4pt]
  -- (frame.north west) ;
  \filldraw[fill=gray,draw=gray] (frame.south west) -- ++(0,#3*3pt) arc [start angle=90, end angle=0, radius=#3*3pt] -- cycle ;
  \filldraw[fill=gray,draw=gray] (frame.south east) -- ++(-#3*3pt,0) arc [start angle=180, end angle=90, radius=#3*3pt] -- cycle ;
},
IfValueTF={#1}{title=【#2】〈#1〉}{title=【#2】},#4}

%% ミニタイトル：左の矢じり＋灰色の下線
\tcbset{ascbox@ascolorbox/.style={after skip=1.5mm, before skip=1.5mm},
ascboxsizeset@ascolorbox/.style={top=0mm,bottom=0mm,right=-1mm,left=4.3mm,fontupper=\TRgt},
titleunderline@ascolorbox/.style={underlay pre={\draw[very thick,draw=gray] ([yshift=.7mm,xshift=4mm]frame.south west) -- ([yshift=.7mm]frame.south east);}}}
\DeclareTCBox{\ascboxZ}{ O{dart} s O{.6} s }%
{empty,ascbox@ascolorbox,ascboxsizeset@ascolorbox,tcbox width=auto limited,% 長い見出しは段幅で折り返す
IfBooleanTF={#4}{}{titleunderline@ascolorbox},
IfBooleanTF={#2}%
{underlay={\node[#1,thick,draw=black!40!white,fill=black!70!white,draw,inner sep=#3mm,anchor=west] at ([xshift=1.86pt,yshift=-.5\baselineskip]frame.north west) {};}}%
{underlay={\node[#1,thick,draw=black!70!white,fill=black!40!white,draw,inner sep=#3mm,anchor=west] at ([xshift=1.86pt,yshift=-.5\baselineskip]frame.north west) {};}}%
}

%% Data：黒地のラベル＋罫枠（テーマ・難易度など）
\newlength\TR@datalabelwd
\NewDocumentEnvironment{Data}{O{Data}}{%
  \settowidth\TR@datalabelwd{{\TRgt\normalsize #1}}%
  \begin{tcolorbox}[enhanced,sharp corners,colback=white,colframe=black,
    boxrule=.45pt,left=\dimexpr\TR@datalabelwd+5.2mm\relax,right=1.4mm,
    top=.8mm,bottom=.8mm,before skip=.4mm,after skip=1.6mm,fontupper=\TRgt,
    overlay={\fill[black!72] (frame.north west) rectangle
      ([xshift=\dimexpr\TR@datalabelwd+3.8mm\relax]frame.south west);
      \node[anchor=west,inner sep=0pt,text=white,font=\TRgt\normalsize]
        at ([xshift=1.9mm]frame.west) {#1};}]
}{\end{tcolorbox}}

%% 小問帯：斜線入りの灰色の帯に (1) などを置く
\NewDocumentCommand{\Shomon}{m}{%
  \par\addvspace{.6\zw}\noindent\begin{tikzpicture}[x=1mm,y=1mm]
    \path[use as bounding box] (0,0) rectangle (\linewidth,6.2mm);
    \shade[left color=black!54,right color=black!40] (0,.65) rectangle (9.2,5.55);
    \begin{scope}
      \clip (0,.65) rectangle (9.2,5.55);
      \fill[black!12] (.7,.65)--(1.8,.65)--(5.1,5.55)--(4.0,5.55)--cycle;
      \fill[black!45] (1.8,.65)--(2.7,.65)--(6.0,5.55)--(5.1,5.55)--cycle;
      \fill[black!88] (2.7,.65)--(4.0,.65)--(7.3,5.55)--(6.0,5.55)--cycle;
    \end{scope}
    \node[anchor=center,inner sep=0pt,font=\TRgt\bfseries\large]
      at (11.95,3.1) {\makebox[5.5mm][c]{(#1)}};
    \coordinate (TRribbonE) at (current bounding box.east);
    \shade[left color=black!28,right color=black!14]
      (14.7,.65)--($(TRribbonE)+(-3.5mm,-2.45mm)$)--(TRribbonE)--
      ($(TRribbonE)+(-3.5mm,2.45mm)$)--(14.7,5.55)--cycle;
    \fill[black!18] ($(TRribbonE)+(-5.2mm,-2.45mm)$)--($(TRribbonE)+(-1.7mm,0)$)--
      ($(TRribbonE)+(-5.2mm,2.45mm)$)--($(TRribbonE)+(-6.05mm,2.45mm)$)--
      ($(TRribbonE)+(-2.55mm,0)$)--($(TRribbonE)+(-6.05mm,-2.45mm)$)--cycle;
    \fill[black!47] ($(TRribbonE)+(-4.35mm,-2.45mm)$)--($(TRribbonE)+(-.85mm,0)$)--
      ($(TRribbonE)+(-4.35mm,2.45mm)$)--($(TRribbonE)+(-5.2mm,2.45mm)$)--
      ($(TRribbonE)+(-1.7mm,0)$)--($(TRribbonE)+(-5.2mm,-2.45mm)$)--cycle;
    \fill[black!88] ($(TRribbonE)+(-3.5mm,-2.45mm)$)--(TRribbonE)--
      ($(TRribbonE)+(-3.5mm,2.45mm)$)--($(TRribbonE)+(-4.35mm,2.45mm)$)--
      ($(TRribbonE)+(-.85mm,0)$)--($(TRribbonE)+(-4.35mm,-2.45mm)$)--cycle;
  \end{tikzpicture}\par\nobreak\vspace{.2mm}}

%% 公式ボックス：灰色のタイトル帯
\DeclareTColorBox{koushiki}{ m O{} }{enhanced,breakable,arc=.625mm,auto outer arc,
  colback=white,colframe=black!65,boxrule=.5pt,left=2mm,right=2mm,top=1.5mm,bottom=1.3mm,
  before skip=1.2mm,after skip=1.2mm,title={#1},titlerule=.35pt,
  fonttitle=\TRgt,coltitle=white,colbacktitle=black!47,#2}

%% メモ：角丸の淡い灰色（補足・確認事項）
\DeclareTColorBox{memo}{ O{} }{enhanced,breakable,boxrule=0pt,arc=1.6mm,
  colback=black!9,colframe=black!9,left=2.8mm,right=2.8mm,top=1.7mm,bottom=1.6mm,
  before skip=1.2mm,after skip=1.2mm,fontupper=\TRgt,#1}

%% 非公式の注記：最後の段の \notefill の直前に置く（ノート罫と重ならないよう下に空きを取る）
\NewDocumentCommand{\Hikoushiki}{}{{\footnotesize\HiraginoKakuGothic{3}本資料は鉄緑会公式のものではない。\par}\vspace{.6\zw}}

%% 柱（上罫＋左右の柱文字）とノンブル（―n―）
\newdimen\top@geom@TETSUMANE    \top@geom@TETSUMANE=20mm
\newdimen\bottom@geom@TETSUMANE \bottom@geom@TETSUMANE=20mm
\newdimen\left@geom@TETSUMANE   \left@geom@TETSUMANE=16mm
\newdimen\right@geom@TETSUMANE  \right@geom@TETSUMANE=16mm
\newcommand{\Rhead@TETSUMANE}[1]{\begin{tikzpicture}[remember picture, overlay]
  \draw[line width=.5pt] ([yshift=-\top@geom@TETSUMANE+3mm, xshift=-\left@geom@TETSUMANE] current page.north east) -- ([yshift=-\top@geom@TETSUMANE+3mm, xshift=\right@geom@TETSUMANE] current page.north west);
  \node[anchor=east,inner sep=0pt,yshift=-\top@geom@TETSUMANE+6mm, xshift=-\left@geom@TETSUMANE] at (current page.north east) {\footnotesize #1};
  \draw[line width=.5pt] ([yshift=\bottom@geom@TETSUMANE-3mm, xshift=-\left@geom@TETSUMANE] current page.south east) -- ([yshift=\bottom@geom@TETSUMANE-3mm, xshift=\right@geom@TETSUMANE] current page.south west);
  \node[anchor=center] at ([yshift=\bottom@geom@TETSUMANE-6mm, xshift=0.5*(-\left@geom@TETSUMANE+\right@geom@TETSUMANE)]current page.south) {\footnotesize\TRgt\symbol{"2015}\hspace*{1ex}\thepage\hspace*{1ex}\symbol{"2015}};
\end{tikzpicture}}
\newcommand{\Lhead@TETSUMANE}[1]{\begin{tikzpicture}[remember picture, overlay]
  \draw[line width=.5pt] ([yshift=-\top@geom@TETSUMANE+3mm, xshift=\left@geom@TETSUMANE] current page.north west) -- ([yshift=-\top@geom@TETSUMANE+3mm, xshift=-\right@geom@TETSUMANE] current page.north east);
  \node[anchor=west,inner sep=0pt,yshift=-\top@geom@TETSUMANE+6mm, xshift=\left@geom@TETSUMANE] at (current page.north west){\footnotesize #1};
  \draw[line width=.5pt] ([yshift=\bottom@geom@TETSUMANE-3mm, xshift=-\right@geom@TETSUMANE] current page.south east) -- ([yshift=\bottom@geom@TETSUMANE-3mm, xshift=\left@geom@TETSUMANE] current page.south west);
  \node[anchor=center] at ([yshift=\bottom@geom@TETSUMANE-6mm,xshift=0.5*(\left@geom@TETSUMANE-\right@geom@TETSUMANE)]current page.south) {\footnotesize\TRgt\symbol{"2015}\hspace*{1ex}\thepage\hspace*{1ex}\symbol{"2015}};
\end{tikzpicture}}
\usepackage[top=\top@geom@TETSUMANE,bottom=\bottom@geom@TETSUMANE,left=\left@geom@TETSUMANE,right=\right@geom@TETSUMANE,headheight=14pt]{geometry}
\newcommand{\RunningHead}[2]{% #1 = 偶数ページ（右柱）, #2 = 奇数ページ（左柱）
  \fancyhead[ER]{\Rhead@TETSUMANE{\TRgt #1}}%
  \fancyhead[OL]{\Lhead@TETSUMANE{\TRgt #2}}}
\pagestyle{fancy}
\renewcommand{\headrule}{}
\fancyhead{}
\fancyfoot{}
\makeatother

\RunningHead{数学　第0講　本書の使い方}{数学　第0講　本書の使い方}

\begin{document}

%% ---- 本文で使う短縮命令（本文の中で定義する） ----
\newcommand{\Kotae}[1]{\namiKKAuto{#1}\hspace{\zw}\textgt{……\kakko{答}}}
\newcommand{\Kou}[2]{\clearpage\RunningHead{数学　第#1講　#2}{数学　第#1講　#2}\label{kou#1}}
\newcommand{\Kai}{\ascboxZ{解答}}
\newcommand{\vv}[1]{\overrightarrow{\mathrm{#1}}}
\newcommand{\Pt}[1]{\mathrm{#1}}
\newcommand{\Mae}[1]{\textgt{第#1講}}
\newcommand{\QED}{{\unskip\nobreak\hfil\penalty50\hskip1em\null\nobreak\hfil$\blacksquare$\parfillskip=0pt\finalhyphendemerits=0\par}}

%% ==================== 表紙 ====================
\thispagestyle{empty}
\begin{center}
\vspace*{6mm}
\begin{tikzpicture}[x=1mm,y=1mm]
  \fill[black!80] (0,0) rectangle (148,3);
  \fill[black!45] (0,4.2) rectangle (148,5);
  \fill[black!20] (0,6.2) rectangle (148,6.6);
\end{tikzpicture}

\vspace{14mm}
{\HiraginoKakuGothic{4}\fontsize{13pt}{16pt}\selectfont 高校数学の入口から東大入試まで、一本の道で}

\vspace{6mm}
{\HiraginoKakuGothic{8}\fontsize{40pt}{48pt}\selectfont 東大理系数学}

\vspace{5mm}
{\HiraginoKakuGothic{8}\fontsize{40pt}{48pt}\selectfont 完全マスター}

\vspace{10mm}
{\HiraginoKakuGothic{3}\fontsize{11pt}{15pt}\selectfont
数学Ⅰ・Ⅱ・Ⅲ・A・B（数列）・C（ベクトル、平面上の曲線と複素数平面）\par
全43講＋導入講\quad 例題・参考問題すべてに詳しい解説つき\par}

\vspace{14mm}
\begin{tikzpicture}[x=1mm,y=1mm]
  \foreach \i/\t/\s/\g/\c in {0/第1部/数と式・関数の土台/12/black,1/第2部/数学Ⅱ・B/24/black,2/第3部/数学C/36/black,3/第4部/数学Ⅲ/60/white,4/第5部/東大実戦/80/white}{
    \fill[black!\g] (\i*30,0) rectangle (\i*30+28,22);
    \node[font=\TRgt\large,text=\c] at (\i*30+14,15) {\t};
    \node[font=\TRgt\footnotesize,text=\c,align=center,text width=26mm] at (\i*30+14,6.5) {\s};
  }
  \draw[line width=.8pt,-{Stealth[length=3mm]}] (0,-5) -- (148,-5);
  \node[font=\TRgt\footnotesize,anchor=north west] at (0,-6.5) {高校初歩};
  \node[font=\TRgt\footnotesize,anchor=north east] at (148,-6.5) {東大入試本番レベル};
\end{tikzpicture}

\vfill
{\HiraginoKakuGothic{3}\footnotesize
理由を読む\quad →\quad 解説を隠して思い出す\quad →\quad 問題で使う\par}
\vspace{4mm}
\begin{tikzpicture}[x=1mm,y=1mm]
  \fill[black!20] (0,0) rectangle (148,.4);
  \fill[black!45] (0,1.6) rectangle (148,2.4);
  \fill[black!80] (0,3.6) rectangle (148,6.6);
\end{tikzpicture}
\end{center}

%% ==================== 第0講 ====================
\clearpage
\label{kou0}
\begin{ascolorbox4A}[出典：オリジナル]{例題0　出発点を決める診断テスト}
  次の\kakko{1}〜\kakko{10}を、何も見ずに解け。1問2分以内を目安とし、解けなかった問題・2分を超えた問題に印を付けよ。\par
  \kakko{1}　$(x+2)(x-3)$ を展開せよ。\qquad
  \kakko{2}　2次方程式 $x^2-5x+6=0$ を解け。\par
  \kakko{3}　2次関数 $y=x^2-4x+1$ の最小値を求めよ。\qquad
  \kakko{4}　$\sin 30^\circ$ と $\cos 60^\circ$ の値を答えよ。\par
  \kakko{5}　異なる5人から3人を選ぶ選び方は何通りか。\qquad
  \kakko{6}　$\log_2 8$ の値を求めよ。\par
  \kakko{7}　$x^3$ を $x$ で微分せよ。\qquad
  \kakko{8}　定積分 $\displaystyle\int_0^1 x^2\,dx$ を計算せよ。\par
  \kakko{9}　$\vec{a}=(1,\ 2)$、$\vec{b}=(3,\ -1)$ の内積 $\vec{a}\cdot\vec{b}$ を求めよ。\qquad
  \kakko{10}　$\displaystyle\lim_{n\to\infty}\frac{1}{n}$ を求めよ。
\end{ascolorbox4A}

\begin{Data}
  テーマ：自分がどの講から始めればよいかを、5〜20分で決める\par
  難易度：A（教科書の基本）。全部できなくてよい。できない所が「始める場所」になる\par
  この講のゴール：本書の進め方・記号・言葉の約束を知り、迷わず第1講以降に進めるようになる
\end{Data}

\begin{multicols*}{2}
  \ascboxZ{診断テストの答えと、始める講}
  まず答え合わせをする。印の付いた問題のうち、\textgt{いちばん番号の小さい問題}の「始める講」から本書を始める。全問できた人は第12講から始めてよい。

  \begin{zuhyo}
  \begin{tblr}{
    colspec = {c c c},
    hline{1,Z} = {0.08em}, hline{2} = {0.05em},
    row{1} = {font=\TRgt},
  }
    問題 & 答え & 始める講 \\
    (1) & $x^2-x-6$ & 第1講 \\
    (2) & $x=2,\ 3$ & 第3講 \\
    (3) & $-3$ & 第4講 \\
    (4) & $\frac12,\ \frac12$ & 第7講 \\
    (5) & $10$ 通り & 第8講 \\
    (6) & $3$ & 第18講 \\
    (7) & $3x^2$ & 第19講 \\
    (8) & $\frac13$ & 第20講 \\
    (9) & $1$ & 第25講 \\
    (10) & $0$ & 第30講 \\
  \end{tblr}
  \end{zuhyo}

  \begin{ptbs}{注意}[診断に関係なく、必ず読む講]
    第6講（論理と証明の型）、第10・11講（整数）、第39講（東大答案の書き方）は、計算ができる人ほど読み飛ばして失点する。診断の結果に関係なく必ず読む。
  \end{ptbs}

  \ascboxZ{1つの講の進め方（3周法）}
  本書の各講は「例題 → Data → 解説 → 参考問題」の順に並んでいる。次の3周で進めると、読み直しの回数が最も少なくて済む。

  \textgt{1周目（理由を読む）}：例題を5分だけ自力で考える。わからなくてよい。次に解説を読み、各行について「なぜこの式変形をしてよいのか」を確認する。KEY ボックスはノートに書き写す。

  \textgt{2周目（隠して思い出す、翌日）}：解説を紙で隠し、例題を見て「最初の一手」と「その理由」を1行で書く。それから解き直す。書けなかった行だけ解説を見る。最後に、各講の終わりにある CHECK ボックスの問いに、声に出すか書いて答える。

  \textgt{3周目（問題で使う、3〜7日後）}：参考問題を解く。解けなければ、解答の中で使っている例題の小問（帯に書いた (1)(2) の番号）に戻る。

  \begin{memo}
    「わかった」と「解ける」は別物である。解説を読んで納得した直後は、誰でも解ける気がする。翌日に白紙から再現できたときに初めて「解ける」と判定する。
  \end{memo}

  \ascboxZ{難易度の記号}
  例題の小問と参考問題には、Data に次の難易度を書いた。

  \begin{zuhyo}
  \begin{tblr}{
    colspec = {c X[l]}, width=\linewidth,
    hline{1,Z} = {0.08em}, hline{2} = {0.05em},
    row{1} = {font=\TRgt},
  }
    記号 & 目安 \\
    A & 教科書の例題。ここが入口 \\
    B & 入試の基本〜標準。共通テストで必要 \\
    C & 東大の標準問題。合格にはここを確実に取る \\
    D & 東大のやや難〜難。部分点を取りにいく \\
  \end{tblr}
  \end{zuhyo}

  \ascboxZ{本書の範囲}
  本書は、数学Ⅰ・Ⅱ・Ⅲ・A、数学B の「数列」、数学C の「ベクトル」「平面上の曲線と複素数平面」を扱う。東京大学（理科）の出題範囲・試験時間・配点は、受験する年度の募集要項で必ず自分で確認すること。

  \ascboxZ{記号の約束}
  数学の記号は「短く書いた日本語」である。読めない記号があると、そこで思考が止まる。次の表の記号は、本書で説明なしに使う。

  \begin{zuhyo}
  \begin{tblr}{
    colspec = {c X[l]}, width=\linewidth,
    hline{1,Z} = {0.08em}, hline{2} = {0.05em},
    row{1} = {font=\TRgt},
    cells = {font=\small},
  }
    記号 & 読み方と意味 \\
    $x\in A$ & $x$ は集合 $A$ の要素である（$A$ に属する） \\
    $A\subset B$ & $A$ の要素はすべて $B$ の要素である \\
    $A\cap B$ & $A$ と $B$ の両方に属するもの（共通部分） \\
    $A\cup B$ & $A$ と $B$ の少なくとも一方に属するもの \\
    $\overline{A}$ & $A$ に属さないもの（補集合） \\
    $p\Longrightarrow q$ & $p$ ならば $q$ \\
    $p\Longleftrightarrow q$ & $p$ と $q$ は同値（行きも帰りも成り立つ） \\
    $\therefore$ & ゆえに（前の行から導かれる） \\
    $\because$ & なぜならば（理由を書く） \\
    $a\leqq b$ & $a$ は $b$ 以下（$a=b$ も含む） \\
    $a<b$ & $a$ は $b$ より小さい（$a=b$ は含まない） \\
    $|x|$ & $x$ の絶対値（数直線上の $0$ からの距離） \\
    $x\fallingdotseq 1.41$ & $x$ はおよそ $1.41$（近似値） \\
    $n!$ & $n$ の階乗 $n(n-1)\cdots 2\cdot 1$。$0!=1$ と約束する \\
    $\displaystyle\sum_{k=1}^{n}a_k$ & $a_1+a_2+\cdots+a_n$（第21講） \\
    $\displaystyle\lim_{n\to\infty}a_n$ & $n$ を限りなく大きくしたときに $a_n$ が近づく値（第30講） \\
  \end{tblr}
  \end{zuhyo}

  \ascboxZ{言葉の約束（暗黙の了解をなくす）}
  教科書や解答が説明なしに使う言葉を、ここで全部決めておく。

  \textgt{数の名前}：自然数は $1,\ 2,\ 3,\ \cdots$（$0$ を含めない）。整数は $\cdots,\ -1,\ 0,\ 1,\ \cdots$。有理数は「整数$\div$0でない整数」で表せる数。実数は数直線上の点で表せる数全体。有理数でない実数を無理数という（$\sqrt2$ や $\pi$）。

  \textgt{以上・以下・より大きい・未満}：「$3$ 以上」「$3$ 以下」は $3$ を含む。「$3$ より大きい」「$3$ 未満」は $3$ を含まない。

  \textgt{「〜とおく」}：新しい文字に名前を付ける宣言である。おいた文字が動ける範囲（例：$t=\sin\theta$ なら $-1\leqq t\leqq 1$）を必ず同時に書く。範囲を書き忘れるのが、入試で最も多い減点理由の1つである。

  \textgt{同値変形}：変形した後の式から、元の式に戻れる変形を同値変形という。両辺に同じ数を足す、$0$ でない数を掛ける、は同値変形である。両辺を2乗する、$0$ かもしれない式で両辺を割る、は同値変形とは限らない（第2講・第3講で詳しく扱う）。

  \textgt{必要条件・十分条件}：「必要条件で答えの候補を絞り、十分性（その候補が本当に条件を満たすこと）を確かめる」という流れは、東大の答案の背骨である（第6講・第39講）。

  \textgt{「求めよ」と「示せ」}：「求めよ」は答えの値や式を出し、どうやって出したかを書く。「示せ（証明せよ）」は、結論が正しい理由を、誰が読んでも納得できる順に書く。答えを知っていても、理由が書けなければ点にならない。

  \textgt{「一般性を失わない」}：文字の入れ替えで同じ議論になるとき、1つの場合に絞ってよい、という宣言。入れ替えても条件が変わらないことを確かめてから使う。

  \begin{memo}
    解説の中で「なぜ？」と思ったら、その行に印を付けて先に進んでよい。多くの「なぜ」は、同じ講の KEY かメモで答えている。それでも解決しない印だけを、後でまとめて質問すればよい。
  \end{memo}

  \notefill
\end{multicols*}

%% ==================== 目次 ====================
\clearpage
\ascboxZ{目次}

本書は5部・43講からなる。各講は、それより前の講の内容だけを使って書いてある。Data の「前提」に書いた講を飛ばしていないかを確かめてから進むこと。表の横線は部の切れ目である（第1部 数と式・関数の土台、第2部 数学Ⅱ・B、第3部 数学C、第4部 数学Ⅲ、第5部 東大実戦）。

\begin{zuhyo}
\begin{tblr}{
  colspec = {Q[c,wd=5mm] Q[l,wd=44mm] X[l] Q[r,wd=9mm]}, width=\linewidth,
  hline{1,Z} = {0.08em}, hline{2} = {0.05em},
  hline{13} = {0.03em},
  row{1} = {font=\TRgt},
  cells = {font=\footnotesize},
  rowsep = .9pt,
}
講 & 単元 & この講で身につけること & ページ \\
1 & 式の展開と因数分解 & たすき掛け、対称式を $x+y,\ xy$ で表す & \pageref{kou1} \\
2 & 実数・平方根・絶対値 & $\sqrt{a^2}=|a|$、絶対値の場合分け、二重根号 & \pageref{kou2} \\
3 & 2次方程式と判別式 & 解の公式の導出、判別式で解の個数を判定 & \pageref{kou3} \\
4 & 2次関数と最大・最小 & 平方完成、軸と定義域の位置で場合分け & \pageref{kou4} \\
5 & 2次不等式と解の配置 & 判別式・軸・端点の3点セット & \pageref{kou5} \\
6 & 集合と論理・証明の型 & 必要条件と十分条件、対偶、背理法 & \pageref{kou6} \\
7 & 三角比と図形の計量 & 正弦定理・余弦定理・面積・内接円 & \pageref{kou7} \\
8 & 場合の数 & 区別の有無の判断、順列・組合せ・仕切り & \pageref{kou8} \\
9 & 確率 & 余事象、反復試行、条件付き確率、最大値の確率 & \pageref{kou9} \\
10 & 整数(1) 約数・倍数・余り & 余りによる分類、合同式 & \pageref{kou10} \\
11 & 整数(2) 互除法と不定方程式 & 積の形・不等式で絞る・一次不定方程式 & \pageref{kou11} \\
12 & 式と証明・不等式 & 恒等式、相加平均と相乗平均、等号成立 & \pageref{kou12} \\
13 & 複素数と高次方程式 & 解と係数の関係、剰余の定理、$\omega$ & \pageref{kou13} \\
14 & 図形と方程式 & 直線・円・点と直線の距離・接線 & \pageref{kou14} \\
15 & 軌跡と領域・通過領域 & 逆像法（実数解をもつ条件）で通過範囲 & \pageref{kou15} \\
16 & 三角関数(1) 定義と加法定理 & 単位円による定義、加法定理の証明 & \pageref{kou16} \\
17 & 三角関数(2) 合成と応用 & 合成、$t=\sin\theta+\cos\theta$ の置き換え & \pageref{kou17} \\
18 & 指数関数・対数関数 & 指数の拡張、底の変換、桁数 & \pageref{kou18} \\
19 & 微分法（数学Ⅱ） & 導関数の定義、増減表、解の個数 & \pageref{kou19} \\
20 & 積分法（数学Ⅱ） & 面積、6分の1公式、絶対値の積分 & \pageref{kou20} \\
21 & 数列と和 & $\Sigma$ の計算、階差数列、差に分ける和 & \pageref{kou21} \\
22 & 漸化式 & 型の見分け方と解き方 & \pageref{kou22} \\
23 & 数学的帰納法 & 帰納法の型、2段の帰納法 & \pageref{kou23} \\
24 & 確率と漸化式 & 状態を決めて確率の漸化式を立てる & \pageref{kou24} \\
\end{tblr}
\end{zuhyo}

\begin{zuhyo}
\begin{tblr}{
  colspec = {Q[c,wd=5mm] Q[l,wd=44mm] X[l] Q[r,wd=9mm]}, width=\linewidth,
  hline{1,Z} = {0.08em}, hline{2} = {0.05em},
  hline{7,16} = {0.03em},
  row{1} = {font=\TRgt},
  cells = {font=\footnotesize},
  rowsep = .9pt,
}
講 & 単元 & この講で身につけること & ページ \\
25 & 平面ベクトル & 内積、一次独立と係数比較、交点の位置 & \pageref{kou25} \\
26 & 空間ベクトル & 空間の面積・平面・垂線の足 & \pageref{kou26} \\
27 & 複素数平面(1) & 極形式、回転、ド・モアブルの定理 & \pageref{kou27} \\
28 & 複素数平面(2) & 図形の式、$z+1/z$ が実数になる条件 & \pageref{kou28} \\
29 & 2次曲線・媒介変数・極座標 & 楕円・双曲線・放物線、極方程式 & \pageref{kou29} \\
30 & 数列の極限 & 不定形の処理、はさみうち、無限級数 & \pageref{kou30} \\
31 & 関数の極限と連続性 & $(\sin x)/x$、$e$、中間値の定理 & \pageref{kou31} \\
32 & 微分法（数学Ⅲ） & 積・商・合成関数、対数微分、陰関数 & \pageref{kou32} \\
33 & 微分の応用(1) グラフ & 凹凸・変曲点・漸近線、$e^\pi$ と $\pi^e$ & \pageref{kou33} \\
34 & 微分の応用(2) 不等式と方程式 & 差の関数、定数分離、平均値の定理 & \pageref{kou34} \\
35 & 積分法（数学Ⅲ）の計算 & 置換積分・部分積分・$\sin^n x$ の積分 & \pageref{kou35} \\
36 & 区分求積法と積分の評価 & 和の極限を積分に、積分で和を評価 & \pageref{kou36} \\
37 & 面積と体積 & 回転体、$y$ 軸まわりの回転 & \pageref{kou37} \\
38 & 体積の応用・曲線の長さ & 断面積で体積、サイクロイドの長さ & \pageref{kou38} \\
39 & 東大答案の書き方 & 必要十分、実数条件、答案の型 & \pageref{kou39} \\
40 & 東大型総合(1) 空間と通過範囲 & 空間図形を平面の存在条件に直す & \pageref{kou40} \\
41 & 東大型総合(2) 整数と論証 & 二項係数と素数、余りの周期 & \pageref{kou41} \\
42 & 東大型総合(3) 評価と極限 & 円周率の評価、積分の極限 & \pageref{kou42} \\
43 & 仕上げ：最終確認と本番戦略 & 総点検テスト、150分の使い方 & \pageref{kou43} \\
\end{tblr}

\end{zuhyo}



\ascboxZ{目的別の最短ルート}

\begin{zuhyo}
\begin{tblr}{
  colspec = {X[l] X[l]}, width=\linewidth,
  hline{1,Z} = {0.08em}, hline{2} = {0.05em},
  row{1} = {font=\TRgt},
  cells = {font=\footnotesize},
}
今の状態・目的 & 進め方 \\
高校数学をほぼ習っていない & 第1講から順にすべて。1日1講、3周法で進める \\
数学Ⅰ・A・Ⅱ・B は習ったが穴が多い & 診断テストで始める講を決める。第6・10・11講は必ず読む \\
数学Ⅲ を最短で身につけたい & 第16〜20講（三角・指数対数・微積の基本）を確認してから第30〜38講 \\
東大の過去問に入る前の総点検 & 全講の例題と KEY だけを1周 → 第39〜43講 → 参考問題 \\
入試まで3か月を切っている & 第43講の総点検テストで弱点の講を特定し、その講の例題だけを解き直す \\
\end{tblr}
\end{zuhyo}

\begin{memo}
  どのルートでも、1講を読み終えたら翌日に「例題を見て最初の一手を1行で言う」ことだけは省かない。読み返す時間より、思い出す時間のほうが定着に効く。
\end{memo}
%% ==================== 第1講 ====================
\Kou{1}{式の展開と因数分解}
\begin{ascolorbox4A}[出典：オリジナル]{例題1}
  \kakko{1}　$(2x-3)(x+4)$ を展開せよ。\par
  \kakko{2}　$6x^2+x-12$ を因数分解せよ。\par
  \kakko{3}　$x+y=3$、$xy=1$ のとき、$x^2+y^2$、$x^3+y^3$、$x^5+y^5$ の値をそれぞれ求めよ。\par
  \kakko{4}　$a^3+b^3+c^3-3abc$ を因数分解せよ。
\end{ascolorbox4A}

\begin{Data}
  テーマ：展開（かっこを外す）と因数分解（かけ算の形に戻す）、対称式\par
  難易度：(1)(2) A　(3) B　(4) B\quad 前提：なし（中学の文字式の計算）\par
  この講のゴール：どんな式でも迷わず展開・因数分解でき、$x+y$ と $xy$ から $x^n+y^n$ を作れる
\end{Data}

\begin{multicols*}{2}
  \begin{koushiki}{この講の言葉}
    \textgt{項}：$+$ や $-$ で区切られた1つ1つのかたまり。$2x^2-3x+1$ の項は $2x^2,\ -3x,\ 1$。\par
    \textgt{係数}：文字に掛かっている数。$-3x$ の係数は $-3$。\par
    \textgt{次数}：掛かっている文字の個数。$2x^2$ は2次。式の次数は、項の次数の最大値。\par
    \textgt{展開}：積（かけ算）の形を、和（足し算）の形に直すこと。\par
    \textgt{因数分解}：和の形を、積の形に直すこと。展開の逆。
  \end{koushiki}

  \Shomon{1}
  \ascboxZ{分配法則を2回使う}
  展開の道具は分配法則 $A(B+C)=AB+AC$ だけである。$(2x-3)$ を1つのかたまり $A$ と見て、$(x+4)$ に分配する。
  \[
  \begin{aligned}
  (2x-3)(x+4)&=(2x-3)x+(2x-3)\cdot 4\\
  &=2x^2-3x+8x-12
  \end{aligned}
  \]
  最後に、文字の部分が同じ項（同類項）をまとめる。$-3x+8x=5x$ だから、
  \[(2x-3)(x+4)=\Kotae{2x^2+5x-12}\]

  \begin{memo}
    慣れたら「前×前、外×外、内×内、後×後」の4つの積を書いて足す。$2x\cdot x=2x^2$、$2x\cdot4=8x$、$-3\cdot x=-3x$、$-3\cdot4=-12$。何をしているかは分配法則そのものである。
  \end{memo}

  \begin{koushiki}{乗法公式（展開の近道）}
    $(a+b)^2=a^2+2ab+b^2$\par
    $(a-b)^2=a^2-2ab+b^2$\par
    $(a+b)(a-b)=a^2-b^2$\par
    $(x+a)(x+b)=x^2+(a+b)x+ab$\par
    $(ax+b)(cx+d)=acx^2+(ad+bc)x+bd$\par
    $(a+b)^3=a^3+3a^2b+3ab^2+b^3$\par
    $(a+b)(a^2-ab+b^2)=a^3+b^3$
  \end{koushiki}
  どれも分配法則で展開すれば確かめられる。覚えていない公式は、その場で展開して作ればよい。

  \Shomon{2}
  \ascboxZ{たすき掛けで因数分解する}
  $6x^2+x-12$ を $(ax+b)(cx+d)$ の形にしたい。乗法公式の最後の行から、
  \[ac=6,\quad ad+bc=1,\quad bd=-12\]
  となる整数 $a,\ b,\ c,\ d$ を探せばよい。

  \textgt{手順1}：$ac=6$ となる組 $(a,\ c)$ を書く。$(1,\ 6)$、$(2,\ 3)$。\par
  \textgt{手順2}：$bd=-12$ となる組 $(b,\ d)$ を書く。\par
  \textgt{手順3}：$ad+bc$ が $1$ になる組を探す。下のように縦に並べ、斜めに掛けて足す（これが「たすき」の名前の由来）。

  \begin{zuhyo}
  \begin{tikzpicture}[x=1cm,y=.6cm,line width=.5pt]
    \node at (0,1) {$2$}; \node at (1.4,1) {$3$};
    \node at (0,0) {$3$}; \node at (1.4,0) {$-4$};
    \draw (.25,.85)--(1.1,.15); \draw (.25,.15)--(1.1,.85);
    \node[anchor=west] at (2,1) {$\to\ 3\cdot 3=9$};
    \node[anchor=west] at (2,0) {$\to\ 2\cdot(-4)=-8$};
    \draw (2.1,-.5)--(4.3,-.5);
    \node[anchor=west] at (2,-1) {和 $9+(-8)=1$};
  \end{tikzpicture}
  \end{zuhyo}

  左の列が $a=2,\ c=3$、右の列が $b=3,\ d=-4$ である。$ad+bc=2\cdot(-4)+3\cdot3=1$ となって条件を満たす。よって、
  \[6x^2+x-12=\Kotae{(2x+3)(3x-4)}\]

  \begin{memo}
    答えが出たら必ず展開して確かめる。$(2x+3)(3x-4)=6x^2-8x+9x-12=6x^2+x-12$。この検算は10秒で済み、失点を確実に防ぐ。
  \end{memo}

  \Shomon{3}
  \ascboxZ{対称式は $x+y$ と $xy$ で書ける}
  $x$ と $y$ を入れ替えても変わらない式を\textgt{対称式}という。$x^2+y^2$ や $x^3+y^3$ は対称式である。対称式は、必ず $x+y$ と $xy$ だけで表せることが知られている。そこで、$(x+y)^2=x^2+2xy+y^2$ を変形して、
  \[
  \begin{aligned}
  x^2+y^2&=(x+y)^2-2xy\\
  &=3^2-2\cdot1=\Kotae{7}
  \end{aligned}
  \]
  次に $(x+y)^3=x^3+3x^2y+3xy^2+y^3=x^3+y^3+3xy(x+y)$ を変形すると、
  \[
  \begin{aligned}
  x^3+y^3&=(x+y)^3-3xy(x+y)\\
  &=27-9=\Kotae{18}
  \end{aligned}
  \]

  \ascboxZ{$x^5+y^5$ は「掛けて余分を引く」}
  $x^2+y^2$ と $x^3+y^3$ を掛けると、
  \[(x^2+y^2)(x^3+y^3)=x^5+y^5+x^2y^3+x^3y^2\]
  余分な2項は $x^2y^2(x+y)$ とまとめられる。したがって、
  \[
  \begin{aligned}
  x^5+y^5&=(x^2+y^2)(x^3+y^3)-x^2y^2(x+y)\\
  &=7\cdot18-1^2\cdot3=\Kotae{123}
  \end{aligned}
  \]

  \begin{ptbs}{KEY}[$s_n=x^n+y^n$ の漸化式]
    $s_n=x^n+y^n$ とおくと、$s_{n+2}=(x+y)s_{n+1}-xy\,s_n$ が成り立つ。$x+y,\ xy$ が整数なら、すべての $s_n$ が整数になる。
  \end{ptbs}

  この式は $x^{n+2}+y^{n+2}=(x+y)(x^{n+1}+y^{n+1})-xy(x^n+y^n)$ の右辺を展開すれば確かめられる。本問では $s_{n+2}=3s_{n+1}-s_n$ で、$s_0=2,\ s_1=3$ から、
  \[s_2=7,\ s_3=18,\ s_4=47,\ s_5=123\]
  となり、上の答えと一致する。この「整数性」は東大の整数問題で何度も使われる（第23講）。

  \Shomon{4}
  \ascboxZ{3乗の和の公式を作る}
  $a^3+b^3=(a+b)^3-3ab(a+b)$ を使い、$a+b$ をかたまりで扱う。
  \[
  \begin{aligned}
  &a^3+b^3+c^3-3abc\\
  &=(a+b)^3-3ab(a+b)+c^3-3abc\\
  &=\{(a+b)^3+c^3\}-3ab(a+b+c)
  \end{aligned}
  \]
  $\{\ \}$ の中は $A^3+c^3$（$A=a+b$）の形なので、乗法公式の最後の行より $(A+c)(A^2-Ac+c^2)$ と因数分解できる。どちらの部分にも $a+b+c$ が現れるので、それでくくる。
  \[=(a+b+c)\{(a+b)^2-(a+b)c+c^2-3ab\}\]
  中かっこの中を展開して整理すると、
  \[=\Kotae{(a+b+c)(a^2+b^2+c^2-ab-bc-ca)}\]

  \begin{memo}
    因数分解の方針は次の順に試す。\par
    \kakko{1} 共通因数でくくる\quad \kakko{2} 公式の形を探す\quad \kakko{3} 1つの文字について整理する（次数の最も低い文字を選ぶ）\quad \kakko{4} かたまりを1文字とみる
  \end{memo}

  \begin{simple}[出典：オリジナル]{参考問題1}
    \kakko{1}　$x^4+x^2+1$ を因数分解せよ。\par
    \kakko{2}　$x^2+3xy+2y^2-x-3y-2$ を因数分解せよ。
  \end{simple}
  \Kai
  \kakko{1}　$x^4+x^2+1$ は、そのままでは公式の形にならない。$x^2$ を足して引き、平方の差 $A^2-B^2$ を作る。
  \[
  \begin{aligned}
  x^4+x^2+1&=(x^4+2x^2+1)-x^2\\
  &=(x^2+1)^2-x^2\\
  &=\Kotae{(x^2+x+1)(x^2-x+1)}
  \end{aligned}
  \]
  \kakko{2}　$x$ について整理し（$x$ の2次式とみる）、定数項にあたる部分を先に因数分解する。
  \[
  \begin{aligned}
  &x^2+(3y-1)x+(2y^2-3y-2)\\
  &=x^2+(3y-1)x+(2y+1)(y-2)
  \end{aligned}
  \]
  和が $3y-1$、積が $(2y+1)(y-2)$ となる2つの式は $2y+1$ と $y-2$ である（足すと $3y-1$）。よって、
  \[=\Kotae{(x+2y+1)(x+y-2)}\]

  \begin{ptbs}{CHECK}[翌日、解説を隠して答える（答えは本講の中にある）]
    \kakko{1}　たすき掛けで探す $a,\ b,\ c,\ d$ の3つの条件は何か。\par
    \kakko{2}　$x^3+y^3$ を $x+y$ と $xy$ で表せ。\par
    \kakko{3}　$s_n=x^n+y^n$ の漸化式を書け。
  \end{ptbs}

  \notefill
\end{multicols*}

%% ==================== 第2講 ====================
\Kou{2}{実数・平方根・絶対値}
\begin{ascolorbox4A}[出典：オリジナル]{例題2}
  \kakko{1}　$\dfrac{1}{\sqrt3-\sqrt2}$ の分母を有理化せよ。\par
  \kakko{2}　$\sqrt{(a-2)^2}$ を、$a$ の値で場合分けして根号を使わずに表せ。\par
  \kakko{3}　方程式 $|x-1|+|x-3|=4$ を解け。\par
  \kakko{4}　$\sqrt{7+4\sqrt3}$ を簡単にせよ。
\end{ascolorbox4A}

\begin{Data}
  テーマ：平方根の定義、絶対値の場合分け、有理化と二重根号\par
  難易度：(1) A　(2) A　(3) B　(4) B\quad 前提：第1講\par
  この講のゴール：「$\sqrt{\ }$ は0以上」「絶対値は距離」という定義から、すべての計算規則を説明できる
\end{Data}

\begin{multicols*}{2}
  \begin{koushiki}{平方根と絶対値の定義}
    \textgt{平方根}：$a\geqq0$ のとき、2乗して $a$ になる数のうち\textgt{0以上のもの}を $\sqrt a$ と書く。例えば $\sqrt9=3$ であり、$-3$ ではない。\par
    \textgt{絶対値}：数直線上で、$x$ と $0$ の距離を $|x|$ と書く。
    \[|x|=\begin{cases}x&(x\geqq0)\\-x&(x<0)\end{cases}\]
  \end{koushiki}
  「$x<0$ のとき $|x|=-x$」は、マイナスが付いているのに正の数である。例えば $x=-5$ なら $-x=5$。ここでつまずく人が多いので、必ず数を入れて確かめる。

  \Shomon{1}
  \ascboxZ{和と差の積で根号を消す}
  分母の $\sqrt3-\sqrt2$ に $\sqrt3+\sqrt2$ を掛けると、$(a-b)(a+b)=a^2-b^2$ より根号が消える。分母と分子に同じ数を掛けても値は変わらないので、
  \[
  \begin{aligned}
  \frac{1}{\sqrt3-\sqrt2}&=\frac{\sqrt3+\sqrt2}{(\sqrt3-\sqrt2)(\sqrt3+\sqrt2)}\\
  &=\frac{\sqrt3+\sqrt2}{3-2}=\Kotae{\sqrt3+\sqrt2}
  \end{aligned}
  \]

  \begin{memo}
    有理化する理由は、値の大きさが見えるようにするためである。$\sqrt3+\sqrt2\fallingdotseq1.732+1.414=3.146$ とすぐにわかる。
  \end{memo}

  \Shomon{2}
  \ascboxZ{$\sqrt{A^2}=|A|$ である}
  「$\sqrt{(a-2)^2}=a-2$」は誤りである。$a=0$ を入れると左辺は $\sqrt4=2$、右辺は $-2$ になる。$\sqrt{\ }$ は0以上という定義から、正しくは $\sqrt{A^2}=|A|$ である。
  \[\sqrt{(a-2)^2}=|a-2|\]
  あとは絶対値の定義に従い、中身 $a-2$ の符号で分ける。
  \[\sqrt{(a-2)^2}=\Kotae{\begin{cases}a-2&(a\geqq2)\\2-a&(a<2)\end{cases}}\]

  \begin{ptbs}{KEY}[$\sqrt{A^2}=|A|$。$A$ の符号がわからなければ場合分け]
    根号の中が2乗のときは、まず絶対値にする。符号が決まっているときだけ絶対値を外せる。
  \end{ptbs}

  \Shomon{3}
  \ascboxZ{絶対値の中身の符号で区間を分ける}
  $|x-1|$ は $x=1$ で、$|x-3|$ は $x=3$ で、中身の符号が変わる。そこで数直線を $x<1$、$1\leqq x<3$、$3\leqq x$ の3つに分ける。

  \begin{zuhyo}
  \begin{tblr}{
    colspec = {c c c c},
    hline{1,Z} = {0.08em}, hline{2} = {0.05em},
    row{1} = {font=\TRgt},
  }
    区間 & $|x-1|$ & $|x-3|$ & 左辺 \\
    $x<1$ & $1-x$ & $3-x$ & $4-2x$ \\
    $1\leqq x<3$ & $x-1$ & $3-x$ & $2$ \\
    $3\leqq x$ & $x-1$ & $x-3$ & $2x-4$ \\
  \end{tblr}
  \end{zuhyo}

  各区間で方程式を解き、\textgt{その区間に入っているか}を確かめる。\par
  $x<1$ のとき：$4-2x=4$ より $x=0$。$0<1$ なので適する。\par
  $1\leqq x<3$ のとき：$2=4$ となり、解なし。\par
  $3\leqq x$ のとき：$2x-4=4$ より $x=4$。$4\geqq3$ なので適する。\par
  以上より、
  \[x=\Kotae{0,\ 4}\]

  \begin{memo}
    図形的には、$|x-1|+|x-3|$ は「点 $x$ から点 $1$ までの距離と、点 $3$ までの距離の和」である。$1$ と $3$ の間にいれば和は常に $2$、外に出ると増える。和が $4$ になるのは両側に $1$ ずつはみ出た $x=0,\ 4$ で、計算と一致する。
  \end{memo}

  \Shomon{4}
  \ascboxZ{二重根号は $(\sqrt p+\sqrt q)^2$ の形を探す}
  $(\sqrt p+\sqrt q)^2=p+q+2\sqrt{pq}$ である（$p,\ q>0$）。これと $7+4\sqrt3=7+2\sqrt{12}$ を見比べ、
  \[p+q=7,\quad pq=12\]
  となる $p,\ q$ を探すと $p=4,\ q=3$ である。よって $7+4\sqrt3=(\sqrt4+\sqrt3)^2=(2+\sqrt3)^2$ となり、$2+\sqrt3>0$ だから、
  \[\sqrt{7+4\sqrt3}=\Kotae{2+\sqrt3}\]

  \begin{memo}
    $\sqrt{p+q-2\sqrt{pq}}$ のときは $\sqrt p-\sqrt q$ になるが、$p>q$ の順にしておく。$\sqrt{\ }$ の値は0以上だから、$\sqrt q-\sqrt p$（負の数）を答えてはいけない。
  \end{memo}

  \begin{simple}[出典：オリジナル]{参考問題2}
    $\sqrt5$ の整数部分を $a$、小数部分を $b$ とする。$a$、$b$ と $b^2+4b$ の値を求めよ。
  \end{simple}
  \Kai
  \textgt{整数部分}とは、その数を超えない最大の整数、\textgt{小数部分}とは「元の数 $-$ 整数部分」のことである。$4<5<9$ の各辺の正の平方根をとると $2<\sqrt5<3$ なので、
  \[a=\Kotae{2},\quad b=\sqrt5-a=\Kotae{\sqrt5-2}\]
  $b^2+4b$ は、$b+2=\sqrt5$ を使えるように変形すると計算が軽くなる。
  \[b^2+4b=(b+2)^2-4=5-4=\Kotae{1}\]
  \begin{memo}
    $b$ をそのまま代入しても答えは出るが、「目的の式を、わかっている形（ここでは $b+2$）で書き直す」という発想は東大でも繰り返し使う。
  \end{memo}

  \begin{ptbs}{CHECK}[翌日、解説を隠して答える（答えは本講の中にある）]
    \kakko{1}　$\sqrt{a^2}$ はいつも何に等しいか。\par
    \kakko{2}　絶対値の方程式で、各区間の解に必ず行う確認は何か。\par
    \kakko{3}　二重根号を外すとき探す $p,\ q$ の条件は何か。
  \end{ptbs}

  \notefill
\end{multicols*}

%% ==================== 第3講 ====================
\Kou{3}{2次方程式と判別式}
\begin{ascolorbox4A}[出典：オリジナル]{例題3}
  \kakko{1}　2次方程式 $2x^2-3x-1=0$ を解け。\par
  \kakko{2}　2次方程式 $ax^2+bx+c=0\ (a\neq0)$ の解の公式を、平方完成によって導け。\par
  \kakko{3}　$k$ を実数の定数とする。$x$ の2次方程式 $x^2+2(k-1)x+k+5=0$ が重解をもつような $k$ の値と、そのときの重解を求めよ。\par
  \kakko{4}　\kakko{3}の方程式が実数解をもたないような $k$ の値の範囲を求めよ。
\end{ascolorbox4A}

\begin{Data}
  テーマ：2次方程式の解法（因数分解・解の公式）と、判別式による実数解の個数の判定\par
  難易度：(1) A　(2) A　(3) B　(4) B\quad 前提：第1・2講\par
  この講のゴール：解の公式を自分で導け、判別式の符号から実数解の個数を即答できる
\end{Data}

\begin{multicols*}{2}
  \begin{koushiki}{方程式を「解く」とは}
    方程式を成り立たせる $x$ の値（解）を\textgt{すべて}求め、それ以外にないことも示すこと。そのために、同値変形（戻れる変形）だけで $x=\cdots$ の形まで変形する。
  \end{koushiki}

  因数分解できる方程式は、次の性質で解く。
  \[AB=0\iff A=0\ \text{または}\ B=0\]
  例えば $x^2-5x+6=0$ は $(x-2)(x-3)=0$ だから $x=2,\ 3$。

  \Shomon{1}
  \ascboxZ{因数分解できなければ解の公式}
  $2x^2-3x-1$ はたすき掛けで整数の組が見つからない。そこで解の公式（\kakko{2}で導く）を使う。$a=2,\ b=-3,\ c=-1$ を代入して、
  \[
  \begin{aligned}
  x&=\frac{-(-3)\pm\sqrt{(-3)^2-4\cdot2\cdot(-1)}}{2\cdot2}\\
  &=\Kotae{\frac{3\pm\sqrt{17}}{4}}
  \end{aligned}
  \]

  \begin{memo}
    $b$ が負のとき、$-b$ は正になる。$-(-3)=3$ を $-3$ と書く誤りが非常に多い。代入するときは必ずかっこを付ける。
  \end{memo}

  \Shomon{2}
  \ascboxZ{平方完成：$(x+\square)^2$ の形を作る}
  方針は、$x$ を1か所にまとめることである。$x^2$ と $x$ の2か所にある $x$ を、$(x+p)^2$ という1つのかたまりにする。\par
  \textgt{手順1}：両辺を $a\ (\neq0)$ で割る。
  \[x^2+\frac ba x+\frac ca=0\]
  \textgt{手順2}：$(x+p)^2=x^2+2px+p^2$ と見比べると、$2p=\frac ba$、つまり $p=\frac{b}{2a}$ にすればよい。$x^2+\frac bax=\left(x+\frac b{2a}\right)^2-\frac{b^2}{4a^2}$ だから、
  \[\left(x+\frac{b}{2a}\right)^2=\frac{b^2}{4a^2}-\frac ca=\frac{b^2-4ac}{4a^2}\]
  \textgt{手順3}：2乗して $\frac{b^2-4ac}{4a^2}$ になる数は $\pm\frac{\sqrt{b^2-4ac}}{2a}$ である（ただし $b^2-4ac\geqq0$ のとき）。よって、
  \[x+\frac b{2a}=\pm\frac{\sqrt{b^2-4ac}}{2a}\]
  したがって、
  \[x=\Kotae{\dfrac{-b\pm\sqrt{b^2-4ac}}{2a}}\]

  \begin{memo}
    $\sqrt{4a^2}=|2a|$ であり、$a<0$ なら $-2a$ になる。しかし前に $\pm$ が付いているので、$+$ と $-$ が入れ替わるだけで結果は同じになる。だから $2a$ と書いてよい。
  \end{memo}

  \Shomon{3}
  \ascboxZ{判別式 $D$ の符号で解の個数がわかる}
  解の公式の根号の中 $D=b^2-4ac$ を\textgt{判別式}という。

  \begin{koushiki}{判別式と実数解の個数}
    $D>0$ $\iff$ 異なる2つの実数解\par
    $D=0$ $\iff$ 1つの実数解（重解）$x=-\frac{b}{2a}$\par
    $D<0$ $\iff$ 実数解なし\par
    $b=2b'$ のときは $\frac D4=b'^2-ac$ で判定してよい。
  \end{koushiki}

  本問は $x$ の係数が $2(k-1)$ なので $b'=k-1$ として、
  \[
  \begin{aligned}
  \frac D4&=(k-1)^2-1\cdot(k+5)\\
  &=k^2-3k-4=(k-4)(k+1)
  \end{aligned}
  \]
  重解をもつ条件は $D=0$ だから $k=4$ または $k=-1$。重解は $x=-\frac{b}{2a}=-(k-1)$ である。\par
  $k=4$ のとき：方程式は $x^2+6x+9=(x+3)^2=0$。\par
  $k=-1$ のとき：方程式は $x^2-4x+4=(x-2)^2=0$。
  \[\Kotae{k=4\ \text{のとき}\ x=-3,\quad k=-1\ \text{のとき}\ x=2}\]

  \Shomon{4}
  \ascboxZ{$D<0$ を解く}
  実数解をもたない条件は $D<0$、つまり $(k-4)(k+1)<0$ である。積が負になるのは、2つの因数の符号が異なるときである。$k+1>k-4$ だから、$k+1>0$ かつ $k-4<0$ のときに限る（2次不等式の一般的な解き方は第5講）。
  \[\Kotae{-1<k<4}\]

  \begin{ptbs}{KEY}[「2次方程式」と書いてなければ $x^2$ の係数0を疑う]
    $x^2$ の係数に文字が入っていたら、その係数が0の場合（1次方程式）を別に調べる。判別式は2次方程式のときしか使えない。
  \end{ptbs}

  \begin{simple}[出典：オリジナル]{参考問題3}
    $x$ の方程式 $(k-1)x^2+2x+1=0$ がただ1つの実数解をもつような実数 $k$ の値を求めよ。
  \end{simple}
  \Kai
  「2次方程式」とは書いていないので、$x^2$ の係数 $k-1$ が $0$ の場合を分ける。\par
  $k=1$ のとき：方程式は $2x+1=0$ となり、ただ1つの解 $x=-\frac12$ をもつ。よって適する。\par
  $k\neq1$ のとき：2次方程式なので、ただ1つの実数解 $\iff$ 重解 $\iff D=0$ である。
  \[\frac D4=1^2-(k-1)\cdot1=2-k=0\]
  より $k=2$（このとき $x^2+2x+1=(x+1)^2=0$ で $x=-1$）。以上より、
  \[k=\Kotae{1,\ 2}\]

  \begin{memo}
    $k=1$ を落とすと、答えは $k=2$ だけになって減点される。東大でも「最高次の係数が0になる場合」の見落としはよく狙われる。
  \end{memo}

  \begin{ptbs}{CHECK}[翌日、解説を隠して答える（答えは本講の中にある）]
    \kakko{1}　解の公式を平方完成で導く3つの手順を言え。\par
    \kakko{2}　判別式の符号と実数解の個数の対応を言え。\par
    \kakko{3}　$x^2$ の係数に文字があるとき、最初に何を確かめるか。
  \end{ptbs}

  \notefill
\end{multicols*}
%% ==================== 第4講 ====================
\Kou{4}{2次関数と最大・最小}
\begin{ascolorbox4A}[出典：オリジナル]{例題4}
  \kakko{1}　2次関数 $y=2x^2-8x+5$ のグラフの頂点と軸を求めよ。\par
  $a$ を実数とする。$0\leqq x\leqq 2$ における $f(x)=x^2-2ax+3$ について、次の問いに答えよ。\par
  \kakko{2}　最小値 $m(a)$ を求めよ。\qquad
  \kakko{3}　最大値 $M(a)$ を求めよ。\par
  \kakko{4}　$0\leqq x\leqq2$ を満たすすべての $x$ に対して $f(x)>0$ となるような $a$ の値の範囲を求めよ。
\end{ascolorbox4A}

\begin{Data}
  テーマ：平方完成とグラフ、定義域が決まった2次関数の最大・最小（軸の位置で場合分け）\par
  難易度：(1) A　(2) B　(3) B　(4) C\quad 前提：第1〜3講\par
  この講のゴール：「グラフをかいて、軸と定義域の位置関係で分ける」を、どんな文字定数でも実行できる
\end{Data}

\begin{multicols*}{2}
  \begin{koushiki}{関数・定義域・最大値の意味}
    \textgt{関数}：$x$ の値を1つ決めると、$y$ の値がただ1つ決まる対応。$y=f(x)$ と書く。\par
    \textgt{定義域}：$x$ が動ける範囲。\par
    \textgt{最大値}：$y$ がとる値のうち最大のもの。実際にとる値でなければならない（「近づくが届かない値」は最大値ではない）。
  \end{koushiki}

  \Shomon{1}
  \ascboxZ{平方完成で頂点を読む}
  $y=a(x-p)^2+q$ の形にすると、グラフは $y=ax^2$ を $x$ 方向に $p$、$y$ 方向に $q$ だけ平行移動した放物線である。頂点は $(p,\ q)$、軸は直線 $x=p$ である。\par
  $x^2$ の係数2で、$x$ を含む項だけをくくる。
  \[
  \begin{aligned}
  y&=2(x^2-4x)+5\\
  &=2\{(x-2)^2-4\}+5\\
  &=2(x-2)^2-3
  \end{aligned}
  \]
  2行目は $x^2-4x=(x-2)^2-4$（$x$ の係数 $-4$ の半分 $-2$ を使い、その2乗の $4$ を引く）。よって、
  \[\text{頂点}\ \Kotae{(2,\ -3)},\quad \text{軸}\ \Kotae{\text{直線}\ x=2}\]

  \begin{memo}
    $(x-2)^2\geqq0$ なので、$y=2(x-2)^2-3\geqq-3$。等号は $x=2$ のとき。最小値 $-3$ がグラフを見なくても式から言える。これが平方完成の本当の目的である。
  \end{memo}

  \Shomon{2}
  \ascboxZ{軸 $x=a$ と定義域の位置で3通り}
  $f(x)=(x-a)^2-a^2+3$ より、グラフは下に凸で、軸は直線 $x=a$ である。下に凸の放物線は、軸から離れるほど値が大きい。したがって、定義域の中で\textgt{軸に最も近い $x$} で最小になる。軸が定義域の左・中・右のどこにあるかで、その $x$ が変わる。

  \begin{zuhyo}
  \begin{tikzpicture}[x=.38cm,y=.28cm,line width=.5pt]
    \foreach \sx/\ax/\lo/\hi/\lab in {0/-0.8/-1.5/1.8/{$a<0$},6.2/1.1/-1.5/3.7/{$0\leqq a\leqq2$},12.4/2.9/0.3/3.7/{$2<a$}}{
      \begin{scope}[shift={(\sx,0)}]
      \fill[black!15] (0,-.35) rectangle (2,0);
      \draw[->] (-1.6,0) -- (3.9,0);
      \draw[domain=\lo:\hi,samples=40] plot (\x,{.55*(\x-\ax)^2+.6});
      \draw[dashed] (\ax,0) -- (\ax,4.2);
      \node[font=\footnotesize] at (\ax,4.9) {軸};
      \node[font=\footnotesize,below] at (0,-.3) {$0$};
      \node[font=\footnotesize,below] at (2,-.3) {$2$};
      \node[font=\footnotesize] at (1.1,-2.6) {\lab};
      \end{scope}}
    \fill (0,{.55*(0+0.8)^2+.6}) circle (2.2pt);
    \fill (7.3,.6) circle (2.2pt);
    \fill (14.4,{.55*(2-2.9)^2+.6}) circle (2.2pt);
  \end{tikzpicture}
  \end{zuhyo}

  黒丸が最小値をとる点である。灰色の帯が定義域 $0\leqq x\leqq2$ を表す。

  \begin{zuhyo}
  \begin{tblr}{
    colspec = {c c c},
    hline{1,Z} = {0.08em}, hline{2} = {0.05em},
    row{1} = {font=\TRgt},
  }
    $a$ の範囲 & 最小をとる $x$ & $m(a)$ \\
    $a<0$ & $0$ & $f(0)=3$ \\
    $0\leqq a\leqq 2$ & $a$ & $f(a)=-a^2+3$ \\
    $2<a$ & $2$ & $f(2)=-4a+7$ \\
  \end{tblr}
  \end{zuhyo}

  以上より、
  \[m(a)=\Kotae{\begin{cases}3&(a<0)\\-a^2+3&(0\leqq a\leqq 2)\\-4a+7&(2<a)\end{cases}}\]

  \begin{ptbs}{KEY}[2次関数の最小値は「軸と定義域の位置関係」]
    下に凸なら、軸が定義域の左・中・右のどこにあるかで3通りに分ける。境目（$a=0,\ 2$）はどちらに含めてもよいが、全体で漏れと重なりが出ないようにする。
  \end{ptbs}

  \Shomon{3}
  \ascboxZ{最大値は「軸から遠い端」}
  下に凸の放物線では、最大値は定義域の両端 $x=0,\ 2$ のどちらかでとる。どちらが大きいかは、軸からの距離で決まる。軸 $x=a$ が定義域の中央 $x=1$ より左なら右端 $x=2$ が遠く、右なら左端 $x=0$ が遠い。

  \begin{koushiki}{下に凸の2次関数の最大値}
    定義域の両端のうち、軸から遠い方で最大になる。軸が定義域の中央にあるときは、両端で等しい。
  \end{koushiki}

  $f(0)=3$、$f(2)=-4a+7$ なので、
  \[M(a)=\Kotae{\begin{cases}-4a+7&(a\leqq 1)\\3&(1<a)\end{cases}}\]

  \begin{memo}
    $a=1$ のとき $f(0)=f(2)=3$ で両端が等しい。$-4\cdot1+7=3$ なので、どちらの式に含めても同じ値になる。
  \end{memo}

  \Shomon{4}
  \ascboxZ{「すべての $x$ で $>0$」は「最小値 $>0$」}
  $0\leqq x\leqq2$ のすべての $x$ で $f(x)>0$ となることは、その範囲での最小値 $m(a)$ が正であることと同値である。最小値が正なら他の値はそれ以上なので正、逆にすべてが正なら最小値も正だからである。\kakko{2}の場合分けごとに $m(a)>0$ を解く。\par
  $a<0$ のとき：$m(a)=3>0$ で常に成り立つ。よって $a<0$ はすべて適する。\par
  $0\leqq a\leqq2$ のとき：$-a^2+3>0\iff a^2<3\iff-\sqrt3<a<\sqrt3$。$0\leqq a\leqq2$ と合わせて $0\leqq a<\sqrt3$。\par
  $2<a$ のとき：$-4a+7>0\iff a<\frac74$。これは $2<a$ と同時に成り立たないので、適する $a$ はない。\par
  3つの場合を合わせて（「または」でつなぐ）、
  \[\Kotae{a<\sqrt3}\]

  \begin{memo}
    場合分けの各結果は、必ず「その場合の条件」と共通部分をとる。最後に全体を合わせるときは和集合（または）をとる。この2つの操作を取り違えると答えが壊れる。
  \end{memo}

  \begin{simple}[出典：オリジナル]{参考問題4}
    $a$ を正の定数とする。$0\leqq x\leqq a$ における $g(x)=-x^2+4x$ の最大値を求めよ。
  \end{simple}
  \Kai
  $g(x)=-(x-2)^2+4$ で、上に凸、軸は $x=2$。今度は定義域の右端 $a$ が動く。上に凸の放物線は軸に近いほど大きい。\par
  $0<a<2$ のとき：定義域全体が軸より左にあり、$x$ が増えると $g(x)$ も増える。最大値は $g(a)=-a^2+4a$。\par
  $2\leqq a$ のとき：軸 $x=2$ が定義域に入るので、最大値は $g(2)=4$。
  \[\text{最大値}\ \Kotae{\begin{cases}-a^2+4a&(0<a<2)\\4&(2\leqq a)\end{cases}}\]

  \begin{ptbs}{CHECK}[翌日、解説を隠して答える（答えは本講の中にある）]
    \kakko{1}　平方完成の本当の目的は何か。\par
    \kakko{2}　下に凸の2次関数の最小値・最大値をとる $x$ の決め方を言え。\par
    \kakko{3}　「すべての $x$ で $f(x)>0$」を何に言い換えるか。
  \end{ptbs}

  \notefill
\end{multicols*}

%% ==================== 第5講 ====================
\Kou{5}{2次不等式と解の配置}
\begin{ascolorbox4A}[出典：オリジナル]{例題5}
  \kakko{1}　2次不等式 $x^2-x-6<0$ を解け。\par
  \kakko{2}　すべての実数 $x$ に対して $x^2+ax+a+3>0$ が成り立つような実数 $a$ の値の範囲を求めよ。\par
  $a$ を実数とし、$x$ の2次方程式 $x^2-2ax+a+2=0$ を考える。\par
  \kakko{3}　この方程式が異なる2つの正の解をもつような $a$ の値の範囲を求めよ。\par
  \kakko{4}　この方程式の異なる2つの解が、ともに $0<x<3$ の範囲にあるような $a$ の値の範囲を求めよ。
\end{ascolorbox4A}

\begin{Data}
  テーマ：2次不等式をグラフで解く、解の配置（判別式・軸・端点の値）\par
  難易度：(1) A　(2) B　(3) B　(4) C\quad 前提：第3・4講\par
  この講のゴール：方程式の解の条件を「放物線と $x$ 軸の交わり方」に翻訳し、3つの条件で過不足なく表せる
\end{Data}

\begin{multicols*}{2}
  \Shomon{1}
  \ascboxZ{不等式は「グラフが $x$ 軸より下」と読む}
  $y=x^2-x-6=(x-3)(x+2)$ のグラフは下に凸で、$x$ 軸と $x=-2,\ 3$ で交わる。$x^2-x-6<0$ は「グラフが $x$ 軸より下にある $x$」の範囲である。下に凸なので、それは2つの交点の間である。

  \begin{zuhyo}
  \begin{tikzpicture}[x=.55cm,y=.17cm,line width=.5pt]
    \draw[->] (-3.4,0) -- (4.4,0) node[right] {$x$};
    \draw[domain=-3:4,samples=40] plot (\x,{(\x-3)*(\x+2)});
    \draw[line width=2pt,black!45] (-2,0) -- (3,0);
    \node[below left,font=\footnotesize] at (-2,0) {$-2$};
    \node[below right,font=\footnotesize] at (3,0) {$3$};
  \end{tikzpicture}
  \end{zuhyo}

  よって、
  \[\Kotae{-2<x<3}\]

  \begin{koushiki}{2次不等式の答え（$\alpha<\beta$、$x^2$ の係数が正）}
    $(x-\alpha)(x-\beta)<0\iff\alpha<x<\beta$\par
    $(x-\alpha)(x-\beta)>0\iff x<\alpha,\ \beta<x$\par
    $x^2$ の係数が負なら、両辺に $-1$ を掛けて正にしてから使う（不等号の向きが逆になる）。
  \end{koushiki}

  \Shomon{2}
  \ascboxZ{常に正 $\iff$ $x$ 軸と共有点なし}
  $y=x^2+ax+a+3$ は下に凸である。すべての $x$ で $y>0$ になるのは、グラフ全体が $x$ 軸より上にあるとき、つまり $x$ 軸と共有点をもたないときである。それは判別式 $D<0$ と同値である。
  \[D=a^2-4(a+3)=a^2-4a-12=(a-6)(a+2)\]
  $D<0$ を解いて、
  \[\Kotae{-2<a<6}\]

  \begin{memo}
    上に凸の放物線は、$x$ を大きくすればいつか負になるので「常に正」にはなりえない。$x^2$ の係数が文字のときは、正・0・負で場合分けが必要である（参考問題5）。
  \end{memo}

  \Shomon{3}
  \ascboxZ{解の配置は3点セット}
  $f(x)=x^2-2ax+a+2$ とおく。方程式 $f(x)=0$ の解は、放物線 $y=f(x)$ と $x$ 軸の交点の $x$ 座標である。「異なる2つの正の解」は、「放物線が $x>0$ の部分で $x$ 軸と2回交わる」と言い換えられる。

  \begin{zuhyo}
  \begin{tikzpicture}[x=.55cm,y=.3cm,line width=.5pt]
    \draw[->] (-1,0) -- (5,0) node[right] {$x$};
    \draw[->] (0,-2.6) -- (0,4.5) node[above] {$y$};
    \draw[domain=0:4.6,samples=40] plot (\x,{(\x-2.3)^2-2});
    \draw[dashed] (2.3,0) -- (2.3,-2);
    \node[below,font=\footnotesize] at (2.3,-2) {軸 $x=a$};
    \fill (0,{(2.3)^2-2}) circle (2pt);
    \node[left,font=\footnotesize] at (0,{(2.3)^2-2}) {$f(0)$};
    \node[below left,font=\footnotesize] at (0,0) {$\mathrm O$};
  \end{tikzpicture}
  \end{zuhyo}

  図のように交わるための条件は、次の3つが同時に成り立つことである。\par
  \textgt{(ア) 判別式}：$x$ 軸と2回交わるので $\frac D4>0$。\par
  \textgt{(イ) 軸}：2つの交点の真ん中に軸があるので、軸 $x=a$ は $x>0$ の部分にある。\par
  \textgt{(ウ) 端点の値}：$x=0$ で放物線は $x$ 軸より上なので $f(0)>0$。

  (ア) $\frac D4=a^2-(a+2)=(a-2)(a+1)>0$ より $a<-1,\ 2<a$。\par
  (イ) $a>0$。\par
  (ウ) $f(0)=a+2>0$ より $a>-2$。\par
  3つの共通部分をとって、
  \[\Kotae{a>2}\]

  \begin{ptbs}{KEY}[解の配置は「判別式・軸・端点」の3点セット]
    放物線の絵をかき、交点の位置を決めるのに必要な条件を3つ書く。1つでも欠けると、絵と違う形の放物線が条件を満たしてしまう。
  \end{ptbs}

  \begin{memo}
    3つのうち1つでも抜くとどうなるか確かめる。(イ)を抜くと $a=-1.5$ が条件(ア)(ウ)を満たすが、このとき軸 $x=-1.5$ は負で、2解とも負になる。だから3つすべて必要である。
  \end{memo}

  \Shomon{4}
  \ascboxZ{端点が2つになれば条件も1つ増える}
  2つの解がともに $0<x<3$ にあるには、放物線が $0<x<3$ の中で $x$ 軸と2回交わればよい。今度は両端 $x=0$ と $x=3$ の両方で、放物線が $x$ 軸より上にある必要がある。\par
  (ア) $\frac D4>0$：$a<-1,\ 2<a$\par
  (イ) 軸：$0<a<3$\par
  (ウ) $f(0)=a+2>0$：$a>-2$\par
  (エ) $f(3)=9-6a+a+2=11-5a>0$：$a<\frac{11}5$\par
  4つの共通部分は、$2<a$ と $a<\frac{11}5$ の重なりである。よって、
  \[\Kotae{2<a<\frac{11}{5}}\]

  \begin{simple}[出典：オリジナル]{参考問題5}
    すべての実数 $x$ に対して $ax^2+2x+a>0$ が成り立つような実数 $a$ の値の範囲を求めよ。
  \end{simple}
  \Kai
  $x^2$ の係数 $a$ に文字が入っているので、$a=0$ を分ける。\par
  $a=0$ のとき：不等式は $2x>0$ となり、$x=-1$ で成り立たない。不適。\par
  $a\neq0$ のとき：すべての $x$ で正になるには、下に凸（$a>0$）かつ $x$ 軸と共有点なし（$\frac D4=1-a^2<0$）が必要十分である。$1-a^2<0\iff a<-1,\ 1<a$。$a>0$ と合わせて $a>1$。
  \[\Kotae{a>1}\]

  \begin{ptbs}{CHECK}[翌日、解説を隠して答える（答えは本講の中にある）]
    \kakko{1}　解の配置の3点セットとは何か。\par
    \kakko{2}　3点のうち軸の条件を抜くと、どんな反例が出るか。\par
    \kakko{3}　「常に正」で $x^2$ の係数が文字のときの注意は何か。
  \end{ptbs}

  \notefill
\end{multicols*}

%% ==================== 第6講 ====================
\Kou{6}{集合と論理・証明の型}
\begin{ascolorbox4A}[出典：オリジナル]{例題6}
  \kakko{1}　実数 $x$ について、命題「$x=1\Longrightarrow x^2=1$」とその逆の真偽を答えよ。偽ならば反例を挙げよ。\par
  \kakko{2}　実数 $x$ について、「$x^2=x$」は「$x=1$」であるための何条件か。必要条件・十分条件・必要十分条件・どれでもない、から選べ。\par
  \kakko{3}　整数 $n$ について、「$n^2$ が偶数ならば $n$ は偶数である」ことを証明せよ。\par
  \kakko{4}　$\sqrt2$ が無理数であることを証明せよ。\par
  \kakko{5}　整数 $a,\ b,\ c$ が $a^2+b^2=c^2$ を満たすとき、$a,\ b$ の少なくとも一方は偶数であることを証明せよ。
\end{ascolorbox4A}

\begin{Data}
  テーマ：命題と集合、必要条件・十分条件、否定、対偶による証明、背理法\par
  難易度：(1) A　(2) B　(3) B　(4) B　(5) C\quad 前提：第2講\par
  この講のゴール：東大の答案で必要になる「論理の言葉」を、意味から正しく使える
\end{Data}

\begin{multicols*}{2}
  \begin{koushiki}{この講の言葉}
    \textgt{命題}：正しいか正しくないかがはっきり決まる文や式。正しいとき\textgt{真}、正しくないとき\textgt{偽}という。\par
    \textgt{反例}：「$p\Rightarrow q$」が偽であることを示す例。$p$ を満たすのに $q$ を満たさないもの。\par
    \textgt{逆・裏・対偶}：「$p\Rightarrow q$」に対し、逆は「$q\Rightarrow p$」、裏は「$\overline p\Rightarrow\overline q$」、対偶は「$\overline q\Rightarrow\overline p$」。$\overline p$ は「$p$ でない」。
  \end{koushiki}

  \Shomon{1}
  \ascboxZ{「ならば」の真偽}
  「$x=1\Rightarrow x^2=1$」：$x=1$ なら $x^2=1$ は必ず成り立つ。\textgt{真}。\par
  逆「$x^2=1\Rightarrow x=1$」：$x=-1$ は $x^2=1$ を満たすが $x=1$ ではない。\textgt{偽}（反例 $x=-1$）。
  \[\Kotae{\text{命題は真、逆は偽（反例}\ x=-1\text{）}}\]

  \ascboxZ{「ならば」は集合の包含}
  条件 $p$ を満たすものの集合を $P$、$q$ を満たすものの集合を $Q$ とすると、
  \[p\Rightarrow q\ \text{が真}\iff P\subset Q\]
  である。$P$ の要素がすべて $Q$ に入っていれば、$p$ を満たすものは必ず $q$ も満たすからである。

  \begin{zuhyo}
  \begin{tikzpicture}[line width=.5pt]
    \draw (0,0) ellipse (1.6cm and .95cm);
    \fill[black!15] (-.35,0) ellipse (.8cm and .5cm);
    \draw (-.35,0) ellipse (.8cm and .5cm);
    \node[font=\footnotesize] at (-.35,0) {$P$};
    \node[font=\footnotesize] at (1,.1) {$Q$};
    \node[font=\footnotesize,anchor=west] at (1.8,.3) {$P\subset Q$ のとき};
    \node[font=\footnotesize,anchor=west] at (1.8,-.3) {$p$ は十分、$q$ は必要};
  \end{tikzpicture}
  \end{zuhyo}

  \Shomon{2}
  \ascboxZ{必要条件・十分条件の定義}
  「$p\Rightarrow q$」が真のとき、\par
  $p$ は $q$ であるための\textgt{十分条件}（$p$ さえ言えれば $q$ に十分）、\par
  $q$ は $p$ であるための\textgt{必要条件}（$p$ であるためには、少なくとも $q$ が必要）\par
  という。両方向が真のときは\textgt{必要十分条件}（同値）である。

  $p$：「$x^2=x$」、$q$：「$x=1$」とする。$x^2=x\iff x(x-1)=0\iff x=0,\ 1$ である。\par
  $q\Rightarrow p$：$x=1$ なら $x^2=1=x$ で真。\par
  $p\Rightarrow q$：反例 $x=0$ があり偽。\par
  よって「$x^2=x$」は「$x=1$」であるための
  \[\Kotae{\text{必要条件であるが十分条件ではない}}\]

  \begin{ptbs}{KEY}[矢印の先が「必要」、元が「十分」]
    $p\Rightarrow q$ のとき、矢印の元 $p$ が十分条件、矢印の先 $q$ が必要条件。集合で言えば、小さい方が十分、大きい方が必要。
  \end{ptbs}

  \begin{koushiki}{否定のつくり方}
    「$p$ かつ $q$」の否定は「$\overline p$ または $\overline q$」\par
    「$p$ または $q$」の否定は「$\overline p$ かつ $\overline q$」\par
    「すべての $x$ で $p$」の否定は「ある $x$ で $\overline p$」\par
    「ある $x$ で $p$」の否定は「すべての $x$ で $\overline p$」\par
    「少なくとも一方は偶数」の否定は「両方とも奇数」
  \end{koushiki}

  \Shomon{3}
  \ascboxZ{直接が難しければ対偶を示す}
  「$p\Rightarrow q$」と、その対偶「$\overline q\Rightarrow\overline p$」は、真偽が必ず一致する（集合で $P\subset Q$ と $\overline Q\subset\overline P$ は同じことを言っている）。\par
  「$n^2$ が偶数 $\Rightarrow$ $n$ が偶数」を直接示すのは難しい。$n^2$ から $n$ の情報を取り出す方法がないからである。そこで対偶「$n$ が奇数 $\Rightarrow$ $n^2$ が奇数」を示す。

  \textgt{証明}　$n$ が奇数のとき、整数 $k$ を用いて $n=2k+1$ と書ける。
  \[n^2=4k^2+4k+1=2(2k^2+2k)+1\]
  $2k^2+2k$ は整数だから、$n^2$ は奇数である。対偶が真なので、元の命題も真である。\QED

  \Shomon{4}
  \ascboxZ{背理法：否定を仮定して矛盾を導く}
  \textgt{背理法}とは、示したいことの否定を仮定し、そこから矛盾（あり得ないこと）を導いて、否定が誤りだったと結論する方法である。「無理数である」とは「有理数でない」ことなので、「ない」ことを示すのに向いている。

  \textgt{証明}　$\sqrt2$ が有理数であると仮定する。すると、互いに素な（1以外の公約数をもたない）自然数 $p,\ q$ を用いて $\sqrt2=\frac pq$ と書ける。分母を払って2乗すると、
  \[p^2=2q^2\]
  右辺は偶数なので $p^2$ は偶数であり、\kakko{3}より $p$ は偶数である。$p=2m$（$m$ は自然数）とおくと、$4m^2=2q^2$ より $q^2=2m^2$。同様に $q$ も偶数となる。$p,\ q$ がともに偶数なので公約数2をもち、互いに素であることに矛盾する。よって $\sqrt2$ は無理数である。\QED

  \begin{memo}
    「互いに素」と置ける理由：どんな分数も約分を最後までやれば、分母と分子が互いに素になる。この「最後まで約分した形」を最初に置いておくのが、矛盾を作るための仕掛けである。
  \end{memo}

  \Shomon{5}
  \ascboxZ{偶奇は「4で割った余り」で調べる}
  「少なくとも一方は偶数」の否定は「$a,\ b$ がともに奇数」である。これを仮定して矛盾を導く（背理法）。\par
  まず、整数の2乗を4で割った余りを調べる。\par
  偶数 $2k$ のとき：$(2k)^2=4k^2$ で余り $0$。\par
  奇数 $2k+1$ のとき：$4k^2+4k+1$ で余り $1$。\par
  つまり、\textgt{平方数を4で割った余りは0か1}である。

  \textgt{証明}　$a,\ b$ がともに奇数であると仮定する。$a^2,\ b^2$ はどちらも4で割って1余るので、$a^2+b^2$ は4で割って2余る。一方、$c^2$ は4で割って0か1余る。よって $a^2+b^2=c^2$ は成り立たず、矛盾する。ゆえに $a,\ b$ の少なくとも一方は偶数である。\QED

  \begin{memo}
    なぜ「4で」割るのか。2で割った余りだけでは、$a^2+b^2$（奇数＋奇数＝偶数）と $c^2$（偶数でもよい）の間に矛盾が出ない。余りの情報をもっと細かくするために4で割る。どの数で割るかの選び方は第10講で扱う。
  \end{memo}

  \begin{simple}[出典：オリジナル]{参考問題6}
    実数 $x,\ y$ について、次の\kakko{1}\kakko{2}は何条件か答えよ。\par
    \kakko{1}　「$x+y>0$ かつ $xy>0$」は「$x>0$ かつ $y>0$」であるための何条件か。\par
    \kakko{2}　「$x>1$ または $y>1$」は「$x+y>2$」であるための何条件か。
  \end{simple}
  \Kai
  \kakko{1}　$x>0,\ y>0$ なら $x+y>0,\ xy>0$ は明らか。逆に $xy>0$ なら $x,\ y$ は同符号で、さらに $x+y>0$ なら両方正である。両方向が真なので、\Kotae{\text{必要十分条件}}。\par
  \kakko{2}　「$x+y>2\Rightarrow x>1$ または $y>1$」：対偶「$x\leqq1$ かつ $y\leqq1\Rightarrow x+y\leqq2$」は真なので、真。\par
  「$x>1$ または $y>1\Rightarrow x+y>2$」：反例 $x=2,\ y=-5$。偽。\par
  よって、\Kotae{\text{必要条件であるが十分条件ではない}}。

  \begin{memo}
    \kakko{2}で対偶を使ったのは、「または」の結論を直接示しにくいからである。否定すると「かつ」になり、2つの不等式を足すだけで済む。
  \end{memo}

  \begin{ptbs}{CHECK}[翌日、解説を隠して答える（答えは本講の中にある）]
    \kakko{1}　必要条件・十分条件を、矢印と集合で説明せよ。\par
    \kakko{2}　「少なくとも一方は偶数」の否定を言え。\par
    \kakko{3}　$\sqrt2$ の証明で $p,\ q$ を互いに素とおく理由は何か。
  \end{ptbs}

  \notefill
\end{multicols*}
%% ==================== 第7講 ====================
\Kou{7}{三角比と図形の計量}
\begin{ascolorbox4A}[出典：オリジナル]{例題7}
  $\triangle\mathrm{ABC}$ において、$\mathrm{AB}=5$、$\mathrm{AC}=3$、$\angle\mathrm{BAC}=120^\circ$ とする。\par
  \kakko{1}　$\sin120^\circ$、$\cos120^\circ$、$\tan120^\circ$ の値を求めよ。\qquad
  \kakko{2}　辺 BC の長さを求めよ。\par
  \kakko{3}　$\triangle\mathrm{ABC}$ の面積 $S$ を求めよ。\qquad
  \kakko{4}　外接円の半径 $R$ を求めよ。\qquad
  \kakko{5}　内接円の半径 $r$ を求めよ。\par
  \kakko{6}　$\angle\mathrm{BAC}$ の二等分線と辺 BC の交点を D とする。線分 AD の長さを求めよ。
\end{ascolorbox4A}

\begin{Data}
  テーマ：単位円による三角比の定義、余弦定理・正弦定理、面積、内接円、角の二等分線\par
  難易度：(1)〜(3) A　(4)(5) B　(6) B\quad 前提：第2講\par
  この講のゴール：三角形の「辺・角・面積・円」のうち、わかっているものから残りを求める道具を選べる
\end{Data}

\begin{multicols*}{2}
  \begin{koushiki}{三角比の定義（単位円）}
    原点 O を中心とする半径1の円（単位円）の上に、$x$ 軸の正の向きから反時計回りに角 $\theta$ だけ回った点 P をとる。P の座標を $(\cos\theta,\ \sin\theta)$ と\textgt{定義}する。また $\tan\theta=\dfrac{\sin\theta}{\cos\theta}$（直線 OP の傾き）とする。
  \end{koushiki}
  中学で習う直角三角形の定義（$\sin\theta=\frac{\text{高さ}}{\text{斜辺}}$ など）は、$0^\circ<\theta<90^\circ$ でしか使えない。単位円で定義すると、$90^\circ$ 以上の角にも同じ式が使える。

  \Shomon{1}
  \ascboxZ{単位円上の点の座標を読む}
  $120^\circ$ 回った点 P は、$x$ 軸の負の向きから $60^\circ$ 上にある。$30^\circ,\ 60^\circ,\ 90^\circ$ の直角三角形の辺の比 $1:\sqrt3:2$ から、P$\left(-\frac12,\ \frac{\sqrt3}2\right)$ である。

  \begin{zuhyo}
  \begin{tikzpicture}[x=1.2cm,y=1.2cm,line width=.5pt]
    \draw[->] (-1.3,0) -- (1.4,0) node[right] {$x$};
    \draw[->] (0,-.3) -- (0,1.35) node[above] {$y$};
    \draw (1,0) arc (0:180:1);
    \draw (0,0) -- (-.5,.866) node[above left] {$\mathrm P$};
    \draw[dashed] (-.5,.866) -- (-.5,0) node[below,font=\footnotesize] {$-\frac12$};
    \draw[dashed] (-.5,.866) -- (0,.866) node[right,font=\footnotesize] {$\frac{\sqrt3}2$};
    \draw (.25,0) arc (0:120:.25);
    \node[font=\footnotesize] at (.3,.33) {$120^\circ$};
    \node[below right,font=\footnotesize] at (0,0) {$\mathrm O$};
    \node[below,font=\footnotesize] at (1,0) {$1$};
  \end{tikzpicture}
  \end{zuhyo}

  定義より、
  \[\sin120^\circ=\Kotae{\frac{\sqrt3}2},\ \cos120^\circ=\Kotae{-\frac12}\]
  $\tan120^\circ=\dfrac{\sqrt3/2}{-1/2}$ だから、
  \[\tan120^\circ=\Kotae{-\sqrt3}\]

  \begin{memo}
    P は半径1の円周上にあるので、三平方の定理から常に $\sin^2\theta+\cos^2\theta=1$。また、$180^\circ-\theta$ の点は $y$ 軸について対称な位置にあるので $\sin(180^\circ-\theta)=\sin\theta$、$\cos(180^\circ-\theta)=-\cos\theta$。
  \end{memo}

  \Shomon{2}
  \ascboxZ{2辺と間の角 → 余弦定理}

  \begin{koushiki}{余弦定理}
    $\triangle\mathrm{ABC}$ で $a=\mathrm{BC},\ b=\mathrm{CA},\ c=\mathrm{AB}$ とすると、
    \[a^2=b^2+c^2-2bc\cos A\]
  \end{koushiki}
  \textgt{なぜ成り立つか}：A を原点、B を $(c,\ 0)$ に置く。C は A から距離 $b$、角 $A$ の位置なので、単位円の定義を $b$ 倍して C$(b\cos A,\ b\sin A)$。2点間の距離の公式で、
  \[
  \begin{aligned}
  a^2&=(b\cos A-c)^2+(b\sin A)^2\\
  &=b^2(\cos^2A+\sin^2A)-2bc\cos A+c^2
  \end{aligned}
  \]
  $\cos^2A+\sin^2A=1$ より成り立つ。$A$ が鈍角でも同じ計算で済む。

  \begin{zuhyo}
  \begin{tikzpicture}[x=.55cm,y=.55cm,line width=.5pt]
    \coordinate (A) at (0,0); \coordinate (B) at (5,0); \coordinate (C) at (-1.5,2.598);
    \draw (A) -- (B) -- (C) -- cycle;
    \node[below] at (A) {$\mathrm A$}; \node[below] at (B) {$\mathrm B$}; \node[above] at (C) {$\mathrm C$};
    \node[below,font=\footnotesize] at (2.5,0) {$5$};
    \node[left,font=\footnotesize] at (-.75,1.3) {$3$};
    \draw (.5,0) arc (0:120:.5);
    \node[font=\footnotesize] at (.75,.75) {$120^\circ$};
    \draw[dashed] (A) -- (.9375,1.624) node[above right,font=\footnotesize] {$\mathrm D$};
  \end{tikzpicture}
  \end{zuhyo}

  本問では $b=3,\ c=5,\ A=120^\circ$ なので、
  \[a^2=9+25-2\cdot3\cdot5\cdot\left(-\frac12\right)=49\]
  $a>0$ より、
  \[\mathrm{BC}=\Kotae{7}\]

  \Shomon{3}
  \ascboxZ{面積は $\frac12\times$ 2辺 $\times\sin$（間の角）}
  上の座標で、C の高さ（$y$ 座標）は $b\sin A$。底辺 AB $=c$ なので、
  \[S=\frac12bc\sin A\]
  よって、
  \[S=\frac12\cdot3\cdot5\cdot\frac{\sqrt3}2=\Kotae{\frac{15\sqrt3}4}\]

  \Shomon{4}
  \ascboxZ{対辺と対角 → 正弦定理}
  \begin{koushiki}{正弦定理}
    $\triangle\mathrm{ABC}$ の外接円の半径を $R$ とすると、
    \[\frac a{\sin A}=\frac b{\sin B}=\frac c{\sin C}=2R\]
  \end{koushiki}
  \textgt{なぜ成り立つか}：B を通る直径 BA$'$ を引く。円周角の定理より $\angle\mathrm{BA'C}$ は $A$ または $180^\circ-A$ に等しく、$\angle\mathrm{BCA'}=90^\circ$。直角三角形 BCA$'$ で $a=2R\sin\angle\mathrm{BA'C}=2R\sin A$。\par
  本問では $\dfrac{7}{\sin120^\circ}=2R$ なので、
  \[R=\frac{7}{2\cdot\frac{\sqrt3}2}=\frac7{\sqrt3}=\Kotae{\frac{7\sqrt3}3}\]

  \Shomon{5}
  \ascboxZ{内接円は「面積を3つに分ける」}
  内接円の中心 I から3辺までの距離はすべて $r$ である。$\triangle\mathrm{ABC}$ を $\triangle\mathrm{IBC},\ \triangle\mathrm{ICA},\ \triangle\mathrm{IAB}$ に分けると、それぞれ高さ $r$ なので、
  \[S=\frac12r(a+b+c)\]
  $a+b+c=7+3+5=15$ だから $\dfrac{15\sqrt3}4=\dfrac{15}2r$。よって、
  \[r=\Kotae{\frac{\sqrt3}2}\]

  \Shomon{6}
  \ascboxZ{角の二等分線は面積の和で}
  AD は $\angle\mathrm A$ を $60^\circ$ ずつに分ける。$\triangle\mathrm{ABD}+\triangle\mathrm{ACD}=\triangle\mathrm{ABC}$ を、AD $=x$ とおいて面積の公式で書く。
  \[\frac12\cdot5x\sin60^\circ+\frac12\cdot3x\sin60^\circ=\frac12\cdot5\cdot3\sin120^\circ\]
  $\sin60^\circ=\sin120^\circ$ なので両辺をこれで割り、2倍すると $8x=15$。よって、
  \[\mathrm{AD}=\Kotae{\frac{15}8}\]

  \begin{ptbs}{KEY}[わかっているもので道具を選ぶ]
    2辺と間の角 → 余弦定理・面積公式。3辺 → 余弦定理で角。対辺と対角 → 正弦定理。内接円 → 面積を3分割。角の二等分線 → 面積の和。
  \end{ptbs}

  \begin{simple}[出典：オリジナル]{参考問題7}
    円に内接する四角形 ABCD において、$\mathrm{AB}=2$、$\mathrm{BC}=3$、$\mathrm{CD}=1$、$\angle\mathrm{ABC}=60^\circ$ である。対角線 AC の長さ、辺 DA の長さ、四角形 ABCD の面積を求めよ。
  \end{simple}
  \Kai
  $\triangle\mathrm{ABC}$ で余弦定理より、
  \[\mathrm{AC}^2=4+9-2\cdot2\cdot3\cdot\frac12=7\]
  よって $\mathrm{AC}=\Kotae{\sqrt7}$。\par
  円に内接する四角形の向かい合う角の和は $180^\circ$ なので $\angle\mathrm{ADC}=120^\circ$。$\mathrm{DA}=x\ (>0)$ として $\triangle\mathrm{ACD}$ で余弦定理を使う。
  \[7=x^2+1-2x\cos120^\circ=x^2+x+1\]
  $x^2+x-6=(x+3)(x-2)=0$、$x>0$ より $\mathrm{DA}=\Kotae{2}$。\par
  面積は2つの三角形の和である。
  \[\frac12\cdot2\cdot3\sin60^\circ+\frac12\cdot2\cdot1\sin120^\circ=\Kotae{2\sqrt3}\]

  \begin{ptbs}{CHECK}[翌日、解説を隠して答える（答えは本講の中にある）]
    \kakko{1}　余弦定理を座標で証明する手順を言え。\par
    \kakko{2}　内接円の半径を求める式とその理由を言え。\par
    \kakko{3}　角の二等分線の長さを求める方法を言え。
  \end{ptbs}

  \notefill
\end{multicols*}

%% ==================== 第8講 ====================
\Kou{8}{場合の数}
\begin{ascolorbox4A}[出典：オリジナル]{例題8}
  \kakko{1}　0、1、2、3、4 の5個の数字から異なる3個を選んで並べ、3桁の整数をつくる。何個できるか。\par
  \kakko{2}　男子4人と女子3人が1列に並ぶとき、女子どうしが隣り合わない並び方は何通りか。\par
  \kakko{3}　6人を2人ずつ3つの組に分ける方法は何通りか。また、6人を2人ずつ部屋 A、B、C に入れる方法は何通りか。\par
  \kakko{4}　$x+y+z=10$ を満たす0以上の整数の組 $(x,\ y,\ z)$ は何組あるか。また、正の整数の組は何組あるか。
\end{ascolorbox4A}

\begin{Data}
  テーマ：積の法則、順列と組合せ、「区別するか」の判断、仕切りの考え方\par
  難易度：(1) A　(2) B　(3) B　(4) B\quad 前提：なし\par
  この講のゴール：数え上げを「何を区別するか」「何回重複して数えたか」から組み立て、公式に頼らず説明できる
\end{Data}

\begin{multicols*}{2}
  \begin{koushiki}{和の法則・積の法則}
    \textgt{和の法則}：同時には起こらない A、B があり、A が $m$ 通り、B が $n$ 通りなら、A または B は $m+n$ 通り。\par
    \textgt{積の法則}：A の起こり方が $m$ 通りで、そのそれぞれに対して B が $n$ 通りずつ起こるなら、A と B の両方は $mn$ 通り。
  \end{koushiki}
  積の法則の「それぞれに対して $n$ 通りずつ」が重要である。A の結果によって B の通り数が変わるなら、単純に掛けてはいけない。

  \Shomon{1}
  \ascboxZ{制限の強い位置から決める}
  3桁の整数なので、百の位に0は使えない。制限のある百の位から決めると、各段階の通り数が一定になる。\par
  百の位：0以外の4通り。\par
  十の位：残り4個（0を含む）から1つで4通り。\par
  一の位：残り3個から1つで3通り。\par
  積の法則より、
  \[4\times4\times3=\Kotae{48\ \text{個}}\]

  \begin{memo}
    一の位から決めると、百の位の候補数が「一の位に0を使ったかどうか」で変わってしまい、積の法則がそのままでは使えない。「制限の強い所から」は数え上げの基本方針である。
  \end{memo}

  \Shomon{2}
  \ascboxZ{隣り合わない → 先に並べて、すき間に入れる}
  まず男子4人を並べる。並べ方は $4!=24$ 通り。男子の間と両端に、すき間が5か所できる。
  \[\wedge\,\text{男}\,\wedge\,\text{男}\,\wedge\,\text{男}\,\wedge\,\text{男}\,\wedge\]
  女子3人を、異なる3か所のすき間に1人ずつ入れれば、女子どうしは隣り合わない。5か所から3か所を選んで3人を並べる方法は ${}_5\mathrm P_3=5\cdot4\cdot3=60$ 通り。よって、
  \[24\times60=\Kotae{1440\ \text{通り}}\]

  \begin{koushiki}{順列と組合せ}
    \textgt{順列}：異なる $n$ 個から $r$ 個を選んで1列に並べる方法 ${}_n\mathrm P_r=n(n-1)\cdots(n-r+1)=\dfrac{n!}{(n-r)!}$\par
    \textgt{組合せ}：異なる $n$ 個から $r$ 個を選ぶ方法（順番は区別しない）${}_n\mathrm C_r=\dfrac{{}_n\mathrm P_r}{r!}=\dfrac{n!}{r!(n-r)!}$
  \end{koushiki}
  \textgt{なぜ $r!$ で割るか}：選んだ $r$ 個の並べ方は $r!$ 通りある。${}_n\mathrm P_r$ は、同じ $r$ 個の選び方を $r!$ 回ずつ数えている。だから $r!$ で割ると、選び方の数になる。

  \begin{memo}
    「隣り合う」ときは逆に、隣り合うものをひとかたまりにして並べ、最後にかたまりの中の並べ方を掛ける。
  \end{memo}

  \Shomon{3}
  \ascboxZ{組に名前があるか（区別するか）}
  部屋 A、B、C に入れる方が簡単なので先に数える。A に入る2人を ${}_6\mathrm C_2=15$ 通り、B を残り4人から ${}_4\mathrm C_2=6$ 通り、C は残り2人で1通り。
  \[15\times6\times1=\Kotae{90\ \text{通り}}\quad\text{（部屋 A、B、C）}\]
  組に名前がないときは、部屋の名前 A、B、C を消す。例えば $\{1,2\},\{3,4\},\{5,6\}$ という分け方は、部屋への入れ方としては $3!=6$ 通り（どの組をどの部屋にするか）に数えられていた。よって、
  \[\frac{90}{3!}=\Kotae{15\ \text{通り}}\quad\text{（名前のない3組）}\]

  \begin{ptbs}{KEY}[迷ったら全部区別して数え、重複度で割る]
    人・番号付きのもの・名前の付いた箱は区別する。区別しないものは「区別して数えた数 $\div$ 1つのものが何回数えられたか」で求める。
  \end{ptbs}

  \Shomon{4}
  \ascboxZ{○と仕切り｜の並べ方に直す}
  $x+y+z=10$ の0以上の整数解は、○10個と仕切り｜2本の並べ方と1対1に対応する。1本目の仕切りより左の○の数を $x$、2本の間を $y$、右を $z$ とする。例えば
  \[\circ\circ\circ\,|\,\circ\circ\circ\circ\circ\,|\,\circ\circ\]
  は $(x,\ y,\ z)=(3,\ 5,\ 2)$ を表す。12か所のうち仕切りを置く2か所を選べば並びが決まるので、
  \[{}_{12}\mathrm C_2=\Kotae{66\ \text{組}}\]
  正の整数の解は、$x=x'+1$ などとおくと $x'+y'+z'=7$（$x',\ y',\ z'\geqq0$）になる。同様に、
  \[{}_9\mathrm C_2=\Kotae{36\ \text{組}}\]

  \begin{memo}
    正の整数の場合は、○10個の間のすき間9か所から仕切りを入れる2か所を選ぶ（${}_9\mathrm C_2$）と考えてもよい。すき間には仕切りが1本しか入らないので、どの部分も○が1個以上になる。
  \end{memo}

  \begin{koushiki}{同じものを含む順列}
    a が $p$ 個、b が $q$ 個、c が $r$ 個（$p+q+r=n$）の並べ方は $\dfrac{n!}{p!\,q!\,r!}$ 通り。全部区別して並べた $n!$ 通りを、同じ文字どうしの入れ替えの数で割っている。
  \end{koushiki}

  \begin{simple}[出典：オリジナル]{参考問題8}
    両親と子ども4人の6人が円形のテーブルに着席する。回転して一致する座り方は同じとみなすとき、座り方は全部で何通りか。また、両親が隣り合う座り方は何通りか。
  \end{simple}
  \Kai
  回転して一致するものを同じとみなすので、1人（例えば父）の席を固定して、残りを並べる。全部で
  \[(6-1)!=\Kotae{120\ \text{通り}}\]
  両親が隣り合うときは、両親をひとかたまりにして5つのものの円順列 $(5-1)!=24$ 通り。かたまりの中の父母の並び方が2通りあるので、
  \[24\times2=\Kotae{48\ \text{通り}}\]

  \begin{ptbs}{CHECK}[翌日、解説を隠して答える（答えは本講の中にある）]
    \kakko{1}　${}_n\mathrm P_r$ を $r!$ で割ると組合せになる理由は何か。\par
    \kakko{2}　名前のない3組に分けるとき $3!$ で割る理由は何か。\par
    \kakko{3}　$x+y+z=n$ の0以上の整数解の数え方を言え。
  \end{ptbs}

  \notefill
\end{multicols*}

%% ==================== 第9講 ====================
\Kou{9}{確率}
\begin{ascolorbox4A}[出典：オリジナル]{例題9}
  \kakko{1}　2個のさいころを同時に投げるとき、出た目の和が7になる確率を求めよ。\par
  \kakko{2}　3枚の硬貨を同時に投げるとき、少なくとも1枚が表になる確率を求めよ。\par
  \kakko{3}　1個のさいころを5回投げるとき、6の目がちょうど2回出る確率を求めよ。\par
  \kakko{4}　赤玉3個と白玉2個が入った袋から、玉を1個ずつ2回取り出す（取り出した玉は戻さない）。2回目に赤玉が出る確率を求めよ。また、2回目に赤玉が出たとき、1回目も赤玉である条件付き確率を求めよ。\par
  \kakko{5}　1個のさいころを $n$ 回投げるとき、出た目の最大値が4である確率を求めよ。
\end{ascolorbox4A}

\begin{Data}
  テーマ：同様に確からしい、余事象、反復試行、条件付き確率、最大値の確率\par
  難易度：(1)(2) A　(3) B　(4) B　(5) C\quad 前提：第8講\par
  この講のゴール：確率を「根元事象の個数の比」または「確率の積と和」として、根拠をもって計算できる
\end{Data}

\begin{multicols*}{2}
  \begin{koushiki}{確率の定義}
    1回の試行で起こりうる結果（根元事象）が全部で $N$ 通りあり、どれも\textgt{同様に確からしい}（起こりやすさが等しい）とする。事象 A に含まれる結果が $a$ 通りなら、A の確率は $P(\mathrm A)=\dfrac aN$。
  \end{koushiki}
  「同様に確からしい」が成り立つように数え方を選ぶことが重要である。例えば2個のさいころの「目の和」は2〜12の11通りあるが、和7と和2は起こりやすさが違う。だから和ではなく「目の組 36 通り」を根元事象にする。

  \Shomon{1}
  \ascboxZ{2個のさいころは区別して36通り}
  2個のさいころを A、B と区別する（同じに見えても、物としては別である）。目の出方は $6\times6=36$ 通りで、どれも同様に確からしい。和が7になるのは
  \[(1,6),\ (2,5),\ (3,4),\ (4,3),\ (5,2),\ (6,1)\]
  の6通り。よって、
  \[\frac6{36}=\Kotae{\frac16}\]

  \Shomon{2}
  \ascboxZ{「少なくとも」は余事象}
  事象 A に対し、A が起こらないという事象を\textgt{余事象} $\overline{\mathrm A}$ という。$P(\mathrm A)=1-P(\overline{\mathrm A})$ である。\par
  「少なくとも1枚が表」の余事象は「3枚とも裏」で、その確率は $\left(\frac12\right)^3=\frac18$。よって、
  \[1-\frac18=\Kotae{\frac78}\]

  \begin{memo}
    「少なくとも1枚表」を直接数えると、表が1枚・2枚・3枚の3つの場合を足すことになる。余事象は1つの場合で済む。「少なくとも」「〜でない」を見たら余事象を考える。
  \end{memo}

  \Shomon{3}
  \ascboxZ{反復試行：どの回かを選んで掛ける}
  各回で6が出る確率は $\frac16$、出ない確率は $\frac56$ で、各回の結果は互いに影響しない（\textgt{独立}）。例えば「1回目と2回目だけ6」の確率は $\left(\frac16\right)^2\left(\frac56\right)^3$ である。6が出る2回をどの回にするかが ${}_5\mathrm C_2=10$ 通りあり、どれも同じ確率で、同時には起こらない。和の法則で、
  \[
  \begin{aligned}
  {}_5\mathrm C_2\left(\frac16\right)^2\left(\frac56\right)^3&=10\cdot\frac1{36}\cdot\frac{125}{216}\\
  &=\Kotae{\frac{625}{3888}}
  \end{aligned}
  \]

  \begin{koushiki}{反復試行の確率}
    1回で A が起こる確率が $p$ の試行を $n$ 回独立に繰り返すとき、A がちょうど $r$ 回起こる確率は ${}_n\mathrm C_r\,p^r(1-p)^{n-r}$。
  \end{koushiki}

  \Shomon{4}
  \ascboxZ{条件付き確率は「全体を絞る」}
  玉は同じ色でもすべて区別する（赤1、赤2、赤3、白1、白2）。1回目、2回目の順に取り出す方法は $5\times4=20$ 通りで、どれも同様に確からしい。\par
  2回目が赤：1回目が赤のとき $3\times2=6$ 通り、1回目が白のとき $2\times3=6$ 通り。合計12通り。
  \[P(\text{2回目赤})=\frac{12}{20}=\Kotae{\frac35}\]
  2回目が赤だったという情報で、起こりうる結果は上の12通りに絞られる。そのうち1回目も赤は6通り。よって、
  \[P(\text{1回目赤}\mid\text{2回目赤})=\frac6{12}=\Kotae{\frac12}\]

  \begin{koushiki}{条件付き確率}
    事象 A が起こったときに B が起こる確率は
    \[P_{\mathrm A}(\mathrm B)=\frac{P(\mathrm A\cap\mathrm B)}{P(\mathrm A)}\]
    ここから $P(\mathrm A\cap\mathrm B)=P(\mathrm A)P_{\mathrm A}(\mathrm B)$（乗法定理）。
  \end{koushiki}

  \begin{memo}
    2回目に赤が出る確率 $\frac35$ は、1回目に赤が出る確率と同じである。くじ引きは、引く順番によらず当たる確率が等しい。直観では「後から引くと不利」と感じやすいが、計算すると等しい。
  \end{memo}

  \Shomon{5}
  \ascboxZ{最大値は「以下」の確率の差}
  「最大値が4」を直接数えるのは難しい（4が何回出るかで場合が分かれる）。そこで、
  \par A：$n$ 回とも4以下、\quad B：$n$ 回とも3以下\par
  とおく。最大値が4 $\iff$ 「すべて4以下」かつ「すべて3以下ではない」である。B が起これば A も起こる（B $\subset$ A）ので、
  \[P(\text{最大値}=4)=P(\mathrm A)-P(\mathrm B)\]
  $P(\mathrm A)=\left(\frac46\right)^n$、$P(\mathrm B)=\left(\frac36\right)^n$ だから、
  \[\Kotae{\left(\frac23\right)^n-\left(\frac12\right)^n}\]

  \begin{ptbs}{KEY}[最大値・最小値の確率は「以下」「以上」の差で]
    $P(\text{最大値}=k)=P(\text{すべて}\ k\ \text{以下})-P(\text{すべて}\ k-1\ \text{以下})$。東大で繰り返し出る型である。
  \end{ptbs}

  検算：$n=1$ のとき $\frac23-\frac12=\frac16$。1回だけ投げて最大値が4なのは4が出たときなので、確かに $\frac16$ である。

  \begin{simple}[出典：オリジナル]{参考問題9}
    ある病気にかかっている人の割合は1\%である。ある検査は、病気の人には90\%の確率で陽性と出て、病気でない人にも5\%の確率で陽性と出る。この検査で陽性と出た人が、実際に病気である確率を求めよ。
  \end{simple}
  \Kai
  A：病気である、B：陽性と出る、とする。求めるのは $P_{\mathrm B}(\mathrm A)=\dfrac{P(\mathrm A\cap\mathrm B)}{P(\mathrm B)}$。乗法定理より、
  \[
  \begin{aligned}
  P(\mathrm A\cap\mathrm B)&=0.01\times0.9=0.009\\
  P(\overline{\mathrm A}\cap\mathrm B)&=0.99\times0.05=0.0495
  \end{aligned}
  \]
  陽性になるのはこの2つの場合だけなので $P(\mathrm B)=0.0585$。よって、
  \[\frac{0.009}{0.0585}=\frac{90}{585}=\Kotae{\frac2{13}}\]
  \begin{memo}
    陽性でも病気である確率は約15\%しかない。病気でない人の数が圧倒的に多いため、5\%の誤判定でも陽性者の大半を占めるからである。条件付き確率では「分母が何か」を必ず言葉で確認する。
  \end{memo}

  \begin{ptbs}{CHECK}[翌日、解説を隠して答える（答えは本講の中にある）]
    \kakko{1}　2個のさいころを区別して数える理由は何か。\par
    \kakko{2}　「最大値が $k$」の確率の求め方を言え。\par
    \kakko{3}　条件付き確率の分母は何か、言葉で言え。
  \end{ptbs}

  \notefill
\end{multicols*}
%% ==================== 第10講 ====================
\Kou{10}{整数(1) 約数・倍数・余り}
\begin{ascolorbox4A}[出典：オリジナル]{例題10}
  \kakko{1}　360 の正の約数の個数と、その総和を求めよ。\par
  \kakko{2}　$n$ を整数とする。$n^3-n$ は6の倍数であることを証明せよ。\par
  \kakko{3}　$3^{100}$ を7で割った余りを求めよ。\par
  \kakko{4}　$p$ を5以上の素数とする。$p^2-1$ は24の倍数であることを証明せよ。
\end{ascolorbox4A}

\begin{Data}
  テーマ：素因数分解と約数、割り算の余りによる分類、合同式\par
  難易度：(1) A　(2) B　(3) B　(4) C\quad 前提：第6講\par
  この講のゴール：整数の問題を「余りで分類する」「積の形にする」の2つの基本操作で処理できる
\end{Data}

\begin{multicols*}{2}
  \begin{koushiki}{割り算の定義と言葉}
    整数 $a$ と正の整数 $b$ に対し、
    \[a=bq+r,\quad 0\leqq r<b\]
    を満たす整数 $q$（商）と $r$（余り）が\textgt{ただ1組}決まる。$r=0$ のとき、$a$ は $b$ の\textgt{倍数}、$b$ は $a$ の\textgt{約数}という。\par
    \textgt{素数}：1とそれ自身以外に正の約数をもたない2以上の整数（2, 3, 5, 7, 11, …）。
  \end{koushiki}
  余りは「$0$ 以上 $b$ 未満」と決まっている。$-7=3\cdot(-3)+2$ なので、$-7$ を3で割った余りは $2$ である（$-1$ ではない）。

  \Shomon{1}
  \ascboxZ{約数は素因数の指数の組合せ}
  $360=2^3\cdot3^2\cdot5$ である。360 の正の約数は $2^a3^b5^c$ の形で、
  \[0\leqq a\leqq3,\quad 0\leqq b\leqq2,\quad 0\leqq c\leqq1\]
  の整数 $a,\ b,\ c$ の選び方と1対1に対応する。$a$ は4通り、$b$ は3通り、$c$ は2通りなので、積の法則より約数の個数は、
  \[4\times3\times2=\Kotae{24\ \text{個}}\]
  総和は、次の積を展開すると、すべての約数がちょうど1回ずつ項として現れることを使う。
  \[
  \begin{aligned}
  &(1+2+4+8)(1+3+9)(1+5)\\
  &=15\cdot13\cdot6=\Kotae{1170}
  \end{aligned}
  \]

  \begin{koushiki}{約数の個数と総和}
    $N=p^aq^br^c$（$p,\ q,\ r$ は異なる素数）のとき、\par
    個数：$(a+1)(b+1)(c+1)$\par
    総和：$(1+p+\cdots+p^a)(1+q+\cdots+q^b)(1+r+\cdots+r^c)$
  \end{koushiki}

  \Shomon{2}
  \ascboxZ{連続する整数の積を見つける}
  $n^3-n=n(n^2-1)=(n-1)n(n+1)$ と因数分解できる。これは連続する3つの整数の積である。\par
  連続する2つの整数には必ず偶数が含まれるので、積は2の倍数。連続する3つの整数には必ず3の倍数が含まれるので、積は3の倍数。\par
  2の倍数かつ3の倍数で、2と3は互いに素なので、6の倍数である。\QED

  \begin{memo}
    「2の倍数かつ3の倍数なら6の倍数」は、2と3が互いに素だから言える。「2の倍数かつ4の倍数なら8の倍数」は誤り（4が反例）。公約数があるときは、最小公倍数の倍数になるだけである。
  \end{memo}

  \ascboxZ{余りで分類する方法でも示せる}
  $n$ を3で割った余りで、$n=3k,\ 3k+1,\ 3k+2$ の3通りに分ける。どの場合も $n-1,\ n,\ n+1$ のどれかが3の倍数になる。このように、\textgt{整数を有限個の場合に分けて全部調べる}のが、整数問題の最も強力な方法である。

  \Shomon{3}
  \ascboxZ{合同式で余りだけを追う}
  \begin{koushiki}{合同式}
    $a-b$ が $m$ の倍数のとき、$a\equiv b\pmod m$ と書く（$a$ と $b$ は $m$ で割った余りが等しい）。\par
    $a\equiv b,\ c\equiv d$ のとき、\par
    $a+c\equiv b+d,\ \ a-c\equiv b-d,\ \ ac\equiv bd,\ \ a^n\equiv b^n$
  \end{koushiki}
  \textgt{なぜ掛けてよいか}：$a=b+mk,\ c=d+ml$ とおくと $ac=bd+m(bl+dk+mkl)$ なので、$ac-bd$ は $m$ の倍数。\par
  $3$ の累乗を7で割った余りを順に調べる。
  \[3^1\equiv3,\ 3^2\equiv2,\ 3^3\equiv6,\ 3^4\equiv4,\ 3^5\equiv5,\ 3^6\equiv1\]
  （例えば $3^3=27=7\cdot3+6$。次は前の余りに3を掛けて7で割ればよい。）\par
  $3^6\equiv1$ なので、6乗ごとに1に戻る。$100=6\cdot16+4$ より、
  \[3^{100}=(3^6)^{16}\cdot3^4\equiv1^{16}\cdot81\equiv\Kotae{4}\]
  （$81=7\cdot11+4$。）

  \begin{ptbs}{注意}[合同式で「割る」のは危ない]
    $2\cdot1\equiv2\cdot3\pmod4$ だが、$1\equiv3\pmod4$ ではない。両辺を割ってよいのは、割る数と $m$ が互いに素のときだけである。
  \end{ptbs}

  \Shomon{4}
  \ascboxZ{24 $=8\times3$ に分けて示す}
  $24=8\cdot3$ で、8と3は互いに素なので、$p^2-1$ が8の倍数かつ3の倍数であることを示せばよい。$p^2-1=(p-1)(p+1)$ である。

  \textgt{8の倍数であること}：$p$ は5以上の素数なので奇数である。$p-1,\ p+1$ は連続する2つの偶数で、片方は4の倍数である（偶数を2つおきに見ると、4の倍数が交互に現れる）。よって積は $2\times4=8$ の倍数。

  \textgt{3の倍数であること}：$p-1,\ p,\ p+1$ は連続する3つの整数なので、どれかが3の倍数。$p$ は3より大きい素数なので3の倍数ではない。よって $p-1$ か $p+1$ が3の倍数。\par
  以上より $p^2-1$ は24の倍数である。\QED

  \begin{memo}
    確かめ：$p=5$ のとき $24$、$p=7$ のとき $48$、$p=11$ のとき $120$。どれも24の倍数である。証明の前後に小さい数で確かめると、方針の誤りに早く気づける。
  \end{memo}

  \begin{simple}[出典：オリジナル]{参考問題10}
    方程式 $x^2-3y^2=2$ を満たす整数 $x,\ y$ は存在しないことを証明せよ。
  \end{simple}
  \Kai
  両辺を3で割った余りを比べる。$3y^2$ は3の倍数なので、左辺は $x^2$ と同じ余りをもつ。\par
  $x\equiv0,\ 1,\ 2\pmod3$ のとき、$x^2\equiv0,\ 1,\ 4\equiv1$。つまり平方数を3で割った余りは0か1である。\par
  一方、右辺の2を3で割った余りは2である。よって左辺と右辺の余りが一致することはなく、等式を満たす整数 $x,\ y$ は存在しない。\QED
  \begin{ptbs}{KEY}[整数解がないことは「余り」で示す]
    係数に現れる数（ここでは3）で割った余りを比べるのが第一候補。平方数の余りの表（mod 3 で0, 1、mod 4 で0, 1）は覚えておく。
  \end{ptbs}

  \begin{ptbs}{CHECK}[翌日、解説を隠して答える（答えは本講の中にある）]
    \kakko{1}　約数の総和が積の形で書ける理由は何か。\par
    \kakko{2}　合同式で両辺を割ってよいのはどんなときか。\par
    \kakko{3}　平方数を3で割った余り、4で割った余りは何か。
  \end{ptbs}

  \notefill
\end{multicols*}

%% ==================== 第11講 ====================
\Kou{11}{整数(2) 互除法と不定方程式}
\begin{ascolorbox4A}[出典：オリジナル]{例題11}
  \kakko{1}　391 と 299 の最大公約数を求めよ。\par
  \kakko{2}　方程式 $7x+5y=1$ を満たす整数 $x,\ y$ の組をすべて求めよ。\par
  \kakko{3}　方程式 $xy-2x-3y=0$ を満たす正の整数 $x,\ y$ の組をすべて求めよ。\par
  \kakko{4}　$x\leqq y\leqq z$ を満たす正の整数 $x,\ y,\ z$ で、$\dfrac1x+\dfrac1y+\dfrac1z=1$ となるものをすべて求めよ。
\end{ascolorbox4A}

\begin{Data}
  テーマ：ユークリッドの互除法、1次不定方程式、積の形に変形、不等式で範囲を絞る\par
  難易度：(1) A　(2) B　(3) B　(4) C\quad 前提：第10講\par
  この講のゴール：整数解を求める3つの型（1次不定方程式・積の形・大小関係で絞る）を見分けて使える
\end{Data}

\begin{multicols*}{2}
  \begin{koushiki}{この講の言葉}
    \textgt{最大公約数}：2つの整数に共通する約数のうち最大のもの。$\gcd(a,\ b)$ と書くことがある。\par
    \textgt{互いに素}：最大公約数が1であること。\par
    \textgt{不定方程式}：解が無数にあったり、整数という条件で絞ったりする方程式。
  \end{koushiki}

  \Shomon{1}
  \ascboxZ{互除法：割り算を繰り返す}
  \begin{ptbs}{KEY}[$a=bq+r$ のとき $\gcd(a,\ b)=\gcd(b,\ r)$]
    大きい数を小さい数で割った余りに置き換えても、最大公約数は変わらない。
  \end{ptbs}
  \textgt{なぜか}：$a$ と $b$ の公約数 $d$ は、$r=a-bq$ も割り切る。逆に $b$ と $r$ の公約数は $a=bq+r$ も割り切る。つまり「$a,\ b$ の公約数」と「$b,\ r$ の公約数」は全く同じ集まりなので、その最大のものも等しい。\par
  これを繰り返す。
  \[
  \begin{aligned}
  391&=299\cdot1+92\\
  299&=92\cdot3+23\\
  92&=23\cdot4
  \end{aligned}
  \]
  割り切れた所の割る数が最大公約数である。$\gcd(391,299)=\gcd(299,92)=\gcd(92,23)=23$。
  \[\Kotae{23}\]

  \begin{memo}
    素因数分解（$391=17\cdot23$、$299=13\cdot23$）が見つからなくても、互除法なら割り算だけで必ず答えが出る。数が大きいときや、文字式の最大公約数（参考問題11）で威力を発揮する。
  \end{memo}

  \Shomon{2}
  \ascboxZ{1組見つけて、引き算する}
  \textgt{手順1：1組の解を見つける}。$x$ に小さい値を入れて $1-7x$ が5の倍数になるものを探すと、$x=-2$ のとき $1+14=15=5\cdot3$。よって $(x,\ y)=(-2,\ 3)$ は解である。
  \[7\cdot(-2)+5\cdot3=1\]
  \textgt{手順2：辺々引く}。$7x+5y=1$ から上の式を引くと、
  \[7(x+2)=-5(y-3)\]
  \textgt{手順3：互いに素を使う}。右辺は5の倍数なので、$7(x+2)$ は5の倍数。7と5は互いに素なので、$x+2$ が5の倍数である。$x+2=5k$（$k$ は整数）とおくと、$7\cdot5k=-5(y-3)$ より $y-3=-7k$。
  \[\Kotae{(x,\ y)=(5k-2,\ -7k+3)\ \ (k\ \text{は整数})}\]

  \begin{koushiki}{互いに素と割り切り}
    $a$ と $b$ が互いに素で、$bc$ が $a$ の倍数ならば、$c$ が $a$ の倍数である。
  \end{koushiki}
  \begin{memo}
    1組の解が見つけにくいときは、互除法の式を下から逆にたどる。$ax+by=c$ が整数解をもつのは、$c$ が $\gcd(a,\ b)$ の倍数のとき、かつそのときに限る。
  \end{memo}

  \Shomon{3}
  \ascboxZ{積の形 $(\ )(\ )=$ 整数 を作る}
  $xy-2x-3y$ に定数を足して因数分解できる形にする。$(x-3)(y-2)=xy-2x-3y+6$ なので、両辺に6を足して、
  \[(x-3)(y-2)=6\]
  左辺の2つのかっこはどちらも整数で、積が6である。6の約数の組は有限個しかないので、すべて調べられる。\par
  $x\geqq1$ より $x-3\geqq-2$、$y\geqq1$ より $y-2\geqq-1$ である。積が6になる整数の組 $(x-3,\ y-2)$ は、
  \[(1,6),(2,3),(3,2),(6,1),(-1,-6),(-2,-3),\cdots\]
  のうち、負の組は $y-2\geqq-1$ を満たさない（$-6,\ -3$ などが小さすぎる）。正の4組から、
  \[\Kotae{(x,\ y)=(4,8),\ (5,5),\ (6,4),\ (9,3)}\]

  \begin{memo}
    検算：$(4,8)$ なら $32-8-24=0$。全組を元の式に代入して確かめる。積の形にすれば「候補が有限個」になる。これが整数問題で積の形を目指す理由である。
  \end{memo}

  \Shomon{4}
  \ascboxZ{大小関係で最小の文字の範囲を絞る}
  $x\leqq y\leqq z$ なので $\frac1x\geqq\frac1y\geqq\frac1z$。よって、
  \[1=\frac1x+\frac1y+\frac1z\leqq\frac3x\]
  より $x\leqq3$。また $\frac1x<1$（残りの2項が正）より $x\geqq2$。\par
  $x=2$ のとき：$\frac1y+\frac1z=\frac12$。同様に $\frac12\leqq\frac2y$ より $y\leqq4$、$\frac1y<\frac12$ より $y\geqq3$。$y=3$ なら $z=6$、$y=4$ なら $z=4$。\par
  $x=3$ のとき：$\frac1y+\frac1z=\frac23$。$\frac23\leqq\frac2y$ より $y\leqq3$、$y\geqq x=3$ より $y=3$、$z=3$。
  \[\Kotae{(x,y,z)=(2,3,6),\ (2,4,4),\ (3,3,3)}\]

  \begin{ptbs}{KEY}[整数解の3つの型]
    1次式 $ax+by=c$ → 1組見つけて引き算。\par
    積が作れる → 積＝定数にして約数を調べる。\par
    対称な式 → 大小を仮定し、最小の文字の範囲を不等式で絞る。
  \end{ptbs}

  \begin{simple}[出典：オリジナル]{参考問題11}
    $1\leqq n\leqq50$ を満たす整数 $n$ のうち、$2n+1$ と $n+3$ の最大公約数が5となるものは何個あるか。
  \end{simple}
  \Kai
  互除法の考え方は文字式でも使える。$2n+1=2(n+3)-5$ なので、
  \[\gcd(2n+1,\ n+3)=\gcd(n+3,\ 5)\]
  5の正の約数は1と5だけなので、最大公約数が5になるのは $n+3$ が5の倍数のとき、つまり $n$ を5で割った余りが2のときに限る。$1\leqq n\leqq50$ では $n=2,\ 7,\ 12,\ \cdots,\ 47$ で、
  \[\Kotae{10\ \text{個}}\]

  \begin{ptbs}{CHECK}[翌日、解説を隠して答える（答えは本講の中にある）]
    \kakko{1}　互除法で最大公約数が変わらない理由は何か。\par
    \kakko{2}　$ax+by=c$ の整数解を求める3手順を言え。\par
    \kakko{3}　積の形 $(\ )(\ )=$ 定数 にする利点は何か。
  \end{ptbs}

  \notefill
\end{multicols*}
%% ==================== 第12講 ====================
\Kou{12}{式と証明・不等式}
\begin{ascolorbox4A}[出典：オリジナル]{例題12}
  \kakko{1}　等式 $x^2+3x+1=a(x-1)^2+b(x-1)+c$ が $x$ についての恒等式となるように、定数 $a,\ b,\ c$ の値を定めよ。\par
  \kakko{2}　$x>0$ のとき、$x+\dfrac4x$ の最小値を求めよ。\par
  \kakko{3}　実数 $a,\ b,\ c$ に対して $a^2+b^2+c^2\geqq ab+bc+ca$ が成り立つことを証明せよ。また、等号が成り立つ条件を求めよ。\par
  \kakko{4}　$a>0,\ b>0$ のとき、$\left(a+\dfrac1b\right)\left(b+\dfrac4a\right)$ の最小値を求めよ。
\end{ascolorbox4A}

\begin{Data}
  テーマ：恒等式と係数比較、不等式の証明（差をとる・平方の和）、相加平均と相乗平均\par
  難易度：(1) A　(2) B　(3) B　(4) C\quad 前提：第1・6講\par
  この講のゴール：「最小値」と言うために必要な「等号が実際に成り立つこと」の確認を、必ず実行できる
\end{Data}

\begin{multicols*}{2}
  \begin{koushiki}{恒等式}
    文字にどんな値を入れても成り立つ等式を\textgt{恒等式}という。これに対し、特定の値でだけ成り立つ等式が\textgt{方程式}である。\par
    $x$ の多項式 $P(x),\ Q(x)$ について、$P(x)=Q(x)$ が恒等式 $\iff$ 両辺の同じ次数の係数がすべて等しい。
  \end{koushiki}

  \Shomon{1}
  \ascboxZ{係数比較法}
  右辺を展開して $x$ について整理する。
  \[
  \begin{aligned}
  &a(x^2-2x+1)+b(x-1)+c\\
  &=ax^2+(-2a+b)x+(a-b+c)
  \end{aligned}
  \]
  左辺 $x^2+3x+1$ と係数を比べて、
  \[a=1,\quad -2a+b=3,\quad a-b+c=1\]
  順に解いて、
  \[\Kotae{a=1,\ b=5,\ c=5}\]

  \ascboxZ{数値代入法（必要条件として使う）}
  恒等式ならどの $x$ でも成り立つので、計算が楽な値を入れてよい。$x=1$ を入れると $5=c$。ただし「ある値で成り立つ」は恒等式であるための必要条件にすぎない。数値代入で求めたら、最後に展開して確かめる（十分性の確認）。

  \begin{memo}
    2次式どうしの等式なら、異なる3つの値で成り立てば恒等式になることが知られている（差が「3つの解をもつ2次以下の式」になり、そのような式は0しかない）。だから3つの値を代入して求めた場合は、確かめを省略してよい理由がある。
  \end{memo}

  \Shomon{2}
  \ascboxZ{相加平均 $\geqq$ 相乗平均}
  \begin{koushiki}{相加平均と相乗平均の関係}
    $a\geqq0,\ b\geqq0$ のとき、
    \[\frac{a+b}2\geqq\sqrt{ab}\quad\text{つまり}\quad a+b\geqq2\sqrt{ab}\]
    等号は $a=b$ のときに限り成り立つ。
  \end{koushiki}
  \textgt{証明}：$a+b-2\sqrt{ab}=(\sqrt a-\sqrt b)^2\geqq0$。等号は $\sqrt a=\sqrt b$、つまり $a=b$ のとき。\par
  $x>0$ より $x>0,\ \frac4x>0$ なので、
  \[x+\frac4x\geqq2\sqrt{x\cdot\frac4x}=4\]
  等号は $x=\frac4x$、すなわち $x^2=4$、$x>0$ より $x=2$ のとき成り立つ。よって、
  \[\text{最小値}\ \Kotae{4}\ \ (x=2\ \text{のとき})\]

  \begin{ptbs}{KEY}[「$\geqq4$」だけでは最小値4とは言えない]
    最小値とは「実際にとる値の中で最小」。不等式 $\geqq4$ に加えて、等号が成り立つ $x$ が実在することを示して、初めて最小値4と言える。
  \end{ptbs}

  \Shomon{3}
  \ascboxZ{差をとって「平方の和」を作る}
  不等式 $A\geqq B$ を示すには、$A-B\geqq0$ を示すのが基本である。
  \[
  \begin{aligned}
  &(a^2+b^2+c^2)-(ab+bc+ca)\\
  &=\frac12\{(a-b)^2+(b-c)^2+(c-a)^2\}\geqq0
  \end{aligned}
  \]
  2行目は、$\{\ \}$ の中を展開すると $2a^2+2b^2+2c^2-2ab-2bc-2ca$ になることから確かめられる。実数の2乗は0以上なので、その和も0以上である。\QED\par
  等号は $(a-b)^2=(b-c)^2=(c-a)^2=0$、すなわち
  \[\Kotae{a=b=c}\]
  のとき成り立つ。

  \begin{memo}
    「2乗の和 $=0$」なら、各2乗がすべて0である。1つでも正なら和は正になるからである。等号条件はこの事実から求める。
  \end{memo}

  \Shomon{4}
  \ascboxZ{よくある誤り：等号が同時に成り立たない}
  \textgt{誤答}：$a+\frac1b\geqq2\sqrt{\frac ab}$、$b+\frac4a\geqq2\sqrt{\frac{4b}a}$ を掛けて、
  \[\left(a+\frac1b\right)\left(b+\frac4a\right)\geqq4\sqrt{\frac ab\cdot\frac{4b}a}=8\]
  「最小値8」とするのは誤りである。1つ目の等号は $ab=1$、2つ目の等号は $ab=4$ のときで、同時には成り立たない。だから値8は実際にはとらない。

  \textgt{正しい解き方}：先に展開して、変数を1つのかたまり $ab$ にまとめる。
  \[
  \begin{aligned}
  \left(a+\frac1b\right)\left(b+\frac4a\right)&=ab+4+1+\frac4{ab}\\
  &=ab+\frac4{ab}+5
  \end{aligned}
  \]
  $ab>0$ なので相加・相乗平均の関係より $ab+\frac4{ab}\geqq2\sqrt4=4$。等号は $ab=2$ のとき（例えば $a=1,\ b=2$）成り立つ。よって、
  \[\text{最小値}\ 4+5=\Kotae{9}\]

  \begin{simple}[出典：オリジナル]{参考問題12}
    $x>1$ のとき、$x+\dfrac1{x-1}$ の最小値と、そのときの $x$ の値を求めよ。
  \end{simple}
  \Kai
  分母が $x-1$ なので、$x$ も $x-1$ の形にそろえる（積が定数になるようにする）。
  \[x+\frac1{x-1}=(x-1)+\frac1{x-1}+1\]
  $x-1>0$ より、$(x-1)+\frac1{x-1}\geqq2$。等号は $x-1=\frac1{x-1}$、$(x-1)^2=1$、$x-1>0$ より $x=2$ のとき。
  \[\text{最小値}\ \Kotae{3}\ \ (x=2)\]

  \begin{ptbs}{CHECK}[翌日、解説を隠して答える（答えは本講の中にある）]
    \kakko{1}　数値代入法で係数を求めたとき、何を確かめるか。\par
    \kakko{2}　相加・相乗平均で「最小値」と言うために必要なことは何か。\par
    \kakko{3}　不等式を2つ掛けて最小値を誤る原因は何か。
  \end{ptbs}

  \notefill
\end{multicols*}

%% ==================== 第13講 ====================
\Kou{13}{複素数と高次方程式}
\begin{ascolorbox4A}[出典：オリジナル]{例題13}
  \kakko{1}　$(2+3i)(1-i)$ と $\dfrac{1+2i}{2-i}$ を計算し、$a+bi$（$a,\ b$ は実数）の形で表せ。\par
  \kakko{2}　2次方程式 $x^2-3x+5=0$ の2つの解を $\alpha,\ \beta$ とする。$\alpha^2+\beta^2$ と $\alpha^3+\beta^3$ の値を求めよ。\par
  \kakko{3}　整式 $P(x)$ を $x-1$ で割ると余りが2、$x+2$ で割ると余りが $-7$ である。$P(x)$ を $(x-1)(x+2)$ で割ったときの余りを求めよ。\par
  \kakko{4}　3次方程式 $x^3-3x^2+4=0$ を解け。\par
  \kakko{5}　方程式 $x^2+x+1=0$ の解の1つを $\omega$ とする。$\omega^{100}+\omega^{50}+1$ の値を求めよ。
\end{ascolorbox4A}

\begin{Data}
  テーマ：複素数の計算、解と係数の関係、剰余の定理・因数定理、3次方程式、$\omega$\par
  難易度：(1) A　(2) A　(3) B　(4) B　(5) C\quad 前提：第1・3講\par
  この講のゴール：「解を直接求めずに、解の和と積だけで計算する」発想を身につける
\end{Data}

\begin{multicols*}{2}
  \begin{koushiki}{複素数}
    2乗して $-1$ になる数を $i$（虚数単位）と書く：$i^2=-1$。実数 $a,\ b$ を用いて $a+bi$ と表される数を\textgt{複素数}という。\par
    相等：$a+bi=c+di\iff a=c$ かつ $b=d$。\par
    共役：$a+bi$ に対し $a-bi$ を\textgt{共役な複素数}という。$(a+bi)(a-bi)=a^2+b^2$（実数）。
  \end{koushiki}
  計算は、$i$ を文字のように扱って展開し、$i^2$ が出たら $-1$ に置き換えるだけである。

  \Shomon{1}
  \ascboxZ{$i^2=-1$ に置き換える／共役を掛ける}
  \[
  \begin{aligned}
  (2+3i)(1-i)&=2-2i+3i-3i^2\\
  &=2+i+3=\Kotae{5+i}
  \end{aligned}
  \]
  割り算は、分母の共役 $2+i$ を分母と分子に掛けて、分母を実数にする（第2講の有理化と同じ発想）。
  \[
  \begin{aligned}
  \frac{1+2i}{2-i}&=\frac{(1+2i)(2+i)}{(2-i)(2+i)}\\
  &=\frac{2+i+4i+2i^2}{4+1}=\frac{5i}5=\Kotae{i}
  \end{aligned}
  \]

  \Shomon{2}
  \ascboxZ{解と係数の関係}
  \begin{koushiki}{解と係数の関係}
    $ax^2+bx+c=0$ の2解を $\alpha,\ \beta$ とすると、
    \[\alpha+\beta=-\frac ba,\qquad\alpha\beta=\frac ca\]
  \end{koushiki}
  \textgt{なぜか}：2解が $\alpha,\ \beta$ なら $ax^2+bx+c=a(x-\alpha)(x-\beta)$ と因数分解できる。右辺を展開した $ax^2-a(\alpha+\beta)x+a\alpha\beta$ と係数を比べればよい。\par
  本問では $\alpha+\beta=3,\ \alpha\beta=5$。解は虚数（$D=9-20<0$）だが、関係式はそのまま使える。第1講の対称式の変形で、
  \[\alpha^2+\beta^2=(\alpha+\beta)^2-2\alpha\beta=9-10=\Kotae{-1}\]
  \[
  \begin{aligned}
  \alpha^3+\beta^3&=(\alpha+\beta)^3-3\alpha\beta(\alpha+\beta)\\
  &=27-45=\Kotae{-18}
  \end{aligned}
  \]

  \begin{memo}
    $\alpha^2+\beta^2=-1<0$ になったが誤りではない。$\alpha,\ \beta$ が虚数なので、その2乗は負にもなりうる。「2乗の和は0以上」は実数のときだけの性質である。
  \end{memo}

  \Shomon{3}
  \ascboxZ{剰余の定理：割り算の式に値を入れる}
  \begin{koushiki}{剰余の定理・因数定理}
    整式 $P(x)$ を $x-a$ で割った余りは $P(a)$。\par
    特に、$P(a)=0\iff P(x)$ は $x-a$ で割り切れる（因数定理）。
  \end{koushiki}
  \textgt{なぜか}：$P(x)=(x-a)Q(x)+r$ と書ける（$r$ は定数）。$x=a$ を代入すると $P(a)=r$。\par
  本問の条件は $P(1)=2,\ P(-2)=-7$。2次式 $(x-1)(x+2)$ で割った余りは1次以下なので $ax+b$ とおける。
  \[P(x)=(x-1)(x+2)Q(x)+ax+b\]
  $x=1,\ -2$ を代入すると、商の部分が消えて、
  \[a+b=2,\quad -2a+b=-7\]
  これを解いて $a=3,\ b=-1$。余りは、
  \[\Kotae{3x-1}\]

  \begin{memo}
    「割る式の次数より、余りの次数は小さい」という割り算の約束から、余りを $ax+b$ と置ける。この置き方が書けるかどうかが、この型の問題の分かれ目である。
  \end{memo}

  \Shomon{4}
  \ascboxZ{因数定理で1つ解を見つける}
  $P(x)=x^3-3x^2+4$ とおく。整数の解があるなら、それは定数項4の約数 $\pm1,\ \pm2,\ \pm4$ のどれかである（$x$ が整数解なら $x(x^2-3x)=-4$ だから）。$P(-1)=-1-3+4=0$ なので、$P(x)$ は $x+1$ で割り切れる。割り算をすると、
  \[P(x)=(x+1)(x^2-4x+4)=(x+1)(x-2)^2\]
  よって、
  \[x=\Kotae{-1,\ 2}\ \ (2\ \text{は重解})\]

  \Shomon{5}
  \ascboxZ{$\omega^3=1$ を作って次数を下げる}
  $\omega^2+\omega+1=0$ の両辺に $\omega-1$ を掛けると、$(\omega-1)(\omega^2+\omega+1)=\omega^3-1$ なので、
  \[\omega^3=1\]
  よって $\omega$ の指数は3で割った余りだけが効く。$100=3\cdot33+1$、$50=3\cdot16+2$ より、$\omega^{100}=\omega$、$\omega^{50}=\omega^2$。
  \[\omega^{100}+\omega^{50}+1=\omega+\omega^2+1=\Kotae{0}\]

  \begin{ptbs}{KEY}[高次の式は「満たす方程式」で次数を下げる]
    $\omega^2=-\omega-1$ や $\omega^3=1$ のように、解が満たす式を使って、次数の高い式を低い式に書き換える。東大の複素数・整式の問題で頻出の操作である。
  \end{ptbs}

  \begin{simple}[出典：オリジナル]{参考問題13}
    $a,\ b$ を実数とする。3次方程式 $x^3+ax^2+bx+10=0$ が $2-i$ を解にもつとき、$a,\ b$ の値と、他の解を求めよ。
  \end{simple}
  \Kai
  係数がすべて実数なので、$2-i$ が解なら共役な $2+i$ も解である（$P(\overline z)=\overline{P(z)}$ なので、$P(z)=0$ なら $P(\overline z)=\overline0=0$）。この2つを解にもつ2次式は、
  \[\{x-(2-i)\}\{x-(2+i)\}=x^2-4x+5\]
  である。残りの解を $c$ とすると、
  \[x^3+ax^2+bx+10=(x^2-4x+5)(x-c)\]
  定数項を比べて $-5c=10$、$c=-2$。右辺 $=(x^2-4x+5)(x+2)=x^3-2x^2-3x+10$。よって、
  \[\Kotae{a=-2,\ b=-3,\ \text{他の解}\ 2+i,\ -2}\]

  \begin{ptbs}{CHECK}[翌日、解説を隠して答える（答えは本講の中にある）]
    \kakko{1}　解と係数の関係はどう導くか。\par
    \kakko{2}　剰余の定理を1行で証明せよ。\par
    \kakko{3}　$\omega^2+\omega+1=0$ から $\omega^3=1$ を作る方法を言え。
  \end{ptbs}

  \notefill
\end{multicols*}

%% ==================== 第14講 ====================
\Kou{14}{図形と方程式}
\begin{ascolorbox4A}[出典：オリジナル]{例題14}
  \kakko{1}　2点 $\mathrm A(1,\ 2)$、$\mathrm B(4,\ 6)$ について、線分 AB の長さ、中点の座標、AB を $1:2$ に内分する点の座標を求めよ。\par
  \kakko{2}　点 $(1,\ 2)$ を通り、直線 $2x-y+3=0$ に垂直な直線の方程式を求めよ。\par
  \kakko{3}　点 $(3,\ 1)$ と直線 $3x+4y-3=0$ の距離を求めよ。\par
  \kakko{4}　円 $x^2+y^2-4x+6y-3=0$ の中心と半径を求めよ。\par
  \kakko{5}　円 $x^2+y^2=5$ と直線 $y=2x+k$ が接するような定数 $k$ の値と、そのときの接点の座標を求めよ。
\end{ascolorbox4A}

\begin{Data}
  テーマ：座標平面上の距離・分点・直線・円、点と直線の距離、円と直線の位置関係\par
  難易度：(1)〜(4) A　(5) B\quad 前提：第2・3講\par
  この講のゴール：図形の条件（垂直・接する・距離）を式に翻訳する辞書を手に入れる
\end{Data}

\begin{multicols*}{2}
  \Shomon{1}
  \ascboxZ{距離は三平方、分点は「比で寄せる」}
  \begin{koushiki}{2点間の距離と内分点}
    $\mathrm A(x_1,y_1),\ \mathrm B(x_2,y_2)$ のとき\par
    距離：$\mathrm{AB}=\sqrt{(x_2-x_1)^2+(y_2-y_1)^2}$\par
    AB を $m:n$ に内分する点：$\left(\dfrac{nx_1+mx_2}{m+n},\ \dfrac{ny_1+my_2}{m+n}\right)$
  \end{koushiki}
  距離は、AB を斜辺とする直角三角形（横 $x_2-x_1$、縦 $y_2-y_1$）に三平方の定理を使ったものである。内分点の式は、A から B へ全体の $\frac m{m+n}$ だけ進んだ点、つまり $x_1+\frac m{m+n}(x_2-x_1)$ を整理したものである。\par
  本問では、
  \[\mathrm{AB}=\sqrt{3^2+4^2}=\Kotae{5}\]
  中点（$1:1$ の内分点）は、
  \[\left(\frac{1+4}2,\ \frac{2+6}2\right)=\Kotae{\left(\frac52,\ 4\right)}\]
  $1:2$ の内分点は（A に近い方）、
  \[\left(\frac{2\cdot1+1\cdot4}3,\ \frac{2\cdot2+1\cdot6}3\right)=\Kotae{\left(2,\ \frac{10}3\right)}\]

  \begin{memo}
    分子では、比の数字が\textgt{たすき掛け}になる（A の座標に B 側の $n$、B の座標に A 側の $m$）。$1:2$ に内分する点は A に近いので、A の重みが大きい $2$ になる、と覚えれば間違えない。
  \end{memo}

  \Shomon{2}
  \ascboxZ{垂直 $\iff$ 傾きの積 $=-1$}
  $2x-y+3=0$ は $y=2x+3$ なので傾き2。傾き $m$ の直線の方向は「右に1、上に $m$」で、これを $90^\circ$ 回すと「右に $-m$、上に1」（左に $m$、上に1）になる。その傾きは $\frac1{-m}=-\frac1m$ なので、垂直な2直線の傾きの積は $-1$ である。\par
  求める直線の傾きは $-\frac12$。点 $(1,\ 2)$ を通るので、
  \[y-2=-\frac12(x-1)\]
  整理して、
  \[\Kotae{x+2y-5=0}\]

  \begin{koushiki}{直線の式の辞書}
    点 $(x_1,y_1)$ を通り傾き $m$：$y-y_1=m(x-x_1)$\par
    一般形 $ax+by+c=0$ では、ベクトル $(a,\ b)$ が直線に垂直（第25講で使う「法線ベクトル」）。\par
    平行 $\iff$ 傾きが等しい／垂直 $\iff$ 傾きの積が $-1$
  \end{koushiki}

  \Shomon{3}
  \ascboxZ{点と直線の距離の公式}
  \begin{koushiki}{点と直線の距離}
    点 $(x_0,y_0)$ と直線 $ax+by+c=0$ の距離 $d$ は、
    \[d=\frac{|ax_0+by_0+c|}{\sqrt{a^2+b^2}}\]
  \end{koushiki}
  \textgt{なぜか}：点 P$(x_0,y_0)$ から直線へ垂線を下ろす。垂線の向きは $(a,\ b)$ なので、垂線の足 H は $(x_0+ta,\ y_0+tb)$ と書ける。これを直線の式に入れると $ax_0+by_0+c+t(a^2+b^2)=0$ より $t=-\frac{ax_0+by_0+c}{a^2+b^2}$。$d=\mathrm{PH}=|t|\sqrt{a^2+b^2}$ を計算すると公式になる。\par
  本問では、
  \[d=\frac{|3\cdot3+4\cdot1-3|}{\sqrt{9+16}}=\frac{10}5=\Kotae{2}\]

  \Shomon{4}
  \ascboxZ{平方完成で中心と半径を読む}
  中心 $(a,\ b)$、半径 $r$ の円は、「中心からの距離が $r$ の点の集まり」なので、$(x-a)^2+(y-b)^2=r^2$ と表される。$x,\ y$ それぞれで平方完成する。
  \[(x-2)^2-4+(y+3)^2-9-3=0\]
  \[(x-2)^2+(y+3)^2=16\]
  よって、
  \[\text{中心}\ \Kotae{(2,\ -3)},\quad\text{半径}\ \Kotae{4}\]

  \Shomon{5}
  \ascboxZ{接する $\iff$ 中心と直線の距離 $=$ 半径}
  円の中心 $(0,\ 0)$ と直線 $2x-y+k=0$ の距離が半径 $\sqrt5$ に等しいとき、円と直線は接する。
  \[\frac{|k|}{\sqrt{4+1}}=\sqrt5\]
  より $|k|=5$。
  \[k=\Kotae{\pm5}\]
  接点は、中心から直線に下ろした垂線の足である。\kakko{3}の計算と同じく、垂線の足は $(0,0)+t(2,\ -1)$ で、$4t+t+k=0$ より $t=-\frac k5$。\par
  $k=5$ のとき $t=-1$ で、接点 \namiKKAuto{$(-2,\ 1)$}。\par
  $k=-5$ のとき $t=1$ で、接点 \namiKKAuto{$(2,\ -1)$}。\hspace{\zw}\textgt{……\kakko{答}}

  \begin{ptbs}{KEY}[円と直線は「距離」で判定]
    中心と直線の距離 $d$、半径 $r$ について、$d<r$ なら2点で交わる、$d=r$ なら接する、$d>r$ なら共有点なし。判別式を使うより計算が軽い。
  \end{ptbs}

  \begin{memo}
    円 $x^2+y^2=r^2$ 上の点 $(x_1,y_1)$ における接線は $x_1x+y_1y=r^2$ である。接点 $(-2,\ 1)$ では $-2x+y=5$、つまり $y=2x+5$ となり、$k=5$ と一致する。
  \end{memo}

  \begin{simple}[出典：オリジナル]{参考問題14}
    直線 $(k+1)x+(k-1)y-2k=0$ は、定数 $k$ の値によらず定点を通る。その定点の座標を求めよ。
  \end{simple}
  \Kai
  「$k$ によらず」とは「$k$ についての恒等式」ということである。$k$ について整理する。
  \[k(x+y-2)+(x-y)=0\]
  これがすべての $k$ で成り立つ条件は、$k$ の係数と定数項がともに0であること（第12講）。
  \[x+y-2=0,\quad x-y=0\]
  これを解いて、定点は
  \[\Kotae{(1,\ 1)}\]

  \begin{ptbs}{CHECK}[翌日、解説を隠して答える（答えは本講の中にある）]
    \kakko{1}　垂直な2直線の傾きの積が $-1$ になる理由は何か。\par
    \kakko{2}　点と直線の距離の公式の導き方を言え。\par
    \kakko{3}　円と直線が接する条件を言え。
  \end{ptbs}

  \notefill
\end{multicols*}
%% ==================== 第15講 ====================
\Kou{15}{軌跡と領域・通過領域}
\begin{ascolorbox4A}[出典：オリジナル]{例題15}
  \kakko{1}　2点 $\mathrm A(-1,\ 0)$、$\mathrm B(1,\ 0)$ に対し、$\mathrm{PA}:\mathrm{PB}=2:1$ を満たす点 P の軌跡を求めよ。\par
  \kakko{2}　連立不等式 $x^2\leqq y\leqq x+2$ の表す領域 $D$ を図示せよ。\par
  \kakko{3}　点 $(x,\ y)$ が領域 $D$ を動くとき、$x+y$ の最大値と最小値を求めよ。\par
  $t$ を実数とし、直線 $\ell_t:\ y=2tx-t^2$ を考える。\par
  \kakko{4}　$t$ が実数全体を動くとき、$\ell_t$ が通過する範囲を求め、図示せよ。\par
  \kakko{5}　$t$ が $t\geqq0$ の範囲を動くとき、$\ell_t$ が通過する範囲を求め、図示せよ。
\end{ascolorbox4A}

\begin{Data}
  テーマ：軌跡の求め方、不等式の表す領域、領域での最大・最小、直線の通過領域（逆像法）\par
  難易度：(1) B　(2) A　(3) B　(4) C　(5) D\quad 前提：第5・14講\par
  この講のゴール：「その点を通る $t$ が存在する」という存在条件に翻訳して、通過範囲を求められる
\end{Data}

\begin{multicols*}{2}
  \begin{koushiki}{軌跡を求める手順}
    \textgt{軌跡}：条件を満たす点全体の集まり（図形）。\par
    \textgt{手順}：\kakko{1} 動く点を $(x,\ y)$ とおく。\kakko{2} 条件を $x,\ y$ の式で書く。\kakko{3} 同値変形で整理し、図形を読み取る。\kakko{4} 途中で同値でない変形をしたら、除く点がないか確かめる。
  \end{koushiki}

  \Shomon{1}
  \ascboxZ{長さの比は2乗して扱う}
  P$(x,\ y)$ とおく。$\mathrm{PA}:\mathrm{PB}=2:1\iff\mathrm{PA}=2\mathrm{PB}\iff\mathrm{PA}^2=4\mathrm{PB}^2$（長さは0以上なので2乗しても同値）。
  \[(x+1)^2+y^2=4\{(x-1)^2+y^2\}\]
  展開して整理すると $3x^2-10x+3y^2+3=0$、両辺を3で割って平方完成すると、
  \[\left(x-\frac53\right)^2+y^2=\frac{16}9\]
  同値変形だけを使ったので、この円上のすべての点が条件を満たす。
  \[\Kotae{\text{中心}\left(\frac53,\ 0\right)\text{、半径}\ \frac43\ \text{の円}}\]

  \begin{memo}
    2定点からの距離の比が $m:n\ (m\neq n)$ で一定の点の軌跡は円になる（アポロニウスの円）。AB を $2:1$ に内分する点 $\left(\frac13,0\right)$ と外分する点 $(3,0)$ を直径の両端とする円である。
  \end{memo}

  \Shomon{2}
  \ascboxZ{$y\geqq f(x)$ はグラフの上側}
  $y\geqq x^2$ は放物線 $y=x^2$ の上側（境界を含む）、$y\leqq x+2$ は直線 $y=x+2$ の下側（境界を含む）。その共通部分が $D$ である。交点は $x^2=x+2$ より $x=-1,\ 2$。

  \begin{zuhyo}
  \begin{tikzpicture}[x=.62cm,y=.42cm,line width=.5pt]
    \fill[black!15] plot[domain=-1:2,samples=40] (\x,{(\x)^2}) -- cycle;
    \draw[->] (-2.3,0) -- (3,0) node[right] {$x$};
    \draw[->] (0,-1.4) -- (0,5.4) node[above] {$y$};
    \draw[domain=-2.1:2.3,samples=40] plot (\x,{(\x)^2});
    \draw (-2.3,-.3) -- (2.7,4.7) node[right,font=\footnotesize] {$y=x+2$};
    \fill (2,4) circle (2pt) node[left,font=\footnotesize] {$(2,4)$};
    \fill (-1,1) circle (2pt) node[left,font=\footnotesize] {$(-1,1)$};
    \fill (-.5,.25) circle (2pt);
    \draw[dashed] (-1.6,1.35) -- (1.2,-1.45);
    \node[font=\footnotesize,anchor=north west] at (1.2,-.6) {$x+y=-\frac14$};
    \node[below left,font=\footnotesize] at (0,0) {$\mathrm O$};
    \node[font=\footnotesize] at (.9,2.2) {$D$};
  \end{tikzpicture}
  \end{zuhyo}

  図の灰色の部分。ただし境界線を含む。\hfill\textgt{……\kakko{答}}

  \Shomon{3}
  \ascboxZ{$x+y=k$ とおいて直線を動かす}
  $x+y=k$ とおくと $y=-x+k$。これは傾き $-1$、$y$ 切片 $k$ の直線である。$k$ がある値をとる $\iff$ この直線が $D$ と共有点をもつ。そこで直線を平行に動かし、$D$ と共有点をもつ範囲で $y$ 切片 $k$ の最大・最小を探す。\par
  \textgt{最大}：直線を上に動かすと、最後に点 $(2,\ 4)$ を通る。このとき $k=6$。\par
  \textgt{最小}：直線を下に動かすと、放物線 $y=x^2$ に接するところで終わる。接点では放物線の接線の傾き $2x$ が $-1$ に等しく（第19講。ここでは判別式でもよい）、$x=-\frac12$、接点 $\left(-\frac12,\ \frac14\right)$ は $D$ の境界上にある。$k=-\frac12+\frac14=-\frac14$。
  \[\text{最大値}\ \Kotae{6},\quad\text{最小値}\ \Kotae{-\frac14}\]

  \begin{memo}
    判別式で確かめる：$x^2=-x+k$ が重解をもつ条件は $1+4k=0$、$k=-\frac14$。接点の $x=-\frac12$ が $-1\leqq x\leqq2$ に入っていることも確認が必要である（入っていなければ端点で最小）。
  \end{memo}

  \Shomon{4}
  \ascboxZ{逆像法：その点を通る $t$ が存在するか}
  点 $(X,\ Y)$ を1つ固定する。直線 $\ell_t$ がこの点を通るのは $Y=2tX-t^2$ のとき、つまり、
  \[t^2-2Xt+Y=0\quad\cdots\text{①}\]
  のときである。したがって、
  \par\noindent「$(X,Y)$ が通過範囲に入る」\par
  $\iff$「①を満たす実数 $t$ が\textgt{少なくとも1つ}存在する」\par
  $\iff$「$t$ の2次方程式①が実数解をもつ」\par
  $\iff \frac D4=X^2-Y\geqq0$\par
  よって、通過範囲は（$X,\ Y$ を $x,\ y$ に書き直して）
  \[\Kotae{y\leqq x^2\ \text{（放物線}\ y=x^2\ \text{とその下側）}}\]

  \begin{ptbs}{KEY}[通過範囲は「パラメータの存在条件」]
    動く図形の式をパラメータ（$t$）の方程式とみる。点 $(X,Y)$ が範囲に入る $\iff$ その方程式が「許された範囲に」解をもつ。第5講の解の配置がそのまま使える。
  \end{ptbs}

  \begin{memo}
    $y=2tx-t^2$ は、放物線 $y=x^2$ 上の点 $(t,\ t^2)$ における接線である。接線を全部集めると、放物線の下側がすき間なく塗られる。図形的にも答えと合う。
  \end{memo}

  \Shomon{5}
  \ascboxZ{$t\geqq0$ に解をもつ条件（解の配置）}
  今度は「①が $t\geqq0$ の範囲に少なくとも1つ解をもつ」条件を求める。①の2解の積は $Y$、和は $2X$ である（解と係数の関係）。$Y$ の符号で分ける。

  \textgt{$Y\leqq0$ のとき}：$\frac D4=X^2-Y\geqq0$ なので実数解をもつ。2解の積 $Y\leqq0$ なので、2解は異符号か、一方が0。いずれにせよ0以上の解がある。よって $Y\leqq0$ の点はすべて範囲に入る。

  \textgt{$Y>0$ のとき}：2解の積が正なので、実数解があれば2つは同符号。$t\geqq0$ の解をもつには、実数解をもち（$X^2\geqq Y$）、かつ2解が正（和 $2X>0$）であればよい。逆にこの2条件のもとでは正の解がある。

  2つの場合を合わせて、通過範囲は、
  \[\Kotae{y\leqq0\ \text{または}\ \left(x>0\ \text{かつ}\ y\leqq x^2\right)}\]
  図示すると次の灰色の部分（境界を含む）。

  \begin{zuhyo}
  \begin{tikzpicture}[x=.6cm,y=.42cm,line width=.5pt]
    \fill[black!15] (-2.3,-1.8) -- (-2.3,0) -- (0,0) -- plot[domain=0:2,samples=30] (\x,{(\x)^2}) -- (2,-1.8) -- cycle;
    \draw[->] (-2.5,0) -- (2.8,0) node[right] {$x$};
    \draw[->] (0,-1.9) -- (0,4.6) node[above] {$y$};
    \draw[line width=1pt] (-2.3,0) -- (0,0);
    \draw[line width=1pt,domain=0:2.05,samples=30] plot (\x,{(\x)^2});
    \draw[dashed,domain=-2.05:0,samples=30] plot (\x,{(\x)^2});
    \node[below left,font=\footnotesize] at (0,0) {$\mathrm O$};
    \node[font=\footnotesize,anchor=west] at (2.1,3) {$y=x^2$};
  \end{tikzpicture}
  \end{zuhyo}

  \begin{memo}
    \textgt{順像法での確かめ}：$x$ を固定して $y=-(t-x)^2+x^2$ を $t\geqq0$ で動かす。$x\geqq0$ なら $t=x$ で最大 $x^2$ をとり、$y$ は $x^2$ 以下のすべての値をとる。$x<0$ なら $t\geqq0$ で減少し、最大は $t=0$ の $0$。$y$ は $0$ 以下の値をとる。同じ答えになる。
  \end{memo}

  \begin{simple}[出典：オリジナル]{参考問題15}
    $a$ が $0\leqq a\leqq1$ の範囲を動くとき、直線 $y=ax-a^2$ が通過する範囲のうち、$x=1$ 上にある部分（$y$ の範囲）を求めよ。
  \end{simple}
  \Kai
  $x=1$ を固定すると $y=a-a^2=-\left(a-\frac12\right)^2+\frac14$。$0\leqq a\leqq1$ で、$a=\frac12$ のとき最大 $\frac14$、$a=0,\ 1$ のとき最小 $0$。$a$ は連続的に動くので、$y$ はその間のすべての値をとる。
  \[\Kotae{0\leqq y\leqq\frac14}\]
  \begin{memo}
    $x$ を1つ固定して $y$ の範囲を調べ、それを各 $x$ について並べると通過範囲になる。これが順像法（「ファクシミリ論法」と呼ばれることもある）である。
  \end{memo}

  \begin{ptbs}{CHECK}[翌日、解説を隠して答える（答えは本講の中にある）]
    \kakko{1}　軌跡の計算で同値でない変形をしたら何をするか。\par
    \kakko{2}　通過範囲を求めるとき、点 $(X,Y)$ の条件をどう言い換えるか。\par
    \kakko{3}　パラメータに $t\geqq0$ の制限があるときの処理を言え。
  \end{ptbs}

  \notefill
\end{multicols*}

%% ==================== 第16講 ====================
\Kou{16}{三角関数(1) 定義と加法定理}
\begin{ascolorbox4A}[出典：オリジナル]{例題16A}
  \kakko{1}　$150^\circ$ を弧度法で表せ。また、半径2、中心角 $\dfrac\pi3$ の扇形の弧の長さと面積を求めよ。\par
  \kakko{2}　$\sin\dfrac{7\pi}6$、$\cos\left(-\dfrac\pi4\right)$、$\tan\dfrac{5\pi}3$ の値を求めよ。\par
  \kakko{3}　$\sin75^\circ$ と $\cos15^\circ$ の値を求めよ。
\end{ascolorbox4A}

\begin{ascolorbox4A}[出典：東京大学 1999年度]{例題16B}
  \kakko{1}　一般角 $\theta$ に対して $\sin\theta$、$\cos\theta$ の定義を述べよ。\par
  \kakko{2}　\kakko{1}で述べた定義にもとづき、一般角 $\alpha$、$\beta$ に対して
  \[\sin(\alpha+\beta)=\sin\alpha\cos\beta+\cos\alpha\sin\beta,\qquad\cos(\alpha+\beta)=\cos\alpha\cos\beta-\sin\alpha\sin\beta\]
  を証明せよ。
\end{ascolorbox4A}

\begin{Data}
  テーマ：弧度法、一般角と単位円による定義、グラフと周期、加法定理の証明\par
  難易度：例題16A (1)(2) A　(3) B　例題16B C\quad 前提：第7・14講\par
  この講のゴール：三角関数の性質をすべて「単位円上の点の座標」から説明でき、加法定理を自力で証明できる
\end{Data}

\begin{multicols*}{2}
  \Shomon{1}
  \ascboxZ{弧度法：角を「単位円の弧の長さ」で測る}
  単位円で、中心角 $\theta$ に対する弧の長さを、その角の大きさとする。これが\textgt{弧度法}（単位はラジアン、普通は省略）である。一周 $360^\circ$ の弧の長さは $2\pi$ なので、
  \[180^\circ=\pi,\qquad 1^\circ=\frac\pi{180}\]
  よって $150^\circ=150\cdot\frac\pi{180}=\Kotae{\frac{5\pi}6}$。\par
  半径 $r$ の円では弧の長さも $r$ 倍になるので、扇形の弧の長さは $\ell=r\theta$、面積は円全体 $\pi r^2$ の $\frac\theta{2\pi}$ 倍で $S=\frac12r^2\theta$。
  \[\ell=2\cdot\frac\pi3=\Kotae{\frac{2\pi}3},\quad S=\frac12\cdot4\cdot\frac\pi3=\Kotae{\frac{2\pi}3}\]

  \begin{memo}
    なぜ弧度法を使うのか。数学Ⅲで $(\sin x)'=\cos x$ という簡単な公式が成り立つのは、角を弧度法で測ったときだけだからである（第31講）。度数法では余計な定数 $\frac\pi{180}$ が付く。
  \end{memo}

  \Shomon{2}
  \ascboxZ{一般角も単位円の座標で読む}
  第7講の定義をそのまま広げる。点 $(1,\ 0)$ を原点のまわりに角 $\theta$ だけ回した点を P とし、P$(\cos\theta,\ \sin\theta)$ とする。$\theta>0$ なら反時計回り、$\theta<0$ なら時計回りに回す。$2\pi$ より大きい角は、1周以上回る。\par
  $\frac{7\pi}6=\pi+\frac\pi6$：P は第3象限、$x$ 軸の負の向きから $\frac\pi6$ 下。座標 $\left(-\frac{\sqrt3}2,\ -\frac12\right)$。\par
  $-\frac\pi4$：時計回りに $45^\circ$。座標 $\left(\frac{\sqrt2}2,\ -\frac{\sqrt2}2\right)$。\par
  $\frac{5\pi}3=2\pi-\frac\pi3$：座標 $\left(\frac12,\ -\frac{\sqrt3}2\right)$、$\tan$ は傾き $\frac{-\sqrt3/2}{1/2}$。
  \[\Kotae{-\frac12},\quad\Kotae{\frac{\sqrt2}2},\quad\Kotae{-\sqrt3}\]

  \begin{koushiki}{単位円から読める性質}
    周期：$\sin(\theta+2\pi)=\sin\theta,\ \cos(\theta+2\pi)=\cos\theta$\par
    $x$ 軸対称：$\sin(-\theta)=-\sin\theta,\ \cos(-\theta)=\cos\theta$\par
    $\tan(\theta+\pi)=\tan\theta$（周期 $\pi$）\par
    範囲：$-1\leqq\sin\theta\leqq1,\ -1\leqq\cos\theta\leqq1$
  \end{koushiki}

  \begin{zuhyo}
  \begin{tikzpicture}[x=.45cm,y=.8cm,line width=.5pt]
    \draw[->] (-.3,0) -- (13.2,0) node[right] {$\theta$};
    \draw[->] (0,-1.3) -- (0,1.4) node[above] {$y$};
    \draw[domain=0:12.566,samples=80] plot (\x,{sin(\x r)});
    \draw[dashed,domain=0:12.566,samples=80] plot (\x,{cos(\x r)});
    \foreach \t/\l in {3.1416/\pi,6.2832/2\pi,9.4248/3\pi,12.566/4\pi} \node[below,font=\footnotesize] at (\t,0) {$\l$};
    \node[left,font=\footnotesize] at (0,1) {$1$};
    \node[left,font=\footnotesize] at (0,-1) {$-1$};
    \node[font=\footnotesize] at (2.4,1.25) {$y=\sin\theta$};
    \node[font=\footnotesize] at (7.4,1.25) {$y=\cos\theta$（破線）};
  \end{tikzpicture}
  \end{zuhyo}

  $y=\cos\theta$ のグラフは、$y=\sin\theta$ を左に $\frac\pi2$ ずらしたものである：$\cos\theta=\sin\left(\theta+\frac\pi2\right)$。

  \Shomon{3}
  \ascboxZ{知っている角の和・差に分ける}
  $75^\circ=45^\circ+30^\circ$、$15^\circ=45^\circ-30^\circ$ と分け、加法定理（下で証明する）を使う。
  \[
  \begin{aligned}
  \sin75^\circ&=\sin45^\circ\cos30^\circ+\cos45^\circ\sin30^\circ\\
  &=\frac{\sqrt2}2\cdot\frac{\sqrt3}2+\frac{\sqrt2}2\cdot\frac12=\Kotae{\frac{\sqrt6+\sqrt2}4}
  \end{aligned}
  \]
  \[
  \begin{aligned}
  \cos15^\circ&=\cos45^\circ\cos30^\circ+\sin45^\circ\sin30^\circ\\
  &=\Kotae{\frac{\sqrt6+\sqrt2}4}
  \end{aligned}
  \]
  $\cos15^\circ=\sin75^\circ$ になるのは、$15^\circ+75^\circ=90^\circ$ だからである（直角三角形の2つの鋭角）。

  \Shomon{B}
  \ascboxZ{定義を正確に書く}
  \textgt{\kakko{1}の解答}　座標平面上で、原点 O を中心とする半径1の円を考える。点 $(1,\ 0)$ を、原点のまわりに角 $\theta$ だけ回転した点を P とする（$\theta>0$ のときは反時計回り、$\theta<0$ のときは時計回りに $|\theta|$ だけ回転する）。このとき、P の $x$ 座標を $\cos\theta$、$y$ 座標を $\sin\theta$ と定める。

  \ascboxZ{距離を2通りに計算する}
  \textgt{\kakko{2}の解答}　定義より、角 $\alpha,\ \beta$ に対応する単位円上の点は P$(\cos\alpha,\ \sin\alpha)$、Q$(\cos\beta,\ \sin\beta)$ である。\par
  \textgt{手順1}：$\mathrm{PQ}^2$ を2点間の距離の公式で計算する。$\cos^2+\sin^2=1$ を使って、
  \[\mathrm{PQ}^2=2-2(\cos\alpha\cos\beta+\sin\alpha\sin\beta)\]
  \textgt{手順2}：P、Q を原点のまわりに $-\beta$ 回転すると、P$'(\cos(\alpha-\beta),\ \sin(\alpha-\beta))$、Q$'(1,\ 0)$ に移る。回転で2点間の距離は変わらないので、
  \[\mathrm{PQ}^2=\{\cos(\alpha-\beta)-1\}^2+\sin^2(\alpha-\beta)=2-2\cos(\alpha-\beta)\]
  \textgt{手順3}：2つを比べて、
  \[\cos(\alpha-\beta)=\cos\alpha\cos\beta+\sin\alpha\sin\beta\ \cdots\text{①}\]
  ①で $\beta$ を $-\beta$ に置き換える。定義（$x$ 軸に関する対称性）より $\cos(-\beta)=\cos\beta$、$\sin(-\beta)=-\sin\beta$ なので、
  \[\cos(\alpha+\beta)=\cos\alpha\cos\beta-\sin\alpha\sin\beta\]
  \textgt{手順4}：①で $\alpha=\frac\pi2$ とすると、$\cos\frac\pi2=0,\ \sin\frac\pi2=1$ より $\cos\left(\frac\pi2-\beta\right)=\sin\beta$。$\beta$ を $\frac\pi2-\theta$ にすると $\sin\left(\frac\pi2-\theta\right)=\cos\theta$ も得る。これらを使って、
  \[
  \begin{aligned}
  \sin(\alpha+\beta)&=\cos\left\{\left(\frac\pi2-\alpha\right)-\beta\right\}\\
  &=\cos\left(\frac\pi2-\alpha\right)\cos\beta+\sin\left(\frac\pi2-\alpha\right)\sin\beta\\
  &=\sin\alpha\cos\beta+\cos\alpha\sin\beta
  \end{aligned}
  \]
  以上で証明された。\QED

  \begin{ptbs}{KEY}[証明に使ってよいのは「定義」と「それから示したこと」だけ]
    加法定理の証明に、加法定理から導かれる公式（倍角の公式など）を使うと、循環論法になって0点である。何を使ったかを毎行で意識する。
  \end{ptbs}

  \begin{memo}
    手順2の「回転で距離は変わらない」は、三角形 OPQ と OP$'$Q$'$ が、2辺（半径1）とその間の角が等しく合同であることによる。答案ではこの理由を一言書き添える。
  \end{memo}

  \begin{simple}[出典：オリジナル]{参考問題16}
    $\tan\alpha=2$、$\tan\beta=3$（$0<\alpha<\frac\pi2,\ 0<\beta<\frac\pi2$）のとき、$\alpha+\beta$ の値を求めよ。
  \end{simple}
  \Kai
  加法定理を割り算して得られる $\tan(\alpha+\beta)=\dfrac{\tan\alpha+\tan\beta}{1-\tan\alpha\tan\beta}$ を使う。
  \[\tan(\alpha+\beta)=\frac{2+3}{1-6}=-1\]
  $0<\alpha+\beta<\pi$ で $\tan$ の値が $-1$ になるのは $\frac{3\pi}4$ だけである。
  \[\alpha+\beta=\Kotae{\frac{3\pi}4}\]
  \begin{memo}
    「$\tan=-1$ だから $\frac{3\pi}4$」と書くだけでは不十分である。角の範囲を先に決めて、その範囲で値が1つに決まることを確認してから答える。
  \end{memo}

  \begin{ptbs}{CHECK}[翌日、解説を隠して答える（答えは本講の中にある）]
    \kakko{1}　弧度法を使う理由は何か。\par
    \kakko{2}　一般角の $\sin\theta,\ \cos\theta$ の定義を正確に言え。\par
    \kakko{3}　加法定理の証明で、何を2通りに計算するか。
  \end{ptbs}

  \notefill
\end{multicols*}

%% ==================== 第17講 ====================
\Kou{17}{三角関数(2) 合成と応用}
\begin{ascolorbox4A}[出典：オリジナル]{例題17}
  \kakko{1}　$\sin\theta=\dfrac35$（$0<\theta<\frac\pi2$）のとき、$\cos\theta$、$\sin2\theta$、$\cos2\theta$ の値を求めよ。\par
  \kakko{2}　$0\leqq\theta\leqq\pi$ のとき、$\sqrt3\sin\theta+\cos\theta$ の最大値と最小値を求めよ。\par
  \kakko{3}　$0\leqq\theta<2\pi$ のとき、方程式 $\cos2\theta+\sin\theta=0$ を解け。\par
  \kakko{4}　$0\leqq\theta\leqq\pi$ のとき、$y=\sin\theta+\cos\theta+\sin\theta\cos\theta$ の最大値と最小値を求めよ。
\end{ascolorbox4A}

\begin{Data}
  テーマ：2倍角の公式、三角関数の合成、2次方程式に帰着、$t=\sin\theta+\cos\theta$ の置き換え\par
  難易度：(1) A　(2) B　(3) B　(4) C\quad 前提：第4・16講\par
  この講のゴール：三角関数の式を「1種類の三角関数」にそろえて、2次関数・2次方程式の問題に直せる
\end{Data}

\begin{multicols*}{2}
  \begin{koushiki}{2倍角・半角の公式}
    $\sin2\theta=2\sin\theta\cos\theta$\par
    $\cos2\theta=\cos^2\theta-\sin^2\theta=1-2\sin^2\theta=2\cos^2\theta-1$\par
    $\sin^2\theta=\dfrac{1-\cos2\theta}2,\quad\cos^2\theta=\dfrac{1+\cos2\theta}2$
  \end{koushiki}
  どれも加法定理で $\alpha=\beta=\theta$ とした式と、$\sin^2\theta+\cos^2\theta=1$ から作れる。覚えていなくても1行で復元できる。

  \Shomon{1}
  \ascboxZ{$\cos\theta$ の符号は範囲で決める}
  $\cos^2\theta=1-\sin^2\theta=\frac{16}{25}$。$0<\theta<\frac\pi2$ では $\cos\theta>0$ なので $\cos\theta=\frac45$。
  \[\cos\theta=\Kotae{\frac45}\]
  \[\sin2\theta=2\cdot\frac35\cdot\frac45=\Kotae{\frac{24}{25}}\]
  \[\cos2\theta=1-2\cdot\frac9{25}=\Kotae{\frac7{25}}\]

  \Shomon{2}
  \ascboxZ{合成：$a\sin\theta+b\cos\theta=r\sin(\theta+\varphi)$}
  点 $(a,\ b)$ を座標平面にとり、原点からの距離を $r=\sqrt{a^2+b^2}$、$x$ 軸の正の向きとなす角を $\varphi$ とする。$a=r\cos\varphi,\ b=r\sin\varphi$ なので、加法定理より、
  \[
  \begin{aligned}
  a\sin\theta+b\cos\theta&=r(\sin\theta\cos\varphi+\cos\theta\sin\varphi)\\
  &=r\sin(\theta+\varphi)
  \end{aligned}
  \]

  \begin{zuhyo}
  \begin{tikzpicture}[x=.9cm,y=.9cm,line width=.5pt]
    \draw[->] (-.3,0) -- (2.2,0) node[right] {$x$};
    \draw[->] (0,-.3) -- (0,1.3) node[above] {$y$};
    \draw (0,0) -- (1.732,1) node[right,font=\footnotesize] {$(\sqrt3,\ 1)$};
    \draw[dashed] (1.732,1) -- (1.732,0);
    \draw (.45,0) arc (0:30:.45);
    \node[font=\footnotesize] at (.7,.17) {$\varphi$};
    \node[font=\footnotesize,above left] at (.9,.5) {$2$};
  \end{tikzpicture}
  \end{zuhyo}

  本問は $a=\sqrt3,\ b=1$ で、$r=2$、$\varphi=\frac\pi6$。よって $\sqrt3\sin\theta+\cos\theta=2\sin\left(\theta+\frac\pi6\right)$。\par
  $0\leqq\theta\leqq\pi$ のとき $\frac\pi6\leqq\theta+\frac\pi6\leqq\frac{7\pi}6$。この範囲で $\sin$ は、$\frac\pi2$ で最大1、$\frac{7\pi}6$ で最小 $-\frac12$。
  \[\text{最大値}\ \Kotae{2}\ \left(\theta=\frac\pi3\right)\]
  \[\text{最小値}\ \Kotae{-1}\ (\theta=\pi)\]

  \begin{ptbs}{KEY}[合成したら「中身の範囲」を必ず書く]
    $\theta$ の範囲から $\theta+\varphi$ の範囲を作り、単位円上でその弧を見て最大・最小を読む。中身の範囲を書かずに $\pm r$ を答えるのが最も多い誤り。
  \end{ptbs}

  \Shomon{3}
  \ascboxZ{1種類の三角関数にそろえる}
  $\cos2\theta$ と $\sin\theta$ が混じっている。$\cos2\theta=1-2\sin^2\theta$ を使えば $\sin\theta$ だけの式になる。
  \[1-2\sin^2\theta+\sin\theta=0\]
  $s=\sin\theta$ とおくと（$-1\leqq s\leqq1$）、
  \[2s^2-s-1=(2s+1)(s-1)=0\]
  $s=1$ または $s=-\frac12$。どちらも範囲内。\par
  $\sin\theta=1$：$\theta=\frac\pi2$。\par
  $\sin\theta=-\frac12$：単位円で $y$ 座標が $-\frac12$ の点は2つあり、$\theta=\frac{7\pi}6,\ \frac{11\pi}6$。
  \[\theta=\Kotae{\frac\pi2,\ \frac{7\pi}6,\ \frac{11\pi}6}\]

  \Shomon{4}
  \ascboxZ{$t=\sin\theta+\cos\theta$ とおく}
  $(\sin\theta+\cos\theta)^2=1+2\sin\theta\cos\theta$ なので、$t=\sin\theta+\cos\theta$ とおくと
  \[\sin\theta\cos\theta=\frac{t^2-1}2\]
  と $t$ で表せる。よって
  \[y=t+\frac{t^2-1}2=\frac12(t+1)^2-1\]
  \textgt{$t$ の範囲}：合成すると $t=\sqrt2\sin\left(\theta+\frac\pi4\right)$。$0\leqq\theta\leqq\pi$ より $\frac\pi4\leqq\theta+\frac\pi4\leqq\frac{5\pi}4$ で、$\sin$ は $-\frac{\sqrt2}2$ 以上1以下の値をとる。よって $-1\leqq t\leqq\sqrt2$。\par
  この範囲で $y=\frac12(t+1)^2-1$ は増加するので（軸 $t=-1$ が左端）、\par
  $t=\sqrt2$（$\theta=\frac\pi4$）で最大 $\frac12(\sqrt2+1)^2-1=\frac12+\sqrt2$、\par
  $t=-1$（$\theta=\pi$）で最小 $-1$。
  \[\text{最大値}\ \Kotae{\frac12+\sqrt2},\quad\text{最小値}\ \Kotae{-1}\]

  \begin{memo}
    $t$ の範囲を $-\sqrt2\leqq t\leqq\sqrt2$ としてしまうと、最小値は変わらないが、問題によっては誤答になる。「おいた文字の範囲」は $\theta$ の範囲から毎回作り直す（第0講の約束）。
  \end{memo}

  \begin{simple}[出典：オリジナル]{参考問題17}
    $0\leqq\theta<2\pi$ のとき、不等式 $\sin\theta>\cos\theta$ を解け。
  \end{simple}
  \Kai
  $\sin\theta-\cos\theta>0$。合成すると（点 $(1,\ -1)$ で $r=\sqrt2,\ \varphi=-\frac\pi4$）、
  \[\sqrt2\sin\left(\theta-\frac\pi4\right)>0\]
  $-\frac\pi4\leqq\theta-\frac\pi4<\frac{7\pi}4$ の範囲で $\sin>0$ となるのは $0<\theta-\frac\pi4<\pi$。よって、
  \[\Kotae{\frac\pi4<\theta<\frac{5\pi}4}\]
  \begin{memo}
    単位円で「直線 $y=x$ より上にある点」と読んでもよい。$\sin\theta>\cos\theta$ は「$y$ 座標 $>$ $x$ 座標」だからである。
  \end{memo}

  \begin{ptbs}{CHECK}[翌日、解説を隠して答える（答えは本講の中にある）]
    \kakko{1}　合成の $r$ と $\varphi$ はどう決めるか。\par
    \kakko{2}　$t=\sin\theta+\cos\theta$ のとき $\sin\theta\cos\theta$ を $t$ で表せ。\par
    \kakko{3}　置き換えた文字の範囲はどう作るか。
  \end{ptbs}

  \notefill
\end{multicols*}
%% ==================== 第18講 ====================
\Kou{18}{指数関数・対数関数}
\begin{ascolorbox4A}[出典：オリジナル]{例題18}
  \kakko{1}　$8^{\frac23}$ と $27^{-\frac13}$ の値を求めよ。\par
  \kakko{2}　$\log_2 12-\log_2 3$ と $\log_4 8$ の値を求めよ。\par
  \kakko{3}　方程式 $4^x-3\cdot2^{x+1}+8=0$ を解け。\par
  \kakko{4}　不等式 $\log_{\frac12}(x-1)>-2$ を解け。\par
  \kakko{5}　$\log_{10}2=0.3010$ とする。$2^{100}$ は何桁の整数か。また、その最高位の数字を求めよ。
\end{ascolorbox4A}

\begin{Data}
  テーマ：指数の拡張、対数の定義と性質、底の変換、置き換えによる方程式、真数条件、桁数\par
  難易度：(1)(2) A　(3) B　(4) B　(5) C\quad 前提：第3・5講\par
  この講のゴール：対数を「指数を取り出す道具」として理解し、大きな数の大きさを見積もれる
\end{Data}

\begin{multicols*}{2}
  \begin{koushiki}{指数の拡張（$a>0$）}
    $a^0=1,\quad a^{-n}=\dfrac1{a^n},\quad a^{\frac1n}=\sqrt[n]a,\quad a^{\frac mn}=\sqrt[n]{a^m}$\par
    指数法則：$a^pa^q=a^{p+q},\ (a^p)^q=a^{pq},\ (ab)^p=a^pb^p$
  \end{koushiki}
  \textgt{なぜこう決めるか}：指数法則 $a^pa^q=a^{p+q}$ を、0や負の数や分数の指数でも成り立たせたいからである。例えば $a^0a^1=a^{0+1}=a$ となるには $a^0=1$ でなければならない。$a^{\frac12}a^{\frac12}=a^1$ となるには、$a^{\frac12}$ は「2乗して $a$ になる正の数」$\sqrt a$ でなければならない。

  \Shomon{1}
  \ascboxZ{底を素数の累乗にそろえる}
  \[8^{\frac23}=(2^3)^{\frac23}=2^{3\cdot\frac23}=2^2=\Kotae{4}\]
  \[27^{-\frac13}=(3^3)^{-\frac13}=3^{-1}=\Kotae{\frac13}\]

  \Shomon{2}
  \ascboxZ{対数は「指数を答える記号」}
  \begin{koushiki}{対数の定義と性質}
    $a>0,\ a\neq1,\ M>0$ のとき、$a^p=M\iff p=\log_aM$（$a$ を底、$M$ を真数という）。\par
    $\log_aMN=\log_aM+\log_aN$\par
    $\log_a\dfrac MN=\log_aM-\log_aN,\quad\log_aM^r=r\log_aM$\par
    底の変換：$\log_aM=\dfrac{\log_bM}{\log_ba}$
  \end{koushiki}
  $\log_28$ は「2を何乗すると8か」の答えで、$3$ である。性質はすべて指数法則の言い換えである。例えば $M=a^p,\ N=a^q$ とすると $MN=a^{p+q}$ なので $\log_aMN=p+q$。\par
  底の変換は、$p=\log_aM$ とおいて $a^p=M$ の両辺の $\log_b$ をとると $p\log_ba=\log_bM$ となることから導ける。
  \[\log_212-\log_23=\log_2\frac{12}3=\log_24=\Kotae{2}\]
  \[\log_48=\frac{\log_28}{\log_24}=\Kotae{\frac32}\]

  \begin{memo}
    確かめ：$4^{\frac32}=(\sqrt4)^3=8$。対数の値が出たら、指数に戻して確かめると誤りがすぐ見つかる。
  \end{memo}

  \Shomon{3}
  \ascboxZ{$t=2^x$ とおいて2次方程式に}
  $4^x=(2^2)^x=(2^x)^2$、$2^{x+1}=2\cdot2^x$ なので、$t=2^x$ とおく。$2^x$ はどんな $x$ でも正なので、\textgt{$t>0$}。
  \[t^2-6t+8=(t-2)(t-4)=0\]
  $t=2,\ 4$（どちらも $t>0$ を満たす）。$2^x=2$ より $x=1$、$2^x=4$ より $x=2$。
  \[x=\Kotae{1,\ 2}\]

  \Shomon{4}
  \ascboxZ{真数条件、底が1より小さいと向きが逆}
  \begin{zuhyo}
  \begin{tikzpicture}[x=.55cm,y=.55cm,line width=.5pt]
    \draw[->] (-2.3,0) -- (4.6,0) node[right] {$x$};
    \draw[->] (0,-2.3) -- (0,4.4) node[above] {$y$};
    \draw[domain=-2.2:2.05,samples=40] plot (\x,{2^(\x)});
    \draw[domain=0.22:4.2,samples=40] plot (\x,{ln(\x)/ln(2)});
    \draw[dashed] (-2,-2) -- (4.2,4.2);
    \node[font=\footnotesize] at (1.35,4.1) {$y=2^x$};
    \node[font=\footnotesize] at (3.6,1.4) {$y=\log_2x$};
    \node[font=\footnotesize] at (-1.2,-1.8) {$y=x$};
    \node[below left,font=\footnotesize] at (0,0) {$\mathrm O$};
  \end{tikzpicture}
  \end{zuhyo}
  $y=\log_ax$ は $y=a^x$ の逆関数で、グラフは直線 $y=x$ について対称である。$a>1$ なら増加、$0<a<1$ なら減少する。\par
  \textgt{真数条件}：真数は正なので $x-1>0$、つまり $x>1$。\par
  $-2=\log_{\frac12}\left(\frac12\right)^{-2}=\log_{\frac12}4$ と書き直すと、不等式は $\log_{\frac12}(x-1)>\log_{\frac12}4$。底 $\frac12$ は1より小さく、$\log_{\frac12}$ は減少関数なので、真数の大小は逆になる。
  \[x-1<4\]
  真数条件と合わせて、
  \[\Kotae{1<x<5}\]

  \begin{ptbs}{KEY}[対数の不等式は「真数条件」と「底と1の大小」]
    最初に真数条件を書く。底が1より大きければ向きはそのまま、1より小さければ逆。この2点を落とすと答えが変わる。
  \end{ptbs}

  \Shomon{5}
  \ascboxZ{桁数は常用対数の整数部分から}
  正の数 $N$ が $n$ 桁の整数 $\iff 10^{n-1}\leqq N<10^n\iff n-1\leqq\log_{10}N<n$。\par
  $\log_{10}2^{100}=100\times0.3010=30.10$ なので $30\leqq\log_{10}2^{100}<31$。よって
  \[\Kotae{31\ \text{桁}}\]
  最高位は小数部分 $0.10$ で決まる。$2^{100}=10^{30.10}=10^{30}\times10^{0.10}$。$\log_{10}1=0<0.10<0.3010=\log_{10}2$ より $1<10^{0.10}<2$。$10^{0.10}$ の整数部分が最高位なので、
  \[\text{最高位の数字}\ \Kotae{1}\]

  \begin{memo}
    $\log_{10}2=0.3010$ は近似値である。厳密には $0.30100\cdots$ と $0.30105$ の間などの幅があるが、本問の結論（30.10 の整数部分と、0.10 が0と0.301の間にあること）は、その幅では変わらない。東大では近似の幅を自分で確かめさせることがある。
  \end{memo}

  \begin{simple}[出典：オリジナル]{参考問題18}
    $\log_{10}2=0.3010$、$\log_{10}3=0.4771$ とする。$2^{30}$、$3^{20}$、$5^{12}$ を小さい順に並べよ。
  \end{simple}
  \Kai
  常用対数をとって比べる（$\log_{10}$ は増加関数なので、対数の大小と元の数の大小は一致する）。$\log_{10}5=\log_{10}\frac{10}2=1-0.3010=0.6990$ に注意する。
  \[
  \begin{aligned}
  \log_{10}2^{30}&=30\times0.3010=9.030\\
  \log_{10}3^{20}&=20\times0.4771=9.542\\
  \log_{10}5^{12}&=12\times0.6990=8.388
  \end{aligned}
  \]
  よって、
  \[\Kotae{5^{12}<2^{30}<3^{20}}\]

  \begin{ptbs}{CHECK}[翌日、解説を隠して答える（答えは本講の中にある）]
    \kakko{1}　$a^0=1$、$a^{\frac12}=\sqrt a$ と決める理由は何か。\par
    \kakko{2}　対数の不等式を解くときの2つの注意は何か。\par
    \kakko{3}　$N$ が $n$ 桁であることを常用対数で書け。
  \end{ptbs}

  \notefill
\end{multicols*}

%% ==================== 第19講 ====================
\Kou{19}{微分法（数学Ⅱ）}
\begin{ascolorbox4A}[出典：オリジナル]{例題19}
  \kakko{1}　$f(x)=x^2$ の導関数を、定義に従って求めよ。\par
  \kakko{2}　曲線 $y=x^3-3x$ 上の点 $(2,\ 2)$ における接線の方程式を求めよ。\par
  \kakko{3}　関数 $f(x)=x^3-3x^2-9x+5$ の極値を求めよ。\par
  \kakko{4}　$a$ を実数の定数とする。方程式 $x^3-3x-a=0$ の異なる実数解の個数を調べよ。
\end{ascolorbox4A}

\begin{Data}
  テーマ：微分係数と導関数の定義、接線、増減表と極値、定数分離による解の個数\par
  難易度：(1) A　(2) A　(3) B　(4) C\quad 前提：第4・14講\par
  この講のゴール：「導関数＝接線の傾き」から増減表を作り、グラフを使って方程式の解の個数を数えられる
\end{Data}

\begin{multicols*}{2}
  \begin{koushiki}{微分係数と導関数}
    $x$ が $a$ から $a+h$ まで変わるときの平均変化率 $\dfrac{f(a+h)-f(a)}h$ は、2点を結ぶ直線（割線）の傾きである。$h\to0$ としたときの極限
    \[f'(a)=\lim_{h\to0}\frac{f(a+h)-f(a)}h\]
    を\textgt{微分係数}といい、曲線の点 $(a,\ f(a))$ における接線の傾きを表す。$a$ を $x$ に変えた関数 $f'(x)$ を\textgt{導関数}という。
  \end{koushiki}
  「$\lim_{h\to0}$」は、$h$ を0に限りなく近づけたときに近づく値を表す（$h=0$ を代入するのではない。$h=0$ では分母が0になる）。

  \Shomon{1}
  \ascboxZ{定義通りに計算し、$h$ で約分してから近づける}
  \[
  \begin{aligned}
  f'(x)&=\lim_{h\to0}\frac{(x+h)^2-x^2}h\\
  &=\lim_{h\to0}\frac{2xh+h^2}h=\lim_{h\to0}(2x+h)=\Kotae{2x}
  \end{aligned}
  \]
  $h\neq0$ の間は $h$ で約分してよい。約分した後の $2x+h$ は、$h$ を0に近づけると $2x$ に近づく。

  \begin{koushiki}{導関数の公式}
    $(x^n)'=nx^{n-1}$（$n$ は自然数）、$(c)'=0$（$c$ は定数）\par
    $\{kf(x)+lg(x)\}'=kf'(x)+lg'(x)$\par
    接線：点 $(a,\ f(a))$ における接線は $y-f(a)=f'(a)(x-a)$
  \end{koushiki}

  \Shomon{2}
  \ascboxZ{接線は「点と傾き」で決まる}
  $y'=3x^2-3$ なので、$x=2$ での傾きは $3\cdot4-3=9$。点 $(2,\ 2)$ を通り傾き9の直線は、
  \[y-2=9(x-2)\]
  よって、
  \[\Kotae{y=9x-16}\]

  \Shomon{3}
  \ascboxZ{$f'(x)$ の符号で増減表を作る}
  $f'(x)>0$ の区間では接線が右上がりなので $f(x)$ は増加、$f'(x)<0$ の区間では減少する。増加から減少に変わる点の値を\textgt{極大値}、減少から増加に変わる点の値を\textgt{極小値}という。
  \[f'(x)=3x^2-6x-9=3(x+1)(x-3)\]

  \begin{zuhyo}
  \begin{tblr}{colspec={c|ccccc}, hlines, vlines, row{1}={font=\TRgt}}
  $x$ & $\cdots$ & $-1$ & $\cdots$ & $3$ & $\cdots$ \\
  $f'(x)$ & $+$ & $0$ & $-$ & $0$ & $+$ \\
  $f(x)$ & $\nearrow$ & $10$ & $\searrow$ & $-22$ & $\nearrow$ \\
  \end{tblr}
  \end{zuhyo}

  $f(-1)=-1-3+9+5=10$、$f(3)=27-27-27+5=-22$。よって、
  \[\Kotae{\text{極大値}\ 10\ (x=-1),\ \text{極小値}\ -22\ (x=3)}\]

  \begin{ptbs}{注意}[$f'(a)=0$ でも極値とは限らない]
    $f(x)=x^3$ は $f'(0)=0$ だが、$x=0$ の前後でずっと増加しているので極値ではない。極値の判定は「$f'$ の符号が変わるか」で行う。だから増減表を書く。
  \end{ptbs}

  \Shomon{4}
  \ascboxZ{定数分離：$a=g(x)$ とグラフの交点}
  方程式を $x^3-3x=a$ と変形する。実数解の個数は、曲線 $y=x^3-3x$ と直線 $y=a$ の共有点の個数に等しい（共有点の $x$ 座標が解）。文字 $a$ を片側に集めたので、曲線は動かず、水平な直線だけが上下する。\par
  $g(x)=x^3-3x$ とおくと $g'(x)=3(x+1)(x-1)$。$g(-1)=2$（極大）、$g(1)=-2$（極小）。

  \begin{zuhyo}
  \begin{tikzpicture}[x=.75cm,y=.42cm,line width=.5pt]
    \draw[->] (-2.4,0) -- (2.6,0) node[right] {$x$};
    \draw[->] (0,-3.4) -- (0,3.6) node[above] {$y$};
    \draw[domain=-2.15:2.15,samples=60] plot (\x,{(\x)^3-3*(\x)});
    \draw[dashed] (-2.4,2) -- (2.4,2) node[right,font=\footnotesize] {$y=2$};
    \draw[dashed] (-2.4,-2) -- (2.4,-2) node[right,font=\footnotesize] {$y=-2$};
    \draw (-2.4,.8) -- (2.4,.8) node[right,font=\footnotesize] {$y=a$};
    \node[below left,font=\footnotesize] at (0,0) {$\mathrm O$};
  \end{tikzpicture}
  \end{zuhyo}

  直線 $y=a$ を上下に動かして、共有点の数を数える。
  \[\Kotae{\begin{cases}a<-2,\ 2<a\ \text{のとき}\ 1\ \text{個}\\a=\pm2\ \text{のとき}\ 2\ \text{個}\\-2<a<2\ \text{のとき}\ 3\ \text{個}\end{cases}}\]

  \begin{ptbs}{KEY}[解の個数は「文字定数を分離してグラフの交点」]
    $f(x)=a$ の形にすれば、動くのは水平線だけ。極値と端の様子（$x\to\pm\infty$）さえわかれば、個数が読み取れる。
  \end{ptbs}

  \begin{simple}[出典：オリジナル]{参考問題19}
    点 $(2,\ a)$ から曲線 $y=x^3-3x$ に引ける接線の本数を、$a$ の値で場合分けして求めよ。
  \end{simple}
  \Kai
  接点を $(t,\ t^3-3t)$ とおく。接線は $y=(3t^2-3)(x-t)+t^3-3t$、整理して $y=(3t^2-3)x-2t^3$。これが点 $(2,\ a)$ を通る条件は、
  \[a=-2t^3+6t^2-6\]
  3次関数のグラフでは、接点が異なれば接線も異なる（1本の直線が2点で接することはない）ので、接線の本数はこの $t$ の方程式の実数解の個数に等しい。$h(t)=-2t^3+6t^2-6$ とおくと $h'(t)=-6t(t-2)$。$h(0)=-6$（極小）、$h(2)=2$（極大）。$y=a$ との交点を数えて、
  \[\Kotae{\begin{cases}a<-6,\ 2<a:\ 1\ \text{本}\\a=-6,\ 2:\ 2\ \text{本}\\-6<a<2:\ 3\ \text{本}\end{cases}}\]
  \begin{memo}
    「接線の本数」を「接点の個数」に置き換えるには、接点と接線が1対1に対応することが必要である。4次関数などでは1本の直線が2点で接することがあり、この置き換えはそのままでは使えない。
  \end{memo}

  \begin{ptbs}{CHECK}[翌日、解説を隠して答える（答えは本講の中にある）]
    \kakko{1}　微分係数の定義と、その図形的な意味を言え。\par
    \kakko{2}　$f'(a)=0$ でも極値にならない例を挙げよ。\par
    \kakko{3}　方程式の解の個数を定数分離で数える手順を言え。
  \end{ptbs}

  \notefill
\end{multicols*}

%% ==================== 第20講 ====================
\Kou{20}{積分法（数学Ⅱ）}
\begin{ascolorbox4A}[出典：オリジナル]{例題20}
  \kakko{1}　不定積分 $\displaystyle\int(3x^2-4x+1)\,dx$ と、定積分 $\displaystyle\int_1^2(3x^2-4x+1)\,dx$ を求めよ。\par
  \kakko{2}　放物線 $y=x^2-2x$ と $x$ 軸で囲まれた部分の面積を求めよ。\par
  \kakko{3}　放物線 $y=x^2$ と直線 $y=x+2$ で囲まれた部分の面積を求めよ。\par
  \kakko{4}　$0\leqq x\leqq1$ とする。$f(x)=\displaystyle\int_0^1|t-x|\,dt$ を $x$ の式で表し、$f(x)$ の最小値を求めよ。
\end{ascolorbox4A}

\begin{Data}
  テーマ：不定積分と定積分、面積、$\frac16$ 公式、絶対値を含む関数の定積分\par
  難易度：(1) A　(2) A　(3) B　(4) C\quad 前提：第19講\par
  この講のゴール：「積分は微分の逆」と「積分は面積」の2つの顔をつなげ、面積の問題を素早く処理できる
\end{Data}

\begin{multicols*}{2}
  \begin{koushiki}{不定積分と定積分}
    微分すると $f(x)$ になる関数 $F(x)$ を $f(x)$ の\textgt{原始関数}という。原始関数は定数の差を除いて1つに決まり、$\displaystyle\int f(x)\,dx=F(x)+C$（$C$ は積分定数）と書く。\par
    定積分：$\displaystyle\int_a^bf(x)\,dx=\Big[F(x)\Big]_a^b=F(b)-F(a)$\par
    $\displaystyle\int x^n\,dx=\frac{x^{n+1}}{n+1}+C$（$n=0,\ 1,\ 2,\ \cdots$）
  \end{koushiki}

  \Shomon{1}
  \ascboxZ{微分の逆をたどる}
  $x^3,\ x^2,\ x$ を微分すると $3x^2,\ 2x,\ 1$ になることから、
  \[\int(3x^2-4x+1)\,dx=\Kotae{x^3-2x^2+x+C}\]
  （$C$ は積分定数）。定積分は $C$ を省いた原始関数に上端と下端を入れて引く。
  \[\int_1^2(3x^2-4x+1)\,dx=\Big[x^3-2x^2+x\Big]_1^2\]
  \[=(8-8+2)-(1-2+1)=\Kotae{2}\]

  \ascboxZ{なぜ定積分が面積になるのか}
  $f(x)\geqq0$ とし、$a$ から $x$ までのグラフの下の面積を $S(x)$ とする。$x$ を $h$ だけ増やすと、面積は幅 $h$、高さほぼ $f(x)$ の細い帯の分だけ増える。
  \[\frac{S(x+h)-S(x)}h\fallingdotseq f(x)\]
  $h\to0$ とすると $S'(x)=f(x)$、つまり $S(x)$ は $f(x)$ の原始関数である。$S(a)=0$ なので、$S(b)=S(b)-S(a)=F(b)-F(a)$。これが「定積分 $=$ 面積」の理由である。

  \Shomon{2}
  \ascboxZ{面積は「上の式 $-$ 下の式」を積分}
  $x^2-2x=x(x-2)$ なので、放物線は $x=0,\ 2$ で $x$ 軸と交わり、$0<x<2$ では $x$ 軸より下にある。面積は「上（$y=0$）$-$ 下（$y=x^2-2x$）」を積分する。
  \[
  \begin{aligned}
  S&=\int_0^2\{0-(x^2-2x)\}\,dx=\Big[-\frac{x^3}3+x^2\Big]_0^2\\
  &=-\frac83+4=\Kotae{\frac43}
  \end{aligned}
  \]

  \begin{koushiki}{6分の1公式}
    \[\int_\alpha^\beta(x-\alpha)(x-\beta)\,dx=-\frac{(\beta-\alpha)^3}6\]
  \end{koushiki}
  \textgt{証明}：$h=\beta-\alpha$ とし、$u=x-\alpha$ とおきかえると、$x-\beta=u-h$ で、$x$ が $\alpha$ から $\beta$ まで動くとき $u$ は $0$ から $h$ まで動く。
  \[\int_0^hu(u-h)\,du=\Big[\frac{u^3}3-\frac{hu^2}2\Big]_0^h=\frac{h^3}3-\frac{h^3}2=-\frac{h^3}6\]
  （$u$ への置きかえは、グラフを左に $\alpha$ ずらしても面積が変わらないことにあたる。）\par
  検算：\kakko{2}は $\int_0^2(x-0)(x-2)\,dx=-\frac{8}6$ の符号を変えたもので、$\frac43$ と一致する。

  \Shomon{3}
  \ascboxZ{放物線と直線の間は6分の1公式}
  交点は $x^2=x+2$ より $x=-1,\ 2$。$-1<x<2$ では直線が上にある。
  \[
  \begin{aligned}
  S&=\int_{-1}^2\{(x+2)-x^2\}\,dx=-\int_{-1}^2(x+1)(x-2)\,dx\\
  &=\frac{\{2-(-1)\}^3}6=\Kotae{\frac92}
  \end{aligned}
  \]

  \begin{ptbs}{KEY}[放物線と直線（または放物線どうし）の面積]
    交点の $x$ 座標 $\alpha<\beta$ がわかれば、面積は $\dfrac{|a|(\beta-\alpha)^3}6$（$a$ は2つの式の差の $x^2$ の係数）。計算時間を大幅に減らせる。
  \end{ptbs}

  \Shomon{4}
  \ascboxZ{絶対値の中身の符号で積分区間を分ける}
  積分の変数は $t$ で、$x$ は定数として扱う。$|t-x|$ は $t\leqq x$ で $x-t$、$t\geqq x$ で $t-x$。$0\leqq x\leqq1$ なので、区間 $[0,\ 1]$ を $x$ で2つに分ける。
  \[
  \begin{aligned}
  f(x)&=\int_0^x(x-t)\,dt+\int_x^1(t-x)\,dt\\
  &=\Big[xt-\frac{t^2}2\Big]_0^x+\Big[\frac{t^2}2-xt\Big]_x^1\\
  &=\frac{x^2}2+\left(\frac12-x+\frac{x^2}2\right)
  \end{aligned}
  \]
  まとめて、
  \[f(x)=\Kotae{x^2-x+\frac12}\]
  平方完成すると $f(x)=\left(x-\frac12\right)^2+\frac14$。$x=\frac12$ は $0\leqq x\leqq1$ に入るので、
  \[\text{最小値}\ \Kotae{\frac14}\ \left(x=\frac12\right)\]

  \begin{memo}
    図で確かめる：$y=|t-x|$ のグラフは $t=x$ で折れ曲がる V 字で、面積は2つの直角二等辺三角形の和 $\frac{x^2}2+\frac{(1-x)^2}2$。$x=\frac12$（真ん中で折れる）のとき最小になるのは直観とも合う。
  \end{memo}

  \begin{simple}[出典：オリジナル]{参考問題20}
    関数 $f(x)$ が $f(x)=x^2+2\displaystyle\int_0^1f(t)\,dt$ を満たすとき、$f(x)$ を求めよ。
  \end{simple}
  \Kai
  $\int_0^1f(t)\,dt$ は区間が定数なので、ある定数である。これを $c$ とおくと $f(x)=x^2+2c$。$c$ の定義に代入すると、
  \[c=\int_0^1(t^2+2c)\,dt=\frac13+2c\]
  よって $c=-\frac13$。
  \[f(x)=\Kotae{x^2-\frac23}\]
  \begin{memo}
    上端・下端が定数の定積分は「数」、上端に $x$ がある定積分は「$x$ の関数」である。前者は定数とおき、後者は両辺を $x$ で微分する（$\frac{d}{dx}\int_a^xf(t)\,dt=f(x)$）。
  \end{memo}

  \begin{ptbs}{CHECK}[翌日、解説を隠して答える（答えは本講の中にある）]
    \kakko{1}　定積分が面積になる理由を言え。\par
    \kakko{2}　6分の1公式を書き、証明の方針を言え。\par
    \kakko{3}　$\int_0^1|t-x|\,dt$ で積分区間をどこで分けるか。
  \end{ptbs}

  \notefill
\end{multicols*}
%% ==================== 第21講 ====================
\Kou{21}{数列と和}
\begin{ascolorbox4A}[出典：オリジナル]{例題21}
  \kakko{1}　初項3、公差4の等差数列 $\{a_n\}$ の一般項と、初項から第 $n$ 項までの和を求めよ。\par
  \kakko{2}　初項2、公比3の等比数列 $\{b_n\}$ の一般項と、初項から第 $n$ 項までの和を求めよ。\par
  \kakko{3}　$\displaystyle\sum_{k=1}^nk(k+1)$ を求めよ。\par
  \kakko{4}　数列 $1,\ 2,\ 5,\ 10,\ 17,\ \cdots$ の一般項を求めよ。\par
  \kakko{5}　$\displaystyle\sum_{k=1}^n\frac1{k(k+1)}$ と $\displaystyle\sum_{k=1}^nk\cdot2^k$ を求めよ。
\end{ascolorbox4A}

\begin{Data}
  テーマ：等差数列・等比数列、$\Sigma$ の意味と公式、階差数列、差に分ける和、$S-rS$ の和\par
  難易度：(1)(2) A　(3) A　(4) B　(5) C\quad 前提：第1講\par
  この講のゴール：$\Sigma$ を「足し算の省略記号」として読み書きでき、和の公式の出どころを説明できる
\end{Data}

\begin{multicols*}{2}
  \begin{koushiki}{数列の言葉}
    数を1列に並べたものを\textgt{数列}といい、$n$ 番目の数（第 $n$ 項）を $a_n$ と書く。$a_n$ を $n$ の式で表したものを\textgt{一般項}という。\par
    \textgt{等差数列}：隣の項との差が一定（公差 $d$）。$a_n=a_1+(n-1)d$\par
    \textgt{等比数列}：隣の項との比が一定（公比 $r$）。$a_n=a_1r^{n-1}$
  \end{koushiki}
  $n-1$ になる理由：第1項から第 $n$ 項までは、公差を足す（公比を掛ける）操作を $n-1$ 回行う。植木算と同じである。

  \Shomon{1}
  \ascboxZ{等差数列の和は「逆順に並べて足す」}
  $a_n=3+(n-1)\cdot4=\Kotae{4n-1}$。\par
  和 $S_n=a_1+a_2+\cdots+a_n$ を逆順にも書いて縦に足すと、どの列も $a_1+a_n$ になる。それが $n$ 列あるので $2S_n=n(a_1+a_n)$。
  \[S_n=\frac{n(3+4n-1)}2=\Kotae{n(2n+1)}\]

  \Shomon{2}
  \ascboxZ{等比数列の和は「$S-rS$」}
  $b_n=\Kotae{2\cdot3^{n-1}}$。\par
  $S_n=2+2\cdot3+\cdots+2\cdot3^{n-1}$、$3S_n=2\cdot3+\cdots+2\cdot3^{n-1}+2\cdot3^n$。下から上を引くと、間の項がすべて消えて $2S_n=2\cdot3^n-2$。
  \[S_n=\Kotae{3^n-1}\]

  \begin{koushiki}{等比数列の和（$r\neq1$）}
    \[a+ar+\cdots+ar^{n-1}=\frac{a(r^n-1)}{r-1}\]
    $r=1$ のときは、すべての項が $a$ なので和は $na$。
  \end{koushiki}

  \Shomon{3}
  \ascboxZ{$\Sigma$ は分けて公式を使う}
  \begin{koushiki}{$\Sigma$ の意味と公式}
    $\displaystyle\sum_{k=1}^na_k=a_1+a_2+\cdots+a_n$（$k$ に1から $n$ まで入れて全部足す）\par
    $\displaystyle\sum(pa_k+qb_k)=p\sum a_k+q\sum b_k$\par
    $\displaystyle\sum_{k=1}^nc=nc,\quad\sum_{k=1}^nk=\frac{n(n+1)}2$\par
    $\displaystyle\sum_{k=1}^nk^2=\frac{n(n+1)(2n+1)}6,\quad\sum_{k=1}^nk^3=\left\{\frac{n(n+1)}2\right\}^2$
  \end{koushiki}
  \textgt{$\sum k^2$ の出どころ}：$(k+1)^3-k^3=3k^2+3k+1$ を $k=1$ から $n$ まで足すと、左辺は途中が消えて $(n+1)^3-1$ だけ残る。右辺は $3\sum k^2+3\sum k+n$。これを $\sum k^2$ について解くと公式になる。\par
  本問は、
  \[
  \begin{aligned}
  \sum_{k=1}^nk(k+1)&=\sum_{k=1}^nk^2+\sum_{k=1}^nk\\
  &=\frac{n(n+1)(2n+1)}6+\frac{n(n+1)}2\\
  &=\frac{n(n+1)}6(2n+1+3)
  \end{aligned}
  \]
  したがって、
  \[\sum_{k=1}^nk(k+1)=\Kotae{\frac{n(n+1)(n+2)}3}\]

  \begin{memo}
    検算：$n=1$ のとき左辺 $1\cdot2=2$、右辺 $\frac{1\cdot2\cdot3}3=2$。$n=2$ のとき左辺 $2+6=8$、右辺 $\frac{2\cdot3\cdot4}3=8$。和の公式は、必ず $n=1,\ 2$ で確かめる。
  \end{memo}

  \Shomon{4}
  \ascboxZ{差をとって規則を見つける（階差数列）}
  隣どうしの差をとると $1,\ 3,\ 5,\ 7,\ \cdots$ となり、等差数列 $b_k=2k-1$ である。第 $n$ 項は、第1項に差を $n-1$ 個足したものなので、$n\geqq2$ のとき、
  \[
  \begin{aligned}
  a_n&=a_1+\sum_{k=1}^{n-1}b_k=1+\sum_{k=1}^{n-1}(2k-1)\\
  &=1+(n-1)^2
  \end{aligned}
  \]
  （$\sum_{k=1}^{m}(2k-1)=m^2$ を $m=n-1$ で使った。）この式は $n=1$ のときも $a_1=1$ を満たす。
  \[a_n=\Kotae{n^2-2n+2}\]

  \begin{ptbs}{注意}[階差の公式は $n\geqq2$ でしか使えない]
    $\sum_{k=1}^{n-1}$ は $n=1$ のとき意味をもたない。$n\geqq2$ で求めた式が $n=1$ でも正しいかを、必ず別に確かめて書く。
  \end{ptbs}

  \Shomon{5}
  \ascboxZ{差に分けて途中を消す}
  $\dfrac1{k(k+1)}=\dfrac1k-\dfrac1{k+1}$（右辺を通分すれば確かめられる）。足すと、隣どうしが打ち消しあう。
  \[
  \begin{aligned}
  &\left(1-\frac12\right)+\left(\frac12-\frac13\right)+\cdots+\left(\frac1n-\frac1{n+1}\right)\\
  &=1-\frac1{n+1}
  \end{aligned}
  \]
  よって、
  \[\sum_{k=1}^n\frac1{k(k+1)}=\Kotae{\frac n{n+1}}\]

  \ascboxZ{（等差）$\times$（等比）の和は $S-2S$}
  $S=1\cdot2+2\cdot2^2+\cdots+n\cdot2^n$ とおく。2倍して1項ずらして書く。
  \[2S=\quad\ 1\cdot2^2+\cdots+(n-1)2^n+n\cdot2^{n+1}\]
  $S-2S$ を計算すると、係数がそろった所は $2^k$ だけ残る。
  \[
  \begin{aligned}
  -S&=2+2^2+\cdots+2^n-n\cdot2^{n+1}\\
  &=(2^{n+1}-2)-n\cdot2^{n+1}
  \end{aligned}
  \]
  よって、
  \[S=\Kotae{(n-1)2^{n+1}+2}\]
  検算：$n=2$ のとき $2+8=10$、式は $1\cdot8+2=10$。

  \begin{simple}[出典：オリジナル]{参考問題21}
    初項から第 $n$ 項までの和 $S_n$ が次の式で表される数列 $\{a_n\}$ の一般項を求めよ。\par
    \kakko{1}　$S_n=n^2+2n$\qquad\kakko{2}　$S_n=n^2+2n+1$
  \end{simple}
  \Kai
  $n\geqq2$ のとき $a_n=S_n-S_{n-1}$、$n=1$ のとき $a_1=S_1$ である。\par
  \kakko{1}　$n\geqq2$ で $a_n=(n^2+2n)-\{(n-1)^2+2(n-1)\}=2n+1$。$a_1=S_1=3$ はこの式で $n=1$ としたものと一致する。
  \[a_n=\Kotae{2n+1}\]
  \kakko{2}　$n\geqq2$ で同様に $a_n=2n+1$。しかし $a_1=S_1=4$ で、式の値3と一致しない。
  \[\Kotae{a_1=4,\quad a_n=2n+1\ (n\geqq2)}\]

  \begin{ptbs}{CHECK}[翌日、解説を隠して答える（答えは本講の中にある）]
    \kakko{1}　$\sum k^2$ の公式はどこから出てくるか。\par
    \kakko{2}　階差数列から一般項を求めたとき、必ず確かめることは何か。\par
    \kakko{3}　（等差）$\times$（等比）の和の求め方を言え。
  \end{ptbs}

  \notefill
\end{multicols*}

%% ==================== 第22講 ====================
\Kou{22}{漸化式}
\begin{ascolorbox4A}[出典：オリジナル]{例題22}
  次の条件で定められる数列 $\{a_n\}$ の一般項を求めよ。\par
  \kakko{1}　$a_1=1,\quad a_{n+1}=a_n+2n$\par
  \kakko{2}　$a_1=1,\quad a_{n+1}=3a_n+2$\par
  \kakko{3}　$a_1=1,\quad a_2=3,\quad a_{n+2}=3a_{n+1}-2a_n$\par
  \kakko{4}　$a_1=1,\quad a_{n+1}=2a_n+3^n$
\end{ascolorbox4A}

\begin{Data}
  テーマ：漸化式の基本形（階差型・$pa_n+q$ 型・3項間・指数の項を含む型）\par
  難易度：(1) A　(2) B　(3) B　(4) C\quad 前提：第21講\par
  この講のゴール：漸化式を見て「等差・等比・階差のどれに変形するか」を決め、確実に一般項まで運べる
\end{Data}

\begin{multicols*}{2}
  \begin{koushiki}{漸化式とは}
    前の項から次の項を決める規則を\textgt{漸化式}という。初項と漸化式があれば、数列のすべての項が順に決まる。「一般項を求める」とは、$a_n$ を $n$ だけの式で表すことである。
  \end{koushiki}
  解き方の目標は1つだけである。\textgt{等差数列・等比数列・階差数列のどれかの形に変形する}こと。

  \Shomon{1}
  \ascboxZ{$a_{n+1}-a_n=f(n)$ は階差型}
  $a_{n+1}-a_n=2n$ なので、階差数列が $b_n=2n$。$n\geqq2$ のとき、
  \[a_n=1+\sum_{k=1}^{n-1}2k=1+(n-1)n\]
  $n=1$ のときも $a_1=1$ を満たすので、
  \[a_n=\Kotae{n^2-n+1}\]

  \Shomon{2}
  \ascboxZ{$a_{n+1}=pa_n+q$ は「動かない数」を引く}
  もし $a_n$ がずっと同じ値 $\alpha$ なら、$\alpha=3\alpha+2$ を満たすはずで、$\alpha=-1$。この $\alpha$ を使うと、
  \[
  \begin{aligned}
  a_{n+1}&=3a_n+2\\
  \alpha&=3\alpha+2
  \end{aligned}
  \]
  の辺々を引いて、$a_{n+1}-\alpha=3(a_n-\alpha)$、つまり
  \[a_{n+1}+1=3(a_n+1)\]
  数列 $\{a_n+1\}$ は初項 $a_1+1=2$、公比3の等比数列である。$a_n+1=2\cdot3^{n-1}$ なので、
  \[a_n=\Kotae{2\cdot3^{n-1}-1}\]

  \begin{memo}
    $\alpha=3\alpha+2$ を「特性方程式」と呼ぶことがある。これは答案には書かず、$a_{n+1}+1=3(a_n+1)$ という変形だけを書けばよい。この変形は展開すれば元の式に戻るので、誰が読んでも正しいとわかる。
  \end{memo}

  \Shomon{3}
  \ascboxZ{3項間漸化式は2通りに変形する}
  $t^2=3t-2$ の解 $t=1,\ 2$ を使うと、漸化式は次の2通りに書き換えられる（展開すれば元の式になる）。
  \[
  \begin{aligned}
  a_{n+2}-a_{n+1}&=2(a_{n+1}-a_n)\\
  a_{n+2}-2a_{n+1}&=a_{n+1}-2a_n
  \end{aligned}
  \]
  1つ目より $\{a_{n+1}-a_n\}$ は初項 $a_2-a_1=2$、公比2の等比数列なので、$a_{n+1}-a_n=2^n$。\par
  2つ目より $\{a_{n+1}-2a_n\}$ は一定で、$a_{n+1}-2a_n=a_2-2a_1=1$。\par
  2式を辺々引くと $a_{n+1}$ が消えて $a_n=2^n-1$。
  \[a_n=\Kotae{2^n-1}\]

  \begin{memo}
    確かめ：$a_3=3\cdot3-2\cdot1=7=2^3-1$。式から求めた値と、漸化式で直接計算した値を $n=3$ くらいまで比べると、計算ミスにほぼ確実に気づく。
  \end{memo}

  \Shomon{4}
  \ascboxZ{指数の項は $3^{n+1}$ で割って消す}
  両辺を $3^{n+1}$ で割ると、
  \[\frac{a_{n+1}}{3^{n+1}}=\frac23\cdot\frac{a_n}{3^n}+\frac13\]
  $c_n=\dfrac{a_n}{3^n}$ とおくと $c_{n+1}=\frac23c_n+\frac13$ で、\kakko{2}の型になる。$\alpha=\frac23\alpha+\frac13$ より $\alpha=1$。
  \[c_{n+1}-1=\frac23(c_n-1),\quad c_1-1=\frac13-1=-\frac23\]
  よって $c_n-1=-\frac23\left(\frac23\right)^{n-1}=-\left(\frac23\right)^n$。$a_n=3^nc_n=3^n-2^n$。
  \[a_n=\Kotae{3^n-2^n}\]

  \begin{ptbs}{KEY}[漸化式の型と変形]
    $a_{n+1}=a_n+f(n)$ → 階差。\quad$a_{n+1}=pa_n+q$ → $\alpha$ を引いて等比。\par
    3項間 → 2つの解で2通りに変形。\quad$+r^n$ の項 → $r^{n+1}$（または $p^{n+1}$）で割る。
  \end{ptbs}

  \begin{simple}[出典：オリジナル]{参考問題22}
    $a_1=1$、$a_{n+1}=\dfrac{a_n}{2a_n+1}$ で定められる数列 $\{a_n\}$ の一般項を求めよ。
  \end{simple}
  \Kai
  $a_1>0$ で、$a_n>0$ なら $a_{n+1}=\frac{a_n}{2a_n+1}>0$ なので、すべての $n$ で $a_n>0$（第23講の帰納法）。特に $a_n\neq0$ なので逆数をとれる。
  \[\frac1{a_{n+1}}=\frac{2a_n+1}{a_n}=\frac1{a_n}+2\]
  $\left\{\frac1{a_n}\right\}$ は初項1、公差2の等差数列で、$\frac1{a_n}=2n-1$。
  \[a_n=\Kotae{\frac1{2n-1}}\]
  \begin{memo}
    逆数をとる前に「$a_n\neq0$」を確かめた。0で割る操作をするときは、毎回その確認が要る。
  \end{memo}

  \begin{ptbs}{CHECK}[翌日、解説を隠して答える（答えは本講の中にある）]
    \kakko{1}　$a_{n+1}=pa_n+q$ はどう変形するか。\par
    \kakko{2}　3項間漸化式の2通りの変形を言え。\par
    \kakko{3}　$+r^n$ の項がある漸化式の処理を言え。
  \end{ptbs}

  \notefill
  \columnbreak
  \notefill
\end{multicols*}

%% ==================== 第23講 ====================
\Kou{23}{数学的帰納法}
\begin{ascolorbox4A}[出典：オリジナル]{例題23}
  \kakko{1}　すべての自然数 $n$ について、$\displaystyle\sum_{k=1}^nk^3=\left\{\frac{n(n+1)}2\right\}^2$ が成り立つことを証明せよ。\par
  \kakko{2}　5以上のすべての自然数 $n$ について、$2^n>n^2$ が成り立つことを証明せよ。\par
  \kakko{3}　$a_1=\dfrac12$、$a_{n+1}=\dfrac1{2-a_n}$ で定められる数列 $\{a_n\}$ の一般項を推測し、それが正しいことを証明せよ。\par
  \kakko{4}　実数 $x$ が $x+\dfrac1x=3$ を満たすとき、すべての自然数 $n$ について $x^n+\dfrac1{x^n}$ は整数であることを証明せよ。
\end{ascolorbox4A}

\begin{Data}
  テーマ：数学的帰納法の型（基本・出発点の変更・推測して証明・2段の帰納法）\par
  難易度：(1) A　(2) B　(3) B　(4) C\quad 前提：第1・21・22講\par
  この講のゴール：「無限個の主張」を2ステップで証明する仕組みを理解し、仮定をどこで使ったか明示できる
\end{Data}

\begin{multicols*}{2}
  \begin{koushiki}{数学的帰納法}
    自然数 $n$ についての主張 $P(n)$ がすべての $n$ で正しいことを示すには、次の2つを示せばよい。\par
    \textgt{[Ⅰ]}　$P(1)$ が正しい。\par
    \textgt{[Ⅱ]}　ある自然数 $k$ で $P(k)$ が正しいと仮定すると、$P(k+1)$ も正しい。
  \end{koushiki}
  \textgt{なぜこれで全部示せるのか}：[Ⅰ]で $P(1)$。[Ⅱ]で $k=1$ として $P(2)$。もう一度[Ⅱ]で $P(3)$……と、どの自然数にもいつかは届く。ドミノの1枚目が倒れ（[Ⅰ]）、どの1枚も次を倒す（[Ⅱ]）なら、全部倒れる。

  \Shomon{1}
  \ascboxZ{[Ⅱ]では仮定の式を必ず使う}
  示す式を①とする。\par
  [Ⅰ]　$n=1$ のとき、左辺 $=1$、右辺 $=\left(\frac{1\cdot2}2\right)^2=1$ で成り立つ。\par
  [Ⅱ]　$n=k$ のとき①が成り立つ、すなわち $\sum_{i=1}^ki^3=\frac{k^2(k+1)^2}4$ と仮定する。このとき、
  \[
  \begin{aligned}
  \sum_{i=1}^{k+1}i^3&=\frac{k^2(k+1)^2}4+(k+1)^3\\
  &=\frac{(k+1)^2(k^2+4k+4)}4=\frac{(k+1)^2(k+2)^2}4
  \end{aligned}
  \]
  1行目で仮定を使った。これは①で $n=k+1$ とした式である。\par
  [Ⅰ][Ⅱ]より、すべての自然数 $n$ で①が成り立つ。\QED

  \Shomon{2}
  \ascboxZ{出発点を変える・不等式は「間にはさむ」}
  [Ⅰ]　$n=5$ のとき、$2^5=32>25=5^2$。\par
  [Ⅱ]　$k\geqq5$ で $2^k>k^2$ と仮定する。$2^{k+1}=2\cdot2^k>2k^2$ である。あとは $2k^2\geqq(k+1)^2$ を示せば、$2^{k+1}>(k+1)^2$ がつながる。
  \[2k^2-(k+1)^2=k^2-2k-1=(k-1)^2-2\]
  $k\geqq5$ では $(k-1)^2\geqq16$ なので正である。よって $2^{k+1}>2k^2>(k+1)^2$。\par
  [Ⅰ][Ⅱ]より、$n\geqq5$ のすべての自然数で成り立つ。\QED

  \begin{memo}
    不等式の帰納法では「仮定から作れる式（$2k^2$）」と「目標（$(k+1)^2$）」の大小を、別に示すことが多い。$n=4$ では $16=16$ で $>$ が成り立たないので、出発点は5でなければならない。
  \end{memo}

  \Shomon{3}
  \ascboxZ{小さい $n$ で計算して推測する}
  $a_2=\dfrac1{2-\frac12}=\dfrac23$、$a_3=\dfrac1{2-\frac23}=\dfrac34$。そこで $a_n=\dfrac n{n+1}$ と推測する。\par
  [Ⅰ]　$a_1=\frac12$ で成り立つ。\par
  [Ⅱ]　$a_k=\frac k{k+1}$ と仮定すると、
  \[a_{k+1}=\frac1{2-\frac k{k+1}}=\frac{k+1}{2(k+1)-k}=\frac{k+1}{k+2}\]
  となり、$n=k+1$ でも成り立つ。よって、
  \[a_n=\Kotae{\frac n{n+1}}\]

  \begin{ptbs}{KEY}[推測した式は帰納法で証明して初めて答え]
    3項ほど計算して規則を見つけるのは強力な方法だが、推測だけでは答えにならない。東大では「推測 → 帰納法」の流れを答案に書かせる問題が多い。
  \end{ptbs}

  \Shomon{4}
  \ascboxZ{2段の帰納法：2つ前まで仮定する}
  $s_n=x^n+\frac1{x^n}$ とおく。第1講と同じ計算で、
  \[s_n\cdot s_1=x^{n+1}+\frac1{x^{n+1}}+x^{n-1}+\frac1{x^{n-1}}\]
  なので $s_{n+1}=3s_n-s_{n-1}$（$n\geqq2$）。$s_{n+1}$ を決めるのに2つ前までの値が要るので、仮定も2つにする。\par
  [Ⅰ]　$s_1=3$、$s_2=s_1^2-2=7$ はともに整数。\par
  [Ⅱ]　$s_k$ と $s_{k+1}$ がともに整数と仮定すると、$s_{k+2}=3s_{k+1}-s_k$ も整数。\par
  [Ⅰ][Ⅱ]より、すべての $n$ で $s_n$ は整数である。\QED

  \begin{memo}
    帰納法の[Ⅰ]で2つ（$s_1,\ s_2$）を確かめたのは、[Ⅱ]が「連続する2つ」から次を出すからである。1つだけだと、ドミノが倒れ始めない。
  \end{memo}

  \begin{simple}[出典：オリジナル]{参考問題23}
    すべての自然数 $n$ について、$n^3+5n$ は6の倍数であることを、数学的帰納法で証明せよ。
  \end{simple}
  \Kai
  [Ⅰ]　$n=1$ のとき $1+5=6$ で成り立つ。\par
  [Ⅱ]　$k^3+5k$ が6の倍数と仮定する。
  \[(k+1)^3+5(k+1)=(k^3+5k)+3k(k+1)+6\]
  $k(k+1)$ は連続する2整数の積で偶数なので、$3k(k+1)$ は6の倍数。仮定と合わせて右辺は6の倍数である。\par
  [Ⅰ][Ⅱ]より、すべての自然数で成り立つ。\QED

  \begin{ptbs}{CHECK}[翌日、解説を隠して答える（答えは本講の中にある）]
    \kakko{1}　[Ⅰ][Ⅱ]だけで、すべての $n$ について示せる理由は何か。\par
    \kakko{2}　[Ⅱ]の中で必ずしなければならないことは何か。\par
    \kakko{3}　2段の帰納法で[Ⅰ]を2つ確かめる理由は何か。
  \end{ptbs}

  \columnbreak
  \notefill
\end{multicols*}

%% ==================== 第24講 ====================
\Kou{24}{確率と漸化式}
\begin{ascolorbox4A}[出典：オリジナル]{例題24}
  \kakko{1}　1個のさいころを $n$ 回投げるとき、6の目が出た回数が偶数（0回を含む）である確率 $p_n$ を求めよ。\par
  正三角形 ABC の頂点上を動く点 P がある。P は最初 A にあり、1秒ごとに、今いる頂点以外の2つの頂点のどちらかへ、それぞれ確率 $\frac12$ で移る。\par
  \kakko{2}　$n$ 秒後に P が A にある確率 $a_n$ を求めよ。\par
  \kakko{3}　$n$ 秒後に P が B にある確率 $b_n$ を求めよ。
\end{ascolorbox4A}

\begin{Data}
  テーマ：状態を決めて「1回後」の確率の漸化式を立てる、余事象と対称性の利用\par
  難易度：(1) B　(2) C　(3) C\quad 前提：第9・22講\par
  この講のゴール：確率の問題を「状態」と「推移」に分解し、漸化式を立てて解ける（東大の頻出テーマ）
\end{Data}

\begin{multicols*}{2}
  \begin{koushiki}{確率漸化式の立て方}
    \kakko{1} 状態を決める（例：「偶数回」と「奇数回」）。\par
    \kakko{2} $n$ 回後の各状態の確率に文字をおく。全状態の確率の和は1。\par
    \kakko{3} 「$n+1$ 回後に状態 X」となる道を、$n$ 回後の状態ごとに分けて足す。\par
    \kakko{4} 小さい $n$ で検算する。
  \end{koushiki}

  \Shomon{1}
  \ascboxZ{最後の1回で場合分けする}
  $n$ 回後の状態は「偶数回」（確率 $p_n$）か「奇数回」（確率 $1-p_n$）のどちらかである。$n+1$ 回後に偶数回になる道は2つある。\par
  ・$n$ 回後に偶数回で、$n+1$ 回目に6以外（確率 $\frac56$）\par
  ・$n$ 回後に奇数回で、$n+1$ 回目に6（確率 $\frac16$）
  \[p_{n+1}=\frac56p_n+\frac16(1-p_n)=\frac23p_n+\frac16\]
  $\alpha=\frac23\alpha+\frac16$ より $\alpha=\frac12$。$p_{n+1}-\frac12=\frac23\left(p_n-\frac12\right)$。$p_1=\frac56$ なので $p_1-\frac12=\frac13$。
  \[p_n-\frac12=\frac13\left(\frac23\right)^{n-1}=\frac12\left(\frac23\right)^n\]
  よって、
  \[p_n=\Kotae{\frac12+\frac12\left(\frac23\right)^n}\]

  \begin{memo}
    別解（二項定理）：偶数回の確率と奇数回の確率の差は $\left(\frac56-\frac16\right)^n$、和は $\left(\frac56+\frac16\right)^n=1$。足して2で割れば同じ答えになる。2通りで一致すれば安心できる。
  \end{memo}

  \Shomon{2}
  \ascboxZ{推移図をかく}
  \begin{zuhyo}
  \begin{tikzpicture}[line width=.5pt,>={Stealth[length=2mm]}]
    \node[draw,circle,inner sep=1.5pt] (A) at (90:1.2) {A};
    \node[draw,circle,inner sep=1.5pt] (B) at (210:1.2) {B};
    \node[draw,circle,inner sep=1.5pt] (C) at (330:1.2) {C};
    \foreach \u/\v in {A/B,B/A,B/C,C/B,C/A,A/C}{\draw[->,bend left=15] (\u) to (\v);}
    \node[font=\footnotesize] at (-1.25,.45) {$\frac12$};
    \node[font=\footnotesize] at (1.25,.45) {$\frac12$};
    \node[font=\footnotesize] at (0,-1.05) {$\frac12$};
  \end{tikzpicture}
  \end{zuhyo}
  どの矢印も確率 $\frac12$ である。$n+1$ 秒後に A にいるには、$n$ 秒後に A 以外（確率 $1-a_n$）にいて、そこから A へ移る（確率 $\frac12$）しかない。A にいたら必ず A から出ていくからである。
  \[a_{n+1}=\frac12(1-a_n)\]
  $\alpha=\frac12(1-\alpha)$ より $\alpha=\frac13$。$a_{n+1}-\frac13=-\frac12\left(a_n-\frac13\right)$。最初は A にいるので $a_0=1$（0秒後）、$a_0-\frac13=\frac23$。
  \[a_n=\Kotae{\frac13+\frac23\left(-\frac12\right)^n}\]

  \begin{memo}
    検算：$a_1=\frac13-\frac13=0$（1秒後は必ず A を離れている）。$a_2=\frac13+\frac16=\frac12$（B か C に行って、半分の確率で戻る）。どちらも直接の計算と一致する。
  \end{memo}

  \Shomon{3}
  \ascboxZ{対称性で残りを決める}
  最初の位置 A から見て、B と C は対称である（三角形を A を通る軸で折り返すと B と C が入れ替わり、移動の規則は変わらない）。よって $b_n=c_n$。$a_n+b_n+c_n=1$ なので、
  \[b_n=\frac{1-a_n}2=\Kotae{\frac13-\frac13\left(-\frac12\right)^n}\]

  \begin{ptbs}{KEY}[確率漸化式の3つの道具]
    「最後の1回で場合分け」「和が1（余事象）」「対称性」。状態の数を、これらでできるだけ減らしてから漸化式を立てる。
  \end{ptbs}

  \begin{memo}
    $n$ を大きくすると $\left(-\frac12\right)^n\to0$ なので、$a_n,\ b_n$ はどれも $\frac13$ に近づく（第30講）。長い時間がたつと、最初の位置の影響が消えて、3頂点に均等にいるようになる。
  \end{memo}

  \begin{simple}[出典：オリジナル]{参考問題24}
    2つの部屋 A、B を移動する人がいる。A にいるときは次の1分で確率 $\frac13$ で B に移り（確率 $\frac23$ で A にとどまる）、B にいるときは次の1分で確率 $\frac12$ で A に移る。最初 A にいるとき、$n$ 分後に A にいる確率 $q_n$ を求めよ。
  \end{simple}
  \Kai
  $n+1$ 分後に A にいるのは、「A にいてとどまる」か「B にいて A に移る」の2通りである。
  \[q_{n+1}=\frac23q_n+\frac12(1-q_n)=\frac16q_n+\frac12\]
  $\alpha=\frac16\alpha+\frac12$ より $\alpha=\frac35$。$q_0=1$ なので $q_0-\frac35=\frac25$。
  \[q_n=\Kotae{\frac35+\frac25\left(\frac16\right)^n}\]
  検算：$q_1=\frac35+\frac1{15}=\frac23$ で、最初の1分でとどまる確率と一致する。

  \begin{ptbs}{CHECK}[翌日、解説を隠して答える（答えは本講の中にある）]
    \kakko{1}　確率漸化式を立てる4手順を言え。\par
    \kakko{2}　$a_{n+1}=\frac12(1-a_n)$ となる理由を言え。\par
    \kakko{3}　状態の数を減らす3つの道具は何か。
  \end{ptbs}

  \notefill
  \columnbreak
  \notefill
\end{multicols*}
%% ==================== 第25講 ====================
\Kou{25}{平面ベクトル}
\begin{ascolorbox4A}[出典：オリジナル]{例題25}
  \kakko{1}　$\vec a=(2,\ 1)$、$\vec b=(1,\ 3)$ について、$|\vec a|$、内積 $\vec a\cdot\vec b$、$\vec a$ と $\vec b$ のなす角 $\theta$ を求めよ。\par
  $\triangle\mathrm{OAB}$ において、$\vv{OA}=\vec a$、$\vv{OB}=\vec b$ とする。\par
  \kakko{2}　辺 AB を $2:1$ に内分する点を P とする。$\vv{OP}$ を $\vec a,\ \vec b$ で表せ。\par
  \kakko{3}　辺 OA を $2:1$ に内分する点を C、辺 OB を $1:2$ に内分する点を D とし、線分 AD と線分 BC の交点を Q とする。$\vv{OQ}$ を $\vec a,\ \vec b$ で表せ。\par
  \kakko{4}　$|\vec a|=2$、$|\vec b|=3$、$|\vec a-\vec b|=\sqrt7$ のとき、内積 $\vec a\cdot\vec b$ と $\triangle\mathrm{OAB}$ の面積を求めよ。
\end{ascolorbox4A}

\begin{Data}
  テーマ：ベクトルの演算と成分、内積、位置ベクトルと分点、一次独立による係数比較、面積\par
  難易度：(1) A　(2) A　(3) C　(4) B\quad 前提：第7・14講\par
  この講のゴール：図形の問題を「2本の基準ベクトルの組合せ」に直し、係数比較で交点を求められる
\end{Data}

\begin{multicols*}{2}
  \begin{koushiki}{ベクトルの言葉}
    \textgt{ベクトル}：向きと大きさだけをもつ量。矢印で表すが、矢印をどこに置いても（平行移動しても）同じベクトルである。\par
    和 $\vec a+\vec b$：$\vec a$ の終点に $\vec b$ の始点をつないだ矢印。\par
    実数倍 $k\vec a$：向きはそのまま（$k<0$ なら逆向き）、大きさ $|k|$ 倍。\par
    差：$\vv{AB}=\vv{OB}-\vv{OA}$（終点 $-$ 始点）。\par
    成分：$\vec a=(a_1,\ a_2)$ のとき $|\vec a|=\sqrt{a_1^2+a_2^2}$。
  \end{koushiki}

  \Shomon{1}
  \ascboxZ{内積の定義と成分表示}
  \begin{koushiki}{内積}
    $\vec a,\ \vec b$ のなす角を $\theta$（$0\leqq\theta\leqq\pi$）とするとき、
    \[\vec a\cdot\vec b=|\vec a||\vec b|\cos\theta\]
    成分では $\vec a\cdot\vec b=a_1b_1+a_2b_2$。\par
    $\vec a\cdot\vec a=|\vec a|^2$、\quad$\vec a\perp\vec b\iff\vec a\cdot\vec b=0$（$\vec a,\vec b\neq\vec0$）
  \end{koushiki}
  \textgt{成分表示の理由}：$\vec a,\ \vec b,\ \vec a-\vec b$ でできる三角形に余弦定理を使うと、$|\vec a-\vec b|^2=|\vec a|^2+|\vec b|^2-2|\vec a||\vec b|\cos\theta$。左辺を成分で展開して比べると、$|\vec a||\vec b|\cos\theta=a_1b_1+a_2b_2$ が出る。
  \[|\vec a|=\sqrt{4+1}=\Kotae{\sqrt5},\quad\vec a\cdot\vec b=2+3=\Kotae{5}\]
  $|\vec b|=\sqrt{10}$ なので $\cos\theta=\dfrac5{\sqrt5\sqrt{10}}=\dfrac1{\sqrt2}$。$0\leqq\theta\leqq\pi$ より、
  \[\theta=\Kotae{\frac\pi4}\]

  \Shomon{2}
  \ascboxZ{内分点の位置ベクトル}
  $\vv{OP}=\vv{OA}+\frac23\vv{AB}$（A から B へ全体の $\frac23$ 進む）$=\vec a+\frac23(\vec b-\vec a)$。
  \[\vv{OP}=\Kotae{\frac{\vec a+2\vec b}3}\]

  \begin{koushiki}{内分点}
    AB を $m:n$ に内分する点 P：$\vv{OP}=\dfrac{n\vec a+m\vec b}{m+n}$\par
    （分子は「たすき掛け」。座標の内分点と同じ形。）
  \end{koushiki}

  \Shomon{3}
  \ascboxZ{交点は「2通りに表して係数比較」}
  \begin{zuhyo}
  \begin{tikzpicture}[x=1.15cm,y=1.15cm,line width=.5pt]
    \coordinate (O) at (0,0); \coordinate (A) at (4,0); \coordinate (B) at (1,3);
    \coordinate (C) at (8/3,0); \coordinate (D) at (1/3,1); \coordinate (Q) at (17/7,3/7);
    \draw (O) -- (A) -- (B) -- cycle;
    \draw (A) -- (D); \draw (B) -- (C);
    \fill (Q) circle (1.5pt);
    \node[below left] at (O) {$\mathrm O$}; \node[below right] at (A) {$\mathrm A$};
    \node[above] at (B) {$\mathrm B$}; \node[below] at (C) {$\mathrm C$};
    \node[left] at (D) {$\mathrm D$}; \node[above right] at (Q) {$\mathrm Q$};
  \end{tikzpicture}
  \end{zuhyo}
  $\vv{OC}=\frac23\vec a$、$\vv{OD}=\frac13\vec b$ である。\par
  \textgt{Q は AD 上}：実数 $s$ を用いて、
  \[\vv{OQ}=(1-s)\vec a+\frac s3\vec b\ \cdots\text{①}\]
  \textgt{Q は BC 上}：実数 $t$ を用いて、
  \[\vv{OQ}=\frac{2t}3\vec a+(1-t)\vec b\ \cdots\text{②}\]
  $\vec a,\ \vec b$ は $\vec0$ でなく平行でもない（\textgt{一次独立}）ので、同じベクトルの表し方は1通りしかない。①②の係数を比べて、
  \[1-s=\frac{2t}3,\qquad\frac s3=1-t\]
  2つ目から $s=3-3t$。1つ目に入れて $3t-2=\frac{2t}3$、$t=\frac67$、$s=\frac37$。①に代入して、
  \[\vv{OQ}=\Kotae{\frac47\vec a+\frac17\vec b}\]

  \begin{ptbs}{KEY}[一次独立なら係数比較してよい]
    $\vec a,\ \vec b$ が一次独立のとき、$p\vec a+q\vec b=p'\vec a+q'\vec b\iff p=p'$ かつ $q=q'$。答案には「$\vec a,\ \vec b$ は一次独立であるから」と必ず書く。
  \end{ptbs}

  \begin{memo}
    なぜ平行だとだめか：$\vec b=2\vec a$ なら $2\vec a+0\vec b=0\vec a+1\vec b$ のように、係数の違う2通りの表し方ができてしまう。三角形の2辺は平行でないので、この心配はない。
  \end{memo}

  \Shomon{4}
  \ascboxZ{$|\ |^2$ を展開して内積を作る}
  大きさの条件は、2乗して内積の式に直す。
  \[|\vec a-\vec b|^2=|\vec a|^2-2\vec a\cdot\vec b+|\vec b|^2\]
  $7=4-2\vec a\cdot\vec b+9$ より、
  \[\vec a\cdot\vec b=\Kotae{3}\]
  面積は $S=\frac12|\vec a||\vec b|\sin\theta$。$\sin\theta=\sqrt{1-\cos^2\theta}$ を内積で書き直すと次の公式になる。
  \[S=\frac12\sqrt{|\vec a|^2|\vec b|^2-(\vec a\cdot\vec b)^2}\]
  よって $S=\frac12\sqrt{36-9}$ で、
  \[S=\Kotae{\frac{3\sqrt3}2}\]

  \begin{memo}
    成分 $\vec a=(a_1,a_2),\ \vec b=(b_1,b_2)$ がわかっているときは、上の公式は $S=\frac12|a_1b_2-a_2b_1|$ になる。
  \end{memo}

  \begin{simple}[出典：オリジナル]{参考問題25}
    $\triangle\mathrm{ABC}$ と点 P が $3\vv{PA}+4\vv{PB}+5\vv{PC}=\vec0$ を満たす。$\vv{AP}$ を $\vv{AB},\ \vv{AC}$ で表し、$\triangle\mathrm{PBC}:\triangle\mathrm{PCA}:\triangle\mathrm{PAB}$ を求めよ。
  \end{simple}
  \Kai
  始点を A にそろえる。$\vv{PA}=-\vv{AP}$、$\vv{PB}=\vv{AB}-\vv{AP}$、$\vv{PC}=\vv{AC}-\vv{AP}$ を代入して、$12\vv{AP}=4\vv{AB}+5\vv{AC}$。
  \[\vv{AP}=\Kotae{\frac{4\vv{AB}+5\vv{AC}}{12}}\]
  内分点の形にするため、$\vv{AP}=\frac34\cdot\frac{4\vv{AB}+5\vv{AC}}9$ と変形する。$\frac{4\vv{AB}+5\vv{AC}}9$ は BC を $5:4$ に内分する点 E を表すので、P は AE を $3:1$ に内分する点である。$\triangle\mathrm{ABC}$ の面積を $S$ とすると、
  \par $\triangle\mathrm{PBC}=\frac14S$（高さが AE の $\frac14$）、\par
  $\triangle\mathrm{PAB}=\frac34\triangle\mathrm{ABE}=\frac34\cdot\frac59S=\frac5{12}S$、\par
  $\triangle\mathrm{PCA}=\frac34\cdot\frac49S=\frac4{12}S$。
  \[\Kotae{3:4:5}\]
  \begin{memo}
    面積比が係数 $3:4:5$ と一致した。$a\vv{PA}+b\vv{PB}+c\vv{PC}=\vec0$（$a,b,c>0$）なら、面積比は常に $a:b:c$ になる。
  \end{memo}

  \begin{ptbs}{CHECK}[翌日、解説を隠して答える（答えは本講の中にある）]
    \kakko{1}　内積の成分表示はどう導くか。\par
    \kakko{2}　係数比較の前に答案に書くべき一言は何か。\par
    \kakko{3}　$|\vec a-\vec b|$ の値から内積を出す方法を言え。
  \end{ptbs}

  \notefill
\end{multicols*}

%% ==================== 第26講 ====================
\Kou{26}{空間ベクトル}
\begin{ascolorbox4A}[出典：オリジナル]{例題26}
  座標空間に原点 O と3点 $\mathrm A(1,\ 0,\ 0)$、$\mathrm B(0,\ 2,\ 0)$、$\mathrm C(0,\ 0,\ 3)$ がある。\par
  \kakko{1}　$\triangle\mathrm{ABC}$ の面積を求めよ。\par
  \kakko{2}　O から平面 ABC に下ろした垂線の足を H とする。H の座標と、線分 OH の長さを求めよ。\par
  \kakko{3}　四面体 OABC の体積を求め、\kakko{1}\kakko{2}の結果と矛盾しないことを確かめよ。\par
  \kakko{4}　直線 $x=y=z$ と平面 ABC の交点の座標を求めよ。
\end{ascolorbox4A}

\begin{Data}
  テーマ：空間座標とベクトル、空間の三角形の面積、平面上の点の表し方、垂線の足、法線ベクトル\par
  難易度：(1) A　(2) C　(3) B　(4) B\quad 前提：第25講\par
  この講のゴール：平面ベクトルの方法をそのまま空間に持ち上げ、垂直条件を内積で処理できる
\end{Data}

\begin{multicols*}{2}
  \begin{koushiki}{空間ベクトル}
    成分は3つ：$\vec a=(a_1,\ a_2,\ a_3)$、$|\vec a|=\sqrt{a_1^2+a_2^2+a_3^2}$\par
    内積：$\vec a\cdot\vec b=a_1b_1+a_2b_2+a_3b_3$\par
    平面 ABC 上の点 P：$\vv{OP}=\vv{OA}+s\vv{AB}+t\vv{AC}$（$s,\ t$ は実数）\par
    空間では、同じ平面上にない3つのベクトルで、どんなベクトルも1通りに表せる。
  \end{koushiki}

  \begin{zuhyo}
  \begin{tikzpicture}[x={(-.45cm,-.35cm)},y={(.9cm,0cm)},z={(0cm,.75cm)},line width=.5pt]
    \draw[->] (0,0,0) -- (1.9,0,0) node[below left] {$x$};
    \draw[->] (0,0,0) -- (0,2.7,0) node[right] {$y$};
    \draw[->] (0,0,0) -- (0,0,3.6) node[above] {$z$};
    \fill[black!12] (1,0,0) -- (0,2,0) -- (0,0,3) -- cycle;
    \draw (1,0,0) -- (0,2,0) -- (0,0,3) -- cycle;
    \node[left] at (1,0,0) {$\mathrm A$};
    \node[below] at (0,2,0) {$\mathrm B$};
    \node[left] at (0,0,3) {$\mathrm C$};
    \node[left] at (0,0,0) {$\mathrm O$};
    \draw[dashed] (0,0,0) -- (36/49,18/49,12/49);
    \fill (36/49,18/49,12/49) circle (1.3pt) node[right,font=\footnotesize] {$\mathrm H$};
  \end{tikzpicture}
  \end{zuhyo}

  \Shomon{1}
  \ascboxZ{面積公式は空間でもそのまま使える}
  $\vv{AB}=(-1,\ 2,\ 0)$、$\vv{AC}=(-1,\ 0,\ 3)$。
  \[|\vv{AB}|^2=5,\quad|\vv{AC}|^2=10,\quad\vv{AB}\cdot\vv{AC}=1\]
  第25講の面積公式より、
  \[S=\frac12\sqrt{5\cdot10-1^2}=\Kotae{\frac72}\]

  \Shomon{2}
  \ascboxZ{垂直 $\iff$ 平面内の2本と内積0}
  H は平面 ABC 上にあるので、実数 $s,\ t$ で
  \[\vv{OH}=\vv{OA}+s\vv{AB}+t\vv{AC}=(1-s-t,\ 2s,\ 3t)\]
  と書ける。OH が平面 ABC に垂直 $\iff$ OH が平面内の平行でない2本 AB、AC の両方に垂直。
  \[
  \begin{aligned}
  \vv{OH}\cdot\vv{AB}&=-(1-s-t)+4s=-1+5s+t=0\\
  \vv{OH}\cdot\vv{AC}&=-(1-s-t)+9t=-1+s+10t=0
  \end{aligned}
  \]
  1つ目から $t=1-5s$。2つ目に代入して $9-49s=0$、$s=\frac9{49}$、$t=\frac4{49}$。
  \[\mathrm H\Kotae{\left(\frac{36}{49},\ \frac{18}{49},\ \frac{12}{49}\right)}\]
  \[\mathrm{OH}=\frac{\sqrt{36^2+18^2+12^2}}{49}=\frac{42}{49}=\Kotae{\frac67}\]

  \begin{ptbs}{KEY}[平面に垂直 $\iff$ 平面内の平行でない2本に垂直]
    1本に垂直なだけでは足りない（その1本を軸に回転できてしまう）。2本と内積0の連立方程式を解く。
  \end{ptbs}

  \ascboxZ{法線ベクトルを使う方法（検算に便利）}
  $\vec n=(6,\ 3,\ 2)$ とすると、$\vec n\cdot\vv{AB}=-6+6=0$、$\vec n\cdot\vv{AC}=-6+6=0$ なので、$\vec n$ は平面 ABC に垂直（\textgt{法線ベクトル}）。平面上の点 $(x,y,z)$ は $\vec n\cdot(\text{その点}-\mathrm A)=0$ を満たすので、平面の方程式は
  \[6x+3y+2z=6\]
  OH は $\vec n$ に平行なので H$=(6k,\ 3k,\ 2k)$。平面の式に入れて $49k=6$、$k=\frac6{49}$。上と一致する。

  \begin{memo}
    平面の方程式 $ax+by+cz=d$ は教科書では発展扱いのことがあるが、東大の空間図形では検算・見通しに非常に役立つ。答案では、上の「内積0」の形で根拠を示してから使うと安全である。
  \end{memo}

  \Shomon{3}
  \ascboxZ{体積を2通りに計算する}
  四面体 OABC は、底面 $\triangle\mathrm{OAB}$（直角をはさむ2辺が1と2で、面積1）、高さ OC $=3$ の三角すいである。
  \[V=\frac13\cdot1\cdot3=\Kotae{1}\]
  一方、底面を $\triangle\mathrm{ABC}$ とみると、高さは OH である。
  \[\frac13\cdot\frac72\cdot\frac67=1\]
  となり、\kakko{1}\kakko{2}の結果と矛盾しない。

  \begin{memo}
    「体積を2通りに表す」ことで、点と平面の距離を逆に求めることもできる。$\mathrm{OH}=\dfrac{3V}{\triangle\mathrm{ABC}}=\dfrac{3}{7/2}=\dfrac67$。計算量が少ない別解である。
  \end{memo}

  \Shomon{4}
  \ascboxZ{直線上の点をパラメータで表す}
  直線 $x=y=z$ 上の点は $(u,\ u,\ u)$ と書ける。平面 $6x+3y+2z=6$ に入れて $11u=6$。
  \[\Kotae{\left(\frac6{11},\ \frac6{11},\ \frac6{11}\right)}\]
  （平面の方程式を使わない場合は、$(u,u,u)=\vv{OA}+s\vv{AB}+t\vv{AC}$ として成分を比べ、$s,\ t,\ u$ の連立方程式を解けばよい。）

  \begin{simple}[出典：オリジナル]{参考問題26}
    球面 $x^2+y^2+z^2-2x+4y-6z=0$ の中心と半径を求めよ。また、この球面と平面 $z=1$ の交わりはどのような図形か。
  \end{simple}
  \Kai
  平方完成すると $(x-1)^2+(y+2)^2+(z-3)^2=14$。
  \[\Kotae{\text{中心}(1,\ -2,\ 3),\ \text{半径}\ \sqrt{14}}\]
  $z=1$ を代入すると $(x-1)^2+(y+2)^2=14-4=10$。
  よって、交わりは平面 $z=1$ 上の
  \[\Kotae{\text{中心}(1,\ -2,\ 1)\text{、半径}\ \sqrt{10}\ \text{の円}}\]
  \begin{memo}
    断面の円の半径は、三平方の定理（球の半径 $\sqrt{14}$、中心から平面までの距離2）からも $\sqrt{14-4}$ と求められる。空間図形は「断面で平面の問題にする」のが基本である（第38・40講）。
  \end{memo}

  \begin{ptbs}{CHECK}[翌日、解説を隠して答える（答えは本講の中にある）]
    \kakko{1}　直線が平面に垂直であることを示す条件は何か。\par
    \kakko{2}　平面 ABC 上の点の表し方を書け。\par
    \kakko{3}　四面体の体積を2通りに表すと何がわかるか。
  \end{ptbs}

  \notefill
\end{multicols*}
%% ==================== 第27講 ====================
\Kou{27}{複素数平面(1)}
\begin{ascolorbox4A}[出典：オリジナル]{例題27}
  \kakko{1}　$z=1+\sqrt3\,i$ を極形式で表し、$z^6$ を求めよ。\par
  \kakko{2}　複素数平面上で、点 $\beta=3+i$ を点 $\alpha=1+i$ のまわりに $\dfrac\pi3$ だけ回転した点 $\gamma$ を求めよ。\par
  \kakko{3}　方程式 $z^3=8i$ を解け。\par
  \kakko{4}　$(1+i)^n$ が実数となる最小の自然数 $n$ と、そのときの値を求めよ。
\end{ascolorbox4A}

\begin{Data}
  テーマ：複素数平面、絶対値と共役、極形式、積と回転、ド・モアブルの定理、$n$ 乗根\par
  難易度：(1) A　(2) B　(3) C　(4) B\quad 前提：第13・16・25講\par
  この講のゴール：「複素数を掛ける＝回転して拡大」という見方で、計算と図形を行き来できる
\end{Data}

\begin{multicols*}{2}
  \begin{koushiki}{複素数平面}
    複素数 $z=x+yi$ を座標平面の点 $(x,\ y)$ に対応させた平面を\textgt{複素数平面}という（$x$ 軸を実軸、$y$ 軸を虚軸という）。\par
    絶対値：$|z|=\sqrt{x^2+y^2}$（原点からの距離）。$|z|^2=z\overline z$。\par
    $|\alpha-\beta|$：2点 $\alpha,\ \beta$ の距離。\par
    共役 $\overline z=x-yi$：実軸について対称な点。
  \end{koushiki}

  \Shomon{1}
  \ascboxZ{極形式：距離と角で表す}
  $z\neq0$ のとき、$r=|z|$、実軸の正の向きから測った角を $\theta$（\textgt{偏角}、$\arg z$）とすると、単位円の定義を $r$ 倍して、
  \[z=r(\cos\theta+i\sin\theta)\]
  と書ける。これを\textgt{極形式}という。$z=1+\sqrt3\,i$ は $r=\sqrt{1+3}=2$、点 $(1,\ \sqrt3)$ の角は $\frac\pi3$。
  \[z=\Kotae{2\left(\cos\frac\pi3+i\sin\frac\pi3\right)}\]

  \begin{koushiki}{積と商、ド・モアブルの定理}
    $z_1=r_1(\cos\theta_1+i\sin\theta_1)$、$z_2=r_2(\cos\theta_2+i\sin\theta_2)$ のとき\par
    $z_1z_2=r_1r_2\{\cos(\theta_1+\theta_2)+i\sin(\theta_1+\theta_2)\}$\par
    $\dfrac{z_1}{z_2}=\dfrac{r_1}{r_2}\{\cos(\theta_1-\theta_2)+i\sin(\theta_1-\theta_2)\}$\par
    $(\cos\theta+i\sin\theta)^n=\cos n\theta+i\sin n\theta$（$n$ は整数）
  \end{koushiki}
  \textgt{なぜ角が足し算になるか}：$(\cos\theta_1+i\sin\theta_1)(\cos\theta_2+i\sin\theta_2)$ を展開すると、実部は $\cos\theta_1\cos\theta_2-\sin\theta_1\sin\theta_2=\cos(\theta_1+\theta_2)$、虚部は $\sin(\theta_1+\theta_2)$ になる。加法定理そのものである。ド・モアブルの定理は、これを $n$ 回繰り返したもの（帰納法）。
  \[z^6=2^6\left(\cos2\pi+i\sin2\pi\right)=\Kotae{64}\]

  \begin{ptbs}{KEY}[$w$ を掛ける＝$\arg w$ だけ回転し、$|w|$ 倍に拡大]
    $\cos\varphi+i\sin\varphi$ を掛けると、原点のまわりに $\varphi$ 回転する。特に $i$ を掛けると $\frac\pi2$ 回転、$-1$ を掛けると $\pi$ 回転。
  \end{ptbs}

  \Shomon{2}
  \ascboxZ{中心を原点に移してから回す}
  $\alpha$ を中心とする回転は、\kakko{1} $\alpha$ が原点に来るように全体を平行移動（$-\alpha$）、\kakko{2} 原点のまわりに回転、\kakko{3} 元に戻す（$+\alpha$）、の3段階で行う。
  \[\gamma-\alpha=(\beta-\alpha)\left(\cos\frac\pi3+i\sin\frac\pi3\right)\]
  $\beta-\alpha=2$ なので、
  \[\gamma=1+i+2\left(\frac12+\frac{\sqrt3}2i\right)=\Kotae{2+(1+\sqrt3)i}\]

  \begin{zuhyo}
  \begin{tikzpicture}[x=.9cm,y=.9cm,line width=.5pt]
    \draw[->] (-.3,0) -- (3.6,0) node[right,font=\footnotesize] {実軸};
    \draw[->] (0,-.3) -- (0,3) node[above,font=\footnotesize] {虚軸};
    \fill (1,1) circle (1.5pt) node[below left,font=\footnotesize] {$\alpha$};
    \fill (3,1) circle (1.5pt) node[below,font=\footnotesize] {$\beta$};
    \fill (2,2.732) circle (1.5pt) node[right,font=\footnotesize] {$\gamma$};
    \draw (1,1) -- (3,1); \draw (1,1) -- (2,2.732);
    \draw[->] (1.5,1) arc (0:60:.5);
    \node[font=\footnotesize] at (1.75,1.35) {$\frac\pi3$};
  \end{tikzpicture}
  \end{zuhyo}

  \Shomon{3}
  \ascboxZ{$n$ 乗根は極形式で「大きさ」と「角」を別々に}
  $z=r(\cos\theta+i\sin\theta)$（$r>0$）とおく。ド・モアブルの定理より $z^3=r^3(\cos3\theta+i\sin3\theta)$。$8i=8\left(\cos\frac\pi2+i\sin\frac\pi2\right)$ と比べると、
  \[r^3=8,\qquad3\theta=\frac\pi2+2k\pi\ (k\ \text{は整数})\]
  角は $2\pi$ の整数倍のずれを許して等しい、という点に注意する。$r=2$、$\theta=\frac\pi6+\frac{2k\pi}3$。$0\leqq\theta<2\pi$ の範囲では $k=0,\ 1,\ 2$ の3つで、$\theta=\frac\pi6,\ \frac{5\pi}6,\ \frac{3\pi}2$。
  \[z=\Kotae{\sqrt3+i,\ -\sqrt3+i,\ -2i}\]

  \begin{memo}
    3つの解は、原点を中心とする半径2の円周上に $\frac{2\pi}3$ ずつ並び、正三角形の頂点になる。$z^n=w$ の解はいつも、円周上に $\frac{2\pi}n$ ずつ等間隔に並ぶ $n$ 個の点である。
  \end{memo}

  \Shomon{4}
  \ascboxZ{角が $\pi$ の整数倍なら実数}
  $1+i=\sqrt2\left(\cos\frac\pi4+i\sin\frac\pi4\right)$ なので、
  \[(1+i)^n=(\sqrt2)^n\left(\cos\frac{n\pi}4+i\sin\frac{n\pi}4\right)\]
  実数 $\iff$ 虚部 $=0\iff\sin\frac{n\pi}4=0\iff\frac{n\pi}4$ が $\pi$ の整数倍 $\iff n$ が4の倍数。最小の $n$ は4で、
  \[n=\Kotae{4}\]
  このとき、
  \[(1+i)^4=4(\cos\pi+i\sin\pi)=\Kotae{-4}\]

  \begin{simple}[出典：オリジナル]{参考問題27}
    複素数平面上の2点 $\alpha=1+2i$、$\beta=4+3i$ を頂点とする正三角形の、残りの頂点 $\gamma$ を求めよ。
  \end{simple}
  \Kai
  $\gamma$ は、$\beta$ を $\alpha$ のまわりに $\pm\frac\pi3$ 回転した点である（向きが2通りある）。$\beta-\alpha=3+i$。\par
  $+\frac\pi3$ のとき：$(3+i)\left(\frac12+\frac{\sqrt3}2i\right)=\frac{3-\sqrt3}2+\frac{1+3\sqrt3}2i$。\par
  $-\frac\pi3$ のとき：$(3+i)\left(\frac12-\frac{\sqrt3}2i\right)=\frac{3+\sqrt3}2+\frac{1-3\sqrt3}2i$。\par
  それぞれに $\alpha$ を足して、
  \[\gamma=\Kotae{\frac{5-\sqrt3}2+\frac{5+3\sqrt3}2i}\]
  \[\gamma=\Kotae{\frac{5+\sqrt3}2+\frac{5-3\sqrt3}2i}\]

  \begin{ptbs}{CHECK}[翌日、解説を隠して答える（答えは本講の中にある）]
    \kakko{1}　複素数の積で偏角が足し算になる理由は何か。\par
    \kakko{2}　点 $\alpha$ のまわりに $\theta$ 回転する式を書け。\par
    \kakko{3}　$z^n=w$ の解は、複素数平面上でどう並ぶか。
  \end{ptbs}

  \notefill
  \columnbreak
  \notefill
\end{multicols*}

%% ==================== 第28講 ====================
\Kou{28}{複素数平面(2)}
\begin{ascolorbox4A}[出典：オリジナル]{例題28}
  \kakko{1}　複素数平面上で、次の等式を満たす点 $z$ 全体はどのような図形か。\quad(ア)　$|z-2i|=1$\quad(イ)　$|z+1|=|z-i|$\par
  \kakko{2}　点 $z$ が $|z-1|=1$、$z\neq0$ を満たしながら動くとき、$w=\dfrac1z$ が描く図形を求めよ。\par
  \kakko{3}　$z\neq0$ とする。$z+\dfrac1z$ が実数となるような点 $z$ 全体を図示せよ。\par
  \kakko{4}　$|z|=1$ のとき、$|z^2-z+1|$ の最大値を求めよ。
\end{ascolorbox4A}

\begin{Data}
  テーマ：複素数の等式が表す図形、変換 $w=\frac1z$、実数条件 $z=\overline z$、単位円上の点の扱い\par
  難易度：(1) A　(2) C　(3) C　(4) C\quad 前提：第27講\par
  この講のゴール：複素数の条件を「距離」「共役」「回転」のどれで読むかを選び、図形として答えられる
\end{Data}

\begin{multicols*}{2}
  \begin{koushiki}{図形を表す式の辞書}
    $|z-\alpha|=r$：中心 $\alpha$、半径 $r$ の円\par
    $|z-\alpha|=|z-\beta|$：線分 $\alpha\beta$ の垂直二等分線\par
    $z$ が実数 $\iff z=\overline z$\par
    $z$ が純虚数 $\iff z=-\overline z$ かつ $z\neq0$\par
    $\dfrac{\gamma-\alpha}{\beta-\alpha}$ が実数 $\iff$ 3点 $\alpha,\beta,\gamma$ が一直線上\par
    $\dfrac{\gamma-\alpha}{\beta-\alpha}$ が純虚数 $\iff$ $\alpha\beta\perp\alpha\gamma$
  \end{koushiki}
  最後の2つは、$\frac{\gamma-\alpha}{\beta-\alpha}$ の偏角が「$\alpha$ から見た $\beta$ と $\gamma$ の間の角」であることから出る（第27講の商の公式）。

  \Shomon{1}
  \ascboxZ{絶対値は距離と読む}
  (ア)　$|z-2i|$ は点 $z$ と点 $2i$ の距離。それが1なので、
  \[\Kotae{\text{中心}\ 2i\text{、半径1の円}}\]
  (イ)　点 $-1$ からの距離と、点 $i$ からの距離が等しい点の集まりなので、2点を結ぶ線分の垂直二等分線である。$z=x+yi$ として確かめると、$(x+1)^2+y^2=x^2+(y-1)^2$ より $y=-x$。
  \[\Kotae{\text{直線}\ y=-x\ \text{（2点}\ -1,\ i\ \text{の垂直二等分線）}}\]

  \Shomon{2}
  \ascboxZ{$z$ を $w$ で表して条件に代入する}
  $w=\frac1z$ より $z=\frac1w$（$w\neq0$）。条件 $|z-1|=1$ に代入する。
  \[\left|\frac1w-1\right|=1\iff\frac{|1-w|}{|w|}=1\iff|w-1|=|w|\]
  これは、点 $w$ が0と1から等距離にあること、つまり2点 $0,\ 1$ の垂直二等分線上にあることを表す。\par
  逆に、この直線上の点 $w$ は $w\neq0$ なので、$z=\frac1w$ は $z\neq0$ と $|z-1|=1$ を満たす。よって直線全体が描かれる。
  \[\Kotae{\text{直線}\ \mathrm{Re}\,w=\frac12\ \text{（実部が}\ \tfrac12\ \text{の点全体）}}\]

  \begin{memo}
    $\mathrm{Re}\,w$ は $w$ の実部を表す。円 $|z-1|=1$ は原点を通るので、$\frac1z$ で写すと直線になる。原点を通る円は直線に、通らない円は円に写るという性質がある。
  \end{memo}

  \Shomon{3}
  \ascboxZ{実数条件は「共役と等しい」}
  $z+\frac1z$ が実数 $\iff z+\frac1z=\overline z+\frac1{\overline z}$。移項して整理する。$\frac1z-\frac1{\overline z}=\frac{\overline z-z}{z\overline z}$ なので、
  \[(z-\overline z)-\frac{z-\overline z}{|z|^2}=0\iff(z-\overline z)\left(1-\frac1{|z|^2}\right)=0\]
  よって $z=\overline z$（$z$ は実数）または $|z|=1$。\par
  答えは、実軸から原点を除いた部分と、単位円 $|z|=1$ を合わせた図形。

  \begin{zuhyo}
  \begin{tikzpicture}[x=.85cm,y=.85cm,line width=.5pt]
    \draw[line width=1.2pt] (-2.2,0) -- (2.2,0);
    \draw[line width=1.2pt] (0,0) circle (1);
    \draw[->] (0,-1.5) -- (0,1.6) node[above,font=\footnotesize] {虚軸};
    \node[right,font=\footnotesize] at (2.2,0) {実軸};
    \filldraw[fill=white,draw=black] (0,0) circle (2.2pt);
    \node[below left,font=\footnotesize] at (0,0) {$\mathrm O$};
    \node[below right,font=\footnotesize] at (1,0) {$1$};
  \end{tikzpicture}
  \end{zuhyo}
  太線部分。白丸の原点は除く。\hfill\textgt{……\kakko{答}}

  \begin{memo}
    $z=x+yi$ と成分で計算しても答えは出るが、分母が $x^2+y^2$ になり計算が重い。「実数 $\iff$ 共役と等しい」を使うと、因数分解の形が自然に現れる。
  \end{memo}

  \Shomon{4}
  \ascboxZ{$|z|=1$ なら $\overline z=\frac1z$}
  $|z|=1$ のとき $z\overline z=1$ なので $\overline z=\frac1z$。$z$ でくくって、
  \[|z^2-z+1|=|z|\left|z-1+\frac1z\right|=|z+\overline z-1|\]
  $z=\cos\theta+i\sin\theta$ とおくと $z+\overline z=2\cos\theta$（実数）なので、
  \[|z^2-z+1|=|2\cos\theta-1|\]
  $-1\leqq\cos\theta\leqq1$ より $-3\leqq2\cos\theta-1\leqq1$。絶対値が最大になるのは $\cos\theta=-1$（$z=-1$）のときで、
  \[\text{最大値}\ \Kotae{3}\]

  \begin{ptbs}{KEY}[単位円上の $z$ は $\overline z=\frac1z$]
    $|z|=1$ の条件があれば、$\frac1z$ を $\overline z$ に、$z+\frac1z$ を $2\cos\theta$ に書き換えられる。東大の複素数問題で最もよく使う変形である。
  \end{ptbs}

  \begin{simple}[出典：オリジナル]{参考問題28}
    複素数平面上の異なる3点 A$(1)$、B$(z)$、C$(z^2)$ が一直線上にあるような $z$ の条件を求めよ。
  \end{simple}
  \Kai
  3点が異なる条件は $z\neq1$、$z^2\neq1$、$z^2\neq z$、すなわち $z\neq0,\ \pm1$。このとき、一直線上 $\iff\dfrac{z^2-1}{z-1}$ が実数。
  \[\frac{z^2-1}{z-1}=z+1\]
  なので、$z+1$ が実数、つまり $z$ が実数。
  \[\Kotae{z\ \text{は}\ 0,\ \pm1\ \text{以外の実数}}\]

  \begin{ptbs}{CHECK}[翌日、解説を隠して答える（答えは本講の中にある）]
    \kakko{1}　$z$ が実数である条件を共役を使って書け。\par
    \kakko{2}　$|z|=1$ のとき $\overline z$ は何に等しいか。\par
    \kakko{3}　$w=\frac1z$ の描く図形を求める手順を言え。
  \end{ptbs}

  \notefill
  \columnbreak
  \notefill
\end{multicols*}

%% ==================== 第29講 ====================
\Kou{29}{2次曲線・媒介変数・極座標}
\begin{ascolorbox4A}[出典：オリジナル]{例題29}
  \kakko{1}　楕円 $\dfrac{x^2}9+\dfrac{y^2}4=1$ の焦点の座標を求めよ。\qquad
  \kakko{2}　双曲線 $\dfrac{x^2}4-\dfrac{y^2}9=1$ の焦点の座標と漸近線の方程式を求めよ。\par
  \kakko{3}　放物線 $y^2=8x$ の焦点の座標と準線の方程式を求めよ。\qquad
  \kakko{4}　極方程式 $r=2\cos\theta$ を、直交座標 $x,\ y$ の方程式で表せ。\par
  \kakko{5}　楕円 $\dfrac{x^2}9+\dfrac{y^2}4=1$ 上の点 P と、直線 $x+y=6$ の距離の最小値を求めよ。
\end{ascolorbox4A}

\begin{Data}
  テーマ：楕円・双曲線・放物線の定義と標準形、媒介変数表示、極座標と極方程式\par
  難易度：(1)〜(3) A　(4) B　(5) C\quad 前提：第14・17講\par
  この講のゴール：2次曲線を「距離の条件」として理解し、媒介変数で1文字の関数に直して最大・最小を求められる
\end{Data}

\begin{multicols*}{2}
  \begin{koushiki}{2次曲線の定義と標準形}
    \textgt{楕円}：2定点（焦点）からの距離の和が一定な点の軌跡。$\dfrac{x^2}{a^2}+\dfrac{y^2}{b^2}=1\ (a>b>0)$、焦点 $(\pm c,0)$、$c^2=a^2-b^2$、距離の和 $2a$。\par
    \textgt{双曲線}：2定点からの距離の差が一定な点の軌跡。$\dfrac{x^2}{a^2}-\dfrac{y^2}{b^2}=1$、焦点 $(\pm c,0)$、$c^2=a^2+b^2$、漸近線 $y=\pm\dfrac bax$。\par
    \textgt{放物線}：定点（焦点）と定直線（準線）から等距離な点の軌跡。$y^2=4px$、焦点 $(p,0)$、準線 $x=-p$。
  \end{koushiki}

  \Shomon{1}
  \ascboxZ{楕円は $c^2=a^2-b^2$}
  $a=3,\ b=2$ なので $c^2=9-4=5$。
  \[\text{焦点}\ \Kotae{(\sqrt5,\ 0),\ (-\sqrt5,\ 0)}\]

  \begin{zuhyo}
  \begin{tikzpicture}[x=.6cm,y=.6cm,line width=.5pt]
    \draw[->] (-3.6,0) -- (3.8,0) node[right] {$x$};
    \draw[->] (0,-2.5) -- (0,2.7) node[above] {$y$};
    \draw (0,0) ellipse (3 and 2);
    \fill (2.236,0) circle (1.6pt) node[below,font=\footnotesize] {$\mathrm F$};
    \fill (-2.236,0) circle (1.6pt) node[below,font=\footnotesize] {$\mathrm F'$};
    \coordinate (P) at (1.2,1.833);
    \fill (P) circle (1.6pt) node[above right,font=\footnotesize] {$\mathrm P$};
    \draw[dashed] (2.236,0) -- (P) -- (-2.236,0);
    \node[font=\footnotesize] at (1.6,-2.5) {$\mathrm{PF}+\mathrm{PF'}=6$};
  \end{tikzpicture}
  \end{zuhyo}

  \begin{memo}
    なぜ $c^2=a^2-b^2$ か：短軸の端 $(0,\ b)$ も楕円上なので、そこから2焦点までの距離の和も $2a$。対称なので片方は $a$ で、三平方の定理より $b^2+c^2=a^2$。
  \end{memo}

  \Shomon{2}
  \ascboxZ{双曲線は $c^2=a^2+b^2$、漸近線は右辺を0に}
  $a=2,\ b=3$ なので $c^2=4+9=13$。漸近線は $\frac{x^2}4-\frac{y^2}9=0$、つまり $y=\pm\frac32x$ である（$x$ が大きいところでは右辺の1が相対的に無視でき、曲線がこの2直線に近づく）。
  \[\text{焦点}\ \Kotae{(\pm\sqrt{13},\ 0)},\quad\text{漸近線}\ \Kotae{y=\pm\frac32x}\]

  \Shomon{3}
  \ascboxZ{放物線は $4p$ の形にそろえる}
  $y^2=8x=4\cdot2\cdot x$ なので $p=2$。
  \[\text{焦点}\ \Kotae{(2,\ 0)},\quad\text{準線}\ \Kotae{x=-2}\]
  確かめ：点 $(x,y)$ と焦点の距離の2乗 $(x-2)^2+y^2$ と、準線までの距離の2乗 $(x+2)^2$ が等しいとおくと、$y^2=8x$ になる。

  \Shomon{4}
  \ascboxZ{極座標：$x=r\cos\theta,\ y=r\sin\theta$}
  点の位置を、原点からの距離 $r$ と、$x$ 軸の正の向きからの角 $\theta$ で表すのが\textgt{極座標} $(r,\ \theta)$ である。直交座標とは $x=r\cos\theta,\ y=r\sin\theta,\ r^2=x^2+y^2$ で結ばれる。\par
  $r=2\cos\theta$ の両辺に $r$ を掛けると $r^2=2r\cos\theta$、すなわち $x^2+y^2=2x$。
  \[\Kotae{(x-1)^2+y^2=1}\]

  \begin{memo}
    $r$ を掛けると、$r=0$（原点）が新たに加わる可能性がある。しかし元の式でも $\theta=\frac\pi2$ のとき $r=0$ で原点を通るので、点は増えていない。掛け算・割り算のたびに「点が増減していないか」を確認する。
  \end{memo}

  \Shomon{5}
  \ascboxZ{媒介変数で1変数にし、合成する}
  楕円 $\frac{x^2}9+\frac{y^2}4=1$ は、単位円を $x$ 方向に3倍、$y$ 方向に2倍したものなので、P は
  \[\mathrm P(3\cos\theta,\ 2\sin\theta)\quad(0\leqq\theta<2\pi)\]
  と表せる（\textgt{媒介変数表示}）。点と直線の距離の公式より、
  \[d=\frac{|3\cos\theta+2\sin\theta-6|}{\sqrt2}\]
  合成すると $3\cos\theta+2\sin\theta=\sqrt{13}\sin(\theta+\varphi)$（$\varphi$ はある定角）なので、その値は $-\sqrt{13}$ 以上 $\sqrt{13}$ 以下のすべての値をとる。$\sqrt{13}<6$ より絶対値の中は常に負で、
  \[d=\frac{6-(3\cos\theta+2\sin\theta)}{\sqrt2}\geqq\frac{6-\sqrt{13}}{\sqrt2}\]
  等号は $\sin(\theta+\varphi)=1$ となる $\theta$ で成り立つ。よって最小値は、
  \[\Kotae{\frac{6\sqrt2-\sqrt{26}}2}\]

  \begin{ptbs}{KEY}[楕円上の点は $(a\cos\theta,\ b\sin\theta)$]
    2変数 $x,\ y$ の条件つき最大・最小を、1変数 $\theta$ の三角関数の問題に変える。合成の後は「中身の範囲」と「等号成立」を確認する。
  \end{ptbs}

  \begin{simple}[出典：オリジナル]{参考問題29}
    媒介変数 $t\ (t\neq0)$ を用いて $x=t+\dfrac1t$、$y=t-\dfrac1t$ と表される曲線を、$x,\ y$ の方程式で表せ。
  \end{simple}
  \Kai
  $x+y=2t$、$x-y=\frac2t$ なので、掛けて $t$ を消去すると $x^2-y^2=4$。\par
  逆に、$x^2-y^2=4$ を満たす点 $(x,y)$ に対し $t=\frac{x+y}2$ とおくと、$\frac{x-y}2=\frac{x^2-y^2}{2(x+y)}=\frac2{x+y}=\frac1t$（$x+y\neq0$ は $x^2-y^2=4$ から従う）なので、その点は曲線上にある。よって曲線全体は、
  \[\Kotae{\text{双曲線}\ x^2-y^2=4}\]
  \begin{memo}
    パラメータを消去したら、得られた式の点が全部実現されるか（範囲の制限がないか）を必ず確かめる。$x=\cos t,\ y=\cos^2t$ なら $y=x^2$ のうち $-1\leqq x\leqq1$ の部分だけになる、というように、範囲が狭まることがよくある。
  \end{memo}

  \begin{ptbs}{CHECK}[翌日、解説を隠して答える（答えは本講の中にある）]
    \kakko{1}　楕円と双曲線で、$a,\ b,\ c$ の関係式を言え。\par
    \kakko{2}　楕円 $\frac{x^2}{a^2}+\frac{y^2}{b^2}=1$ 上の点の媒介変数表示を書け。\par
    \kakko{3}　パラメータを消去した後に確かめることは何か。
  \end{ptbs}

  \notefill
\end{multicols*}
%% ==================== 第30講 ====================
\Kou{30}{数列の極限}
\begin{ascolorbox4A}[出典：オリジナル]{例題30}
  次の極限を求めよ。\par
  \kakko{1}　$\displaystyle\lim_{n\to\infty}\frac{2n^2+1}{3n^2-n}$\qquad
  \kakko{2}　$\displaystyle\lim_{n\to\infty}\left(\sqrt{n^2+n}-n\right)$\qquad
  \kakko{3}　$\displaystyle\lim_{n\to\infty}\frac{3^n+2^n}{3^{n+1}-2^n}$\par
  \kakko{4}　無限級数 $\displaystyle\sum_{n=1}^\infty\left(\frac13\right)^n$ の和を求めよ。また、$0.999\cdots=1$ であることを無限級数で説明せよ。\par
  \kakko{5}　$\displaystyle\lim_{n\to\infty}\frac n{2^n}=0$ であることを証明せよ。
\end{ascolorbox4A}

\begin{Data}
  テーマ：数列の極限、不定形の処理（最高次でくくる・有理化）、$r^n$ の極限、無限等比級数、はさみうちの原理\par
  難易度：(1) A　(2) B　(3) B　(4) B　(5) C\quad 前提：第21講\par
  この講のゴール：「$\infty-\infty$」「$\frac\infty\infty$」を見たら、形を変えて決着をつける方法を選べる
\end{Data}

\begin{multicols*}{2}
  \begin{koushiki}{極限の言葉}
    $n$ を限りなく大きくするとき、$a_n$ がある値 $\alpha$ に限りなく近づくことを $\displaystyle\lim_{n\to\infty}a_n=\alpha$ と書き、$\{a_n\}$ は $\alpha$ に\textgt{収束}するという。収束しないことを\textgt{発散}という（$+\infty$ に発散、$-\infty$ に発散、振動）。\par
    基本：$\displaystyle\lim_{n\to\infty}\frac1n=0$。収束する数列どうしの和・差・積・商（分母の極限 $\neq0$）の極限は、極限どうしの和・差・積・商。
  \end{koushiki}
  $\infty$ は数ではない。だから $\infty-\infty=0$、$\frac\infty\infty=1$ とはできない。例えば $n^2-n$ も $n-n^2$ も「$\infty-\infty$」の形だが、行き先は $+\infty$ と $-\infty$ で全く違う。このように形だけでは決まらないものを\textgt{不定形}という。

  \Shomon{1}
  \ascboxZ{分母の最高次で分母・分子を割る}
  分母・分子を $n^2$ で割ると、各項が収束する形になる。
  \[\frac{2n^2+1}{3n^2-n}=\frac{2+\frac1{n^2}}{3-\frac1n}\]
  $n\to\infty$ で $\frac1{n^2}\to0,\ \frac1n\to0$ なので、
  \[\lim_{n\to\infty}\frac{2n^2+1}{3n^2-n}=\Kotae{\frac23}\]

  \Shomon{2}
  \ascboxZ{$\sqrt{\ }-\sqrt{\ }$ 型は分子の有理化}
  $\infty-\infty$ の形である。$(\sqrt A-B)(\sqrt A+B)=A-B^2$ を使って根号を外す。
  \[
  \begin{aligned}
  \sqrt{n^2+n}-n&=\frac{(n^2+n)-n^2}{\sqrt{n^2+n}+n}=\frac n{\sqrt{n^2+n}+n}\\
  &=\frac1{\sqrt{1+\frac1n}+1}
  \end{aligned}
  \]
  （最後は分母・分子を $n$ で割った。$n>0$ なので $\frac{\sqrt{n^2+n}}n=\sqrt{1+\frac1n}$。）よって、
  \[\lim_{n\to\infty}\left(\sqrt{n^2+n}-n\right)=\Kotae{\frac12}\]

  \Shomon{3}
  \ascboxZ{いちばん大きい底で割る}
  \begin{koushiki}{$r^n$ の極限}
    $r>1$ のとき $+\infty$、\quad$r=1$ のとき $1$\par
    $-1<r<1$ のとき $0$、\quad$r\leqq-1$ のとき振動（極限なし）
  \end{koushiki}
  分母・分子を $3^n$ で割ると $\left(\frac23\right)^n$ が現れ、これは0に収束する。
  \[\frac{3^n+2^n}{3^{n+1}-2^n}=\frac{1+\left(\frac23\right)^n}{3-\left(\frac23\right)^n}\]
  よって、
  \[\lim_{n\to\infty}\frac{3^n+2^n}{3^{n+1}-2^n}=\Kotae{\frac13}\]

  \Shomon{4}
  \ascboxZ{無限級数は部分和の極限}
  \begin{koushiki}{無限等比級数}
    $\displaystyle\sum_{n=1}^\infty ar^{n-1}=a+ar+ar^2+\cdots$ は、\par
    $a\neq0$ のとき、$|r|<1$ なら収束して和は $\dfrac a{1-r}$、$|r|\geqq1$ なら発散。\par
    $a=0$ のときは0に収束。
  \end{koushiki}
  \textgt{理由}：無限個の和は、まず第 $n$ 項までの和（部分和）$S_n=\frac{a(1-r^n)}{1-r}$ を作り、$n\to\infty$ の極限として定める。$|r|<1$ なら $r^n\to0$ なので $S_n\to\frac a{1-r}$。\par
  本問は初項 $\frac13$、公比 $\frac13$ なので、
  \[\sum_{n=1}^\infty\left(\frac13\right)^n=\frac{\frac13}{1-\frac13}=\Kotae{\frac12}\]
  $0.999\cdots=0.9+0.09+0.009+\cdots$ は初項 $0.9$、公比 $0.1$ の無限等比級数で、和は $\frac{0.9}{1-0.1}=1$。\par
  「$0.999\cdots$ は1に近づくだけで1ではない」という感覚は誤りである。$0.999\cdots$ という記号そのものが「部分和 $0.9,\ 0.99,\ \cdots$ の極限」を表しており、その値はちょうど1である。

  \Shomon{5}
  \ascboxZ{はさみうち：両側から同じ値ではさむ}
  \begin{koushiki}{はさみうちの原理}
    $a_n\leqq c_n\leqq b_n$ がすべての $n$（ある番号以降）で成り立ち、$a_n\to\alpha,\ b_n\to\alpha$ ならば、$c_n\to\alpha$。
  \end{koushiki}
  分母の $2^n$ を、$n$ より速く大きくなる式で下から押さえる。二項定理より、$n\geqq2$ のとき
  \[2^n=(1+1)^n=1+n+\frac{n(n-1)}2+\cdots>\frac{n(n-1)}2\]
  （すべての項が正なので、一部の項だけ残すと小さくなる。）よって、
  \[0<\frac n{2^n}<\frac{n}{\frac{n(n-1)}2}=\frac2{n-1}\]
  $n\to\infty$ で $\frac2{n-1}\to0$ なので、はさみうちの原理より $\displaystyle\lim_{n\to\infty}\frac n{2^n}=0$。\QED

  \begin{ptbs}{KEY}[極限が直接計算できないときは「はさむ」]
    求めたい式を、極限がわかる簡単な式で上下から押さえる。押さえ方には二項定理、$|\sin|\leqq1$、積分（第36講）などを使う。
  \end{ptbs}

  \begin{simple}[出典：オリジナル]{参考問題30}
    無限級数 $\displaystyle\sum_{n=1}^\infty x(1-x)^{n-1}$ が収束するような実数 $x$ の範囲と、そのときの和 $f(x)$ を求めよ。
  \end{simple}
  \Kai
  初項 $x$、公比 $1-x$ の無限等比級数である。初項が0かどうかで分ける。\par
  $x=0$ のとき：すべての項が0なので収束し、和は0。\par
  $x\neq0$ のとき：収束 $\iff|1-x|<1\iff0<x<2$。和は $\frac x{1-(1-x)}=1$。
  \[\Kotae{0\leqq x<2\ ;\quad f(x)=\begin{cases}0&(x=0)\\1&(0<x<2)\end{cases}}\]
  \begin{memo}
    $f(x)$ は $x=0$ で値が飛んでいる。連続な関数を無限個足しても、連続になるとは限らない。初項が0の場合を落とすと、収束範囲を誤る。
  \end{memo}

  \begin{ptbs}{CHECK}[翌日、解説を隠して答える（答えは本講の中にある）]
    \kakko{1}　不定形とは何か。なぜ形だけでは極限が決まらないか。\par
    \kakko{2}　$\sqrt{\ }-\sqrt{\ }$ 型の極限の処理を言え。\par
    \kakko{3}　はさみうちの原理を言え。
  \end{ptbs}

  \notefill
\end{multicols*}

%% ==================== 第31講 ====================
\Kou{31}{関数の極限と連続性}
\begin{ascolorbox4A}[出典：オリジナル]{例題31}
  \kakko{1}　$\displaystyle\lim_{x\to0}\frac{\sin3x}x$ と $\displaystyle\lim_{x\to0}\frac{1-\cos x}{x^2}$ を求めよ。\par
  \kakko{2}　$\displaystyle\lim_{x\to2}\frac{x^2+ax+b}{x-2}=5$ となるような定数 $a,\ b$ の値を求めよ。\par
  \kakko{3}　$e=\displaystyle\lim_{h\to0}(1+h)^{\frac1h}$ とする。$\displaystyle\lim_{x\to0}\frac{\log(1+x)}x=1$ と $\displaystyle\lim_{x\to0}\frac{e^x-1}x=1$ を示せ。ただし $\log$ は $e$ を底とする対数とする。\par
  \kakko{4}　方程式 $x^3+x-1=0$ は、$0<x<1$ の範囲にただ1つの実数解をもつことを示せ。
\end{ascolorbox4A}

\begin{Data}
  テーマ：$\frac{\sin x}x\to1$ の証明と応用、極限から係数を決める、$e$ と自然対数、連続性と中間値の定理\par
  難易度：(1) B　(2) B　(3) C　(4) B\quad 前提：第16・19・30講\par
  この講のゴール：数学Ⅲの微分公式の土台となる3つの極限（$\sin$、$\log$、$e^x$）を自分で導ける
\end{Data}

\begin{multicols*}{2}
  \Shomon{1}
  \ascboxZ{$\dfrac{\sin x}x\to1$ は面積の大小から}
  \begin{zuhyo}
  \begin{tikzpicture}[x=2.2cm,y=2.2cm,line width=.5pt]
    \draw[->] (-.1,0) -- (1.3,0);
    \draw (1,0) arc (0:60:1);
    \fill[black!12] (0,0) -- (1,0) arc (0:40:1) -- cycle;
    \draw (0,0) -- (.766,.643) -- (1,0);
    \draw (0,0) -- (1,.839) -- (1,0);
    \node[below left,font=\footnotesize] at (0,0) {$\mathrm O$};
    \node[below,font=\footnotesize] at (1,0) {$\mathrm A$};
    \node[above left,font=\footnotesize] at (.766,.643) {$\mathrm P$};
    \node[right,font=\footnotesize] at (1,.839) {$\mathrm T$};
    \node[font=\footnotesize] at (.22,.08) {$x$};
  \end{tikzpicture}
  \end{zuhyo}
  $0<x<\frac\pi2$ とし、単位円上に $\angle\mathrm{AOP}=x$ の点 P、A における接線と直線 OP の交点 T をとる。図で
  \[\triangle\mathrm{OAP}<\text{扇形 OAP}<\triangle\mathrm{OAT}\]
  なので $\frac12\sin x<\frac12x<\frac12\tan x$。各辺を $\frac12\sin x\ (>0)$ で割って逆数をとると、
  \[\cos x<\frac{\sin x}x<1\]
  $x\to+0$ で $\cos x\to1$ なので、はさみうちより $\frac{\sin x}x\to1$。$\frac{\sin x}x$ は $x$ を $-x$ にしても値が変わらないので、$x\to-0$ でも同じ。

  \begin{koushiki}{基本の極限}
    \[\lim_{x\to0}\frac{\sin x}x=1\quad\text{（$x$ は弧度法）}\]
    \[\lim_{x\to0}\frac{\log(1+x)}x=1,\quad\lim_{x\to0}\frac{e^x-1}x=1\]
  \end{koushiki}

  $\sin3x$ は「中身」と分母をそろえる。$\theta=3x$ とおくと $x\to0$ のとき $\theta\to0$。
  \[\frac{\sin3x}x=3\cdot\frac{\sin3x}{3x}\to\Kotae{3}\]
  $1-\cos x$ は $1+\cos x$ を掛けて $\sin^2x$ を作る。
  \[\frac{1-\cos x}{x^2}=\frac{\sin^2x}{x^2(1+\cos x)}=\left(\frac{\sin x}x\right)^2\frac1{1+\cos x}\]
  $x\to0$ で $1^2\cdot\frac12$ に近づくので、
  \[\lim_{x\to0}\frac{1-\cos x}{x^2}=\Kotae{\frac12}\]

  \Shomon{2}
  \ascboxZ{分母 $\to0$ で極限が有限なら、分子 $\to0$}
  $x\to2$ で分母 $x-2\to0$。もし分子が0以外の値に近づくなら、商は $\pm\infty$ になってしまう。極限が5（有限）なので、分子も0に近づかなければならない（必要条件）。
  \[4+2a+b=0\quad\text{すなわち}\quad b=-2a-4\]
  このとき分子 $x^2+ax-2a-4=(x-2)(x+a+2)$ と因数分解できるので、
  \[\frac{x^2+ax+b}{x-2}=x+a+2\to a+4\]
  $a+4=5$ より $a=1$、$b=-6$。逆にこのとき極限は5になる（十分性）。
  \[\Kotae{a=1,\ b=-6}\]

  \begin{ptbs}{KEY}[$\lim\dfrac{\text{分子}}{\text{分母}}=\text{有限},\ \text{分母}\to0$ なら 分子 $\to0$]
    必要条件として分子 $\to0$ を使い、係数を1つ減らす。最後に、求めた値で本当に極限が指定の値になるか（十分性）を確かめる。
  \end{ptbs}

  \Shomon{3}
  \ascboxZ{$e$ の定義から $\log$ と $e^x$ の極限へ}
  \textgt{1つ目}：対数の性質より
  \[\frac{\log(1+x)}x=\log(1+x)^{\frac1x}\]
  $x\to0$ のとき $(1+x)^{\frac1x}\to e$ で、$\log$ は連続なので、$\log e=1$ に近づく。\par
  \textgt{2つ目}：$t=e^x-1$ とおくと、$x=\log(1+t)$ で、$x\to0$ のとき $t\to0$。
  \[\frac{e^x-1}x=\frac t{\log(1+t)}=\frac1{\frac{\log(1+t)}t}\to\frac11=1\]
  以上で示された。\QED

  \begin{memo}
    $e=2.71828\cdots$ は、$y=e^x$ のグラフの $x=0$ での接線の傾きがちょうど1になる底である（2つ目の極限がそれを言っている）。この性質のおかげで $(e^x)'=e^x$ という簡単な公式が成り立つ（第32講）。
  \end{memo}

  \Shomon{4}
  \ascboxZ{存在は中間値の定理、ただ1つは単調性}
  \begin{koushiki}{連続と中間値の定理}
    $\displaystyle\lim_{x\to a}f(x)=f(a)$ のとき、$f(x)$ は $x=a$ で\textgt{連続}という（グラフがつながっている）。\par
    \textgt{中間値の定理}：$f(x)$ が $a\leqq x\leqq b$ で連続で、$f(a)$ と $f(b)$ が異符号なら、$f(c)=0$ となる $c$ が $a<c<b$ に少なくとも1つある。
  \end{koushiki}
  $f(x)=x^3+x-1$ は多項式なので連続。$f(0)=-1<0$、$f(1)=1>0$ なので、中間値の定理より $0<x<1$ に解が少なくとも1つある。\par
  さらに $f'(x)=3x^2+1>0$ なので $f(x)$ は増加関数であり、値0をとる $x$ は2つ以上ない。よって解はただ1つ。\QED

  \begin{memo}
    「存在する」と「ただ1つ」は別の主張で、使う道具も違う。存在は中間値の定理、一意性（ただ1つ）は単調性。東大の答案ではこの2つを分けて書く。
  \end{memo}

  \begin{simple}[出典：オリジナル]{参考問題31}
    $\displaystyle\lim_{x\to\infty}\left(\sqrt{x^2+x}-x\right)$ と $\displaystyle\lim_{x\to-\infty}\left(\sqrt{x^2+x}-x\right)$ を求めよ。
  \end{simple}
  \Kai
  $x\to\infty$ のときは第30講\kakko{2}と同じ計算で $\Kotae{\frac12}$。\par
  $x\to-\infty$ のときは、$x=-t$ とおくと $t\to\infty$ で、
  \[\sqrt{x^2+x}-x=\sqrt{t^2-t}+t\]
  これは「$\infty+\infty$」で不定形ではない。$+\infty$ に発散する。
  \[\Kotae{+\infty\ \text{に発散}}\]
  \begin{memo}
    $x\to-\infty$ では $\sqrt{x^2}=|x|=-x$ である。$x=-t$ と置き換えて、正の文字で考えると符号の誤りを防げる。
  \end{memo}

  \begin{ptbs}{CHECK}[翌日、解説を隠して答える（答えは本講の中にある）]
    \kakko{1}　$\frac{\sin x}x\to1$ の証明で使う面積の不等式を書け。\par
    \kakko{2}　分母 $\to0$ で極限が有限のとき、分子について何が言えるか。\par
    \kakko{3}　解の「存在」と「ただ1つ」には、それぞれ何を使うか。
  \end{ptbs}

  \notefill
\end{multicols*}

%% ==================== 第32講 ====================
\Kou{32}{微分法（数学Ⅲ）}
\begin{ascolorbox4A}[出典：オリジナル]{例題32}
  次の関数を微分せよ。\par
  \kakko{1}　$x^2\sin x$、\ $\dfrac x{x^2+1}$\qquad
  \kakko{2}　$\sin^32x$、\ $e^{x^2}$、\ $\log(x^2+1)$\qquad
  \kakko{3}　$y=x^x\ (x>0)$\par
  \kakko{4}　$x^2+y^2=25$ のとき $\dfrac{dy}{dx}$ を $x,\ y$ で表せ。また、$x=\cos^3t$、$y=\sin^3t$ のとき $\dfrac{dy}{dx}$ を $t$ で表せ。\par
  \kakko{5}　$(\sin x)'=\cos x$ を、導関数の定義に従って証明せよ。
\end{ascolorbox4A}

\begin{Data}
  テーマ：積・商・合成関数の微分、三角・指数・対数関数の導関数、対数微分法、陰関数・媒介変数の微分\par
  難易度：(1)(2) A　(3) B　(4) B　(5) B\quad 前提：第19・31講\par
  この講のゴール：どんな関数でも、公式の組合せで機械的に微分でき、公式の由来も説明できる
\end{Data}

\begin{multicols*}{2}
  \begin{koushiki}{微分の公式}
    積：$(fg)'=f'g+fg'$\qquad 商：$\left(\dfrac fg\right)'=\dfrac{f'g-fg'}{g^2}$\par
    合成：$\{f(g(x))\}'=f'(g(x))\,g'(x)$\par
    $(\sin x)'=\cos x,\ (\cos x)'=-\sin x,\ (\tan x)'=\dfrac1{\cos^2x}$\par
    $(e^x)'=e^x,\ (\log|x|)'=\dfrac1x,\ (x^\alpha)'=\alpha x^{\alpha-1}$
  \end{koushiki}

  \Shomon{1}
  \ascboxZ{積と商の微分}
  \textgt{積の公式の理由}：$f(x+h)g(x+h)-f(x)g(x)=\{f(x+h)-f(x)\}g(x+h)+f(x)\{g(x+h)-g(x)\}$ と分けて $h$ で割ると、$h\to0$ で $f'g+fg'$。長方形の面積の増え方（縦の増加×横＋縦×横の増加）と同じである。
  \[(x^2\sin x)'=\Kotae{2x\sin x+x^2\cos x}\]
  \[
  \begin{aligned}
  \left(\frac x{x^2+1}\right)'&=\frac{1\cdot(x^2+1)-x\cdot2x}{(x^2+1)^2}\\
  &=\Kotae{\frac{1-x^2}{(x^2+1)^2}}
  \end{aligned}
  \]

  \Shomon{2}
  \ascboxZ{合成関数：外から順に微分して掛ける}
  $y=f(u),\ u=g(x)$ のとき $\dfrac{dy}{dx}=\dfrac{dy}{du}\cdot\dfrac{du}{dx}$。「$u$ が少し変わると $y$ は何倍変わるか」と「$x$ が少し変わると $u$ は何倍変わるか」を掛ける、という意味である。\par
  $\sin^32x$ は「3乗」→「$\sin$」→「$2x$」の3重の合成。外側から順に微分して掛ける。
  \[(\sin^32x)'=3\sin^22x\cdot\cos2x\cdot2=\Kotae{6\sin^22x\cos2x}\]
  \[(e^{x^2})'=e^{x^2}\cdot2x=\Kotae{2xe^{x^2}}\]
  \[\{\log(x^2+1)\}'=\frac1{x^2+1}\cdot2x=\Kotae{\frac{2x}{x^2+1}}\]

  \Shomon{3}
  \ascboxZ{底にも指数にも $x$ があれば対数微分法}
  $x^x$ は $x^\alpha$ の公式も $a^x$ の公式も使えない（底も指数も動く）。$x>0$ より $y>0$ なので、両辺の自然対数をとる。
  \[\log y=x\log x\]
  両辺を $x$ で微分する。左辺は合成関数の微分で $\frac{y'}y$。
  \[\frac{y'}y=\log x+x\cdot\frac1x=\log x+1\]
  よって、
  \[y'=\Kotae{x^x(\log x+1)}\]

  \Shomon{4}
  \ascboxZ{陰関数は $y$ を $x$ の関数とみて微分}
  $x^2+y^2=25$ の両辺を $x$ で微分する。$y$ は $x$ の関数なので、$y^2$ の微分は合成関数の微分で $2y\cdot\frac{dy}{dx}$。
  \[2x+2y\frac{dy}{dx}=0\]
  $y\neq0$ のとき、
  \[\frac{dy}{dx}=\Kotae{-\frac xy}\]
  （円の接線は半径に垂直なので、傾きが $\frac yx$ の逆数の $-1$ 倍になる。図形的にも合う。）

  \ascboxZ{媒介変数は $\frac{dy}{dt}\div\frac{dx}{dt}$}
  $\dfrac{dx}{dt}=-3\cos^2t\sin t$、$\dfrac{dy}{dt}=3\sin^2t\cos t$。$\frac{dx}{dt}\neq0$ のとき、
  \[\frac{dy}{dx}=\frac{3\sin^2t\cos t}{-3\cos^2t\sin t}=\Kotae{-\tan t}\]

  \Shomon{5}
  \ascboxZ{加法定理で分けて、基本の極限を使う}
  \[
  \begin{aligned}
  &\frac{\sin(x+h)-\sin x}h\\
  &=\frac{\sin x\cos h+\cos x\sin h-\sin x}h\\
  &=\sin x\cdot\frac{\cos h-1}h+\cos x\cdot\frac{\sin h}h
  \end{aligned}
  \]
  第31講より $\frac{\sin h}h\to1$。また $\frac{\cos h-1}h=-\frac{1-\cos h}{h^2}\cdot h\to-\frac12\cdot0=0$。よって $h\to0$ のとき、上の式は $\sin x\cdot0+\cos x\cdot1=\cos x$ に近づく。\QED

  \begin{ptbs}{KEY}[3つの基本極限 → すべての微分公式]
    $\frac{\sin h}h\to1$ から三角関数、$\frac{e^h-1}h\to1$ から $e^x$、$\frac{\log(1+h)}h\to1$ から $\log x$ の導関数が出る。あとは積・商・合成の公式で全部つながる。
  \end{ptbs}

  \begin{memo}
    同じ方法で：$\frac{e^{x+h}-e^x}h=e^x\cdot\frac{e^h-1}h\to e^x$。$\frac{\log(x+h)-\log x}h=\frac1x\cdot\frac{\log(1+\frac hx)}{\frac hx}\to\frac1x$（$x>0$）。
  \end{memo}

  \begin{simple}[出典：オリジナル]{参考問題32}
    $f(x)=x|x|$ は $x=0$ で微分可能か。微分可能ならば $f'(0)$ を求めよ。また、$g(x)=|x|$ についても調べよ。
  \end{simple}
  \Kai
  定義に従い、$h\to0$ のときの $\frac{f(h)-f(0)}h$ の極限を調べる。
  \[\frac{f(h)-f(0)}h=\frac{h|h|}h=|h|\to0\]
  右から近づけても左から近づけても0なので、微分可能で $f'(0)=\Kotae{0}$。\par
  $g$ については $\frac{g(h)-g(0)}h=\frac{|h|}h$ で、$h>0$ なら1、$h<0$ なら $-1$。左右で近づく値が違うので極限はなく、$g(x)=|x|$ は $x=0$ で\Kotae{\text{微分可能でない}}。
  \begin{memo}
    $|x|$ のグラフは原点でとがっていて、接線が1本に決まらない。微分可能とは「グラフがなめらかで、接線が1本に決まる」ことである。
  \end{memo}

  \begin{ptbs}{CHECK}[翌日、解説を隠して答える（答えは本講の中にある）]
    \kakko{1}　積の微分公式が成り立つ理由を言え。\par
    \kakko{2}　対数微分法を使うのはどんなときか。\par
    \kakko{3}　媒介変数表示の $\frac{dy}{dx}$ を書け。
  \end{ptbs}

  \notefill
\end{multicols*}
%% ==================== 第33講 ====================
\Kou{33}{微分の応用(1) グラフ}
\begin{ascolorbox4A}[出典：オリジナル]{例題33}
  \kakko{1}　関数 $y=xe^{-x}$ の増減と凹凸を調べ、グラフの概形をかけ。ただし $\displaystyle\lim_{x\to\infty}xe^{-x}=0$ を用いてよい。\par
  \kakko{2}　曲線 $y=x+\dfrac1x$ の漸近線を求めよ。\par
  \kakko{3}　$x>0$ における関数 $f(x)=\dfrac{\log x}x$ の最大値を求めよ。また、それを用いて $e^\pi$ と $\pi^e$ の大小を比較せよ。
\end{ascolorbox4A}

\begin{Data}
  テーマ：第2次導関数と凹凸・変曲点、増減表、漸近線、グラフを使った大小比較\par
  難易度：(1) B　(2) B　(3) C\quad 前提：第19・31・32講\par
  この講のゴール：増減・凹凸・端の様子・漸近線の4点からグラフの概形を確定し、不等式の証明に使える
\end{Data}

\begin{multicols*}{2}
  \begin{koushiki}{凹凸と変曲点}
    $f''(x)>0$ の区間では、接線の傾き $f'(x)$ が増加するので、グラフは\textgt{下に凸}（$\cup$ 形）。\par
    $f''(x)<0$ の区間では\textgt{上に凸}（$\cap$ 形）。\par
    凹凸が入れ替わる点を\textgt{変曲点}という。
  \end{koushiki}

  \Shomon{1}
  \ascboxZ{$y'$ と $y''$ の符号を1つの表に}
  積の微分で、
  \[
  \begin{aligned}
  y'&=e^{-x}+x\cdot(-e^{-x})=(1-x)e^{-x}\\
  y''&=-e^{-x}-(1-x)e^{-x}=(x-2)e^{-x}
  \end{aligned}
  \]
  $e^{-x}>0$ なので、符号は $1-x$ と $x-2$ で決まる。

  \begin{zuhyo}
  \begin{tblr}{colspec={c|ccccc}, hlines, vlines, row{1}={font=\TRgt}, cells={font=\small}}
  $x$ & $\cdots$ & $1$ & $\cdots$ & $2$ & $\cdots$ \\
  $y'$ & $+$ & $0$ & $-$ & $-$ & $-$ \\
  $y''$ & $-$ & $-$ & $-$ & $0$ & $+$ \\
  凹凸 & $\cap$ & & $\cap$ & & $\cup$ \\
  $y$ & $\nearrow$ & $\frac1e$ & $\searrow$ & $\frac2{e^2}$ & $\searrow$ \\
  \end{tblr}
  \end{zuhyo}

  端の様子：$x\to\infty$ で $y\to0$（与えられた極限）。$x\to-\infty$ で $x\to-\infty,\ e^{-x}\to\infty$ なので $y\to-\infty$。

  \begin{zuhyo}
  \begin{tikzpicture}[x=.8cm,y=3.2cm,line width=.5pt]
    \draw[->] (-.8,0) -- (5.4,0) node[right] {$x$};
    \draw[->] (0,-.5) -- (0,.55) node[above] {$y$};
    \draw[domain=-.38:5.2,samples=80] plot (\x,{(\x)*exp(-(\x))});
    \draw[dashed] (1,0) -- (1,.368) -- (0,.368) node[left,font=\footnotesize] {$\frac1e$};
    \draw[dashed] (2,0) -- (2,.271);
    \fill (2,.271) circle (1.5pt);
    \node[below,font=\footnotesize] at (1,0) {$1$};
    \node[below,font=\footnotesize] at (2,0) {$2$};
    \node[below left,font=\footnotesize] at (0,0) {$\mathrm O$};
    \node[font=\footnotesize] at (3.9,.42) {変曲点 $\left(2,\frac2{e^2}\right)$};
  \end{tikzpicture}
  \end{zuhyo}

  まとめると、
  \[\Kotae{\text{極大値}\ \frac1e\ (x=1),\ \text{変曲点}\left(2,\ \frac2{e^2}\right)}\]
  漸近線は $x$ 軸（$x\to\infty$ のとき）。グラフは上の図。

  \begin{memo}
    $\lim_{x\to\infty}\frac x{e^x}=0$ は「指数関数は多項式より速く大きくなる」ことを表す。証明は第34講\kakko{4}で行う。東大では、使ってよいと書かれていなければ証明が必要になることがある。
  \end{memo}

  \Shomon{2}
  \ascboxZ{漸近線は「縦」と「斜め」を別々に探す}
  \begin{koushiki}{漸近線の求め方}
    \textgt{縦}：$x\to a$ で $y\to\pm\infty$ なら、直線 $x=a$ が漸近線。\par
    \textgt{斜め}：$\displaystyle a=\lim_{x\to\pm\infty}\frac{f(x)}x,\ b=\lim_{x\to\pm\infty}\{f(x)-ax\}$ が有限なら、$y=ax+b$ が漸近線。
  \end{koushiki}
  縦：$x\to+0$ で $\frac1x\to+\infty$、$x\to-0$ で $-\infty$。よって $x=0$。\par
  斜め：$f(x)-x=\frac1x\to0$（$x\to\pm\infty$）なので、$y=x$。
  \[\Kotae{\text{直線}\ x=0\ \text{と直線}\ y=x}\]

  \Shomon{3}
  \ascboxZ{比べたい2数を、1つの関数の2つの値にする}
  $f'(x)=\dfrac{\frac1x\cdot x-\log x}{x^2}=\dfrac{1-\log x}{x^2}$。$0<x<e$ で $f'>0$、$x>e$ で $f'<0$ なので、$x=e$ で最大。
  \[\text{最大値}\ f(e)=\Kotae{\frac1e}\]
  $\pi\neq e$ なので $f(\pi)<f(e)$。
  \[\frac{\log\pi}\pi<\frac1e\iff e\log\pi<\pi\iff\log\pi^e<\log e^\pi\]
  $\log$ は増加関数なので、
  \[\Kotae{\pi^e<e^\pi}\]

  \begin{ptbs}{KEY}[$a^b$ と $b^a$ の比較は $\dfrac{\log x}x$]
    $a^b<b^a\iff b\log a<a\log b\iff\frac{\log a}a<\frac{\log b}b$。2数を同じ関数の値にそろえれば、増減表だけで大小がわかる。
  \end{ptbs}

  \begin{simple}[出典：オリジナル]{参考問題33}
    曲線 $y=e^{-x^2}$ の変曲点の座標を求めよ。
  \end{simple}
  \Kai
  $y'=-2xe^{-x^2}$、$y''=-2e^{-x^2}+4x^2e^{-x^2}=2(2x^2-1)e^{-x^2}$。$y''$ は $x=\pm\frac1{\sqrt2}$ で符号が変わる（$2x^2-1$ の符号）。そのときの $y$ は $e^{-\frac12}=\frac1{\sqrt e}$。
  \[\Kotae{\left(\pm\frac1{\sqrt2},\ \frac1{\sqrt e}\right)}\]
  \begin{memo}
    $y''=0$ となる点でも、前後で符号が変わらなければ変曲点ではない（$y=x^4$ の原点）。必ず符号の変化を確かめる。
  \end{memo}

  \begin{ptbs}{CHECK}[翌日、解説を隠して答える（答えは本講の中にある）]
    \kakko{1}　グラフの凹凸は何の符号で決まるか。\par
    \kakko{2}　斜めの漸近線 $y=ax+b$ の $a,\ b$ の求め方を言え。\par
    \kakko{3}　$e^\pi$ と $\pi^e$ の比較で使う関数と、その最大値は何か。
  \end{ptbs}

  \notefill
  \columnbreak
  \notefill
\end{multicols*}

%% ==================== 第34講 ====================
\Kou{34}{微分の応用(2) 不等式と方程式}
\begin{ascolorbox4A}[出典：オリジナル]{例題34}
  \kakko{1}　$x>0$ のとき、$e^x>1+x+\dfrac{x^2}2$ が成り立つことを証明せよ。\par
  \kakko{2}　$a$ を実数の定数とする。方程式 $e^x=ax$ の異なる実数解の個数を調べよ。\par
  \kakko{3}　$0<a<b$ のとき、$\dfrac{b-a}b<\log b-\log a<\dfrac{b-a}a$ が成り立つことを証明せよ。\par
  \kakko{4}　\kakko{1}を用いて、$\displaystyle\lim_{x\to\infty}\frac x{e^x}=0$ を証明せよ。
\end{ascolorbox4A}

\begin{Data}
  テーマ：差の関数による不等式の証明、定数分離と極限、平均値の定理、不等式からはさみうちへ\par
  難易度：(1) B　(2) C　(3) C　(4) B\quad 前提：第19・30・33講\par
  この講のゴール：不等式は「差の関数の最小値 $\geqq0$」、方程式は「グラフの交点」に直す型を使いこなす
\end{Data}

\begin{multicols*}{2}
  \Shomon{1}
  \ascboxZ{差の関数を作り、符号がわかるまで微分する}
  $f(x)=e^x-1-x-\dfrac{x^2}2$ とおき、$x>0$ で $f(x)>0$ を示す。
  \[f'(x)=e^x-1-x,\quad f''(x)=e^x-1\]
  $x>0$ で $e^x>1$ なので $f''(x)>0$。よって $f'(x)$ は $x\geqq0$ で増加し、$f'(0)=0$ より $x>0$ で $f'(x)>0$。すると $f(x)$ も $x\geqq0$ で増加し、$f(0)=0$ より $x>0$ で $f(x)>0$。\QED

  \begin{ptbs}{KEY}[不等式は「差 $=f(x)$」の増減で]
    $f'(x)$ の符号が見えなければ、もう一度微分する。下から順に「$f''>0$ → $f'$ 増加 → $f'>0$ → $f$ 増加 → $f>0$」と積み上げる。
  \end{ptbs}

  \Shomon{2}
  \ascboxZ{$x=0$ を確認してから定数分離}
  $x=0$ は解でない（左辺1、右辺0）。よって $x\neq0$ として両辺を $x$ で割ってよく、
  \[a=\frac{e^x}x\]
  $g(x)=\dfrac{e^x}x$ とおくと、解の個数は曲線 $y=g(x)$ と直線 $y=a$ の共有点の個数である。
  \[g'(x)=\frac{e^x\cdot x-e^x}{x^2}=\frac{(x-1)e^x}{x^2}\]
  端の様子：$x\to-\infty$ で $g\to-0$、$x\to-0$ で $g\to-\infty$、$x\to+0$ で $g\to+\infty$、$x\to\infty$ で $g\to+\infty$（\kakko{4}より $\frac x{e^x}\to+0$ なので、その逆数）。

  \begin{zuhyo}
  \begin{tblr}{colspec={c|cccccc}, hlines, vlines, row{1}={font=\TRgt}, cells={font=\small}}
  $x$ & $-\infty$ & $\cdots$ & $(0)$ & $\cdots$ & $1$ & $\cdots$ \\
  $g'$ & & $-$ & & $-$ & $0$ & $+$ \\
  $g$ & $-0$ & $\searrow$ & $\pm\infty$ & $\searrow$ & $e$ & $\nearrow$ \\
  \end{tblr}
  \end{zuhyo}

  \begin{zuhyo}
  \begin{tikzpicture}[x=.85cm,y=.42cm,line width=.5pt]
    \draw[->] (-3.2,0) -- (2.9,0) node[right] {$x$};
    \draw[->] (0,-3.4) -- (0,5.8) node[above] {$y$};
    \draw[domain=-3.1:-.25,samples=60] plot (\x,{exp(\x)/(\x)});
    \draw[domain=.3:2.6,samples=60] plot (\x,{exp(\x)/(\x)});
    \draw[dashed] (1,0) -- (1,2.718) -- (0,2.718) node[left,font=\footnotesize] {$e$};
    \node[below,font=\footnotesize] at (1,0) {$1$};
    \node[below left,font=\footnotesize] at (0,0) {$\mathrm O$};
  \end{tikzpicture}
  \end{zuhyo}

  直線 $y=a$ を動かして共有点を数える。
  \[\Kotae{\begin{cases}a<0:\ 1\ \text{個}\\0\leqq a<e:\ 0\ \text{個}\\a=e:\ 1\ \text{個}\\e<a:\ 2\ \text{個}\end{cases}}\]

  \begin{memo}
    $a=0$ のとき、$e^x=0$ は解をもたない。グラフでも $y=0$（$x$ 軸）は曲線と交わらない（$x<0$ の部分は $x$ 軸に近づくだけで届かない）。境目の値は1つずつ確認する。
  \end{memo}

  \Shomon{3}
  \ascboxZ{平均値の定理で「差」を「微分」に変える}
  \begin{koushiki}{平均値の定理}
    $f(x)$ が $a\leqq x\leqq b$ で連続、$a<x<b$ で微分可能ならば、
    \[\frac{f(b)-f(a)}{b-a}=f'(c),\quad a<c<b\]
    を満たす $c$ が存在する。
  \end{koushiki}
  図形的には「2点を結ぶ直線と平行な接線が、間のどこかで引ける」という主張である。\par
  $f(x)=\log x$ に使うと、$\dfrac{\log b-\log a}{b-a}=\dfrac1c$（$a<c<b$）となる $c$ がある。$0<a<c<b$ より $\frac1b<\frac1c<\frac1a$ なので、
  \[\frac1b<\frac{\log b-\log a}{b-a}<\frac1a\]
  各辺に $b-a\ (>0)$ を掛けて示された。\QED

  \Shomon{4}
  \ascboxZ{示した不等式ではさむ}
  \kakko{1}より $x>0$ で $e^x>1+x+\frac{x^2}2>\frac{x^2}2$。よって、
  \[0<\frac x{e^x}<\frac x{\frac{x^2}2}=\frac2x\]
  $x\to\infty$ で $\frac2x\to0$ なので、はさみうちの原理より $\displaystyle\lim_{x\to\infty}\frac x{e^x}=0$。\QED

  \begin{memo}
    \kakko{1}の不等式は、\kakko{4}の極限を示すための「部品」になっている。東大の小問は、前の小問が後の小問の道具になるように作られていることが多い。前の結果をどう使うかを必ず考える。
  \end{memo}

  \begin{simple}[出典：オリジナル]{参考問題34}
    すべての正の実数 $x$ に対して $\log x\leqq ax$ が成り立つような定数 $a$ の値の範囲を求めよ。
  \end{simple}
  \Kai
  $x>0$ で割って定数を分離すると、条件は「すべての $x>0$ で $\frac{\log x}x\leqq a$」。これは「$\frac{\log x}x$ の最大値 $\leqq a$」と同値である。第33講より最大値は $\frac1e$（$x=e$）なので、
  \[\Kotae{a\geqq\frac1e}\]
  \begin{memo}
    「すべての $x$ で $f(x)\leqq a$」$\iff$「$f(x)$ の最大値 $\leqq a$」。ただし最大値が存在しない（上限に届かない）関数では、上限との比較になる。
  \end{memo}

  \begin{ptbs}{CHECK}[翌日、解説を隠して答える（答えは本講の中にある）]
    \kakko{1}　差の関数の符号が見えないときはどうするか。\par
    \kakko{2}　定数分離で割り算をする前に確かめることは何か。\par
    \kakko{3}　平均値の定理の図形的な意味を言え。
  \end{ptbs}

  \notefill
\end{multicols*}

%% ==================== 第35講 ====================
\Kou{35}{積分法（数学Ⅲ）の計算}
\begin{ascolorbox4A}[出典：オリジナル]{例題35}
  \kakko{1}　不定積分 $\displaystyle\int xe^x\,dx$、$\displaystyle\int\log x\,dx$ を求めよ。\par
  \kakko{2}　定積分 $\displaystyle\int_0^1x\sqrt{1-x^2}\,dx$、$\displaystyle\int_0^1\sqrt{1-x^2}\,dx$ を求めよ。\par
  \kakko{3}　$\displaystyle\int_0^{\frac\pi2}\sin^2x\,dx$ と、不定積分 $\displaystyle\int\frac{dx}{x(x+1)}$ を求めよ。\par
  \kakko{4}　$\displaystyle\int_0^1\frac{dx}{1+x^2}$ を求めよ。\par
  \kakko{5}　$I_n=\displaystyle\int_0^{\frac\pi2}\sin^nx\,dx$（$n=0,\ 1,\ 2,\ \cdots$）とする。$n\geqq2$ のとき $I_n=\dfrac{n-1}nI_{n-2}$ を示し、$I_4$ を求めよ。
\end{ascolorbox4A}

\begin{Data}
  テーマ：部分積分、置換積分（$\sin$・$\tan$ 置換を含む）、次数下げ、部分分数分解、積分の漸化式\par
  難易度：(1) A　(2) B　(3) A　(4) B　(5) C\quad 前提：第20・32講\par
  この講のゴール：被積分関数の形から「置換・部分積分・式変形」のどれを使うかを即決できる
\end{Data}

\begin{multicols*}{2}
  \begin{koushiki}{基本の不定積分（$C$ は積分定数）}
    $\displaystyle\int x^\alpha dx=\frac{x^{\alpha+1}}{\alpha+1}+C\ (\alpha\neq-1),\quad\int\frac{dx}x=\log|x|+C$\par
    $\displaystyle\int e^xdx=e^x+C,\quad\int\sin x\,dx=-\cos x+C$\par
    $\displaystyle\int\cos x\,dx=\sin x+C,\quad\int\frac{dx}{\cos^2x}=\tan x+C$\par
    $\displaystyle\int\frac{f'(x)}{f(x)}dx=\log|f(x)|+C$
  \end{koushiki}
  どれも、右辺を微分すると左辺の中身に戻ることで確かめられる。

  \Shomon{1}
  \ascboxZ{部分積分：積の微分の逆}
  \begin{koushiki}{部分積分}
    \[\int f(x)g'(x)\,dx=f(x)g(x)-\int f'(x)g(x)\,dx\]
  \end{koushiki}
  積の微分 $(fg)'=f'g+fg'$ を積分して移項したものである。微分すると簡単になるもの（$x$、$\log x$）を $f$ に、積分しやすいもの（$e^x$、$\sin x$）を $g'$ に選ぶ。
  \[\int xe^x\,dx=xe^x-\int1\cdot e^x\,dx=\Kotae{(x-1)e^x+C}\]
  $\log x$ は $\log x\cdot1$ とみて、$1$ を $g'$（$g=x$）にする。
  \[\int\log x\,dx=x\log x-\int\frac1x\cdot x\,dx\]
  \[=\Kotae{x\log x-x+C}\]

  \Shomon{2}
  \ascboxZ{置換積分：中身を1文字にする}
  \begin{koushiki}{置換積分}
    $x=g(t)$ とおくと $dx=g'(t)\,dt$。定積分では積分区間も $t$ の区間に変える。
  \end{koushiki}
  $\int_0^1x\sqrt{1-x^2}\,dx$ は、$u=1-x^2$ とおくと $du=-2x\,dx$ で、$x$ が $0\to1$ のとき $u$ は $1\to0$。
  \[\int_1^0\sqrt u\cdot\left(-\frac12\right)du=\frac12\int_0^1\sqrt u\,du=\frac12\cdot\frac23=\Kotae{\frac13}\]
  $\int_0^1\sqrt{1-x^2}\,dx$ は $x=\sin\theta$ とおく（$\sqrt{1-x^2}$ を見たら $\sin$ 置換）。$dx=\cos\theta\,d\theta$、$x$ が $0\to1$ のとき $\theta$ は $0\to\frac\pi2$。この範囲で $\cos\theta\geqq0$ なので $\sqrt{1-\sin^2\theta}=\cos\theta$。
  \[\int_0^{\frac\pi2}\cos^2\theta\,d\theta=\Kotae{\frac\pi4}\]
  （最後の計算は\kakko{3}と同じ。）図形的には、半径1の四分円の面積 $\frac\pi4$ である。

  \Shomon{3}
  \ascboxZ{2乗は半角の公式で次数を下げる}
  $\sin^2x=\frac{1-\cos2x}2$ を使う。
  \[\int_0^{\frac\pi2}\frac{1-\cos2x}2\,dx=\Big[\frac x2-\frac{\sin2x}4\Big]_0^{\frac\pi2}=\Kotae{\frac\pi4}\]
  分数式は、分母の因数ごとに分ける（\textgt{部分分数分解}）。$\frac1{x(x+1)}=\frac1x-\frac1{x+1}$ なので、
  \[\int\frac{dx}{x(x+1)}=\log|x|-\log|x+1|+C\]
  \[=\Kotae{\log\left|\frac x{x+1}\right|+C}\]

  \Shomon{4}
  \ascboxZ{$\dfrac1{1+x^2}$ は $\tan$ 置換}
  $x=\tan\theta$ とおくと $dx=\frac{d\theta}{\cos^2\theta}$、$1+\tan^2\theta=\frac1{\cos^2\theta}$。$x$ が $0\to1$ のとき $\theta$ は $0\to\frac\pi4$。
  \[\int_0^{\frac\pi4}\cos^2\theta\cdot\frac{d\theta}{\cos^2\theta}=\int_0^{\frac\pi4}d\theta=\Kotae{\frac\pi4}\]

  \begin{ptbs}{KEY}[形を見て置換を決める]
    $\sqrt{a^2-x^2}$ → $x=a\sin\theta$。\quad$\dfrac1{x^2+a^2}$ → $x=a\tan\theta$。\par
    $f(g(x))g'(x)$ の形 → $u=g(x)$。\quad 積 → 部分積分。
  \end{ptbs}

  \Shomon{5}
  \ascboxZ{部分積分で $I_n$ と $I_{n-2}$ をつなぐ}
  $\sin^nx=\sin^{n-1}x\cdot\sin x$ とみて、$\sin x=(-\cos x)'$ に部分積分を使う。
  \[
  \begin{aligned}
  I_n&=\Big[-\sin^{n-1}x\cos x\Big]_0^{\frac\pi2}\\
  &\quad+(n-1)\int_0^{\frac\pi2}\sin^{n-2}x\cos^2x\,dx
  \end{aligned}
  \]
  第1項は両端で0（$\cos\frac\pi2=0$、$\sin0=0$）。$\cos^2x=1-\sin^2x$ とすると、第2項は $(n-1)(I_{n-2}-I_n)$。よって $I_n=(n-1)I_{n-2}-(n-1)I_n$、すなわち $nI_n=(n-1)I_{n-2}$。\QED\par
  $I_0=\frac\pi2$ から、$I_2=\frac12\cdot\frac\pi2=\frac\pi4$、
  \[I_4=\frac34I_2=\Kotae{\frac{3\pi}{16}}\]

  \begin{memo}
    「$I_n$ を部分積分したら、また $I_n$ が出てきた」ときは、それを左辺に移項して解く。逃げずに同じ記号でまとめるのがコツである。
  \end{memo}

  \begin{simple}[出典：オリジナル]{参考問題35}
    $\displaystyle\int_0^\pi x\sin x\,dx$ と $\displaystyle\int_1^ex\log x\,dx$ を求めよ。
  \end{simple}
  \Kai
  1つ目は $x$ を微分する側にして部分積分する。
  \[\Big[-x\cos x\Big]_0^\pi+\int_0^\pi\cos x\,dx=\pi+0=\Kotae{\pi}\]
  2つ目は $\log x$ を微分する側、$x$ を積分する側（$\frac{x^2}2$）にする。
  \[\Big[\frac{x^2}2\log x\Big]_1^e-\int_1^e\frac x2\,dx=\frac{e^2}2-\frac{e^2-1}4\]
  整理して、
  \[\int_1^ex\log x\,dx=\Kotae{\frac{e^2+1}4}\]

  \begin{ptbs}{CHECK}[翌日、解説を隠して答える（答えは本講の中にある）]
    \kakko{1}　部分積分で、微分する側・積分する側の選び方を言え。\par
    \kakko{2}　$\sqrt{a^2-x^2}$ と $\frac1{x^2+a^2}$ はそれぞれ何で置換するか。\par
    \kakko{3}　$I_n$ を部分積分して再び $I_n$ が出たらどうするか。
  \end{ptbs}

  \notefill
\end{multicols*}
%% ==================== 第36講 ====================
\Kou{36}{区分求積法と積分の評価}
\begin{ascolorbox4A}[出典：オリジナル]{例題36}
  \kakko{1}　$\displaystyle\lim_{n\to\infty}\frac1n\sum_{k=1}^n\left(\frac kn\right)^2$ を求めよ。\qquad
  \kakko{2}　$\displaystyle\lim_{n\to\infty}\sum_{k=1}^n\frac1{n+k}$ を求めよ。\par
  \kakko{3}　自然数 $n$ に対して $\displaystyle\sum_{k=1}^n\frac1k>\log(n+1)$ が成り立つことを示し、$\displaystyle\lim_{n\to\infty}\sum_{k=1}^n\frac1k$ を求めよ。\par
  \kakko{4}　$\dfrac43<\displaystyle\int_0^1e^{x^2}\,dx<e-1$ を証明せよ。
\end{ascolorbox4A}

\begin{Data}
  テーマ：区分求積法（和の極限を積分に）、積分と和の大小比較、積分の評価（不等式ではさむ）\par
  難易度：(1) B　(2) B　(3) C　(4) C\quad 前提：第30・34・35講\par
  この講のゴール：「和 ⇄ 積分」を長方形の面積で行き来し、計算できない積分を不等式ではさめる
\end{Data}

\begin{multicols*}{2}
  \Shomon{1}
  \ascboxZ{長方形の面積の和 → 積分}
  区間 $[0,\ 1]$ を $n$ 等分し、幅 $\frac1n$、高さ $f\left(\frac kn\right)$ の長方形を並べる。$n$ を大きくすると、長方形の面積の和はグラフの下の面積に近づく。

  \begin{zuhyo}
  \begin{tikzpicture}[x=3.6cm,y=2.4cm,line width=.5pt]
    \foreach \k in {1,...,6}{\fill[black!12] ({(\k-1)/6},0) rectangle ({\k/6},{(\k/6)^2}); \draw ({(\k-1)/6},0) rectangle ({\k/6},{(\k/6)^2});}
    \draw[->] (-.05,0) -- (1.15,0) node[right] {$x$};
    \draw[->] (0,-.05) -- (0,1.15) node[above] {$y$};
    \draw[line width=.8pt,domain=0:1.05,samples=40] plot (\x,{(\x)^2});
    \node[below,font=\footnotesize] at (1,0) {$1$};
    \node[below,font=\footnotesize] at (1/6,0) {$\frac1n$};
    \node[font=\footnotesize] at (.62,.85) {$y=x^2$};
  \end{tikzpicture}
  \end{zuhyo}

  \begin{koushiki}{区分求積法}
    $f(x)$ が $0\leqq x\leqq1$ で連続のとき、
    \[\lim_{n\to\infty}\frac1n\sum_{k=1}^nf\left(\frac kn\right)=\int_0^1f(x)\,dx\]
  \end{koushiki}
  本問は $f(x)=x^2$ の場合なので、
  \[\lim_{n\to\infty}\frac1n\sum_{k=1}^n\left(\frac kn\right)^2=\int_0^1x^2\,dx=\Kotae{\frac13}\]

  \begin{memo}
    検算：$\frac1n\sum\frac{k^2}{n^2}=\frac{n(n+1)(2n+1)}{6n^3}\to\frac26=\frac13$。$\Sigma$ の公式で計算できる場合と一致する。区分求積法は、$\Sigma$ が計算できない場合に威力を発揮する。
  \end{memo}

  \Shomon{2}
  \ascboxZ{$\frac1n$ をくくり出し、$\frac kn$ の関数を探す}
  \[\sum_{k=1}^n\frac1{n+k}=\frac1n\sum_{k=1}^n\frac1{1+\frac kn}\]
  $f(x)=\frac1{1+x}$ の区分求積である。
  \[\lim_{n\to\infty}\sum_{k=1}^n\frac1{n+k}=\int_0^1\frac{dx}{1+x}\]
  \[=\Big[\log(1+x)\Big]_0^1=\Kotae{\log2}\]

  \begin{ptbs}{KEY}[区分求積の見つけ方]
    \kakko{1} $\frac1n$ をくくり出す。\kakko{2} 残りを $\frac kn$ だけの式にする。\kakko{3} $\frac kn$ を $x$ に、$\frac1n\sum$ を $\int_0^1dx$ に置き換える。
  \end{ptbs}

  \Shomon{3}
  \ascboxZ{積分と和を長方形で比べる}
  \begin{koushiki}{積分の大小}
    $a\leqq x\leqq b$ で $f(x)\leqq g(x)$ ならば $\displaystyle\int_a^bf(x)\,dx\leqq\int_a^bg(x)\,dx$。\par
    $f,\ g$ が連続で、つねに等しいわけではなければ、不等号 $<$ が成り立つ。
  \end{koushiki}
  $k\leqq x\leqq k+1$ のとき $\frac1x\leqq\frac1k$（等号は $x=k$ のときだけ）。よって、
  \[\int_k^{k+1}\frac{dx}x<\int_k^{k+1}\frac{dx}k=\frac1k\]
  これを $k=1,\ 2,\ \cdots,\ n$ について足すと、左辺は $\int_1^{n+1}\frac{dx}x=\log(n+1)$ になるので、
  \[\sum_{k=1}^n\frac1k>\log(n+1)\]
  $n\to\infty$ で $\log(n+1)\to\infty$ なので、
  \[\lim_{n\to\infty}\sum_{k=1}^n\frac1k=\Kotae{\infty}\]

  \begin{memo}
    各項 $\frac1k$ は0に近づくのに、和は無限大に発散する。「項が0に近づけば和は収束する」は誤りである（逆の「和が収束すれば項は0に近づく」は正しい）。
  \end{memo}

  \Shomon{4}
  \ascboxZ{計算できない積分は、計算できる関数ではさむ}
  $e^{x^2}$ の原始関数は、知っている関数では書けない。そこで、$0\leqq x\leqq1$ で $e^{x^2}$ を上下から押さえる。\par
  \textgt{下から}：すべての実数 $t$ で $e^t\geqq1+t$（等号は $t=0$ のときだけ）。$h(t)=e^t-1-t$ とおくと $h'(t)=e^t-1$ で、$t=0$ で最小値0をとるからである。$t=x^2$ として $e^{x^2}\geqq1+x^2$（等号は $x=0$ だけ）。
  \[\int_0^1e^{x^2}dx>\int_0^1(1+x^2)\,dx=\frac43\]
  \textgt{上から}：$0\leqq x\leqq1$ で $x^2\leqq x$（等号は $x=0,\ 1$ だけ）、$e^t$ は増加関数なので $e^{x^2}\leqq e^x$。
  \[\int_0^1e^{x^2}dx<\int_0^1e^x\,dx=e-1\]
  以上で示された。\QED

  \begin{memo}
    $\frac43\fallingdotseq1.333$、$e-1\fallingdotseq1.718$ で、実際の値は約1.463。はさむ関数の選び方で精度が変わる。東大では「この不等式を示せ」の形で、はさむ関数を自分で選ばせることが多い（第42講）。
  \end{memo}

  \begin{simple}[出典：オリジナル]{参考問題36}
    $\displaystyle\lim_{n\to\infty}\frac1n\log\frac{(n+1)(n+2)\cdots(2n)}{n^n}$ を求めよ。
  \end{simple}
  \Kai
  積は対数で和に変わる。$\frac{(n+1)(n+2)\cdots(2n)}{n^n}=\prod_{k=1}^n\frac{n+k}n$ なので、
  \[\frac1n\log\frac{(n+1)\cdots(2n)}{n^n}=\frac1n\sum_{k=1}^n\log\left(1+\frac kn\right)\]
  区分求積法により、
  \[\int_0^1\log(1+x)\,dx=\Big[(1+x)\log(1+x)-x\Big]_0^1\]
  \[=\Kotae{2\log2-1}\]

  \begin{ptbs}{CHECK}[翌日、解説を隠して答える（答えは本講の中にある）]
    \kakko{1}　区分求積法を見つける3手順を言え。\par
    \kakko{2}　$\sum\frac1k>\log(n+1)$ を示す考え方を言え。\par
    \kakko{3}　$\int_0^1e^{x^2}dx$ をはさむ2つの関数は何か。
  \end{ptbs}

  \notefill
\end{multicols*}

%% ==================== 第37講 ====================
\Kou{37}{面積と体積}
\begin{ascolorbox4A}[出典：オリジナル]{例題37}
  \kakko{1}　曲線 $y=\sin x\ (0\leqq x\leqq\pi)$ と $x$ 軸で囲まれた部分の面積 $S$ と、その部分を $x$ 軸のまわりに1回転してできる立体の体積 $V$ を求めよ。\par
  \kakko{2}　楕円 $\dfrac{x^2}{a^2}+\dfrac{y^2}{b^2}=1\ (a>0,\ b>0)$ で囲まれた部分の面積を求めよ。\par
  \kakko{3}　曲線 $y=\log x$、$x$ 軸、直線 $x=e$ で囲まれた部分を、$y$ 軸のまわりに1回転してできる立体の体積を求めよ。
\end{ascolorbox4A}

\begin{Data}
  テーマ：面積、体積＝断面積の積分、$x$ 軸・$y$ 軸のまわりの回転体\par
  難易度：(1) A　(2) B　(3) C\quad 前提：第20・35講\par
  この講のゴール：「どの方向に切ると断面が簡単か」を決め、断面積を積分して体積を出せる
\end{Data}

\begin{multicols*}{2}
  \begin{koushiki}{体積の基本原理}
    立体を $x$ 軸に垂直な平面で切った断面積を $S(x)$ とすると、$a\leqq x\leqq b$ の部分の体積は
    \[V=\int_a^bS(x)\,dx\]
    $x$ 軸のまわりの回転体では、断面は半径 $|f(x)|$ の円なので $S(x)=\pi\{f(x)\}^2$。
  \end{koushiki}
  \textgt{理由}：立体を厚さ $\Delta x$ の薄い板に切ると、1枚の体積はほぼ $S(x)\Delta x$。これを足し合わせて厚さを0に近づけたものが積分である（第36講の区分求積と同じ考え方）。

  \Shomon{1}
  \ascboxZ{面積と $x$ 軸まわりの回転体}
  \[S=\int_0^\pi\sin x\,dx=\Big[-\cos x\Big]_0^\pi=\Kotae{2}\]
  \[V=\pi\int_0^\pi\sin^2x\,dx=\pi\int_0^\pi\frac{1-\cos2x}2\,dx\]
  $\int_0^\pi\cos2x\,dx=0$ なので、
  \[V=\pi\cdot\frac\pi2=\Kotae{\frac{\pi^2}2}\]

  \Shomon{2}
  \ascboxZ{対称性で4分の1を計算する}
  楕円は $x$ 軸・$y$ 軸について対称なので、第1象限の部分の4倍。第1象限では $y=b\sqrt{1-\frac{x^2}{a^2}}=\frac ba\sqrt{a^2-x^2}$。
  \[\text{面積}=4\cdot\frac ba\int_0^a\sqrt{a^2-x^2}\,dx\]
  $\int_0^a\sqrt{a^2-x^2}\,dx$ は半径 $a$ の四分円の面積 $\frac{\pi a^2}4$（第35講\kakko{2}と同じく $x=a\sin\theta$ で計算できる）。よって、
  \[\text{面積}=4\cdot\frac ba\cdot\frac{\pi a^2}4=\Kotae{\pi ab}\]

  \begin{memo}
    楕円は円 $x^2+y^2=a^2$ を $y$ 方向に $\frac ba$ 倍したもので、面積も $\frac ba$ 倍になる。$\pi a^2\cdot\frac ba=\pi ab$。図形的な見方で答えを予想しておくと、計算ミスに気づける。
  \end{memo}

  \Shomon{3}
  \ascboxZ{$y$ 軸まわりは $y$ で切る}
  \begin{zuhyo}
  \begin{tikzpicture}[x=1.2cm,y=1.4cm,line width=.5pt]
    \fill[black!15] (1,0) -- plot[domain=1:2.718,samples=40] (\x,{ln(\x)}) -- (2.718,0) -- cycle;
    \draw[->] (-.3,0) -- (3.3,0) node[right] {$x$};
    \draw[->] (0,-.3) -- (0,1.4) node[above] {$y$};
    \draw[domain=.6:3.1,samples=40] plot (\x,{ln(\x)});
    \draw (2.718,0) -- (2.718,1.2);
    \draw[dashed] (0,.55) -- (2.718,.55);
    \fill (1.733,.55) circle (1.3pt);
    \node[below,font=\footnotesize] at (1,0) {$1$};
    \node[below,font=\footnotesize] at (2.718,0) {$e$};
    \node[left,font=\footnotesize] at (0,1) {$1$};
    \node[left,font=\footnotesize] at (0,.55) {$y$};
    \node[font=\footnotesize] at (1.2,.85) {$x=e^y$};
    \node[below left,font=\footnotesize] at (0,0) {$\mathrm O$};
  \end{tikzpicture}
  \end{zuhyo}
  灰色の部分を $y$ 軸のまわりに回す。高さ $y$（$0\leqq y\leqq1$）で水平に切ると、灰色の部分は $x=e^y$ から $x=e$ までの線分。これを $y$ 軸のまわりに回すと、外半径 $e$、内半径 $e^y$ の\textgt{円環}（ドーナツ形の板）になる。断面積は $\pi(e^2-e^{2y})$。
  \[
  \begin{aligned}
  V&=\pi\int_0^1(e^2-e^{2y})\,dy=\pi\Big[e^2y-\frac{e^{2y}}2\Big]_0^1\\
  &=\pi\left(e^2-\frac{e^2}2+\frac12\right)
  \end{aligned}
  \]
  よって、
  \[V=\Kotae{\frac{\pi(e^2+1)}2}\]

  \begin{ptbs}{KEY}[回転軸に垂直に切る]
    $y$ 軸まわりなら $y$ で切り、断面（円または円環）の面積を $y$ の式で表す。曲線の式を $x=(y\ \text{の式})$ に直すのが第一歩。
  \end{ptbs}

  \begin{memo}
    検算として「バウムクーヘン型」の公式 $V=2\pi\int_1^ex\log x\,dx$ を使うと、第35講の参考問題より $2\pi\cdot\frac{e^2+1}4$ で一致する（この公式は教科書外なので、答案では断面積による方法を書くのが安全）。
  \end{memo}

  \begin{simple}[出典：オリジナル]{参考問題37}
    放物線 $y=x^2$ と直線 $y=1$ で囲まれた部分を、直線 $y=1$ のまわりに1回転してできる立体の体積を求めよ。
  \end{simple}
  \Kai
  回転軸 $y=1$ に垂直な直線 $x=t$（$-1\leqq t\leqq1$）で切ると、断面は半径 $1-t^2$（軸から曲線までの距離）の円。
  \[
  \begin{aligned}
  V&=\pi\int_{-1}^1(1-x^2)^2\,dx=2\pi\int_0^1(1-2x^2+x^4)\,dx\\
  &=2\pi\left(1-\frac23+\frac15\right)
  \end{aligned}
  \]
  よって、
  \[V=\Kotae{\frac{16\pi}{15}}\]

  \begin{ptbs}{CHECK}[翌日、解説を隠して答える（答えは本講の中にある）]
    \kakko{1}　体積の基本原理（断面積の積分）を言え。\par
    \kakko{2}　$y$ 軸のまわりの回転体はどの方向に切るか。\par
    \kakko{3}　楕円の面積を図形的に予想する方法を言え。
  \end{ptbs}

  \notefill
  \columnbreak
  \notefill
\end{multicols*}

%% ==================== 第38講 ====================
\Kou{38}{体積の応用・曲線の長さ}
\begin{ascolorbox4A}[出典：オリジナル]{例題38}
  \kakko{1}　サイクロイド $x=t-\sin t,\ y=1-\cos t\ (0\leqq t\leqq2\pi)$ の長さを求めよ。\par
  \kakko{2}　曲線 $y=\dfrac{e^x+e^{-x}}2\ (0\leqq x\leqq1)$ の長さを求めよ。\par
  \kakko{3}　座標空間において、2つの円柱の内部 $x^2+y^2\leqq1$、$y^2+z^2\leqq1$ の共通部分の体積を求めよ。\par
  \kakko{4}　円 $x^2+(y-2)^2\leqq1$ を $x$ 軸のまわりに1回転してできる立体の体積を求めよ。
\end{ascolorbox4A}

\begin{Data}
  テーマ：媒介変数表示の曲線の長さ、$y=f(x)$ の長さ、回転体でない立体の体積、円環の回転体\par
  難易度：(1) B　(2) B　(3) C　(4) C\quad 前提：第29・37講\par
  この講のゴール：初めて見る立体でも「切り口が簡単になる方向」を探し、断面積を積分して体積を出せる
\end{Data}

\begin{multicols*}{2}
  \begin{koushiki}{曲線の長さ}
    媒介変数表示 $x=x(t),\ y=y(t)\ (a\leqq t\leqq b)$ の曲線の長さは
    \[L=\int_a^b\sqrt{\left(\frac{dx}{dt}\right)^2+\left(\frac{dy}{dt}\right)^2}\,dt\]
    $y=f(x)\ (a\leqq x\leqq b)$ なら $L=\displaystyle\int_a^b\sqrt{1+\{f'(x)\}^2}\,dx$。
  \end{koushiki}
  \textgt{理由}：$t$ が少し $\Delta t$ 変わると、点は横に約 $\frac{dx}{dt}\Delta t$、縦に約 $\frac{dy}{dt}\Delta t$ 動く。その斜めの距離は三平方の定理で $\sqrt{(\ )^2+(\ )^2}\,\Delta t$。これを足し合わせたものである。

  \Shomon{1}
  \ascboxZ{半角の公式で根号を外す}
  \begin{zuhyo}
  \begin{tikzpicture}[x=.85cm,y=.85cm,line width=.5pt]
    \draw[->] (-.3,0) -- (6.8,0) node[right] {$x$};
    \draw[->] (0,-.2) -- (0,2.5) node[above] {$y$};
    \draw[line width=.8pt,domain=0:6.2832,samples=80,variable=\t] plot ({\t-sin(\t r)},{1-cos(\t r)});
    \draw[dashed] (3.1416,1) circle (1);
    \node[below,font=\footnotesize] at (6.2832,0) {$2\pi$};
    \node[below,font=\footnotesize] at (3.1416,0) {$\pi$};
    \node[left,font=\footnotesize] at (0,2) {$2$};
  \end{tikzpicture}
  \end{zuhyo}
  サイクロイドは、半径1の円が $x$ 軸上を転がるとき、円周上の1点が描く曲線である。
  \[
  \begin{aligned}
  &\left(\frac{dx}{dt}\right)^2+\left(\frac{dy}{dt}\right)^2=(1-\cos t)^2+\sin^2t\\
  &=2-2\cos t=4\sin^2\frac t2
  \end{aligned}
  \]
  $0\leqq t\leqq2\pi$ では $\sin\frac t2\geqq0$ なので、根号を外すと $2\sin\frac t2$。
  \[L=\int_0^{2\pi}2\sin\frac t2\,dt=\Big[-4\cos\frac t2\Big]_0^{2\pi}=\Kotae{8}\]

  \begin{memo}
    $\sqrt{4\sin^2\frac t2}=2\left|\sin\frac t2\right|$ である。範囲によっては絶対値の中が負になるので、根号を外す前に符号を確認する（第2講の $\sqrt{A^2}=|A|$）。
  \end{memo}

  \Shomon{2}
  \ascboxZ{$1+f'^2$ が完全平方になる}
  $f'(x)=\frac{e^x-e^{-x}}2$。$e^x\cdot e^{-x}=1$ を使うと、
  \[1+\{f'(x)\}^2=\frac{4+e^{2x}-2+e^{-2x}}4=\left(\frac{e^x+e^{-x}}2\right)^2\]
  $\frac{e^x+e^{-x}}2>0$ なので、
  \[L=\int_0^1\frac{e^x+e^{-x}}2\,dx=\Big[\frac{e^x-e^{-x}}2\Big]_0^1\]
  \[=\Kotae{\frac{e-e^{-1}}2}\]

  \Shomon{3}
  \ascboxZ{回転体でなければ、断面が簡単な方向を探す}
  $x$ で切ると、断面は「$y^2\leqq1-x^2$ かつ $y^2+z^2\leqq1$」という扱いにくい形になる。そこで、\textgt{2つの条件の両方に出てくる文字 $y$} を固定する。\par
  $y=t$（$-1\leqq t\leqq1$）で切ると、断面は
  \[x^2\leqq1-t^2,\quad z^2\leqq1-t^2\]
  つまり $|x|\leqq\sqrt{1-t^2}$、$|z|\leqq\sqrt{1-t^2}$ の正方形で、1辺は $2\sqrt{1-t^2}$。断面積は $4(1-t^2)$。
  \[V=\int_{-1}^14(1-t^2)\,dt=8\int_0^1(1-t^2)\,dt=\Kotae{\frac{16}3}\]

  \begin{ptbs}{KEY}[共通の文字で切ると、条件が独立になる]
    2つの条件に共通して現れる文字を固定すると、残りの文字の条件がばらばらになり、断面が長方形や円などの簡単な形になる。東大の立体の問題で最も使う発想である。
  \end{ptbs}

  \Shomon{4}
  \ascboxZ{回転すると円環になる部分は「外 $-$ 内」}
  $x$ を固定すると（$-1\leqq x\leqq1$）、円の内部の $y$ の範囲は $2-\sqrt{1-x^2}\leqq y\leqq2+\sqrt{1-x^2}$。これを $x$ 軸のまわりに回すと、外半径 $2+\sqrt{1-x^2}$、内半径 $2-\sqrt{1-x^2}$ の円環で、断面積は
  \[\pi\{(2+\sqrt{1-x^2})^2-(2-\sqrt{1-x^2})^2\}=8\pi\sqrt{1-x^2}\]
  $\int_{-1}^1\sqrt{1-x^2}\,dx$ は半径1の半円の面積 $\frac\pi2$ なので、
  \[V=8\pi\cdot\frac\pi2=\Kotae{4\pi^2}\]

  \begin{memo}
    この立体はドーナツ形（トーラス）である。断面の円の面積 $\pi$ に、円の中心が1周する距離 $2\pi\cdot2$ を掛けても $4\pi^2$ になる（パップス・ギュルダンの定理）。答案には使わず、検算に使う。
  \end{memo}

  \begin{simple}[出典：オリジナル]{参考問題38}
    底面が半径1の円 $x^2+y^2\leqq1$（$z=0$）で、$x$ 軸に垂直な平面による切り口がすべて正三角形となる立体（正三角形は $z\geqq0$ の側にある）の体積を求めよ。
  \end{simple}
  \Kai
  $x=t$ で切ると、底面の円の切り口は長さ $2\sqrt{1-t^2}$ の線分で、これを1辺とする正三角形の面積は $\frac{\sqrt3}4\cdot4(1-t^2)=\sqrt3(1-t^2)$。
  \[V=\int_{-1}^1\sqrt3(1-t^2)\,dt=\sqrt3\cdot\frac43=\Kotae{\frac{4\sqrt3}3}\]

  \begin{ptbs}{CHECK}[翌日、解説を隠して答える（答えは本講の中にある）]
    \kakko{1}　曲線の長さの公式が成り立つ理由を言え。\par
    \kakko{2}　根号 $\sqrt{A^2}$ を外す前に何を確かめるか。\par
    \kakko{3}　2つの円柱の共通部分で、$y$ で切る理由は何か。
  \end{ptbs}

  \notefill
  \columnbreak
  \notefill
\end{multicols*}
%% ==================== 第39講 ====================
\Kou{39}{東大答案の書き方}
\begin{ascolorbox4A}[出典：オリジナル]{例題39}
  \kakko{1}　正の実数 $x,\ y$ が $x+y=2$ を満たすとき、積 $xy$ の取り得る値の範囲を求めよ。\par
  \kakko{2}　実数 $x,\ y$ が $x^2+y^2=1$ を満たしながら動く。$u=x+y$、$v=xy$ とおくとき、点 $(u,\ v)$ の動く範囲を図示せよ。\par
  \kakko{3}　$x$ の方程式 $x^2-2ax+a=0$ が、$0\leqq x\leqq1$ の範囲に少なくとも1つの実数解をもつような実数 $a$ の値の範囲を求めよ。
\end{ascolorbox4A}

\begin{Data}
  テーマ：必要条件と十分条件を分けて書く答案、実数条件、定数分離と除外点\par
  難易度：(1) B　(2) C　(3) C\quad 前提：第5・6・12・19講（第34講があればなおよい）\par
  この講のゴール：答えが合っていても減点される原因（十分性の欠落・範囲の書き忘れ）を、自分の答案から消せる
\end{Data}

\begin{multicols*}{2}
  \begin{koushiki}{東大答案の5つの型}
    \kakko{1} \textgt{文字の宣言}：おいた文字の意味と範囲を書く。\par
    \kakko{2} \textgt{方針の宣言}：「〜の存在条件を求める」など、何をするかを1行で書く。\par
    \kakko{3} \textgt{論理の矢印}：同値（$\iff$）か、一方向（$\Rightarrow$）かを区別する。\par
    \kakko{4} \textgt{場合分けの宣言と回収}：分けた理由を書き、最後に全部を合わせる。\par
    \kakko{5} \textgt{結論の明示}：問われた形で答えを書く。
  \end{koushiki}

  \Shomon{1}
  \ascboxZ{候補を絞る方向と、実現する方向}
  \textgt{よくある不完全な答案}：「$(x-y)^2\geqq0$ より $xy\leqq1$。よって $0\leqq xy\leqq1$。」\par
  これには2つの不足がある。$x,\ y>0$ なので $xy=0$ は起こらない。また、不等式を導いただけでは、その範囲の値が\textgt{全部実際にとれる}ことは示されていない。

  \textgt{必要条件}：$x,\ y>0$ より $xy>0$。$(x-y)^2\geqq0$ と $x+y=2$ より $4-4xy\geqq0$、$xy\leqq1$。よって $xy$ は $0<xy\leqq1$ の中にある。

  \textgt{十分条件}：$0<t\leqq1$ を満たす任意の $t$ に対し、
  \[x=1+\sqrt{1-t},\quad y=1-\sqrt{1-t}\]
  とおく。$0\leqq\sqrt{1-t}<1$ なので $x,\ y>0$。$x+y=2$、$xy=1-(1-t)=t$。よって $t$ は実際に $xy$ の値としてとれる。
  \[\Kotae{0<xy\leqq1}\]

  \begin{ptbs}{KEY}[「範囲を求めよ」は2方向で完成する]
    条件から候補を絞る（必要条件）だけでなく、候補の各値を実現する例を作る（十分条件）。または、最初から同値変形だけでつなぐ。
  \end{ptbs}

  \Shomon{2}
  \ascboxZ{置き換えたら「実数条件」を忘れない}
  $x^2+y^2=(x+y)^2-2xy$ なので、条件は $u^2-2v=1$、すなわち $v=\frac{u^2-1}2$。\par
  しかし、放物線 $v=\frac{u^2-1}2$ 上のすべての点が実現されるわけではない。$u,\ v$ が与えられたとき、$x,\ y$ は $t$ の2次方程式
  \[t^2-ut+v=0\]
  の2解であり（解と係数の関係）、$x,\ y$ が\textgt{実数}であるためには判別式 $u^2-4v\geqq0$ が必要十分である。$v=\frac{u^2-1}2$ を代入すると $u^2-2u^2+2\geqq0$、$u^2\leqq2$。\par
  逆に、$-\sqrt2\leqq u\leqq\sqrt2$、$v=\frac{u^2-1}2$ のとき、上の2次方程式は実数解 $x,\ y$ をもち、$x^2+y^2=u^2-2v=1$ を満たす。よって求める範囲は、
  \[\Kotae{\text{放物線}\ v=\frac{u^2-1}2\ \text{の}\ -\sqrt2\leqq u\leqq\sqrt2\ \text{の部分}}\]

  \begin{zuhyo}
  \begin{tikzpicture}[x=1.1cm,y=1.1cm,line width=.5pt]
    \draw[->] (-1.9,0) -- (1.9,0) node[right] {$u$};
    \draw[->] (0,-.75) -- (0,.9) node[above] {$v$};
    \draw[dashed,domain=-1.8:1.8,samples=40] plot (\x,{((\x)^2-1)/2});
    \draw[line width=1.2pt,domain=-1.4142:1.4142,samples=40] plot (\x,{((\x)^2-1)/2});
    \fill (1.4142,.5) circle (1.5pt); \fill (-1.4142,.5) circle (1.5pt);
    \node[right,font=\footnotesize] at (1.4142,.45) {$\left(\sqrt2,\frac12\right)$};
    \node[left,font=\footnotesize] at (-1.4142,.45) {$\left(-\sqrt2,\frac12\right)$};
    \node[below right,font=\footnotesize] at (0,-.5) {$-\frac12$};
    \node[below left,font=\footnotesize] at (0,0) {$\mathrm O$};
  \end{tikzpicture}
  \end{zuhyo}
  太線部分（両端を含む）。

  \begin{memo}
    $x=\cos\theta,\ y=\sin\theta$ とおけば、$u=\sqrt2\sin\left(\theta+\frac\pi4\right)$ なので $|u|\leqq\sqrt2$ が直接わかる。2通りで同じ範囲が出ることを確認すると、答案に自信がもてる。
  \end{memo}

  \Shomon{3}
  \ascboxZ{定数分離は「割れるか」を先に確かめる}
  方程式を $x^2=a(2x-1)$ と変形する。$a$ を分離するには $2x-1$ で割りたいが、$x=\frac12$ では割れない。$x=\frac12$ を入れると $\frac14=0$ となり成り立たないので、$x=\frac12$ は解にならない。よって $x\neq\frac12$ として、
  \[a=\frac{x^2}{2x-1}\quad\left(0\leqq x\leqq1,\ x\neq\frac12\right)\]
  求める条件は「この範囲の $x$ で、$h(x)=\frac{x^2}{2x-1}$ が値 $a$ をとる」ことである。
  \[h'(x)=\frac{2x(2x-1)-2x^2}{(2x-1)^2}=\frac{2x(x-1)}{(2x-1)^2}\leqq0\]
  なので、$h$ は $0\leqq x<\frac12$ でも $\frac12<x\leqq1$ でも減少する。\par
  $0\leqq x<\frac12$：$h(0)=0$ から減少し、$x\to\frac12-0$ で $-\infty$。値の範囲は $a\leqq0$。\par
  $\frac12<x\leqq1$：$x\to\frac12+0$ で $+\infty$ から減少し、$h(1)=1$。値の範囲は $a\geqq1$。
  \[\Kotae{a\leqq0,\ 1\leqq a}\]

  \begin{memo}
    検算：$a=\frac12$ では $x^2-x+\frac12=0$ の判別式が負で、確かに解がない。$a=0$ では $x=0$、$a=1$ では $x=1$ が解になり、境目が含まれることも確かめられる。
  \end{memo}

  \begin{ptbs}{注意}[減点されやすい書き方]
    「$\Rightarrow$」しか言えていないのに「$\iff$」と書く。おいた文字の範囲を書かない。0で割る可能性を確かめない。場合分けの最後に「以上より」で合わせない。
  \end{ptbs}

  \begin{simple}[出典：オリジナル]{参考問題39}
    正の実数 $x,\ y$ が $x+y=1$ を満たすとき、$x^2+y^2$ の取り得る値の範囲を求めよ。
  \end{simple}
  \Kai
  $x^2+y^2=(x+y)^2-2xy=1-2xy$。$t=xy$ の範囲を\kakko{1}と同じ方法で求める。\par
  必要：$t>0$、$(x-y)^2\geqq0$ より $1-4t\geqq0$、$t\leqq\frac14$。\par
  十分：$0<t\leqq\frac14$ なら、$s^2-s+t=0$ の2解 $\frac{1\pm\sqrt{1-4t}}2$ はともに正の実数で、和1、積 $t$。\par
  よって $0<t\leqq\frac14$ で、$1-2t$ の範囲は、
  \[\Kotae{\frac12\leqq x^2+y^2<1}\]

  \begin{ptbs}{CHECK}[翌日、解説を隠して答える（答えは本講の中にある）]
    \kakko{1}　「範囲を求めよ」で必要な2つの方向は何か。\par
    \kakko{2}　$u=x+y,\ v=xy$ とおいたときに忘れやすい条件は何か。\par
    \kakko{3}　東大答案の5つの型を言え。
  \end{ptbs}

  \notefill
\end{multicols*}

%% ==================== 第40講 ====================
\Kou{40}{東大型総合(1) 空間と通過範囲}
\begin{ascolorbox4A}[出典：東京大学 2026年度（条件を要約・設問を再構成）]{例題40}
  座標空間において、原点 O を中心とする半径5の球面上に、相異なる3点 P、Q、R がある。P と Q は $xy$ 平面上にあり、三角形 PQR の重心は $\mathrm G(2,\ 0,\ 1)$ である。\par
  \kakko{1}　線分 PQ の中点 M の軌跡を求めよ。\par
  \kakko{2}　線分 PQ が通過する範囲を求め、$xy$ 平面上に図示せよ。
\end{ascolorbox4A}

\begin{Data}
  テーマ：空間の条件から平面の条件へ、軌跡の必要十分、通過範囲＝「中点が存在する条件」\par
  難易度：(1) C　(2) D\quad 前提：第14・15・26・39講\par
  この講のゴール：空間図形を「固定できる量」から順に決め、最後は平面上の円と直線の交点の存在条件に帰着できる
\end{Data}

\begin{multicols*}{2}
  \begin{memo}
    この例題は、東京大学が公表した2026年度数学（理科）第3問の条件を要約し、設問を本書の形に再構成したものである。問題文の正確な表現は大学の公表資料で確認すること。以下の解説・図は本書が独自に作成したもので、大学の公式解答ではない。
  \end{memo}

  \Shomon{1}
  \ascboxZ{重心の $z$ 座標から、R の高さが決まる}
  \textgt{固定できるものから決める}。P、Q の $z$ 座標は0。R の $z$ 座標を $r$ とすると、重心の $z$ 座標は3頂点の平均なので $\frac{0+0+r}3=1$、$r=3$。\par
  R は球面上にあるので、R$(a,\ b,\ 3)$ とおくと $a^2+b^2+9=25$、$a^2+b^2=16$。よって R は高さ3の平面上の半径4の円を動き、
  \[\mathrm R(4\cos\theta,\ 4\sin\theta,\ 3)\]
  と表せる。

  \ascboxZ{中点 M を R で表す}
  重心の定義より $\vv{OP}+\vv{OQ}+\vv{OR}=3\vv{OG}$。M は PQ の中点なので、
  \[\vv{OM}=\frac{\vv{OP}+\vv{OQ}}2=\frac{3\vv{OG}-\vv{OR}}2\]
  成分で計算すると、
  \[\mathrm M(3-2\cos\theta,\ -2\sin\theta,\ 0)\]
  M$(u,\ v,\ 0)$ と書くと $(u-3)^2+v^2=4$。ここまでで言えたのは「M ならばこの円上にある」（必要条件）だけである。

  \ascboxZ{P$\neq$Q による除外点}
  P、Q は $xy$ 平面上の円 $x^2+y^2=25$ 上の異なる2点なので、弦 PQ の中点 M はこの円の内部にある：$u^2+v^2<25$。円 $(u-3)^2+v^2=4$ 上で原点からの距離が5になるのは $(5,\ 0)$ だけ（それ以外の点は距離5未満）なので、除くのは $(5,\ 0)$ のみ。

  \ascboxZ{逆に、残りの点は実現できるか（十分性）}
  円 $(u-3)^2+v^2=4$ 上の $(5,\ 0)$ 以外の点 M をとる。M は原点ではなく（中心 $(3,0)$ から原点までの距離は3で、半径2ではない）、半径5の円の内部にある。M を通り OM に垂直な直線は、円 $x^2+y^2=25$ と異なる2点で交わる。それを P、Q とすれば中点は M。さらに $\vv{OR}=3\vv{OG}-2\vv{OM}=(6-2u,\ -2v,\ 3)$ とおくと、
  \[(6-2u)^2+4v^2+9=4\{(u-3)^2+v^2\}+9=25\]
  で R は球面上にあり、$z$ 座標が3なので P、Q と異なる。よって条件を満たす3点が作れる。
  よって M の軌跡は、$xy$ 平面上の
  \[\Kotae{\text{円}\ (x-3)^2+y^2=4\ \text{（点}\ (5,0)\ \text{を除く）}}\]

  \begin{ptbs}{KEY}[空間図形は「固定できる量」から順に決める]
    重心の $z$ 座標 → R の高さ → R の円 → M の円。変数を増やさず、決まるものから決める。P、Q を別々に文字で置かないのがコツ。
  \end{ptbs}

  \Shomon{2}
  \ascboxZ{点 T を通る弦が「存在する条件」}
  $xy$ 平面上の点 T$(x,\ y)$ を固定する。T が線分 PQ 上にあるのは、
  \par\kakko{ア} T が直線 PQ 上にあり、かつ \kakko{イ} T が円板 $x^2+y^2\leqq25$ 内にある\par
  ときである（直線と閉じた円板の共通部分が、ちょうど弦 PQ だから）。\par
  直線 PQ は M を通り OM に垂直なので、\kakko{ア}は内積で
  \[(x-u)u+(y-v)v=0\iff ux+vy=u^2+v^2\]
  M の円の式から $u^2+v^2=6u-5$ なので、
  \[(x-6)u+yv+5=0\quad\cdots\text{①}\]
  T を固定すると、①は $uv$ 平面の\textgt{直線}である。よって、
  \par「T が通過範囲に入る」$\iff$「\kakko{イ} かつ、直線①と円 $(u-3)^2+v^2=4$ が $(5,0)$ 以外の点を共有する」\par
  となる（第15講の逆像法）。

  \ascboxZ{中心と直線の距離 $\leqq$ 半径}
  中心 $(3,\ 0)$ と直線①の距離が2以下：
  \[\frac{|3(x-6)+5|}{\sqrt{(x-6)^2+y^2}}\leqq2\]
  （T$=(6,0)$ では分母が0になるが、この点は円板の外なので考えなくてよい。）両辺とも0以上なので、2乗しても同値である。
  \[(3x-13)^2\leqq4\{(x-6)^2+y^2\}\]
  展開して整理すると $5x^2-30x+25-4y^2\leqq0$、すなわち
  \[\frac{(x-3)^2}4-\frac{y^2}5\leqq1\quad\cdots\text{②}\]

  \ascboxZ{除外点の処理}
  直線①と円の共有点が M$=(5,0)$ \textgt{だけ}になる T を除く必要がある。M$=(5,0)$ を①に入れると $5x-30+5=0$、$x=5$。円板内でこれを満たすのは T$=(5,0)$ だけで、このとき①は $u=5$ となり、円と $(5,0)$ でのみ接する。よって除く点は $(5,\ 0)$ のみである。それ以外の、\kakko{イ}と②を満たす点では、利用できる M があり、その M に対する弦が T を通る（十分性）。

  \ascboxZ{図示のための交点}
  円 $x^2+y^2=25$ と双曲線②の境界を連立すると、$3x^2-10x-25=0$、$x=5,\ -\frac53$。交点は $(5,\ 0)$、$\left(-\frac53,\ \pm\frac{10\sqrt2}3\right)$。

  \begin{zuhyo}
  \begin{tikzpicture}[x=.42cm,y=.42cm,line width=.5pt]
    \fill[black!15] (-1.667,4.714) arc[start angle=109.47,end angle=-109.47,radius=5] -- plot[domain=-4.714:4.714,samples=60] ({3-2*sqrt(1+(\x)^2/5)},{\x}) -- cycle;
    \draw[->] (-5.8,0) -- (6.4,0) node[right] {$x$};
    \draw[->] (0,-5.6) -- (0,5.9) node[above] {$y$};
    \draw (0,0) circle (5);
    \draw[line width=.8pt,domain=-4.714:4.714,samples=60] plot ({3-2*sqrt(1+(\x)^2/5)},{\x});
    \draw[dashed] (3,0) circle (2);
    \draw (4.778,1.475) -- (-.778,-4.939);
    \fill (2,-1.732) circle (1.6pt) node[below right,font=\footnotesize] {$\mathrm M$};
    \filldraw[fill=white,draw=black] (5,0) circle (2.4pt);
    \node[below left,font=\footnotesize] at (0,0) {$\mathrm O$};
    \node[above right,font=\footnotesize] at (5,0) {$(5,0)$};
    \node[font=\footnotesize] at (-2.6,5.3) {$\left(-\frac53,\frac{10\sqrt2}3\right)$};
  \end{tikzpicture}
  \end{zuhyo}
  点線の円は M の軌跡、実線の線分は弦 PQ の一例（$\theta=\frac\pi3$）。
  以上より、通過範囲は
  \[\Kotae{x^2+y^2\leqq25\ \text{かつ}\ \frac{(x-3)^2}4-\frac{y^2}5\leqq1}\]
  を満たす点から、点 $(5,\ 0)$ を除いたもの。図の灰色部分（境界を含む。ただし白丸の $(5,0)$ は除く）。

  \begin{memo}
    塗る側の確認：原点 $(0,0)$ を②に入れると $\frac94>1$ で範囲外、$(1,\ 0)$ は $1\leqq1$ で範囲内。簡単な点を1つ2つ入れて、塗る側が逆でないかを必ず確かめる。
  \end{memo}

  \begin{ptbs}{KEY}[通過範囲の3段階]
    \kakko{1} 動く図形を、1つのパラメータ（ここでは M）で表す。\kakko{2} 点 T を固定し、パラメータの存在条件にする。\kakko{3} 線分なら「直線上＋範囲内」、除外されるパラメータに対応する点を最後に処理する。
  \end{ptbs}

  \begin{simple}[出典：オリジナル]{参考問題40}
    $xy$ 平面上で、原点を中心とする半径2の円周上の異なる2点 P、Q が、線分 PQ の中点が直線 $x=1$ 上にあるように動く。線分 PQ が通過する範囲を求めよ。
  \end{simple}
  \Kai
  中点を M$(1,\ m)$ とおく。P$\neq$Q より M は円の内部にあり、$1+m^2<4$、$m^2<3$。直線 PQ は M を通り OM に垂直なので
  \[x+my=1+m^2\iff m^2-ym+(1-x)=0\]
  T$(x,y)$ が範囲に入る $\iff$ $x^2+y^2\leqq4$ かつ、この $m$ の2次方程式が $m^2<3$ の実数解をもつ。\par
  ここで、もし実数解 $m$ が $m^2\geqq3$ なら、直線と原点の距離は $\frac{1+m^2}{\sqrt{1+m^2}}=\sqrt{1+m^2}\geqq2$ で、直線は円板と高々1点（接点 $(1,\pm\sqrt3)$）でしか交わらない。接点 $(1,\ \pm\sqrt3)$ は $m=0$ の弦（$x=1$）の端点として通過されるので、結局「円板内で実数解をもつ」ことが条件になる。判別式 $y^2-4(1-x)\geqq0$ より、
  \[\Kotae{x^2+y^2\leqq4\ \text{かつ}\ x\geqq1-\frac{y^2}4}\]
  （円の内部で、放物線 $x=1-\frac{y^2}4$ の右側。境界を含む。）

  \begin{ptbs}{CHECK}[翌日、解説を隠して答える（答えは本講の中にある）]
    \kakko{1}　R の $z$ 座標が3に決まる理由は何か。\par
    \kakko{2}　M の軌跡から除く点と、その理由を言え。\par
    \kakko{3}　通過範囲を求める3段階を言え。
  \end{ptbs}

  \notefill
  \columnbreak
  \notefill
\end{multicols*}
%% ==================== 第41講 ====================
\Kou{41}{東大型総合(2) 整数と論証}
\begin{ascolorbox4A}[出典：オリジナル]{例題41}
  \kakko{1}　すべての整数 $n$ について、$n^5-n$ は30の倍数であることを証明せよ。\par
  \kakko{2}　$p$ を素数とする。$1\leqq k\leqq p-1$ を満たす整数 $k$ について、二項係数 ${}_p\mathrm C_k$ は $p$ の倍数であることを証明せよ。\par
  \kakko{3}　$p$ を素数とする。すべての自然数 $a$ について、$a^p-a$ は $p$ の倍数であることを証明せよ。\par
  \kakko{4}　$3^n+4^n$ が5の倍数となるような自然数 $n$ をすべて求めよ。
\end{ascolorbox4A}

\begin{Data}
  テーマ：素因数ごとに分けて示す、素数の性質と二項係数、帰納法との組合せ、余りの周期\par
  難易度：(1) C　(2) C　(3) D　(4) C\quad 前提：第8・10・11・23講\par
  この講のゴール：整数の論証を「分解する・余りで分類する・帰納法」の組合せで組み立てられる
\end{Data}

\begin{multicols*}{2}
  \Shomon{1}
  \ascboxZ{30 $=2\cdot3\cdot5$ に分けて、1つずつ示す}
  $2,\ 3,\ 5$ はどの2つも互いに素なので、$n^5-n$ が2、3、5のそれぞれの倍数であることを示せばよい。因数分解しておく。
  \[n^5-n=(n-1)n(n+1)(n^2+1)\]
  \textgt{2と3}：$(n-1)n(n+1)$ は連続する3整数の積なので、6の倍数（第10講）。\par
  \textgt{5}：$n$ を5で割った余りで分ける。\par
  $n\equiv0,\ 1,\ 4\pmod5$ のとき、$n,\ n-1,\ n+1$ のどれかが5の倍数。\par
  $n\equiv2,\ 3\pmod5$ のとき、$n^2\equiv4$ なので $n^2+1\equiv0$。\par
  いずれの場合も5の倍数。以上より30の倍数である。\QED

  \begin{memo}
    $n\equiv3$ は $n\equiv-2$ と同じなので、$n\equiv\pm2$ とまとめてもよい。余りを $0,\ \pm1,\ \pm2$ と対称にとると、2乗の計算が半分で済む。
  \end{memo}

  \Shomon{2}
  \ascboxZ{素数 $p$ は、$p$ 未満の数の積を割り切らない}
  \begin{koushiki}{素数の基本性質}
    $p$ が素数で、積 $ab$ が $p$ の倍数ならば、$a$ または $b$ が $p$ の倍数である。
  \end{koushiki}
  ${}_p\mathrm C_k=\dfrac{p!}{k!(p-k)!}$ を分母を払った形で書く。
  \[p!=k!\,(p-k)!\,{}_p\mathrm C_k\]
  左辺は $p$ の倍数。右辺の $k!\,(p-k)!$ は、$p$ より小さい正の整数（$1\leqq k\leqq p-1$ より、$k$ も $p-k$ も $p-1$ 以下）の積であり、どの因数も $p$ の倍数ではない。上の性質を繰り返し使うと、$k!\,(p-k)!$ は $p$ の倍数ではない。よって ${}_p\mathrm C_k$ が $p$ の倍数である。\QED

  \begin{memo}
    $p$ が素数でないと成り立たない。${}_4\mathrm C_2=6$ は4の倍数ではない。証明のどこで「素数」を使ったか（上の基本性質）を言えることが大切である。
  \end{memo}

  \Shomon{3}
  \ascboxZ{$a$ についての帰納法＋二項定理}
  $a$ についての数学的帰納法で示す。\par
  [Ⅰ]　$a=1$ のとき、$1-1=0$ は $p$ の倍数。\par
  [Ⅱ]　$a^p-a$ が $p$ の倍数と仮定する。二項定理より、
  \[(a+1)^p=a^p+\sum_{k=1}^{p-1}{}_p\mathrm C_ka^k+1\]
  よって、
  \[(a+1)^p-(a+1)=(a^p-a)+\sum_{k=1}^{p-1}{}_p\mathrm C_ka^k\]
  第1項は仮定より $p$ の倍数、$\Sigma$ の各項は\kakko{2}より $p$ の倍数。よって $(a+1)^p-(a+1)$ も $p$ の倍数。\par
  [Ⅰ][Ⅱ]より、すべての自然数 $a$ で成り立つ。\QED

  \begin{ptbs}{KEY}[前の小問は次の小問の部品]
    \kakko{2}の「${}_p\mathrm C_k$ は $p$ の倍数」は、\kakko{3}の帰納法の[Ⅱ]で二項展開の中間項を消すための部品である。東大の整数問題は、このように小問が積み上がる構成が多い。
  \end{ptbs}

  \begin{memo}
    \kakko{3}はフェルマーの小定理と呼ばれる。$p=5$ とすると $a^5-a$ は5の倍数となり、\kakko{1}の「5の倍数」の部分の別証明になる。
  \end{memo}

  \Shomon{4}
  \ascboxZ{余りは周期的にくり返す}
  5で割った余りで考える。$4\equiv-1\pmod5$ なので $4^n\equiv(-1)^n$。$3^n$ の余りは
  \[3,\ 4,\ 2,\ 1,\ 3,\ 4,\ \cdots\]
  と周期4でくり返す（$3^4=81\equiv1$ だから）。$n$ を4で割った余りで分ける。

  \begin{zuhyo}
  \begin{tblr}{
    colspec = {c c c c},
    hline{1,Z} = {0.08em}, hline{2} = {0.05em},
    row{1} = {font=\TRgt},
  }
    $n$ を4で割った余り & $3^n$ & $(-1)^n$ & 和 \\
    1 & 3 & $-1$ & 2 \\
    2 & 4 & 1 & $5\equiv0$ \\
    3 & 2 & $-1$ & 1 \\
    0 & 1 & 1 & 2 \\
  \end{tblr}
  \end{zuhyo}

  よって、5の倍数になるのは $n$ を4で割った余りが2のときに限る。
  \[\Kotae{n=4m+2\ \ (m=0,\ 1,\ 2,\ \cdots)}\]

  \begin{simple}[出典：オリジナル]{参考問題41}
    $n^2+n+1$ が7の倍数となるような整数 $n$ の条件を求めよ。
  \end{simple}
  \Kai
  $n$ を7で割った余り $r=0,\ 1,\ \cdots,\ 6$ について $r^2+r+1$ を7で割った余りを調べると、
  \[1,\ 3,\ 0,\ 6,\ 0,\ 3,\ 1\]
  （例えば $r=4$ のとき $21\equiv0$。）0になるのは $r=2,\ 4$。
  \[\Kotae{n\ \text{を7で割った余りが2または4}}\]
  \begin{memo}
    余りが $r$ と $6-r$ で同じ値になっている（1と1、3と3、0と0）。$n\to-1-n$ で $n^2+n+1$ が変わらないからである。表の対称性に気づくと、検算になる。
  \end{memo}

  \begin{ptbs}{CHECK}[翌日、解説を隠して答える（答えは本講の中にある）]
    \kakko{1}　30の倍数であることを示すとき、どう分けるか。\par
    \kakko{2}　${}_p\mathrm C_k$ が $p$ の倍数である証明のどこで素数を使うか。\par
    \kakko{3}　$3^n$ を5で割った余りが周期的になる理由は何か。
  \end{ptbs}

  \notefill
\end{multicols*}

%% ==================== 第42講 ====================
\Kou{42}{東大型総合(3) 評価と極限}
\begin{ascolorbox4A}[出典：東京大学 2003年度]{例題42A}
  円周率が3.05より大きいことを証明せよ。
\end{ascolorbox4A}

\begin{ascolorbox4A}[出典：オリジナル]{例題42B}
  $a_n=\displaystyle\int_0^1\frac{x^n}{1+x}\,dx$（$n=0,\ 1,\ 2,\ \cdots$）とする。\par
  \kakko{1}　$\displaystyle\lim_{n\to\infty}a_n=0$ を示せ。\qquad
  \kakko{2}　$a_n+a_{n+1}=\dfrac1{n+1}$ を示せ。\qquad
  \kakko{3}　$\displaystyle\lim_{n\to\infty}na_n=\frac12$ を示せ。\par
  \kakko{4}　$\displaystyle\lim_{n\to\infty}\sum_{k=1}^n\frac{(-1)^{k-1}}k=\log2$ を示せ。
\end{ascolorbox4A}

\begin{Data}
  テーマ：図形を使った評価、根号の評価、積分の漸化式、単調性とはさみうち、交代級数\par
  難易度：例題42A C　例題42B (1) B　(2) B　(3) D　(4) D\quad 前提：第7・30・36講\par
  この講のゴール：「値を直接は求められない量」を、計算できる量で上下からはさむ技術を完成させる
\end{Data}

\begin{multicols*}{2}
  \Shomon{A}
  \ascboxZ{円に内接する正多角形の周の長さと比べる}
  \textgt{方針}：半径1の円の周の長さは $2\pi$。内接する正多角形の周の長さは、円周より短い（各辺は弦で、対応する弧より短い。2点を結ぶ最短の線は線分だから）。周の長さが6.1を超える正多角形を見つければ、$2\pi>6.1$、つまり $\pi>3.05$ が言える。\par
  正六角形では周の長さが6で足りない（$\pi>3$ しか言えない）。辺を増やして\textgt{正十二角形}を使う。

  \begin{zuhyo}
  \begin{tikzpicture}[x=1.3cm,y=1.3cm,line width=.5pt]
    \draw (0,0) circle (1);
    \draw (0:1) \foreach \k in {1,...,12}{ -- ({30*\k}:1)};
    \draw[dashed] (0,0) -- (0:1); \draw[dashed] (0,0) -- (30:1);
    \draw (0:.3) arc (0:30:.3);
    \node[font=\footnotesize] at (15:.48) {$30^\circ$};
    \node[font=\footnotesize,below] at (.5,0) {$1$};
    \fill (0,0) circle (1pt);
  \end{tikzpicture}
  \end{zuhyo}

  \textgt{1辺の長さ}：中心角 $30^\circ$、2辺が1の二等辺三角形に余弦定理を使うと、1辺の長さ $s$ は
  \[s^2=1+1-2\cos30^\circ=2-\sqrt3\]
  \textgt{評価}：示したいのは $12s>6.1$、つまり $s^2>\left(\frac{6.1}{12}\right)^2=\frac{37.21}{144}$。\par
  $1.74^2=3.0276>3$ なので $\sqrt3<1.74$。よって
  \[s^2=2-\sqrt3>0.26=\frac{37.44}{144}>\frac{37.21}{144}\]
  したがって $12s>6.1$。円周は内接正十二角形の周より長いので、
  \[2\pi>12s>6.1\]
  ゆえに $\pi>3.05$ である。\QED

  \begin{ptbs}{KEY}[評価は「目標から逆算して、必要な精度を決める」]
    目標 $6.1$ から、$s^2$ が $0.2584$ を超えればよいと先に計算した。だから $\sqrt3$ は $1.74$ 程度の粗い評価で足りる。必要以上に細かい近似値を使わない。
  \end{ptbs}

  \begin{memo}
    正八角形でも示せる：1辺 $\sqrt{2-\sqrt2}$、周 $8\sqrt{2-\sqrt2}>6.1\iff2-\sqrt2>0.5814\cdots\iff\sqrt2<1.4185\cdots$。いずれにしても、根号の値の評価は「2乗して整数と比べる」で確かめる。
  \end{memo}

  \Shomon{1}
  \ascboxZ{分母を1で押さえて、はさむ}
  $0\leqq x\leqq1$ で $1\leqq1+x$ なので、$0\leqq\frac{x^n}{1+x}\leqq x^n$。積分して、
  \[0\leqq a_n\leqq\int_0^1x^n\,dx=\frac1{n+1}\]
  $n\to\infty$ で右辺は0なので、はさみうちより $\displaystyle\lim_{n\to\infty}a_n=0$。\QED

  \Shomon{2}
  \ascboxZ{足すと分母が消える}
  \[a_n+a_{n+1}=\int_0^1\frac{x^n(1+x)}{1+x}\,dx=\int_0^1x^n\,dx=\frac1{n+1}\]
  \QED

  \Shomon{3}
  \ascboxZ{単調性と\kakko{2}で上下からはさむ}
  $0\leqq x\leqq1$ で $x^{n+1}\leqq x^n$ なので $a_{n+1}\leqq a_n$（$\{a_n\}$ は減少）。\kakko{2}と組み合わせる。\par
  \textgt{下から}：$\frac1{n+1}=a_n+a_{n+1}\leqq2a_n$ より $a_n\geqq\frac1{2(n+1)}$。\par
  \textgt{上から}：\kakko{2}で $n$ を $n-1$ にすると（$n\geqq1$）、$\frac1n=a_{n-1}+a_n\geqq2a_n$ より $a_n\leqq\frac1{2n}$。\par
  よって
  \[\frac n{2(n+1)}\leqq na_n\leqq\frac12\]
  左辺は $n\to\infty$ で $\frac12$ に近づくので、はさみうちより $\displaystyle\lim_{n\to\infty}na_n=\frac12$。\QED

  \begin{memo}
    \kakko{1}では $a_n\leqq\frac1{n+1}$ という粗い評価で足りたが、\kakko{3}では $n$ 倍するので、2倍の精度が必要になる。「減少列なので、2つの和は片方の2倍ではさめる」という発想が鍵である。
  \end{memo}

  \Shomon{4}
  \ascboxZ{符号を交互に付けて足すと、途中が消える}
  \kakko{2}より $\frac1k=a_{k-1}+a_k$。符号 $(-1)^{k-1}$ を付けて $k=1$ から $n$ まで足すと、
  \[
  \begin{aligned}
  \sum_{k=1}^n\frac{(-1)^{k-1}}k&=(a_0+a_1)-(a_1+a_2)+\cdots\\
  &\qquad+(-1)^{n-1}(a_{n-1}+a_n)\\
  &=a_0+(-1)^{n-1}a_n
  \end{aligned}
  \]
  $a_0=\int_0^1\frac{dx}{1+x}=\log2$、\kakko{1}より $a_n\to0$ なので、
  \[\lim_{n\to\infty}\sum_{k=1}^n\frac{(-1)^{k-1}}k=\log2\]
  \QED

  \begin{ptbs}{KEY}[積分の漸化式は「級数の和」の道具]
    $1-\frac12+\frac13-\cdots$ のように直接は和が求まらない級数も、各項を積分で表すと、極限が積分の値として求まる。東大の積分と極限の融合問題の典型である。
  \end{ptbs}

  \begin{simple}[出典：オリジナル]{参考問題42}
    $\sqrt2$、$\sqrt3$ の近似値を用いずに、$\sqrt2+\sqrt3>3.14$ であることを証明せよ。
  \end{simple}
  \Kai
  両辺とも正なので、2乗して比べてよい。$(\sqrt2+\sqrt3)^2=5+2\sqrt6$、$3.14^2=9.8596$。示すことは $2\sqrt6>4.8596$、さらに両辺正なので2乗して $24>4.8596^2$。$4.8596<4.86$ で $4.86^2=23.6196<24$ なので成り立つ。\QED
  \begin{memo}
    「2乗して比べる」は、両辺が正（または0以上）のときにだけ同値変形になる。根号の大小比較では、毎回この確認を書く。
  \end{memo}

  \begin{ptbs}{CHECK}[翌日、解説を隠して答える（答えは本講の中にある）]
    \kakko{1}　円周率の評価で、正六角形では足りない理由は何か。\par
    \kakko{2}　評価で必要な精度を先に決める方法を言え。\par
    \kakko{3}　$na_n\to\frac12$ の証明で使った2つの事実は何か。
  \end{ptbs}

  \notefill
\end{multicols*}

%% ==================== 第43講 ====================
\Kou{43}{仕上げ：最終確認と本番戦略}
\begin{ascolorbox4A}[出典：オリジナル]{例題43　総点検テスト}
  次の\kakko{1}〜\kakko{12}を、1問3分以内で解け。\par
  \kakko{1}　$x+y=4,\ xy=2$ のとき $x^3+y^3$\qquad
  \kakko{2}　不等式 $|2x-3|<5$ を解け\qquad
  \kakko{3}　$x^2-kx+k+3=0$ が異なる2つの正の解をもつ $k$ の範囲\par
  \kakko{4}　異なる5人を、空の部屋ができないように2つの部屋 A、B に分ける方法の数\par
  \kakko{5}　さいころを3回投げるとき、出た目の最大値が5である確率\qquad
  \kakko{6}　$3^{50}$ を8で割った余り\par
  \kakko{7}　$a_1=2,\ a_{n+1}=2a_n-1$ の一般項\qquad
  \kakko{8}　$|\vec a|=1,\ |\vec b|=2,\ \vec a\cdot\vec b=1$ のとき $|\vec a+2\vec b|$\par
  \kakko{9}　$\left(\cos\dfrac\pi{12}+i\sin\dfrac\pi{12}\right)^{12}$\qquad
  \kakko{10}　$\displaystyle\lim_{x\to0}\frac{1-\cos2x}{x^2}$\qquad
  \kakko{11}　$\displaystyle\int_0^1xe^x\,dx$\par
  \kakko{12}　曲線 $y=x^2\ (0\leqq x\leqq1)$ と $x$ 軸、直線 $x=1$ で囲まれた部分を、$x$ 軸のまわりに1回転してできる立体の体積
\end{ascolorbox4A}

\begin{Data}
  テーマ：全範囲の総点検と、試験本番（150分・6題が目安）の時間の使い方\par
  難易度：A〜B（全問正解が目標）\quad 前提：第1〜42講\par
  この講のゴール：弱点の講を特定して復習の優先順位を決め、本番で得点を最大にする進め方を身につける
\end{Data}

\begin{multicols*}{2}
  \ascboxZ{答えと、間違えたときに戻る講}
  \begin{zuhyo}
  \begin{tblr}{
    colspec = {c c c},
    hline{1,Z} = {0.08em}, hline{2} = {0.05em},
    row{1} = {font=\TRgt},
    cells = {font=\small},
  }
    問題 & 答え & 戻る講 \\
    (1) & $40$ & 第1講 \\
    (2) & $-1<x<4$ & 第2講 \\
    (3) & $k>6$ & 第5講 \\
    (4) & $30$ 通り & 第8講 \\
    (5) & $\frac{61}{216}$ & 第9講 \\
    (6) & $1$ & 第10講 \\
    (7) & $a_n=2^{n-1}+1$ & 第22講 \\
    (8) & $\sqrt{21}$ & 第25講 \\
    (9) & $-1$ & 第27講 \\
    (10) & $2$ & 第31講 \\
    (11) & $1$ & 第35講 \\
    (12) & $\frac\pi5$ & 第37講 \\
  \end{tblr}
  \end{zuhyo}

  \ascboxZ{要点だけの解説}
  \kakko{1}　$(x+y)^3-3xy(x+y)=64-24$。\par
  \kakko{2}　$-5<2x-3<5$ より $-1<x<4$。\par
  \kakko{3}　判別式 $k^2-4k-12>0$、和 $k>0$、積 $k+3>0$ をすべて満たす（2解の和と積が正 $\iff$ 2実数解がともに正）。$k<-2,\ 6<k$ と合わせて $k>6$。\par
  \kakko{4}　1人ずつ A か B を選ぶ $2^5=32$ 通りから、全員が同じ部屋の2通りを除く。\par
  \kakko{5}　$\left(\frac56\right)^3-\left(\frac46\right)^3=\frac{125-64}{216}$。\par
  \kakko{6}　$3^2=9\equiv1\pmod8$ より $3^{50}=(3^2)^{25}\equiv1$。\par
  \kakko{7}　$a_{n+1}-1=2(a_n-1)$、$a_1-1=1$。\par
  \kakko{8}　$|\vec a+2\vec b|^2=1+4\cdot1+4\cdot4=21$。\par
  \kakko{9}　ド・モアブルの定理より $\cos\pi+i\sin\pi$。\par
  \kakko{10}　$\frac{1-\cos2x}{x^2}=\frac{2\sin^2x}{x^2}\to2$。\par
  \kakko{11}　$\Big[(x-1)e^x\Big]_0^1=0-(-1)$。\par
  \kakko{12}　$\pi\int_0^1x^4\,dx=\frac\pi5$。

  \begin{ptbs}{KEY}[間違えた問題は「講」に戻って例題から]
    総点検で間違えた問題は、答えを覚え直すのではなく、戻る講の例題を白紙から解き直す。例題が解ければ、同じ型の入試問題にも対応できる。
  \end{ptbs}

  \ascboxZ{本番の時間の使い方（目安）}
  東京大学（理科）の数学は、例年150分で6題という形式で行われてきた。試験時間・題数・配点は、受験する年度の募集要項で必ず確認すること。以下は150分・6題を前提とした、時間配分の一つの目安である。

  \begin{zuhyo}
  \begin{tblr}{
    colspec = {c X[l]}, width=\linewidth,
    hline{1,Z} = {0.08em}, hline{2} = {0.05em},
    row{1} = {font=\TRgt},
    cells = {font=\small},
  }
    時間 & すること \\
    最初の10分 & 全6題の問題文と小問を読み、方針が立つ順に順番を決める \\
    次の100分 & 1題あたり約25分を上限に、解ける問題から答案を完成させる \\
    次の25分 & 止まった問題の小問(1)・(2)だけを確実に書く \\
    最後の15分 & 見直し（条件の使い忘れ・場合分けの漏れ・範囲・計算） \\
  \end{tblr}
  \end{zuhyo}

  \begin{koushiki}{手が止まったときの5つの問い}
    \kakko{1} 使っていない条件はないか。\par
    \kakko{2} 小さい値（$n=1,\ 2,\ 3$ など）で実験したか。\par
    \kakko{3} 図をかいたか。対称性はないか。\par
    \kakko{4} 前の小問の結果を使えないか。\par
    \kakko{5} 求めるものを別の言葉（存在条件・個数・最大値）に言い換えられないか。
  \end{koushiki}

  \begin{memo}
    部分点の付け方の詳細は公表されていない。だからこそ、途中の論理を省かず、第39講の「答案の5つの型」に沿って書く。途中までしか解けなくても、方針と、そこまでに確定したことを正しく書いておくことが大切である。
  \end{memo}

  \ascboxZ{本書を終えた後の進め方}
  \textgt{第1段階}：全講の例題を、解説を隠して解き直す。各講で「最初の一手」が1行で言えない講に印を付け、その講だけ参考問題まで戻る。\par
  \textgt{第2段階}：東大の過去問に取り組む。1題ごとに、どの講の例題の考え方を使ったかを書き込む。結びつかない問題は、新しい型としてノートにまとめる。\par
  \textgt{第3段階}：本番と同じ150分で6題を解く演習を繰り返し、上の時間配分を自分に合うように調整する。

  \begin{memo}
    第0講で書いた3周法は、過去問でも同じである。解説を読んで「わかった」で終わらせず、翌日に白紙から再現し、数日後に類題で使う。この繰り返しが、合格に必要な「初見で方針を立てる力」になる。
  \end{memo}

  \begin{ptbs}{CHECK}[翌日、解説を隠して答える（答えは本講の中にある）]
    \kakko{1}　手が止まったときの5つの問いを言え。\par
    \kakko{2}　本番の最初の10分ですることは何か。\par
    \kakko{3}　本書を終えた後の3段階の進め方を言え。
  \end{ptbs}

  \notefill
  \columnbreak
  \ascboxZ{学習記録（3周法の日付を書く）}
  \begin{zuhyo}
  \begin{tblr}{
    colspec = {Q[c,wd=5mm] X[c] X[c] X[c] Q[c,wd=5mm] X[c] X[c] X[c]}, width=\linewidth,
    hlines, vlines, vline{5} = {1pt},
    row{1} = {font=\TRgt},
    cells = {font=\scriptsize},
    rowsep = 1.2pt, colsep = 1.5pt,
  }
    講 & 1周 & 2周 & 3周 & 講 & 1周 & 2周 & 3周 \\
    1 & & & & 23 & & & \\
    2 & & & & 24 & & & \\
    3 & & & & 25 & & & \\
    4 & & & & 26 & & & \\
    5 & & & & 27 & & & \\
    6 & & & & 28 & & & \\
    7 & & & & 29 & & & \\
    8 & & & & 30 & & & \\
    9 & & & & 31 & & & \\
    10 & & & & 32 & & & \\
    11 & & & & 33 & & & \\
    12 & & & & 34 & & & \\
    13 & & & & 35 & & & \\
    14 & & & & 36 & & & \\
    15 & & & & 37 & & & \\
    16 & & & & 38 & & & \\
    17 & & & & 39 & & & \\
    18 & & & & 40 & & & \\
    19 & & & & 41 & & & \\
    20 & & & & 42 & & & \\
    21 & & & & 43 & & & \\
    22 & & & &  & & & \\
  \end{tblr}
  \end{zuhyo}

  \Hikoushiki
  \notefill
\end{multicols*}

\end{document}
